% concatenated from arXivVersion_3.tex
\documentclass[11pt]{amsart}

% --- Core Encoding & Fonts ---
\usepackage[T1]{fontenc}
\usepackage[utf8]{inputenc}
\usepackage[mathscr]{eucal}
\usepackage{dsfont}
\usepackage{textcomp}

% --- Math Support ---
\usepackage{amssymb, amsfonts, amsthm}
\usepackage{mathtools} % Loads amsmath automatically

% --- Graphics & Colors ---
\usepackage[dvipsnames]{xcolor} % Loaded once, before TikZ
\usepackage{tikz}
\usetikzlibrary{positioning}

% --- Layout & Tools ---
\usepackage[a4paper, twoside=false, vmargin={2cm,3cm}, includehead]{geometry}
\usepackage{comment}
\usepackage{mdwlist}
\usepackage{enumitem} % Prefer enumitem over 'enumerate' for better control

% --- Hyperref (Load last, setup once) ---
\usepackage{hyperref}
\hypersetup{
    colorlinks=true,
    linktocpage=true,
    linkcolor=red,
    filecolor=blue,
    citecolor=blue,
    urlcolor=cyan,
}

% --- Theorem Environments ---
\newtheorem{lemma}{Lemma}[section]
\newtheorem{theorem}[lemma]{Theorem}
\newtheorem*{theorem*}{Theorem}
\newtheorem{corollary}[lemma]{Corollary}
\newtheorem{question}{Question}
\newtheorem{proposition}[lemma]{Proposition}
\newtheorem*{proposition*}{Proposition}
\newtheorem*{MCP}{Multiplicative correspondence principle}
\newtheorem{conjecture}[lemma]{Conjecture}
\newtheorem{problem}[lemma]{Problem}
\newtheorem*{problem*}{Problem}

\theoremstyle{definition}
\newtheorem{claim}{Claim}
\newtheorem*{claim*}{Claim}
\newtheorem*{convention}{Convention}
\newtheorem*{notation}{Notation}
\newtheorem{definition}[lemma]{Definition}
\newtheorem*{equivdefinition}{Equivalent Definition}
\newtheorem{example}[lemma]{Example}
\newtheorem{examples}[lemma]{Examples}
\newtheorem*{remark}{Remark}
\newtheorem*{remarks}{Remarks}

% --- Custom Commands & Operators ---

% Enumerations (Using enumitem syntax is better, but these overrides work)
\renewcommand{\theenumi}{\roman{enumi}}
\renewcommand{\labelenumi}{{\rm (\roman{enumi})}}
\renewcommand{\theenumii}{{\rm \alph{enumii}}}

% Blackboard letters
\DeclareMathOperator*{\oE}{\overline{\mathbb{E}}}
\DeclareMathOperator*{\E}{\mathbb{E}}
\newcommand{\lE}{\mathbb{E}^{\log}}
\newcommand{\C}{{\mathbb C}}
\newcommand{\D}{{\mathbb D}}
\newcommand{\F}{{\mathbb F}}
\newcommand{\N}{{\mathbb N}}
\newcommand{\NN}{{\mathcal N}}
\renewcommand{\P}{{\mathbb P}}
\newcommand{\Q}{{\mathbb Q}}
\newcommand{\R}{{\mathbb R}}
\renewcommand{\S}{\mathbb{S}}
\newcommand{\T}{{\mathbb T}}
\newcommand{\Z}{{\mathbb Z}}
\newcommand{\U}{{\mathbb U}}

% Calligraphic letters
\newcommand{\CA}{{\mathcal A}} \newcommand{\CB}{{\mathcal B}}
\newcommand{\CC}{{\mathcal C}} \newcommand{\CD}{{\mathcal D}}
\newcommand{\CE}{{\mathcal E}} \newcommand{\CF}{{\mathcal F}}
\newcommand{\CH}{{\mathcal H}} \newcommand{\CI}{{\mathcal I}}
\newcommand{\CK}{{\mathcal K}} \newcommand{\CM}{{\mathcal M}}
\newcommand{\CN}{{\mathcal N}} \newcommand{\CP}{{\mathcal P}}
\newcommand{\CQ}{{\mathcal Q}} \newcommand{\CX}{{\mathcal X}}
\newcommand{\CY}{{\mathcal Y}} \newcommand{\CV}{{\mathcal V}}
\newcommand{\CU}{{\mathcal U}} \newcommand{\CW}{{\mathcal W}}
\newcommand{\CZ}{{\mathcal Z}}
\newcommand{\FC}{{\mathfrak C}} \newcommand{\FD}{{\mathfrak D}}
\newcommand{\FL}{{\mathfrak L}} \newcommand{\FH}{{\mathfrak H}}
\newcommand{\FS}{{\mathfrak S}} \newcommand{\OK}{{\mathcal{O}_K}}

% Bold letters
\renewcommand{\a}{\mathbf{a}}
\newcommand{\ba}{\mathbf{a}}
\renewcommand{\b}{\mathbf{b}}
\newcommand{\bb}{\mathbf{b}}
\newcommand{\bc}{\mathbf{c}}
\newcommand{\bm}{\mathbf{m}}
\newcommand{\bn}{\mathbf{n}}
\newcommand{\be}{\mathbf{e}}
\newcommand{\bg}{\mathbf{g}}
\newcommand{\bh}{\mathbf{h}}
\newcommand{\bk}{\mathbf{k}}
\newcommand{\bp}{\mathbf{p}}
\newcommand{\bq}{\mathbf{q}}
\newcommand{\mf}{\mathfrak}

% Handle special accents
\let\ringaccent\r
\newcommand{\aaaccent}{\ringaccent a} 

\renewcommand{\r}{{\mathbf{r}}}
\newcommand{\br}{{\mathbf{r}}}
\newcommand{\bs}{{\mathbf{s}}}
\newcommand{\bt}{{\mathbf{t}}}
\newcommand{\bu}{{\mathbf{u}}}
\newcommand{\bv}{{\mathbf{v}}}
\newcommand{\x}{{\mathbf{x}}}
\newcommand{\bx}{{\mathbf{x}}}
\newcommand{\by}{{\mathbf{y}}}
\newcommand{\bz}{{\mathbf{z}}}
\newcommand{\bI}{{\mathbf{I}}}
\newcommand{\bM}{{\mathbf{M}}}

\newcommand{\ube}{{\underline{\mathbf{e}}}}
\newcommand{\ubg}{{\underline{\mathbf{g}}}}
\newcommand{\ubh}{{\underline{\mathbf{h}}}}
\newcommand{\ubk}{{\underline{\mathbf{k}}}}
\newcommand{\ubm}{{\underline{\mathbf{m}}}}
\newcommand{\ubr}{{\underline{\mathbf{r}}}}
\newcommand{\bN}{{\mathbf{N}}}

\newcommand{\balpha}{{\boldsymbol{\alpha}}}
\newcommand{\bepsilon}{{\boldsymbol{\epsilon}}}
\newcommand{\bgamma}{{\boldsymbol{\gamma}}}
\newcommand{\bdelta}{{\boldsymbol{\delta}}}
\newcommand{\bzero}{{\boldsymbol{0}}}
\newcommand{\btheta}{{\boldsymbol{\theta}}}
\newcommand{\beeta}{{\boldsymbol{\eta}}}
\newcommand{\bbeta}{{\boldsymbol{\beta}}}

\newcommand{\ub}{{\underline{b}}}
\newcommand{\ug}{{\underline{g}}}
\newcommand{\uh}{{\underline{h}}}
\newcommand{\uk}{{\underline{k}}}
\newcommand{\ul}{{\underline{l}}}
\newcommand{\un}{{\underline{n}}}
\newcommand{\um}{{\underline{m}}}
\newcommand{\ur}{{\underline{r}}}
\newcommand{\uu}{{\underline{u}}}
\newcommand{\uv}{{\underline{v}}}

% Various abbreviations
\newcommand{\ve}{\varepsilon}
\newcommand{\veps}{\varepsilon}
\newcommand{\wh}{\widehat}
\newcommand{\wt}{\widetilde}
\newcommand{\e}{\epsilon}
\newcommand{\one}{\mathbf{1}}
\newcommand{\eps}{\epsilon}
\newcommand{\ueps}{{\underline{\epsilon}}}
\newcommand{\ue}{{\underline{\epsilon}}}

\DeclareMathOperator{\supp}{supp}

% Brackets and Norms
\renewcommand{\angle}[1]{\mathopen{}\left\langle #1\mathclose{}\right\rangle}
\newcommand{\ZN}{\Z_{\widetilde N}}
\newcommand{\tN}{{\widetilde N}}
\newcommand{\tM}{{\widetilde M}}
\newcommand{\CCM}{{\CM_1^c}}
\newcommand{\tZN}{\Z_\tN}
\newcommand{\norm}[1]{\left\Vert #1\right\Vert}
\newcommand{\fnnorm}[1]{|\!|\!| #1|\!|\!|}
\newcommand{\nnorm}[1]{\left|\!\left|\!\left| #1\right|\!\right|\!\right|}
\newcommand{\bignnorm}[1]{\Big|\!\Big|\!\Big| #1\Big|\!\Big|\!\Big|}
\newcommand{\lip}{{\text{\rm Lip}}}
\newcommand{\inv}{^{-1}}
\newcommand{\st}{{\text{\rm st}}}
\newcommand{\er}{{\text{\rm er}}}

\DeclareMathOperator{\spec}{Spec}
\DeclareMathOperator{\id}{id}
\DeclareMathOperator{\Span}{Span}
\newcommand{\dmult}{d_\textrm{mult}}
\DeclareMathOperator{\reel}{Re}
\DeclareMathOperator{\Imag}{Im}
\renewcommand{\Im}{\Imag}
\DeclareMathOperator{\nonpol}{nonpol}
\DeclareMathOperator{\pol}{pol}
\DeclareMathOperator{\poly}{poly}
\DeclareMathOperator{\fracdeg}{frac\; deg}
\newcommand{\Krat}{{\CK_{\text{\rm rat}}}}

\newcommand{\abs}[1]{\mathopen{}\left| #1\mathclose{}\right|}
\newcommand{\bigabs}[1]{\bigl| #1 \bigr|}
\newcommand{\Bigabs}[1]{\Bigl| #1 \Bigr|}
\newcommand{\brac}[1]{\mathopen{}\left( #1 \mathclose{}\right)}
\newcommand{\bigbrac}[1]{\bigl( #1 \bigr)}
\newcommand{\Bigbrac}[1]{\Bigl( #1 \Bigr)}
\newcommand{\sfloor}[1]{{\lfloor #1 \rfloor}}
\newcommand{\floor}[1]{{\left \lfloor #1 \right \rfloor}}
\newcommand{\ceil}[1]{{\left \lceil #1 \right \rceil}}
\newcommand{\rem}[1]{\left \{ #1 \right \}}

% Author-specific colors
\newcommand{\FM}[1]{{\color{violet}{{#1}}}}
\newcommand{\WS}[1]{{\color{blue}{{#1}}}}
\newcommand{\SD}[1]{{\color{green!60!black}{{#1}}}}
\newcommand{\AK}[1]{{\color{ForestGreen}{{#1}}}}

\author{Sebasti\'an Donoso, Andreu Ferr\'e Moragues, Andreas Koutsogiannis, and Wenbo Sun}

\address[Sebasti{\'a}n Donoso]{Departamento de Ingenier\'{\i}a Matem\'atica and Centro de Modelamiento Matem{\'a}tico, Universidad de Chile \& IRL 2807 - CNRS, Beauchef 851, Santiago, Chile} \email{sdonoso@dim.uchile.cl}
\address[Andreu Ferr\'e Moragues]{Department of Mathematical Analysis and Applied Mathematics, Faculty of Mathematics, Complutense University of Madrid, 28040 Madrid, Spain
}
\email{anferre@ucm.es}

\address[Andreas Koutsogiannis]{
Department of Mathematics, Aristotle University of Thessaloniki, Thessaloniki 54124, Greece}
\email{akoutsogiannis@math.auth.gr}

\address[Wenbo Sun]{Department of Mathematics, Virginia Tech, 225 Stanger Street, Blacksburg, VA, 24061, USA}
\email{swenbo@vt.edu}


\thanks{The first author was partially funded by ANID/Fondecyt/1241346 and Centro de Modelamiento Matemático (CMM) FB210005, BASAL funds for centers of excellence from ANID-Chile. This project was implemented in the framework “3rd Call for H.F.R.I.’s Research Projects to Support Faculty Members \& Researchers” (H.F.R.I. Project Number: 24979). The fourth author was supported by the NSF Grant DMS-2247331.}

\subjclass[2020]{Primary: 05D10, Secondary:11N37, 11B30, 37A44}

\keywords{Partition regularity, multiplicative functions, number fields, pythagorean equations.} 
\title{Partition regularity in imaginary quadratic rings of integers}
%\date{January 2025}

\begin{document}
\begin{abstract}
%We prove a version of Hal\'asz's theorem for number fields. This, together with some newly developed concentration estimates \AK{``developed by'' OR ``following the work of FKM''},  allows us to obtain partition regularity results for certain homogeneous quadratic equations over quadratic imaginary rings of integers.

 
We obtain partition regularity results for homogeneous quadratic equations whose parametrized solutions admit nice factorizations into linear forms over rings of integers of imaginary quadratic fields. To do so, we develop number-theoretic results of independent interest on such fields, such as  a characterization for aperiodic completely multiplicative functions, the Tur\'an-Kubilius inequality, and a new concentration estimate for multiplicative functions. 


 



%\SD{We obtain partition regularity results for certain homogeneous quadratic equations over rings of integers of imaginary quadratic fields. To do so, we develop number-theoretic results of independent interest, such as a version of Halász's theorem for number fields and new concentration estimates for multiplicative functions.}
\end{abstract}
\maketitle
\tableofcontents

%\section{Introduction and main results}

%\subsection{TODO LIST}
%\begin{enumerate}
    %\item Cleanup the appendix and merge it into only one section at the end. Needs: proof of $B_N$ is a good set and fixing the equivalence of APs at the end.
 %   \item Double check that the proofs in the appendix are fine.
 %   \item Clean up the number theory background section: point of the section should be to give some of the basic number theory stuff that we need and also to prove the extension to ideal valued functions that we will use in the proof of our version of Hal\'asz's theorem.
 %   \item Put in remarks that explain how the Gaussian integer (more generally UFD) case(s) is easier and why some of the things we do would simplify there.
  %  \item Adapt the proofs of concentration estimates to the new averaging scheme we will need to use.
  %  \item Fix the pythagorean pairs for square coefficients section (the proofs) to take into account the new averaging scheme.
 %   \item Future directions: the real case and quadratic concentration estimates for rings of integers. Obstacles, conjectures and potential approaches for both of these questions -> see the Future discussions section.
  %  \item Clear all the questions that are in colors and proofread
  %  \item Write an introduction: mostly historical interest, context of the problem, what we do, extensions, obstacles, conjectures, sell the paper a bit, etc. - Also fill in the Notation subsection.
%    \item Write an abstract.
 %   \item References.


%Give examples of generalization of Prop 10.4 in [Sun1]

%Consider the examples $(a,b,c,d,e,f)=(k^2b,b,1,(k+1)^2f^2-5b,kf,f)$ for $k \in \N\setminus \{1\}$, $b, f \in \Z$ arbitrary. (We simplified and took $c=1$ to get many examples easily).

%Remark about how if we insist on only using $x^2,y^2,z^2$, a ``diagonal'' quadratic equation, then, necessarily linear factorization of all terms implies Gaussian integers (it is a very special case!)
%\[
%\begin{cases}
%    x=(m+\ell_1n)(m+\ell_2n) \\
%    y= mn \\
%    z=(m+\ell_3n)(m+\ell_4n).
%\end{cases}
%\]
% This is why we go to the given examples above.

%Give the example $x^2+y^2+z^2=0$ as an application in $\Z[i]$. Full partition regularity for all possible pairs. Similarly for $a,b,c \in \pm \square$ in $\Z[i]$.

%Finally, also give the examples of the form $ax^2+dby^2=cz^2$, where $a,b,c$ are squares in $\Z[\tau_d]$ and we only get partition regularity for one pair, because for the other one we would need quadratic concentration estimates: due to length this is a potential future project.

%Part 2 of the intro: discuss how Hal\'asz's theorem is an important ingredient in our proofs, what it is for $\Z$, how we extend it to $\Z[i]$ and other UFDs, and the approach we take for general $\mathcal{O}$.


%\end{enumerate}
%\subsection{Provisional stuff - probably to be deleted}
%Questions to think about:

%(1) Extend the approach of FKM2 to show that $ax^{2}+by^{2}=cz^{2}$ is partition regular over $\mathbb{Z}[\sqrt{ac}]$. 

%{\color{red}{This is more or less done here. So, for example, in $\mathbb{Z}[i]$ we can find monochromatic pairs of the equation $x^2+y^2=-z^2,$ which doesn't make sense in $\mathbb{Z}.$ Also, by working in $\mathbb{Z}[\sqrt{2}]$ which has some additional issues to deal with, we can say something about $x^2+y^2=2z^2$ unconditionally. FKM2 deals with this under conjecture 2.}}

%For the problem, the main difficulty is the concentration result for other number fields (see Proposition 2.11 of FKM)

%(2) Check the open questions left from FKM2 on partition regularity over $\mathbb{Z}$.

%For this problem, the main difficulty is to show that the average converges to 0 if one of the multiplicative functions is aperiodic. In the cases where FKM2 left open, a major issue is that we don't know how to control the average under consideration by some Gowers norms.

%The parametrization of the equation $ax^2+by^2=cz^2$ that we would need would follow from the one used in [FKM2]. Namely, we put
%\[
%x=k c(am^2-bn^2), \quad y=k2acmn, \quad z = k\sqrt{ac}(am^2+bn^2),
%\]
%and we take advantage of the fact that in our number field, $ac$ is a perfect square without any conditions on the pythagorean equation we started with.

%We could also deal with the case where $bc$ is a square from [FKM2, Theorem 1.2] if we go to the number field $\Z[\sqrt{bc}]$, again unconditionally.

%Next, we use the spectral theorem to translate the measure question into a question of multiplicative functions. We are interested in the averages


%\[
%\liminf_{N \to \infty} \frac{1}{N^4}\sum_{n, m \in \Z[\sqrt{ac}] \cap [-N,N]^2} \mu( T_{kc(am^2-bn^2)}A \cap T_{k2acmn}A)>0
%\]
%Using the spectral theorem, this is going to become (after cancelling $T_k$ as it preserves the measure) into 
%\begin{equation}\label{345}
%    \liminf_{N \to \infty} \frac{1}{N^4}\sum_{1 \leq n, m \in \Z[\sqrt{ac}] \cap [-N,N]^2} \int_{\mathcal{M}} f(c(am^2-bn^2))\overline{f(2acmn)} \ d\sigma(f)>0,
%\end{equation}
%where $\mathcal{M}:= \{ f: \Z[\sqrt{ac}] \to \mathbb{S}^1 : f(xy) = f(x)f(y), \ x, y \in \Z[\sqrt{ac}]\}$.


%{\color{red} the map $f\mapsto\limsup_{N\to\infty}\Vert 1_{B_{N}}f\Vert_{U^{s}([N]^{2})}$ is not a norm, but only a seminorm. Does the open ampping theorem apply?}
%Idea: Split $\mathcal{M}$ into ``aperiodic'', ``pretentious'' (need to define equivalent of $f \sim n^{it}\chi$) and then the atom at $\{1\}.$

%Idea 2: Pretentious functions on $\Z[\sqrt{ac}]$ may be defined via $f \sim n^{it_1}m^{it_2}\chi,$ where $\chi$ takes values on roots of unity and $z=n+m\sqrt{ac}$.






%Idea 3: there should be a definition of pretentious by using Hecke characters. That is, $f$ is pretentious if it has a finite {\em pretentious distance} to a character. 


%NEED: To provide an adequate translation of the following result from number theory. Namely, [FKM, Lemma 3.5] which states that a multiplicative function is aperiodic if and only if it is non-pretentious.

%To have such a clasification for a number field $\mathcal{O}_{K}$ with $D$-dimensions. A function is \emph{aperiodic} if its average along any $D$-dimensional arithmetic progression is zero. Let $f$ be a multiplicative function on $\mathcal{O}_{K}$. Then by Theorem 1.5 of [Sun2], we have that
%$$\mathbb{E}_{n_{1},\dots,n_{D}}f(n_{1},\dots,n_{D})e^{2\pi i(\alpha_{1}n_{1}+\dots+\alpha_{D}n_{D})}=0$$
%if one of $\alpha_{1},\dots,\alpha_{D}$ is irrational. What we need to check is what the limit looks like when all of $\alpha_{1},\dots,\alpha_{D}$ are rational. This will lead to the definition of pretentious functions in our context. To do so, check Theorem 2 in DaboussiDelange.

%We would like to define a distance function for multiplicative functions on $\Z[\sqrt{ac}]$. Could the following be useful?
%\[
%\mathbb{D}(f,g):= \sum_{p \text{ prime in } \Z[\sqrt{ac}]} \frac{1}{|N(p)|}(1-\text{Re}(f(p)\overline{g(p)})).
%\]
%Comment for what we discussed Wenbo: I hope this definition might work for the following reason: even if our ring is no longer a UFD, you can still factor each element into primes, so this distance, for functions that take values in the unit circle, should still be telling us, with weights, how far the two functions are when evaluated at the primes, which is why I suspect it may still be relevant for us, even though uniqueness of factorization is lost.

%\medskip

%Define the notions of ``pretentious'' functions and, their complement, ``aperiodic'' functions. Here, we have to find a reasonable proof of [Lemma 3.5, FKM].

%Hálasz's theorem is proved in Section 6 of Elliot's book from 79 in the bibliography file. Need to check what can be adapted for number fields in order to define aperiodic and pretentious functions so that we can have the dychotomy we desire between these two notions (splitting the Borel set into two disjoint parts we can deal with to control the averages or dilated averages).
%\medskip

%To deal with the pretentious part, find the analogue of the ``concentration estimate'' (this expected to be hard as you need some analytic tricks). 

%\medskip

%If we are in the aperiodic case, we are expecting to be easier because of known Gowers-seminorms bounds. 

%Let us briefly describe the proof strategy in the paper. We aim to prove a partition regularity result for pairs of elements that satisfy a Pythagorean equation. Our starting point is the Gaussian integers because they enjoy many properties that are very useful to us: they are a UFD, they have only four units, they are a ring of integers.

%The proof strategy is as follows. First, instead of proving a partition regularity statement, we are actually able to show a stronger density result (using multiplicative density). A quick use of Furstenberg's correspondence theorem transforms the statement to prove into a recurrence property for averages of certain shifted sets under multiplicative actions.

%One more maneuver is then performed: the spectral theorem transforms the statement one more time so it suffices to look at averages of certain multiplicative functions evaluated at certain binary quadratic forms.

%After this reduction, we split the problem into two main parts: functions that are ``random'', called aperiodic, for which the averages will be $0$ because we can factor the polynomials of the parametrization into linear terms. This makes it so we can apply a Theorem from Sun2 which guarantees that said averages converge to $0$ provided the function is aperiodic.

%Next, we need to deal with the complement of aperiodic functions; (for us, they will become pretentious functions). Now, in order to be able to do that, we first need to understand what this complement is. This is where it is crucial that we develop an analog of Hal\'asz's theorem for Gaussian integers: it allows us to see that the non-aperiodic functions are non-other than those that correlate with a special type of multiplicative function: these are of the form $u \mapsto \NN(u)^{i\tau}$, for $\tau \in \R$.

%Once this is properly understood, we deal with this part developing some new linear concentration estimates for this type of multiplicative functions. Using some weights is also necessary, but this is fine.

%The fact that the aforementioned sets are measurable, I believe could also be done in a different way (as in [FKM2]), but in order to understand and treat the pretentious part properly, there is always some Hal\'asz in the background, it seems very hard to avoid.


%From what I have looked at FKM2, it seems possible, but definitely non-trivial to develop quadratic estimates (at least for Gaussian integers). This, together with some conjecture would allow one to deal with much more complicated Pyth equations that don't factor nicely.


%Finally, as for other rings of integers, it is definitely possible that some of the results we have can be translated, but to me the biggest roadblock is Hal\'asz's theorem. This is what is requiring UFD and finitely many units to do. Perhaps there's a way to get around finitely many units, but I do think UFD is being used non-trivially in several parts of the proof. That already reduces the ability for us to extend the result in this direction.


%It's also not clear what the analog of $u \mapsto \NN(u)^{i\tau}$ should be for other fields of integers (especially if the norm is non-negative). My interpretation of the $n \mapsto n^{it}$ map is that it transforms $(\N,\times)$ into an additive group via $n \mapsto \log n$, and then uses composition with characters $t \mapsto e^{it}$ to create multiplicative functions whose averages do not converge. How to do this in non-imaginary quadratic number fields is also an issue.




%\subsection{Partition regularity of Pythagorean pairs}
%The main result we will obtain is the following partition regularity result for Pythagorean pairs, which was missed in \cite{S23}, but that we can obtain in the next:
%\begin{theorem}\label{mainthm1}
%Let $a, b, c \in \Z[i]^{\times}$ be squares. Then, for every finite coloring of $\Z[i]^{\times}$, we can find
%\begin{itemize}
%    \item[(i)] distinct $x, y \in \Z[i]^{\times}$ with the same color, and $z \in \Z[i]^{\times}$ such that $ax^2+by^2=cz^2$.
%      \item[(ii)] distinct $y, z \in \Z[i]^{\times}$ with the same color, and $x \in \Z[i]^{\times}$ such that $ax^2+by^2=cz^2$.
%\end{itemize}
%\end{theorem}

%We will, in fact, prove a stronger version than %Theorem~\ref{mainthm1} 
%using densities, instead of colorings.


%Indeed, in order to fully take advantage of the invariance under dilations (due to homogeneity) of the standard Pythagorean equation: $x^2+y^2=z^2$, we will make use of the following notion of density that is invariant under dilations.

%Recall that a \emph{multiplicative F\o lner sequence in } $\Z[i]^{\times}$ is a sequence $\Phi:=(\Phi_K)_{K \in \N}$ of finite subsets of $\Z[i]^{\times}$ such that
%\[
%\lim_{K \to \infty} \frac{|\Phi_K \cap (x \cdot \Phi_K)|}{|\Phi_K|}=1, \text{ for all } x \in \Z[i]^{\times}.
%\]
%A prototypical example of such a F\o lner sequence is that of \eqref{folnerMult}. With each F\o lner sequence comes a notion of largeness for subsets $\Lambda \subset \Z[i]^{\times}$ called \emph{multiplicative upper density} with respect to $\Phi$, defined via
%\[
%\overline{d}_{\Phi}(\Lambda) := \limsup_{K \to \infty} \frac{|\Lambda \cap \Phi_K|}{|\Phi_K|},
%\]
%where the upper bar is dropped if the limit actually exists. We say that $\Lambda \subset \Z[i]^{\times}$ has positive multiplicative upper Banach density if there exists some multiplicative F\o lner sequence $\Phi=(\Phi_K)_{K \in \N}$ along which $\overline{d}_{\Phi}(\Lambda)>0$. In view of the finite subadditivity of such densities, the next result is, in fact, stronger than Theorem~\ref{mainthm1}:
%\begin{theorem}\label{mainthm2}
%Let $a, b, c \in \Z[i]^{\times}$ be squares. Then, for every $\Lambda \subset \Z[i]^{\times}$ with positive multiplicative density, we can find
%\begin{itemize}
%    \item[(i)] distinct $x, y \in \Lambda$ such that $ax^2+by^2=cz^2$.
%    \item[(ii)] distinct $y, z \in \Lambda$ such that $ax^2+by^2=cz^2$.
%\end{itemize}
%\end{theorem}


%\subsection{Parametric reformulation}
%In order to be able to apply the tools we have from Ergodic Theory, we need to write the solutions to the generalized Pythagorean equation
%\begin{equation}\label{pythequation}
%ax^2+by^2=cz^2
%\end{equation}
%using a convenient parametrization. By assumption, $a=a_0^2$, $b=b_0^2$ and $c=c_0^2$, for some $a,b,c, \in \Z[i]^{\times}$. We will utilize the following parametrization::
%\[
%x= k\ell_1(m^2-n^2), \quad  y=k\ell_2mn, \quad z= k\ell_3(m^2+n^2), \quad m, n \in \Z[i]^{\times},
%\]
%where $\ell_1=a_0bc$, $\ell_2=ab_0c$ and $\ell_3=abc_0$. Therefore, Theorem~\ref{mainthm2} can be established via the following
%\begin{theorem}\label{mainthm3}
%Suppose that $\Lambda \subset \Z[i]^{\times}$ is such that $\overline{d}_{\Phi}(\Lambda)>0$ for some multiplicative F\o lner sequence $(\Phi)$. Then, for every $\ell,\ell' \in \Z[i]^{\times}$ there exist
%\begin{itemize}
%    \item[(i)] Elements $m, n \in \Z[i]^{\times}$ with $\NN(m)\neq\NN(n)$ such that $\ell(m^2-n^2)$ and $\ell'mn$ are distinct and
%    \[
%    \overline{d}_{\Phi}((\ell(m^2-n^2))^{-1}\Lambda \cap (\ell'mn)^{-1} \Lambda)>0.
%    \]
%    \item[(ii)]  Elements $m, n \in \Z[i]^{\times}$ with $\NN(m)\neq\NN(n)$ such that $\ell(m^2+n^2)$ and $\ell' mn$ are distinct and
%    \[
%    \overline{d}_{\Phi}((\ell(m^2+n^2))^{-1}\Lambda \cap (\ell'mn)^{-1} \Lambda)>0.
%    \]
%\end{itemize}
%\end{theorem}
%\begin{remark}
%Given that $2(m^2+n^2)=(m+n)^2+(m-n)^2$ and $4mn=(m+n)^2-(m-n)^2$, applying (ii) with $2\ell$ instead of $\ell$ and $4\ell'$ instead of $\ell'$ we can get
%\begin{itemize}
%    \item[(iii)] Elements $m,n \in \Z[i]^{\times}$ with $\NN(m)\neq\NN(n)$, such that $\ell(m^2+n^2)$ and $\ell'(m^2-n^2)$ are distinct and
%    \[
%    \overline{d}_{\Phi}((\ell(m^2+n^2))^{-1}\Lambda \cap (\ell'(m^2-n^2))^{-1} \Lambda)>0.
%    \]
%\end{itemize}
%\end{remark}



\section{Introduction}\label{s1}
\subsection{Partition and density regularity}
A central problem in Ramsey theory is to determine when the existence of solutions of a polynomial equation is preserved under finitely many partitions. For example, let $q(x,y,z)=0$ be a polynomial equation with 3 variables. Is it true that for any finite coloring of the domain $D$ of the polynomial $q$, there exist $x,y,z$ of the same color with $q(x,y,z)=0$? We say that the equation $q(x,y,z)=0$ is \emph{partition regular over $D$} if the answer to this question is affirmative. A famous example due to Schur \cite{Schur} is that the linear equation $x+y=z$ is partition regular over $\N$. This result was extended by Rado in \cite{Rado} where he characterized all the linear equations that are partition regular.

On the other hand, the partition regularity problem for non-linear equations is significantly harder. In \cite{Gr07} and \cite{Gr08}, Erd\H{o}s and Graham asked whether the Pythagorean equation $x^{2}+y^{2}=z^{2}$ is partition regular over $\N$. 
Although this question remains open, significant progress has been made in the last decade. Specifically, for the Pythagorean equation and variations of it, one can find pairs from the set $\{x,y,z\}$ that are of the same color. We formalize this relaxation of the problem with the following definition. 

\begin{definition}
    Let $D \subset \C$ be a set and $q\colon D^{3}\to \C$ be a function. We say that $q(x,y,z)=0$ is \emph{partition regular over $D$ in $(x,y)$} if for every finite coloring of $D$, there exist distinct nonzero $x,y,z\in D$ with $x,y$ having the same color  such that $q(x,y,z)=0$. Partition regularity over $D$ in  $(x,z)$ and $(y,z)$ is defined analogously.
\end{definition} %\SD{should we say $D\subseteq \C$, as we use that and here we talk about $x,y,z$ being equal to 0.} 
Pioneer work achieving quadratic partition regularity results includes that of S\'ark\"{o}zy in \cite{sarkozy_1978a}, who showed that the equation $x-y=z^{2}$ is partition regular over $\N$ in $(x,y)$, as well as the work of Khalfalah and Szemer\'{e}di (see \cite{KS06}), who showed that the equation $x+y=z^{2}$ is partition regular over $\N$ in  $(x,y)$. 

Using a decomposition result for multiplicative functions in terms of Gowers norms (introduced in \cite{Go01}, and which have been successfully employed in the study of patterns and very general partition regularity results, e.g., \cite{GT08b}, \cite{BMR20}, \cite{Fr21}, \cite{Fr22}, \cite{FHo17}, among many others), Frantzikinakis and Host showed in \cite{FHo17} that given $a,b,c\in\Z$, if $\sqrt{-ac}, \sqrt{-bc}, \sqrt{-(a+b)c}\in\Z$, then the equation $ax^{2}+by^{2}+cz^{2}=0$ is partition regular over $\N$ in $(x,y)$ (such an example is the equation $9x^{2}+16y^{2}-z^{2}=0$). 
This result was extended in \cite{S18, S23}, where it was shown that if $\sqrt{-ac}, \sqrt{-bc}, \sqrt{-(a+b)c}$ belong to some number field $K$, then
the equation $ax^{2}+by^{2}+cz^{2}=0$ is partition regular over $\mathcal{O}_{K}$ (the ring of integers of $K$) in $(x,y)$  (for example $x^{2}+y^{2}-z^{2}=0$ is partition regular over $\Z[\sqrt{2}]$ in $(x,y)$, and is partition regular over $\Z[i]$ in $(x,z)$ and $(y,z)$).

The result of \cite{FHo17} was improved recently by Frantzikinakis, Klurman and Moreira \cite{FKM1} where the original assumption on the square roots was weakened to $\sqrt{-ac}, \sqrt{-bc}\in\Z$. This was further improved to $\sqrt{-ac}$ or $\sqrt{-bc}\in\Z$ by the same authors in \cite{FKM2}. As a special case, they showed that $x^{2}+y^{2}-z^{2}=0$ is partition regular over $\N$ in any two of the variables. One of the key innovations of \cite{FKM2,FKM1} is that the authors replaced the usage of a decomposition result for multiplicative functions introduced in \cite{FHo17} by a convenient splitting of the set of multiplicative functions into aperiodic and non-aperiodic, and then studied each of the components separately (as opposed to the previous approach, which decomposed all functions simultaneously).
Indeed, using this new methodology, it was shown in \cite{FKM1} that $x^{2}+y^{2}-z^{2}=0$ is also partition regular over $\N$ in $(x,z)$ and $(y,z)$. It is important to highlight that, in order to do so, a quadratic concentration estimate for the non-aperiodic functions was developed and used. %\AK{Don't they have this in [4] too? If yes, we have to restate things.} \FM{[4] is the second paper, so even if it follows, I think it's best to stick to the original source.}
%\SD{this last sentence feels disconnected} 



In this paper, we extend the splitting approach in \cite{FKM2,FKM1} in the setting of imaginary quadratic number fields (i.e., $K=\Q(\sqrt{-d})$ for some squarefree $d\in\N$) and  provide novel applications to various partition regularity problems, extending many results from \cite{FHo17,FKM2,FKM1,S18,S23}. 
  In this setting, the ring of integers $\OK$ of $K$ is given by $\OK=\Z[\tau_{d}]$, where $\tau_{d}=\sqrt{-d}$ if $d \equiv 2,3 \pmod 4$, and $\tau_d=\frac{1+\sqrt{-d}}{2}$ if $d \equiv 1 \pmod 4$.
  The next result is an extension we obtain of the aforementioned partition result \cite[Corollary~1.3]{FKM1} to imaginary quadratic number fields.

%\subsubsection{Partition reularity over Gaussian integers}

 



\begin{theorem}\label{th1}
    Let  $d\in\N$ be squarefree. For $a,b,c\in\Z\backslash\{0\}$, if $\sqrt{-ac}, \sqrt{-bc}\in \Z[\tau_{d}]$, then the equation $ax^{2}+by^{2}+cz^{2}=0$ is partition regular over $\Z[\tau_{d}]$ in $(x,y)$.
\end{theorem}





In fact, we obtain a density version of Theorem \ref{th1}.
Recall that a \emph{multiplicative F\o lner sequence in} $\OK$ is a sequence of finite subsets of $\mathcal{O}_K^{\times}$ such that
\[
\lim_{M \to \infty} \frac{|\Phi_M \cap (x^{-1} \cdot \Phi_M)|}{|\Phi_M|} = 1, \text{ for all } x \in \mathcal{O}_K^{\times}.
\]
%As was discussed, we will work with rings $R$ that are of the form $\OK$ for some given number field $K$. 
The precise F\o lner sequences that we will use (cf. \eqref{folnerMult}) are a bit intricate to define, so for now, the reader can assume that F\o lner sequences look like the following example when $\OK=\Z$, which contains all the main ingredients for the generalization that appears in \eqref{folnerMult}:
\[
\Phi_M:=\{p_1^{a_1}\dots p_k^{a_k} : M \leq a_i \leq 2M, 1 \leq i \leq M\}.
\]
where $(p_i)_{i\in \N}$ are the prime numbers. 

For each F\o lner sequence we have a notion of largeness for subsets $\Lambda \subseteq R^{\times}$, defined via
\begin{equation}\label{eq: largeness}
\overline{d}_{\Phi}(\Lambda):=\limsup_{M \to \infty} \frac{|\Lambda \cap \Phi_M|}{|\Phi_M|},
\end{equation}
where the upper bar is dropped if the limit does actually exist. We say that $\Lambda \subseteq \mathcal{O}_K^{\times}$ has positive \emph{multiplicative upper Banach density} if there exists some multiplicative F\o lner sequence $\Phi=(\Phi_M)_{M}$ along which we have $\overline{d}_{\Phi}(\Lambda)>0$. 

%The density defined in \eqref{eq: largeness} is the notion of largeness we will use in Definition~\ref{def: density reg}, which allows us to formulate the stronger versions of Theorems~\ref{th1}, \ref{th2} in terms of density regularity.


%stronger conlcusion in the sense of the following definition:
\begin{definition}[Density regularity]\label{def: density reg}
    Let $q\colon D^{3}\to \C$ be a function. We say that $q(x,y,z)=0$ is \emph{density regular over $D$ in $(x,y)$} if for every subset $\Lambda \subseteq D$ with positive multiplicative upper Banach density, there exist distinct nonzero $x,y,z\in D$ with $x,y\in \Lambda$ such that $q(x,y,z)=0$. Density regularity over $D$ in $(x,z)$ and $(y,z)$ is defined analogously.
\end{definition}

 


We have the following result, which is stronger than (and easily implies) Theorem \ref{th1}, given the subadditivity of the density we just introduced.

\begin{theorem}\label{th1density}
    Let  $d\in\N$ be squarefree. For $a,b,c\in\Z\backslash\{0\}$, if $\sqrt{-ac}, \sqrt{-bc}\in \Z[\tau_{d}]$, then the equation $ax^{2}+by^{2}+cz^{2}=0$ is density regular over $\Z[\tau_{d}]$ in $(x,y)$.
\end{theorem}

It is important to note that more standard additive versions of density will not be helpful to deal with the equation $x^2+y^2=z^2$, since if, for example, we look at $\Lambda=\{ u \in \Z[i]^{\times} : u \equiv 1 \pmod{2}\}$, then having $x, y \in \Lambda$, implies that $x^2+y^2$ can never be a square (we can still look at it modulo $4$ and argue similarly as in the case of $\N$).
It turns out that the more convenient form of density will be one that is multiplicative instead of additive, given the nature of the quadratic equations under consideration. 



%Indeed, so that we can fully take advantage of the invariance under dilations (due to homogeneity) of the quadratic equations that we will consider, we will make use of the following notion of density (which is invariant under dilations):


We remark that Theorems \ref{th1} and \ref{th1density} improve the results in \cite{S18,S23} by dropping the requirement that $\sqrt{-(a+b)c}\in \Z[\tau_{d}]$. 
For example, Theorem \ref{th1} implies that if $d\in\N$ is squarefree, then the equation $x^{2}+dy^{2}-z^{2}=0$ is partition regular over the ring of integers of $\Z[\tau_d]$. This result was previously known only over the ring of integers coming from the number field $\Q(\sqrt{-d},\sqrt{-(d+1)})$ (see \cite{S23}).

 
Our method applies to not only quadratic forms of the form $ax^{2}+by^{2}+cz^{2}$, but to any quadratic forms $q(x,y,z)=0$ for which the parametrized solutions for $x$ and $y$ factorize linearly in  $\Z[\tau_{d}]$. We refer the reader to Theorem \ref{thmain} for details.



It is natural to ask whether Theorems \ref{th1} and \ref{th1density} hold when $d<0$, or whether they hold when only one of $\sqrt{-ac}, \sqrt{-bc}$ belongs to $\Z[\tau_{d}]$.
The discussion of these natural questions is postponed to Section~\ref{SEC: further}.
%\AK{What about the other pairs? Say that we have some results --corollaries--below over $\Z[i].$}



 

\subsection{Number-Theoretic tools} In order to prove the above-mentioned partition regularity results, we obtain a series of results in analytic number theory for imaginary quadratic number fields. To the best of our knowledge, these results were previously unknown, and many of them require substantial new ideas compared to the classical results over integers. We believe these results are of independent interest, and may have potential applications in other number-theoretic problems.

Here is a summary of the number-theoretic tools we develop in this paper. We postpone their precise statements to the next Section.

\subsubsection{Characterizations for completely multiplicative functions}
 It is a classical result of Daboussi and Delange (see \cite[Corollary~1]{daboussidelange82}) that a completely multiplicative function $f\colon \Z^{\times}\to\U$ is either \emph{aperiodic}, meaning that its average along every arithmetic progression is zero, or \emph{pretentious}, meaning that it behaves similarly to the product of a Dirichlet character and an Archimedean character. 
This is a crucial tool for the splitting method introduced in \cite{FKM2,FKM1}.
In this paper, we obtain an analog of \cite[Corollary~1]{daboussidelange82} in the setting of imaginary quadratic number fields (i.e., Theorem \ref{characterization}) by providing a list of equivalent definitions of aperiodic completely multiplicative functions (see Proposition \ref{prop: dircharequivalence}). 
An interesting feature is that to define a pretentious multiplicative function, one needs to consider its extension to ideal domained functions. We refer the readers to Section \ref{ss5} for details.
 

%{\color{red} Should we just delete this subsubsection altogether?}  \WS{I think so}
\iffalse
\subsubsection{Hal\'asz's theorem over number fields}
The proof of the aforementioned result of Daboussi and Delange relies on Hal\'asz's theorem, which is a profound result in analytic number theory that characterizes the asymptotic behavior of the average of a multiplicative function. In this paper, we use an analog of Hal\'asz's theorem over arbitrary number fields that can be found in \cite[Theorem~6.1]{lucht2001} (not just imaginary quadratic ones) for multiplicative functions (see Theorem \ref{C22intro} for the precise statement). 
\fi
%The essential obstacles to carry out the proof of the  classical Hal\'asz's theorem are that number fields may not be uniquely factorizable, and that there may be infinitely many units. To overcome these difficulties, we consider functions whose domains are ideals instead of numbers, and this is good enough for the purpose of this paper. We refer the readers to Section \ref{SEChalasz} for details.  

\subsubsection{The Tur\'an-Kubilius inequality}
Another important step in proving the partition regularity results of this paper is obtaining an analog of \cite[Proposition~2.5]{FKM1} (namely, Theorem \ref{p25}), which is a linear concentration estimate that allows us to easily treat the pretentious multiplicative functions. As concentration estimates are based on the Tur\'an-Kubilius inequality, in order to prove  Theorem \ref{p25}, we obtain an analog of the inequality for imaginary quadratic fields in Theorem \ref{TK}.
We refer the readers to Section \ref{SECconcest} for all the details.

\subsection{Organization of the paper}
%The structure of the paper is as follows. 
We state all the main results in density regularity and number theory in Section \ref{spps}.
Then in Section~\ref{sec: proofstrat} we explain the strategy  for the proof  of the main density regularity result (i.e., Theorem \ref{thmain}) %and explain the tools that we develop and/or use. 
and break it down into several positivity results. A lot of the key reductions that help properly frame the problem are also explained in detail.

Next, in Section~\ref{SEC: numthybackg} we provide some necessary number theoretical background. %that is needed in what follows, and develop a new extension theorem for a certain class of multiplicative functions. 
%This helps us obtain in turn a theorem characterizing aperiodicity of multiplicative functions that we then use together with the main result of Section~\ref{SEChalasz}: an extension of Hal\'asz's theorem using the pretentious approach, which allows us to easily get two required measurability results.
We then prove the characterization result for completely multiplicative functions in Section \ref{ss5}, and the Tur\'an-Kubilius inequality as well as the main concentration estimate we need in Section \ref{SECconcest} respectively.
Finally, in  Section~\ref{sec: proofmain} we combine all  the tools developed up to that point to prove the positivity results that are needed to complete the proof of Theorem \ref{thmain}.

%Next, we prove an analog of a concentration estimate (namely \cite[Lemma~2.5]{KMPT}) in our setting in Section~\ref{SECconcest}, which constitutes all the main ingredients necessary for the positivity results, contained in Proposition~\ref{fffss} in Section~\ref{sec: proofmain} . 

%The proofs of said positivity results are given in Section~\ref{sec: proofmain} with all the tools developed up to that point. 
We finish the paper with some discussion on possible future avenues in Section~\ref{SEC: further}. 
Lastly, the Appendix is used to prove a result that allows us to reconcile ball averages with box averages.
%{\color{red} TBC.}3
\\ \\
\noindent {\bf{Acknowledgements.}} We thank Diego C\'espedes for bringing the paper \cite{lucht2001} to our attention. This simplified a previous version of our work, in which we had proved a version of Hal\'asz's theorem from scratch, a result already established in \cite{lucht2001}.
 



\section{Precise statements of the main results}\label{spps}
In this section we will introduce the relevant technical definitions so that the main results can be precisely stated.
\subsection{The main density regularity result and applications}

As was stated in the Introduction, Theorem \ref{th1density} can be generalized to any quadratic equation $q(x,y,z)=0$ for which the parametrized solutions for $x$ and $y$ factorizes linearly in  $\Z[\tau_{d}]$. In order to make this more precise, we introduce the following definition.
 
\begin{definition}
    Let $K$ be a number field and $q\colon\mathcal{O}_{K}^{3}\to\C$ be a map. We say that $q$ admits an \emph{$\mathcal{O}_{K}$-factorization in $(x,y)$} if there exist some map $P\colon \mathcal{O}_K^{2}\to \OK$ and maps $L_{i}\colon \mathcal{O}_{K}^{2}\to \OK, 1\leq i\leq 4,$ of the form $L_{i}(m,n)=a_{i,1}m+a_{i,2}n$ for some $a_{i,1},a_{i,2}\in\mathcal{O}_{K}$ with $L_{1}, L_2, L_3, L_{4}$ being pairwise linearly independent such that for all $m,n,k\in \mathcal{O}_{K}$ the following is a solution to the equation $q(x,y,z)=0$:
    $$x=kL_{1}(m,n)L_{2}(m,n), y=kL_{3}(m,n)L_{4}(m,n) \text{ and } z=kP(m,n).$$
    We say that $p$ admits a \emph{strong $\mathcal{O}_{K}$-factorization in $(x,y)$} if we may further require $a_{1,1}=a_{2,1}=a_{3,1}=a_{4,1}=1$.
\end{definition}

We will proceed by proving the following result which implies Theorems \ref{th1} and ~\ref{th1density}. 
 \begin{theorem}\label{thmain}
     Let $K=\Q(\sqrt{-d})$ with $d\in\N$ squarefree and $q\colon\mathcal{O}_{K}^{3}\to\C$ be a map admitting an $\mathcal{O}_{K}$-factorization in $(x,y)$.
     Then $q(x,y,z)=0$ is density regular over $\OK$ in $(x,y)$.    
     %Then, for every $\Lambda\subseteq \mathcal{O}_{K}^{\times}$ with positive multiplicative density, there exists distinct nonzero $x,y,z\in \mathcal{O}_{K}$ with $x,y\in\Lambda$    such that 
    % $q(x,y,z)=0$.
 \end{theorem}

It was essentially proved in \cite{S23} that Theorem \ref{thmain} holds for all number fields $K$ if $q$ admits a strong $\mathcal{O}_{K}$-factorization in $(x,y)$. In the present paper, we relax this condition by $\mathcal{O}_{K}$-factorization.
\begin{proof}[Proof of Theorems \ref{th1} and ~\ref{th1density} assuming Theorem~\ref{thmain}]
    Note that the solutions of the equation $ax^{2}+by^{2}+cz^{2}=0$ admit the following parametrization:
    $$x=k\sqrt{-bc^{3}}(m-n)(m+n),\;\; y=2k\sqrt{-ac^{3}}mn,\;\; z=k\sqrt{abc^{2}}(m^{2}+n^{2}),$$
    where $\sqrt{-ac^{3}},\sqrt{-bc^{3}}\in\Z[\tau_{d}]$ and $\sqrt{abc^{2}}=\sqrt{-ac}\cdot\sqrt{-bc}\in\Z[\tau_{d}]$.
   The conclusion follows from Theorem \ref{thmain} and the Pigeonhole Principle.
\end{proof}


Next we present some applications of Theorem \ref{thmain}.
In \cite{FKM2,FKM1}, it was proven that $x^{2}+y^{2}-z^{2}=0$ is partition regular over $\Z$ in any of the three pairs of variables. 
Using Theorem \ref{th1}, in the case where $d=1$, i.e., the Gaussian integers, we obtain examples where partition regularity holds in any of the three pairs of variables.

\begin{corollary}\label{cor: gaussian integers nice}
    The equation $ax^{2}+by^{2}+cz^{2}=0$ is partition regular over $\Z[i]$ in $(x,y)$, $(x,z)$ and $(y,z)$ if $\sqrt{-ac}, \sqrt{-bc},\sqrt{-ab}\in \Z[i]$ (or equivalently if $\sqrt{a}, \sqrt{b},\sqrt{c}\in \Z[i]$).
\end{corollary}
 
For example, we have that the equation $x^{2}+y^{2}+z^{2}=0$ is partition regular over $\Z[i]$ in any of the three pairs of variables. (Note that this equation cannot be partition regular over $\Z$.)

\begin{remark}
Full partition regularity is impossible for every generalized Pythagorean equation (of the form  $ax^2+by^2+cz^2=0$) with squares over a ring of integers $\OK$. Indeed, if we hope to have full partition regularity of the equation $ax^2+bx^2=cz^2$ over $\OK$, then, at the very least, we will have full partition regularity for the equation $ax+by=cz$ over $\OK$. Rado proved in \cite{rado43} an extension of his result over $\Z$ for subrings of $\C$ (in particular, any $\OK$ fits the bill), which implies that what is now known as Rado's condition is necessary. Whether it is sufficient is wide open for any $\OK$, including, of course, the case where $\OK=\Z$.
%\begin{itemize}
 %   \item We will see that, in fact, the partition can be combined into one color; but of course, the $y$ cannot be expected to be the same.
  %  \item We can also cover more general homogeneous equations of the form $q(x,y,z)=ax^2+by^2+cz^2+dxy+exz+fyz = 0$, for some $a,b,c,d,e,f \in \Z[i]$. {\color{red} Discuss here how far we can push our method and what it depends on}
  %  \item As a Corollary to Theorem~\ref{mainthm1}, it is interesting to note that the following equations are pair-partition regular over $\Z[i]^{\times}$:
   % \[
   % a^2x^2 \pm b^2y^2=\pm c^2z^2
    %\]
    %for $a, b, c \in \N$.
 %   \item 
%\end{itemize}
\end{remark}

 

%For example, {\color{red} say what we can get for $x^{2}+dy^{2}=z^{2}$, what we can get for $x^{2}+y^{2}=2z^{2}$ (in any two variables) and for $x^{2}+2y^{2}=z^{2}$ (in $(x,z)$).}

%In the case of Gaussian integers, there are examples where  partition regular holds for all the three pairs.

%\begin{corollary}
%    The equation $ax^{2}+by^{2}+cz^{2}=0$ is partition regular over $\Z[i]$ in $(x,y)$, $(x,z)$ and $(y,z)$ if $\sqrt{-ac}, \sqrt{-bc},\sqrt{-ab}\in \Z[i]$ (or equivalently if $\sqrt{a}, \sqrt{b},\sqrt{c}\in \Z[i]$).
%\end{corollary}

%For example, we have that $x^{2}+y^{2}+z^{2}=0$ is partition regular over $\Z[i]$ in any of the three pairs of variables.
   
 It is natural to ask for examples for which partition regularity holds in any of the three pairs of variables over other number fields. To do so we need to include terms $xy,yz, zx$ (or we can only hope to obtain a similar result in the Gaussian integers; see the remark after Theorem~\ref{th2density} below). Consider the more general quadratic polynomials of the form
\begin{equation}\label{fp}
    q(x,y,z)=ax^{2}+by^{2}+cz^{2}+exy+fxz+gyz
\end{equation}
for some $a,b,c,e,f,g\in\Z$. Denote 
$$\Delta_{1}(q):=f^{2}-4ac,\; \Delta_{2}(q):=g^{2}-4bc,\;\text{and}\; \Delta_{3}(q):=(f+g)^{2}-4(a+b+e)c.$$
It was shown in \cite[Proposition~10.4]{S23} that $q(x,y,z)=0$ is partition regular in $(x,y)$ over $\mathcal{O}_{K}$, where $K=\Q(\sqrt{\Delta_{1}(q)},\sqrt{\Delta_{2}(q)},\sqrt{\Delta_{3}(q)})$. We remark that
\cite[Proposition~10.4]{S23} can be recovered using Theorem~\ref{thmain}, which we will prove below. In addition,
  we show that the equation is partition regular with some other choices of number fields, leading to new applications.
 Let $$\Delta_{4}(q)=c(ce^{2}+bf^{2}+ag^{2}-efg-4abc).$$
% \AK{Just a feeling: I am not sure I like all these results in the introduction. Should we split it into more sections or subsections? E.g., ``Main results''.}
We have the following.

%that if $\Delta_{3}(Q)\neq 0$, 
%then 
%$q(x,y,z)=0$ is partition regular in $(x,z)$ over $\mathcal{O}_{K}$, where $K=\Q(\sqrt{\Delta_{2}(Q)},\sqrt{\Delta_{3}(Q)})$. 
%When $\Delta_{3}=0$, it was shown in  \cite[Proposition~10.4]{S23} that 
%$q(x,y,z)=0$ is partition regular in $(x,y)$ over $\mathcal{O}_{K}$, where $K=\Q(\sqrt{\Delta_{1}(Q)},\sqrt{\Delta_{2}(Q)})$.
%Howoever, the method used in \cite{S23} can not cover partition regularity result in $(x,z)$ when $\Delta_{3}(Q)=0$.

%In this paper, we fill the gap for this missing case when $\Delta_{2}(Q)<0$:
%\begin{theorem}\label{th2}
%    Let  $p$ be given by (\ref{fp}) and $d\in\N$. Suppose that $\sqrt{\Delta_{2}(Q)}\in\Z[\tau_{d}]$. If $\Delta_{3}(Q)=0$ and $\Delta_{1}(Q)\neq \Delta_{2}(Q)$,\footnote{The assumption that $\Delta_{1}(Q)\neq \Delta_{2}(Q)$ is to ensure the equation to be non-degenerate. See \cite[Remark~10.5]{S23} for details (on the type (ii) equations).} then the equation $q(x,y,z)=0$ is partition regular in $(x,z)$ over $\Z[\tau_{d}]$.
%\end{theorem}

%Combining Theorem \ref{th2} with the aforementioned partition regularity  in $(x,y)$ from \cite[Proposition~10.4]{S23}, we have 
%\begin{corollary}
%     Let  $p$ be given by (\ref{fp}) and $d\in\N$. Suppose that   $\Delta_{3}(Q)=0$ and $\sqrt{\Delta_{1}(Q)}, \sqrt{\Delta_{2}(Q)}\in\Z[\tau_{d}]$ for some $d\in\N$. Then the equation $q(x,y,z)=0$ is partition regular  over $\Z[\tau_{d}]$ in $(x,y)$, $(x,z)$ and $(y,z)$.
%\end{corollary}

% As an application, we have that the equation
%$$4x^{2}+y^{2}+z^{2}+31xy+4xz+2yz=0$$
%is partition regular over $\Z[\sqrt{-3}]$ in any of the three pairs of variables. 
% {\color{red} Check the parametrization.}

\iffalse

\begin{theorem}\label{th2}
    Let $q$ be given by (\ref{fp}) and $d\in\N$ be squarefree. 
     Suppose that $a,b,c\neq 0$ and $\Delta_{4}(Q)\neq 0.$\footnote{If $\Delta_{4}(Q)=0$, then the proof of Theorem \ref{th2} implies that $Q$ can be factorized into the product of two linear equations with $\Q(\sqrt{-d})$-coefficients, and the problem is reduced to the linear case.} Then the   equation $q(x,y,z)=0$ is  
\begin{itemize}
    \item partition regular in $(x,z)$ if $\sqrt{\Delta_{2}(Q)},\sqrt{\Delta_{4}(Q)}\in\Z[\tau_{d}]$;
    \item partition regular in $(y,z)$ if $\sqrt{\Delta_{1}(Q)},\sqrt{\Delta_{4}(Q)}\in\Z[\tau_{d}]$; and
    \item partition regular in $(x,y)$ if $\sqrt{\Delta_{1}(Q)},\sqrt{\Delta_{2}(Q)},\sqrt{\Delta_{4}(Q)}\in\Z[\tau_{d}]$.
\end{itemize}
%\AK{$\sqrt{\Delta_{3}(Q)}$ doesn't appear anywhere?}
\end{theorem}

%Combining Theorem \ref{th2} with the aforementioned partition regularity  in $(x,y)$ from \cite[Proposition~10.4]{S23}, we have 

%\begin{corollary}
%     Let  $p$ be given by (\ref{fp}) and $d\in\N$. Suppose that   $\sqrt{\Delta_{1}(Q)}, \sqrt{\Delta_{2}(Q)},\sqrt{\Delta_{4}(Q)}\in\Z[\tau_{d}]$ for some $d\in\N$. Then the equation $q(x,y,z)=0$ is partition regular  over $\Z[\tau_{d}]$ in $(x,y)$, $(x,z)$ and $(y,z)$.
%\end{corollary}

\fi


 
 


\begin{theorem}\label{th2density}
    Let $q$ be given by \eqref{fp} and $d\in\N$ be squarefree. 
     Suppose that $a,b,c\neq 0$ and $\Delta_{4}(q)\neq 0$.\footnote{If $\Delta_{4}(q)=0$, then the proof of Theorem \ref{th2density} implies that $q$ can be factorized into the product of two linear equations with $\Q(\sqrt{-d})$-coefficients, and the problem is reduced to the linear case.} Then the equation $q(x,y,z)=0$ is  
\begin{itemize}
    \item density regular in $(x,z)$ if $\sqrt{\Delta_{2}(q)},\sqrt{\Delta_{4}(q)}\in\Z[\tau_{d}]$;
    \item density regular in $(y,z)$ if $\sqrt{\Delta_{1}(q)},\sqrt{\Delta_{4}(q)}\in\Z[\tau_{d}]$; and
    \item density regular in $(x,y)$ if $\sqrt{\Delta_{1}(q)},\sqrt{\Delta_{2}(q)},\sqrt{\Delta_{4}(q)}\in\Z[\tau_{d}]$.
\end{itemize}
\end{theorem}

 As an application of Theorem \ref{th2density}, we have that the equation
$$x^{2}+y^{2}+3z^{2}+xy=0$$
is partition regular over $\Z[\sqrt{-3}]$ in any of the three pairs of variables, which is a result that cannot be obtained with the methods of \cite{S23}. 
 %{\color{red} Check the parametrization.}

\begin{proof}[Proof of Theorem \ref{th2density} assuming Theorem \ref{thmain}]
 %   The conclusion follows from the following parametrization of the solutions to $q(x,y,z)=0$ (when $\Delta_{3}(Q)=0$):
 %   $$x=k(m-\sqrt{\Delta_{2}(Q)}n)(m+\sqrt{\Delta_{2}(Q)}n), y=k(m^{2}-\Delta_{1}(Q)n^{2}), z=\pm k(\Delta_{1}(Q)-\Delta_{2}(Q))mn.$$
 %   {\color{red} this case is already covered by [FKM2]!!!}

    
     We first consider the case when $f=g=0$. In this case, $\Delta_{1}(q), \Delta_{2}(q)$ and $\Delta_{4}(q)$ are equal to $-4ac$, $-4bc$ and $c^2(e^{2}-4ab)$ respectively. %So we have $\sqrt{-bc}, \sqrt{d^{2}-4ab}\in\Z[\tau_{d}]$. 
    We may rewrite $q(x,y,z)=0$ as 
    $$a(x-\lambda_{-}y)(x-\lambda_{+}y)+cz^{2}=0$$
    where
    $$\lambda_{\pm}=\frac{-e\pm\sqrt{e^{2}-4ab}}{2a}.$$
    Since $\Delta_{4}(q)\neq 0$, we have $\lambda_{+}\neq \lambda_{-}$. Since $a,b\neq 0$, we have $\lambda_{\pm}\neq 0$.
    So we have a solution if 
    $$x-\lambda_{-}y=kcm^{2},\;\; x-\lambda_{+}y=-kan^{2}, \;\; \text{and}\;\; z=kamn,$$
    or equivalently,
     $$x=\frac{k}{\lambda_{+}-\lambda_{-}}(\lambda_{-}cm^{2}+\lambda_{+}an^{2}),\;\; y=\frac{k}{\lambda_{+}-\lambda_{-}}(cm^{2}+an^{2}),\;\;  \text{and}\;\; z=kamn.$$

     Note that $\lambda_{+}-\lambda_{-}=\frac{\sqrt{e^{2}-4ab}}{a}=\frac{\sqrt{\Delta_{4}(q)}}{a}$, $\sqrt{-ac}=\frac{1}{1}\sqrt{\Delta_{2}(q)}$ and  $\sqrt{-\lambda_{+}\lambda_{-}ac}=\sqrt{-bc}=\frac{1}{2}\sqrt{\Delta_{2}(q)}$. It follows that $q$ admits an 
     \begin{itemize}
    \item $\OK$-factorization in $(x,z)$ if $\sqrt{\Delta_{2}(q)},\sqrt{\Delta_{4}(q)}\in\Z[\tau_{d}]$.
    \item $\OK$-factorization in $(y,z)$ if $\sqrt{\Delta_{1}(q)},\sqrt{\Delta_{4}(q)}\in\Z[\tau_{d}]$.
    \item $\OK$-factorization in $(x,y)$ if $\sqrt{\Delta_{1}(q)},\sqrt{\Delta_{2}(q)},\sqrt{\Delta_{4}(q)}\in\Z[\tau_{d}]$.
\end{itemize}
%{\color{red} here we need to check that the linear forms are linearly independent}
 We are done by Theorem \ref{thmain}.
 
\iffalse

 
    Since $\sqrt{-bc}, \sqrt{d^{2}-4ab}\in\Z[\tau_{d}]$, we have that $\lambda_{\pm}\in \Q(\sqrt{-d})$.  
   Moreover, $\sqrt{-\lambda_{+}\lambda_{-}ac}=\sqrt{-bc}$. Finally, since $a,b\neq 0$, we have $\lambda_{\pm}\neq 0$. From this we deduce that
     $$x=kL_{1}(m,n)L_{2}(m,n), y=kP(m,n)  \text{ and } z=kL_{3}(m,n)L_{4}(m,n)$$
     where $L_{1},\dots,L_{4}$ are pairwise linearly independent linear forms in $\Q(\sqrt{-d})$ and $P$ is a quadratic equation in  $\Q(\sqrt{-d})$. {\color{red} check that they are linearly independent}
   So
   $Q$ amdits an $\OK$-factorization in $(x,z)$ with $K=\Q(\sqrt{-d})$. We are done by Theorem \ref{thmain} and the Pigeonhole Principle.

\fi
     

    We now consider the general case.
    Let 
    \begin{eqnarray*}
        p'(x,y,z) & = & p(2cx,2cy,z-fx-gy) \\
        & = &-c\left(\Delta_{1}(q)x^{2}+\Delta_{2}(q)y^{2}-z^{2}+(\Delta_{3}(q)-\Delta_{1}(q)-\Delta_{2}(q))xy\right).
    \end{eqnarray*}
    From the previous case and the identity 
    $$4\Delta_{4}(q)=(\Delta_{3}(q)-\Delta_{1}(q)-\Delta_{2}(q))^{2}-4\Delta_{1}(q)\Delta_{2}(q),$$
    we have the conclusion.
\end{proof}
\begin{remark}
It is interesting to note that if one insists on having a diagonal equation (of the form $ax^2+by^2=cz^2$) and wants to obtain full partition regularity for any of the three possible pairs using Theorem~\ref{thmain}, then after some checking which is left to the interested reader, one can show that, necessarily, the quadratic field extension must be $\Z[i]$, the Gaussian integers, which highlights the special properties that this ring of integers enjoys.
%\AK{More details here? E.g., a footnote?}
\end{remark}

\iffalse


\begin{proof}[Proof of Theorem \ref{th2} assuming Theorem \ref{thmain}]
 %   The conclusion follows from the following parametrization of the solutions to $q(x,y,z)=0$ (when $\Delta_{3}(Q)=0$):
 %   $$x=k(m-\sqrt{\Delta_{2}(Q)}n)(m+\sqrt{\Delta_{2}(Q)}n), y=k(m^{2}-\Delta_{1}(Q)n^{2}), z=\pm k(\Delta_{1}(Q)-\Delta_{2}(Q))mn.$$
 %   {\color{red} this case is already covered by [FKM2]!!!}

    
{\color{red} needs more work}
    We first consider the case when $e=f=0$. In this case, $\Delta_{2}(Q)$ and $\Delta_{4}(Q)$ are $-4bc$ and $c^{2}(d^{2}-4ab)$ respectively. So we have $\sqrt{-bc}, \sqrt{d^{2}-4ab}\in\Z[\tau_{d}]$. {\color{red} notation conflict }
    We may rewrite $q(x,y,z)=0$ as 
    $$a(x-\lambda_{-}y)(x-\lambda_{+}y)+cz^{2}=0$$
    where
    $$\lambda_{\pm}=\frac{-d\pm\sqrt{d^{2}-4ab}}{2a}.$$
    Since $\Delta_{4}(Q)\neq 0$, we have $\lambda_{+}\neq \lambda_{-}$.
    So we have a solution if 
    $$x-\lambda_{-}y=kcm^{2}, x-\lambda_{+}y=-kan^{2} \text{ and } z=kamn,$$
    or equvialently,
     $$x=\frac{k}{\lambda_{+}-\lambda_{-}}(\lambda_{-}cm^{2}+\lambda_{+}an^{2}), y=\frac{k}{\lambda_{+}-\lambda_{-}}(cm^{2}+an^{2})  \text{ and } z=kamn.$$
    Since $\sqrt{-bc}, \sqrt{d^{2}-4ab}\in\Z[\tau_{d}]$, we have that $\lambda_{\pm}\in \Q(\sqrt{-d})$.  
   Moreover, $\sqrt{-\lambda_{+}\lambda_{-}ac}=\sqrt{-bc}$. Finally, since $a,b\neq 0$, we have $\lambda_{\pm}\neq 0$. From this we deduce that
     $$x=kL_{1}(m,n)L_{2}(m,n), y=kL_{3}(m,n)L_{4}(m,n)  \text{ and } z=kL_{5}(m,n)L_{6}(m,n)$$
     where $L_{1},\dots,L_{6}$ are pairwise linearly independent linear forms {\color{red} check that they are linearly independent}
   So
   $Q$ amdits an $\OK$-factorization in $(x,y)$ with $K=\Q(\sqrt{-d})$. We are done by Theorem \ref{thmain} and the Pigeonhole Principle.

     

    We now consider the general case.
    Let 
    \begin{eqnarray*}
        p'(x,y,z) & = & p(2cx,2cy,z-ex-fy) \\
        & = &-c\left(\Delta_{1}(Q)x^{2}+\Delta_{2}(Q)y^{2}-z^{2}+(\Delta_{3}(Q)-\Delta_{1}(Q)-\Delta_{2}(Q))xy\right).
    \end{eqnarray*}
    The conclusion follows from the previous case and the identity 
    $$4\Delta_{4}(Q)=(\Delta_{3}(Q)-\Delta_{1}(Q)-\Delta_{2}(Q))^{2}-4\Delta_{1}(Q)\Delta_{2}(Q).$$
\end{proof}

\fi

 
 \subsection{Characterizations for completely multiplicative functions}

% \WS{The deifnition in this section has many overlap with Section 3.2. We need to think of a way to present the definitions}

The strategy of the proofs for the partition regularity results of this paper is to replace the decomposition method used in \cite{S18,S23} with the splitting method introduced in \cite{FKM2,FKM1}. %We first recall the definition of multiplicative functions.
%
%\begin{definition}
%Let $f:\mathcal{O}_K^{\times} \to \U$. We say that $f$ is \emph{completely multiplicative} if
%\[
%f(mn)=f(m)\cdot f(n), \quad \text{ for all } m,n \in \OK.
%\]
%\end{definition}
Recall that a function $f\colon \N\to\U$, where $\U:=\{z \in\C: |z| \leq 1\}$ is \emph{multiplicative} if $f(mn)=f(m)\cdot f(n)$ whenever $(m,n)=1$, and is \emph{completely multiplicative} if $f(mn)=f(m)\cdot f(n)$ for all $m,n\in\N$.
For functions $f,g: \N\to\U$, define
$$\D(f,g):=\sum_{p \text{ is a prime number in } \N} \frac{1}{p}(1-\text{Re}f(p)\overline{g}(p)).$$
We say that $\chi : \Z \to \U$ is a \emph{Dirichlet character of modulus $m$} if $\chi$ is completely multiplicative, satisfies $\chi(a)=0$ if and only if $\gcd(a,m)>1$, and is periodic with period $m$. It is a classical result (see \cite[Corollary~1]{daboussidelange82}) that a completely multiplicative function $f\colon \Z^{\times}\to\U$ is either \emph{aperiodic}, meaning that
$$\lim_{N\to\infty}\frac{1}{N}\sum_{n=1}^{N}f(a+bn)=0$$
for all $a,b\in\N$, or \emph{pretentious}, meaning that $\D(f\chi,n^{i\tau})<\infty$ for some Dirichlet character $\chi$ and $\tau\in\R$.



%One of the key innovation of the paper is that we obtain a classification result for completely multiplicative functions into aperiodic and pertentious ones, depending on the asyptotic behaviors of their averages along arithmatic progressions. Such a result is classical for completely multiplicative functions in $\Z$. 

In this paper, we develop an analog of this characterization for completely multiplicative functions over imaginary quadratic number fields.
  We say that $f\colon \mathcal{O}_K^{\times}\to\U$ is \emph{completely multiplicative} if $f(mn)=f(m)\cdot f(n)$ for all $m,n\in \mathcal{O}_K^{\times}$. Let $d\in\N$ be squarefree and $\NN\colon \Q(\sqrt{-d})\to \Q$ be the norm of the field extension (see Section~\ref{SEC: numthybackg} for more details). %We say that $f\colon \Z[\tau_{d}]^{\times}\to\U$ is \emph{aperiodic} if for all $a_{1},a_{2}\in\N$ and $b_{1},b_{2}\in\Z$ we have that 
%$$\lim_{N\to\infty} \E_{m_{1}, m_{2}\in\Z\colon \NN\left((a_{1}m_{1}+b_{1})+(a_{2}m_{2}+b_{2})\tau_{d}\right)\leq N}f\left((a_{1}m_{1}+b_{1})+(a_{2}m_{2}+b_{2})\tau_{d}\right)=0$$
%for all $a_{1},a_{2}\in\N$ and $b_{1},b_{2}\in\Z$. 
%We remark this notion of aperiodicity is only well defined when $d>0$, as the set $\{n\in\Z[\tau_{d}]\colon \NN(n)\leq N\}$ is infinite when $d<0$.
We define the notion of aperiodic functions in this context in the following way.


\begin{definition}[Aperiodic functions]\label{aperiodicdef}
   Let $d\in\N$ be squarefree and let $\mathcal{AP}[\tau_{d}]$ denote the collection  of all  2-dimensional grids in $\Z[\tau_{d}]$, i.e., the collection of sets of the form
   \begin{equation}\label{eq:adcd} 
\{(a_{1}m+b_{1})+(a_{2}n+b_{2})\tau_{d}\in \Z[\tau_{d}]\colon m,n\in\Z\}
\end{equation}
        for some $a_{1},b_{1},a_{2},b_{2}\in\Z$ with $a_{1},a_{2}\in\N$.
   We say that a function $f\colon\Z[\tau_{d}]^{\times}\to\C$ is \emph{aperiodic} if the limit
\begin{equation}\label{eq: aperiodicdef}
\lim_{N\to\infty}\E_{u\in P,\NN(u) \leq N}f(u)
\end{equation}
%where $\{x_1,x_2\}$ is an integral basis for $\OK$, 
exists and equals to 0 for all $P\in\mathcal{AP}[\tau_{d}]$. 
%{\color{red}
%If I am understanding the definition of progressions that is proposed here (it is slightly different from the one in Wenbo's paper), then, for even for general number fields, using an integral basis it would look like
%\[
%\lim_{N \to \infty} \E_{\NN(u=\sum_{i=1}^d (a_in_i+b_i)x_i) \leq N} g( (a_1n_1+b_1)x_1+\dots+(a_dn_d+b_d)x_d) .
%\]
%Then, we can take a common factor for all the $a_i$'s, and consider multiple characteristic functions along APs in that coordinate, say of the form $\prod_{i=1}^d\mathds{1}_{\alpha_i\Z+\beta_i}(n_i)$, so that the average of interest is a linear combination of averages of the form
%\[
%\lim_{N \to \infty} \E_{\NN(u) \leq N} g( nu+b)
%\]
%for $n=\text{lcm }(a_1,\dots,a_d)$ and some $b \in \mathcal{O}_K$. Finally, in order to deal with this, I found that, even though the quadratic imaginary fields are not UFDs, they do admit factoring into irreducibles (an argument based on induction and the norms being non-negative). This would allow us to write $b$ and $n$ as a product of irreducibles in $\mathcal{O}_K$ and finally be able to have the above average as
%\[
%\lim_{N \to \infty} \E_{\NN(u) \leq N} g( Qu+d)
%\]
%where $Q$ is a non-zero element of $\mathcal{O}_K$, and $d \in (\mathcal{O}_K/\langle Q \rangle)^{\times}$. Thus, we can see that such averages will be zero provided that the averages of $g$ times a Dirichlet character (on a principal ideal of the form $\langle Q\rangle$ is zero or not. This seems to be able to expand the argument even if we lose the  UFD property. However, for infinitely many units, the fact we are using norms and balls is a big issue.
%}
%We say that $g$ is \emph{pretentious} if $\D(f,\chi\cdot \NN(u)^{i\tau})<\infty$ for some Dirichlet character $\chi$ and some $\tau\in\R$. We say that $g$ is \emph{super pretentious} if $\D(f,\NN(u)^{i\tau})<\infty$ for some $\tau\in\R$.
\end{definition}

\begin{remark}
    We remark this notion of aperiodicity is only well defined when $d>0$, as the set $\{n\in\Z[\tau_{d}]\colon \NN(n)\leq N\}$ is infinite when $d<0$.
%We also caution the reader that Definition~\ref{aperiodicdef} is slightly different from the one that appeared in \cite{S23}, but this does not have any significant impact on the present work.
%{\color{red} say that it differs from the definition in \cite{S23}}
We also caution the reader that Definition~\ref{aperiodicdef} is different from the one that appeared in \cite{S23}, where the averages are taken over boxes instead of balls. 
It is an interesting question to ask if these two definitions are equivalent. While aperiodic functions defined in both ways share similar properties (for example, they admit similar seminorm estimates for multiple ergodic averages, and their Gowers norms both vanish, see Appendix \ref{appapp}), in this paper it is crucial that we work with ball averages. We refer the readers to Section \ref{s:as} for details on the discussion regarding the choice of averging scheme.

\end{remark}


 

\iffalse
 
Based on the classification result for completely multiplicative functions over integers,
it is natural to ask the following question:

\begin{question}
    Let  $d\in\N$ be squarefree and $\NN\colon \Q(\sqrt{-d})\to \Q$ be the norm of the field extension. Let $f\colon \Z[\tau_{d}]^{\times}\to\U$  be completely multiplicative. %with $f(\e)=1$ for all units $\e$\footnote{If $f$ does not take value 1 on the units, then it is known to be aperiodic by  Lemma \ref{tv}.}. 
    Is it true that either $f$ is aperiodic or $\D_{0}(f\chi,\NN(n)^{i\tau})<\infty$ for some Dirichlet character $\chi$ {\color{red} define it} and $\tau\in\R$, where we denote  
    $$\D_{0}(f,g):=\sum_{p \text{ is a prime element of } \Z[\tau_{d}]}\frac{1}{\NN(p)}(1-\text{Re}f(p)\overline{g}(p)).$$
    Here we say that $p\in\OK$ is a \emph{prime element} if it is not the products of two non-unit\footnote{Recall that an element of $\OK$ is a unit if and only if its norm is equal to $\pm 1$} elements in $\OK$.
\end{question}
As a special case of Theorem \ref{characterization}, we show that the answer to this question is affirmative when $\Z[\tau_{d}]$ is a principal ideal domain  (or equivalently a unique factorization domain, since $\Z[\tau_d]$ is a Dedekind domain).

\fi

It is then natural to ask if such a classification result holds for completely multiplicative functions over other number fields. While such a result seems feasible for principal ideal domains (or equivalently, a unique factorization domain, since $\Z[\tau_d]$ is a Dedekind domain), for non-principal ideal domains it is a harder task to achieve. In the latter, we need to consider the extensions of the function into ideal-domained functions.
%
%For non-pincipal ideal domains, we do not know the answer to this question. However, we provide an alterative characterization by considering the extensions of the function into ideal-domained functions.


 
Let $\mathfrak{I}(K)$ be the set of all ideals of $\mathcal{O}_K^{\times}$. We say that a map $g\colon \mathfrak{I}(K)\to \U$ is \emph{multiplicative} if $g(IJ)=g(I)\cdot g(J)$ for all coprime  $I,J\in\mathfrak{I}(K)$,\footnote{See Section \ref{SEC: numthybackg} for definition.} and is  \emph{completely multiplicative} if $g(IJ)=g(I)\cdot g(J)$ for all  $I,J\in\mathfrak{I}(K)$.

This allows us to define extensions of completely multiplicative functions.
\begin{definition}\label{def: multfunction extension}
Let $f\colon \mathcal{O}_K^{\times}\to\U$ be a completely multiplicative function with $f(\e)=1$ for all units $\e$. We say that a map $\tilde{f}\colon \mathfrak{I}(K)\to \U$ is an \emph{extension} of $f$ if $\tilde{f}$ is completely multiplicative and $\tilde{f}((n))=f(n)$ for all $n\in \mathcal{O}_K^{\times}.$\footnote{Since $(\e n)=(n)$ for all units $\e$ and $n\in\mathcal{O}_K^{\times}$, $f$ admits an extension only if $f(\e)=1$ for all units $\e$.} 
\end{definition}
For  $f,g\colon\mathfrak{I}(K)\to\C$, define
\begin{equation}\label{ddf}
    \D(f,g):=\sum_{\mathfrak{p} \text{ is a prime ideal of } \mathcal{O}_{K}}\frac{1}{\NN(\mathfrak{p})}(1-\text{Re}f(\mathfrak{p})\overline{g}(\mathfrak{p})).
\end{equation}
It follows from Proposition~\ref{prop: extensions} (whose proof is deferred) that, when $K=\Q(\sqrt{-d}), d\in\N$, the number of extensions of every completely multiplicative function is equal to the ideal class number of $K$, and furthermore, they can be written down explicitly.


\begin{definition}[Dirichlet characters]
     Let $I$ be a non-trivial ideal of $\OK$. We say that a completely multiplicative function $\chi\colon \mathcal{O}_K^{\times}\to\C$ is a \emph{Dirichlet character} over $\OK$ with \emph{period $I$} if $\chi(x+y)=\chi(x)$ for all $x\in \OK, y\in I$ and if $\chi(x)=0$ if and only if $x \mod I\notin (\OK/I)^{\times}$.  
     Given a Dirichlet character $\chi$, we will call the function
     $\chi'\colon \mathcal{O}_K^{\times}\to\C$ a \emph{modified Dirichlet character} over $\OK$ with \emph{period $I$} if $\chi'(x)=\chi(x)$ for $x \mod I\notin (\OK/I)^{\times}$ and $\chi'(x)=1$ otherwise.  
\end{definition}
It is clear from the definition, but important in the sequel, that modified Dirichlet characters are $\S$-valued completely multiplicative functions.
%Let $I$ be a finite index ideal of $\Z[i]$ (this is the same as saying that $I \neq (0)$). Let $\mathcal{B}_{I}$ denote the set of all completely multiplicative function $\chi\colon \Z[i]\to\C$ with $\chi(\cdot+I)=\chi$ and with $\chi(x)=0$ if and only if $x \mod I\in (\Z[i]/I)^{\times}$ (i.e. there exists $x'\in \Z[i]$ with $xx'-1\in I$. Let $\mathcal{B}:=\cup_{I}\mathcal{B}_{I}$. We call the elements in $\mathcal{B}$ \emph{Dirichlet characters} over $\Z[i]$.

\begin{definition}[Pretentious functions]%\label{aperiodicdef}
    We say that a function $f\colon\mathcal{O}_K^{\times}\to\S$  
  is \emph{pretentious} if $f(\e)=1$ for all units $\e$ and there exist some modified Dirichlet character $\chi$, some $\tau\in\R$, and some extension $\tilde{f'}$ of the function
$$f'(u):=f(u)\overline{\chi}(u)\NN(u)^{-i\tau}\footnote{Note that $f'$ is a multiplicative function taking values in $\S$ that is trivial on units.}$$
such that 
  $\D(\tilde{f'}, 1)<\infty$. We say that $f$ is \emph{super pretentious} if we can further require $\chi=1$ in the construction of $f'$.  
\end{definition}

 We have the following classification theorem.

%\AK{About the following: (2) It is not ideal to have footnotes inside statements, in particular inside theorems. (3) After making the ``extension'' a central thing, refer to the relation and say ``some extension $\widetilde{f\chi}$ of $f\chi$ as in (???)''.}

\begin{theorem}\label{characterization}
     Let  $d\in\N$ be squarefree and $\NN\colon \Q(\sqrt{-d})\to \Q$ be the norm of the field extension. Let $f\colon \Z[\tau_{d}]^{\times}\to\S$ be completely multiplicative. Then either $f$ is aperiodic or  pretentious. %$\D(\widetilde{f\chi},\NN(n)^{i\tau})<\infty$ for some Dirichlet character $\chi$,\footnote{In this context a Dirichlet character means a multiplicative function $\chi: \Z[\tau_d]\to \U$ such that $\chi(\cdot)=\chi(\cdot+I)$ for some ideal $I \subseteq \Z[\tau_d]$ and such that $\chi(u)=0$ if and only if $u \in I$.} some extension $\widetilde{f\chi}$ as in Definition~\ref{def: multfunction extension} of $f\chi$, and some $\tau\in\R$.
\end{theorem}



\begin{remark}
Note that in Theorem \ref{characterization} we only classify $\S$-valued functions. We will study $\U$-valued functions in Section 4, but not in Theorem \ref{characterization}. The reason is that in Proposition \ref{prop: dircharequivalence}, we will only characterize $\S$-valued aperiodic functions, to ensure the existence of extensions of $f\chi$.
\end{remark}

%{\color{red} Question: can we actually use this theorem to find an example so that the previous conjecture fails?
%One can show that if $\D(\tilde{f\chi},\NN(n)^{i\tau})<\infty$, then $\D_{0}(f\chi,\NN(n)^{i\tau})<\infty$. Can we find some example where $\D_{0}(f\chi,\NN(n)^{i\tau})<\infty$ but $\D(\tilde{f\chi},\NN(n)^{i\tau})=\infty$? (find an example in $\Z[\sqrt{-5}]$)

%}

As a consequence, we have the following corollary when the domain is a principal ideal domain (which we will shorten to PID). In order to state it we define a notion of distance between two multiplicative functions. We need to recall that $p \in \OK$ is a \emph{prime element} if it is not the product of two non-unit elements in $\OK$.\footnote{Recall that an element of $\OK$ is a unit if and only if its norm is equal to $\pm 1$.}

\begin{definition}
Let $f,g: \Z[\tau_d]\to \U$ be completely multiplicative functions. Then, the pretentious distance between $f$ and $g$ is given by
\[
\D(f,g):= \sum_{p \text{ is a prime element of } \Z[\tau_{d}]}\frac{1}{\NN(p)}(1-\text{Re}f(p)\overline{g}(p))
\]
\end{definition}
We now state a corollary of Theorem~\ref{characterization}.
\begin{corollary} 
     Let  $d\in\N$ be squarefree and $\NN\colon \Q(\sqrt{-d})\to \Q$ be the norm of the field extension. Let $f\colon \Z[\tau_{d}]^{\times}\to\S$ be completely multiplicative. If $\Z[\tau_{d}]$ is a PID, then either $f$ is aperiodic or  $\D(f\chi,\NN(n)^{i\tau})<\infty$ for some Dirichlet character $\chi$, and some $\tau\in\R$.
     %where  $$\D(f,g):=\sum_{p \text{ is a prime element of } \Z[\tau_{d}]}\frac{1}{\NN(p)}(1-\text{Re}f(p)\overline{g}(p))$$
     %for $f,g:\Z[\tau_{d}]^{\times}\to\C$. \AK{(1) Define $\D(f,g)$ in this setting before?}
   % Here we say that $p\in\OK$ is a \emph{prime element} if it is not the product of two non-unit\footnote{Recall that an element of $\OK$ is a unit if and only if its norm is equal to $\pm 1$} elements in $\OK$. \AK{put the last phrase before the corollary as well. The corellary will look beautiful afterwards.}
\end{corollary}

Theorem~\ref{characterization} is one of the major number theoretical inputs of the paper, which allows us to execute the machinery of \cite{FKM1} for imaginary quadratic number fields, especially when $\Z[\tau_{d}]$ is not a principal ideal domain. We believe that Theorem~\ref{characterization} is interesting in its own right, and may potentially have further applications in number theory.

%\WS{Section 2.3 should no longer be an inedependent section since it is not a main result.
%My suggestion is the remove the entire section 2.3, and add the following at the end of Section 2.2:


The proof of Theorem~\ref{characterization} is based on an analog of Hal\'asz's theorem over number fields. The classical Hal\'asz's Theorem is a profound result in analytic number theory that characterizes the asymptotic behavior of the average of a multiplicative function (cf. \cite[Satz~1 and Satz~1$'$]{halasz68}). As it turns out, while this theorem can be relatively easily extended to PIDs such as $\Z[i]$ (see \cite[Thoerems A and 1.2]{dlms24}, which proved a special case of Hal\'asz's theorem over the Gaussian integers for real valued bounded completely multiplicative functions), difficulties arise when the elements in the domain cannot be factorized in a unique way. 
 Nevertheless, by passing to $\mf I(K)$-domained functions instead of $O_K^{\times}$ ones, the following extension of Hal\'asz's theorem is known to hold:
%While the method of the proof of Hal\'asz's theorem can be extended to principal ideal  domains such as $\Z[i]$ (see \cite[Thoerems A and 1.2]{dlms24}), which proved a special case of Hal\'asz's theorem over the Gaussian integers for real valued bounded completely multiplicative functions, difficulties arise when the elements in the domain cannot be factorized in a unique way. To resolve this issue, our strategy is to consider $\mathfrak{I}(K)$-domained functions instead of $\mathcal{O}_K^{\times}$ ones.This leads us to the following result.


\begin{theorem}[Theorem 6.1, \cite{lucht2001}]\label{C22intro}
     Let $K$ be a number field and $g\colon\mf{I}(K)\to\U$ be a multiplicative function. Then, we have the following.
     \begin{enumerate}
         \item If $\inf_{\tau\in\R}\D(g,\NN(\mf u)^{i\tau})=\infty$, then $\lim_{x \to \infty} \E_{\NN(\mf u)\leq x}g(\mf u)=0$.
         \item If $\D(g,\NN(\mf u)^{i\tau})<\infty$ for some $\tau\in\R$, then
         $$\sum_{1 \leq \NN(\mf u)\leq x}g(\mf u) = \gamma_{K}\cdot\frac{x^{1+i\tau}}{1+i\tau}\prod_{\mf p\colon \text{prime ideal in } \mathcal{O}_K,\NN(\mf p)\leq x}(1-1/\NN(\mf p))(1+h_{1+i\tau}(\mf p))+o(x),$$
         where $h_{s}(\mf p):=\sum_{k=1}^{\infty}g(\mf p^{k})\NN(\mf p^{k})^{-s}$, and $\gamma_{K}$ %is the residue of the Dedekind zeta function of our number field at its simple pole at $s=1$.
         is the universal constant defined in Corollary \ref{lem: bounds for number of ideals of norm at most x}. 
     \end{enumerate}
\end{theorem}

%{\color{red}

%Here I completely removed the notion $R_{1}$, since we no longer need to define the zeta function.
%}

We remark that Theorem \ref{C22intro} applies to any number field not just the quadratic imaginary ones. When $K=\Q$, Theorem \ref{C22intro} is the standard version of Hal\'asz's theorem.
Although Theorem \ref{C22intro} is only valid for $\mf I(K)$-domained functions, it turns out that this is good enough for the proof of Theorem~\ref{characterization} (for $\OK$-domained functions).






%}

\iffalse  
\subsection{Hal\'asz's theorem over number fields}
The first step in the study of the aperiodicity properties of a multiplicative function $f$ on $\OK=\Z[\tau_{d}]$ (for some squarefree $d\in\N$) is to study its average along the entire space instead of along arithmetic progressions, i.e., we want to focus on expressions of the form
$$\lim_{x\to\infty}\E_{\NN(u)\leq x}f(u).$$
When $f$ is a multiplicative function on $\Z$, the behavior of these averages is well understood via Hal\'asz's theorem (cf. \cite[Satz~1 and Satz~1$'$]{halasz68}). As it turns out, the theorem can be extended to PIDs such as $\Z[i]$ (see \cite[Thoerems A and 1.2]{dlms24}, which proved a special case of Hal\'asz's theorem over the Gaussian integers for real valued bounded completely multiplicative functions, difficulties arise when the elements in the domain cannot be factorized in a unique way. Nevertheless, by passing to $\mf I(K)$-domained functions instead of $O_K^{\times}$ ones, the following extension of Hal\'asz's theorem is known to hold (cf. \cite[Theorem 6.1]{lucht2001}):
%While the method of the proof of Hal\'asz's theorem can be extended to principal ideal  domains such as $\Z[i]$ (see \cite[Thoerems A and 1.2]{dlms24}), which proved a special case of Hal\'asz's theorem over the Gaussian integers for real valued bounded completely multiplicative functions, difficulties arise when the elements in the domain cannot be factorized in a unique way. To resolve this issue, our strategy is to consider $\mathfrak{I}(K)$-domained functions instead of $\mathcal{O}_K^{\times}$ ones.This leads us to the following result.


\begin{theorem}[Theorem 6.1, \cite{lucht2001}]\label{C22intro}
     Let $K$ be a number field and $g\colon\mf{I}(K)\to\U$ be a multiplicative function. Then, we have the following.
     \begin{enumerate}
         \item If $\inf_{\tau\in\R}\D(g,\NN(\mf u)^{i\tau})=\infty$, then $\lim_{x \to \infty} \E_{\NN(\mf u)\leq x}g(\mf u)=0$.
         \item If $\D(g,\NN(\mf u)^{i\tau})<\infty$ for some $\tau\in\R$, then
         $$\sum_{1 \leq \NN(\mf u)\leq x}g(\mf u) = R_1\cdot\frac{x^{1+i\tau}}{1+i\tau}\prod_{\mf p\colon \text{prime ideal in } \mathcal{O}_K,\NN(\mf p)\leq x}(1-1/\NN(\mf p))(1+h_{1+i\tau}(\mf p))+o(x),$$
         where $h_{s}(\mf p):=\sum_{k=1}^{\infty}g(\mf p^{k})\NN(\mf p^{k})^{-s}$, and $R_1$ is the residue of the Dedekind zeta function of our number field at its simple pole at $s=1$.
     \end{enumerate}
\end{theorem}
We remark that Theorem \ref{C22intro} applies to any number field not just the imaginary quadratic ones. When $K=\Q$, Theorem \ref{C22intro} is the standard version of Hal\'asz's theorem.
The major application of Theorem \ref{C22intro} in this paper is that it allows us to prove the classification theorem (Theorem \ref{characterization}). 
%We believe that Theorem \ref{C22intro} could lead to other applications in number theory in the future.
\fi
\subsection{The Tur\'an-Kubilius inequality}
%\WS{I moved the statement of  Tur\'an-Kubilius ineqaulity here, as I think this is also a major number theorical input for the paper} \AK{OK}


The Tur\'an-Kubilius inequality is an important tool in analytic number theory (see for example \cite[Lemma~4.1]{E79book}). In this paper, we need an analog of it for additive functions over imaginary quadratic fields whose domains are ideals.
We say that $h: \mf I(K) \to \C$ is \emph{additive} if   $h(\mf u \mf v)=h(\mf u)+h(\mf v)$ for any coprime ideals $\mf u, \mf v$.
We have the following Tur\'an-Kubilius inequality for additive functions. 
It is worth pointing out that a different version of this result is proved in \cite[Chapter~10]{kubiliusbook} for $\Z[i]$.

\begin{theorem}[Tur\'an-Kubilius inequality]\label{TK}
    Let $K=\Q(\sqrt{-d})$ for some squarefree $d\in\N$, $h\colon\mf I(K)\to\C$ be an additive function, $a,Q\in\mathcal{O}_K^{\times}$ with $\NN(a)\leq \NN(Q)$ and $(a)$ coprime to $(Q)$, and also fix $M,N\in\N$. Put $N':=2\NN(Q)(N+1)$. Suppose that $Q\in \mf p$ for all prime ideal $\mf p$ with $\NN(\mf p)\leq M$.
    Then \begin{equation}\label{ltk}
    \begin{split}
       \E_{u\in\OK\colon\NN(u) \leq N}\left| h((Qu+a))-A_{Q,N}\right|^{2}
       \leq %(1+o_{N\to\infty}(1)) B_{N'},
       %B_{N'}(O(\sqrt{Q})+o_{Q;N\to\infty}(1))
       (2+o_{Q;N\to\infty}(1))B_{N'}+O(C_{N}),
    \end{split}
\end{equation}
where the implicit constant depends only on $d$,
$$A_{Q,N}:=\sum_{\NN(\mf p^{k})\leq N,\mf p\nmid Q}h(\mf p^{k})\NN(\mf p)^{-k}(1-\NN(\mf p)^{-1}),$$
$$B_{N}:=\sum_{\NN(\mf p^{k})\leq N, \NN(\mf p)>M, k\geq 1}\vert h(\mf p^{k})\vert^{2}\NN(\mf p^{k})^{-1},$$
and
$$C_{N}:=\frac{1}{\sqrt{N}}\sum_{\NN(\mf p^{k})\leq N', \NN(\mf p)>M}\vert h(\mf p^{k})\vert^{2}\cdot \NN(\mf p^{k})^{-1/2}.$$
\end{theorem}


We remark that Theorem \ref{TK} differs from the classical  Tur\'an-Kubilius inequality in the following ways. The first is that in Theorem~\ref{TK} we deal with averages along arithmetic progressions $Qu+a$ instead of along intervals. The second is that in Theorem~\ref{TK} there is an extra error term $O(C_{K})$ which does not exist in the case when $K=\Q$. We do not know how to remove this term in Proposition~\ref{TK}. However, this term is harmless in our applications.



\subsection{Notation}
%{\color{red} Fill as needed.}
Given a finite set $A$, and a function $f$ on $A$ we use the expectation operator $\E$ to denote its averages:
\[
\E_{u \in A} f(u):=\frac{1}{|A|} \sum_{u \in A} f(u).
\]
For $N \in \N$ we will use the notation $[N]$ to mean the set of integers $\{1,\dots,N\}$.

Given $t \in \R$, we put $e(t):=e^{2\pi i t}$.
For $z \in \C$,
we use $\text{Re}(z)$ and $\text{Im}(z)$ for its real and imaginary parts respectively.   We also write $\exp(z)$ for its series $\sum_{n=0}^{\infty} z^n/n!$.  
%\WS{Maybe we should only keep one of $\exp$ and $e$}

We use standard number theory notation, so for example, $a(n) \ll b(n)$, or equivalently, $a(n)=O(b(n)),$ if there exists some constant $C>0$ such that $|a(n)| \leq C|b(n)|$ for large enough $n \in \N$. The little $o$ notation is also used: we write $a(n)=o(b(n))$ if $\lim_{n \to \infty} \frac{a(n)}{b(n)}=0$.


We use $\U$ to denote the closed unit ball in $\C$, so $\U=\{ z \in \C : |z| \leq 1\}$, and $\S^1$ for its boundary, so $\S^1=\{ z \in \C : |z|=1\}$.

We will typically fix a number field $K$, and $\OK$ will stand for its ring of integers. We will always use parentheses for the field and square brackets for the ring of integers, so typical notations will be $\Q(i)$ and $\Z[\sqrt{-2}]$, for example. Since we usually work with the non-zero elements of the ring of integers $\OK$, we often write $\mathcal{O}_K^{\times}$  to denote the set $\OK \setminus \{0\}$. 

To distinguish when we are working with elements $u,v \in \OK$ as opposed to ideals of $\OK$, we use in the latter gothic letters like $\mf u, \mf v$ or $\mf p$ (which we reserve for prime ideals). To denote the set of all ideals of $\OK$ we use $\mf I(K),$ taking only one representative modulo units.

We let $\NN: K \to K$ denote the norm associated to the number field $K$ (see the number theory section below for the relevant definitions). 

Most of our work applies to quadratic imaginary fields. In the case where $K$ is such a field, we will use $\tau_d$ to denote a generator of its ring of integers, so $\OK=\Z[\tau_d]$.


 




\section{An outline for the proof of the density regularity result}\label{sec: proofstrat}


%In this section we will explain how, using the parametrization that was discussed in the introduction, which conveniently exploits the invariance under dilations of the Pythagorean equation, we will obtain our main results: Theorems~\ref{th1} and \ref{th2} (via their density versions Theorems~\ref{th1density} and \ref{th2density}). 

In this section, we explain the outline of the proof of Theorem \ref{thmain}. For the rest of this Section \ref{sec: proofstrat}, we assume that $K=\Q(\sqrt{-d})$ with $d\in\N$ squarefree.

\subsection{Reductions}




 
%\subsection{Reductions}\WS{My suggestion is to remove the entire Section 2.2, and replace it with what I wrote at the beginning of Section 2.3. The reason is that half of this section is identical from Section 10.3 of \cite{S23} and so we can just cite \cite{S23}. This will also save us from many definitions.}


 
In order to ease the notation for the proof of Theorem \ref{thmain}, we apply some standard reductions on the linear forms. First, by a change of variables, we may assume without loss of generality that $a_{4,2} \neq 0$.  The linear substitution $(m,n) \mapsto (a_{4,2}m,n-a_{4,1}m)$ shows that we may further assume that $a_{4,1}=0$ without loss of generality and without changing the assumptions on our linear forms. 

Next, it follows that $a_{1,1},a_{2,1},a_{3,1}$ are all non-zero so if we now make the change $n \mapsto a_{1,1}a_{2,1}a_{3,1} n$ and factor each $a_{i,1}$ from $L_i$ for $i=1,2,3$, we can further assume that $a_{1,1}=a_{2,1}=a_{3,1}=1$ modulo multiples by elements of $\mathcal{O}_K^{\times}$. Lastly, the change $m \mapsto m-a_{3,2}n$ allows us to reduce Theorem \ref{thmain} to the special case when
$$L_{1}(m,n)L_{2}(m,n)=\ell(m+\alpha n)(m+\beta n) \text{ and } L_{3}(m,n)L_{4}(m,n)=\ell'mn$$
for some $\ell, \ell' \in \mathcal{O}_K^{\times}$ and $\alpha,\beta \in \mathcal{O}_K$.


We first use a variation of the Furstenberg's correspondence principle, that applies to actions of $(\mathcal{O}_K^{\times}, \times)$ (see for example \cite[Theorem~2.8]{bf21b}). This allows us to translate Theorem~\ref{thmain} into the following.

\iffalse
\begin{theorem}\label{firstreduction}
Let $(T_u)_{u \in\mathcal{O}_K^{\times}}$ be a measure preserving multiplicative action of $(\mathcal{O}_K^{\times},\times)$ on a probability space $(X,\mathcal{B},\mu)$,\footnote{Meaning that $T_{u}$ is a measure preserving transformation for all $u \in\mathcal{O}_K^{\times}$ with $T_{1}=id$ and $T_{u}\circ T_{v}=T_{uv}$ for all $u,v\in\mathcal{O}_K^{\times}$.} and let $A \in \mathcal{B}$ with $\mu(A)>0$. Suppose that $q: \mathcal{O}_K^3 \to \C$ admits an $\OK$-factorization in $(x,y)$. Then, there exists a set of positive lower density of pairs $(m,n)\in (\mathcal{O}_K^{\times})^2$ \WS{not defined} so that the elements $ L_1(m,n)\cdot L_2(m,n)$ and $L_3(m,n)\cdot L_4(m,n)$ are distinct, and
%\begin{itemize}
%\item[(i)] The elements $\ell(m^2-n^2)$ and $\ell'mn$ are distinct, and
\begin{equation}\label{measurerec1}
\mu(T_{L_1(m,n)\cdot L_2(m,n)}^{-1}A \cap T_{L_3(m,n)\cdot L_4(m,n)}^{-1}A)>0.
\end{equation}
%\item[(ii)] The elements $\ell(m^2+n^2)$ and $\ell'mn$ are distinct, and
%\begin{equation}\label{measurerec2}
%\mu(T_{\ell(m^2+n^2)}^{-1}A \cap T_{\ell'mn}^{-1}A)>0.
%\end{equation}
%\end{itemize}
% By ``many'', we mean that the set of such pairs $(m, n)$ has positive lower density. \AK{Also, the way that it is written is not good because it is not about (4) only but that the $L_i$'s are distinct.}
\end{theorem}

 
In order to ease the notation and some of the proofs, we apply some standard reductions on the linear forms. Notice that if we want to show that
\[
\overline{d}_{\Phi}((L_1(m,n)\cdot L_2(m,n))^{-1}\Lambda \cap (L_3(m,n)\cdot L_4(m,n))^{-1} \Lambda)>0
\]
when at least one of the linear forms is independent from the others over the number field $K$, we can assume first, without loss of generality that $a_{4,2} \neq 0$.  The linear substitution $(m,n) \mapsto (a_{4,2}m,n-a_{4,1}m)$ shows that we may assume that $a_{4,1}=0$ without loss of generality, and without changing the assumptions on our linear forms. 

Next, it follows that $a_{1,1},a_{2,1},a_{3,2}$ are all non-zero so if we now make the change $n \mapsto a_{1,1}a_{2,1}a_{3,1} n$ and factor each $a_{i,1}$ from $L_i$ for $i=1,2,3$, we can further assume that $a_{1,1}=a_{2,1}=a_{3,1}=1$ modulo multiples by elements of $\mathcal{O}_K^{\times}$. Lastly, the change $m \mapsto m-a_{3,2}n$ allows us to only look at expressions of the form
\[
\overline{d}_{\Phi}((\ell(m+\alpha n)(m+\beta n))^{-1}\Lambda \cap (\ell'mn)^{-1} \Lambda)>0
\]
for $\ell, \ell' \in \mathcal{O}_K^{\times}$ and $\alpha,\beta \in \mathcal{O}_K$.

Thus, with these reductions, it is actually enough to show the following version of Theorem~\ref{firstreduction}:
\fi


\begin{theorem}\label{firstreductionplus}
Let $(T_u)_{u \in\mathcal{O}_K^{\times}}$ be a measure preserving multiplicative action of $(\mathcal{O}_K^{\times},\times)$ on a probability space $(X,\mathcal{B},\mu)$,\footnote{Meaning that $T_{u}$ is a measure preserving transformation for all $u \in\mathcal{O}_K^{\times}$ with $T_{1}=id$ and $T_{u}\circ T_{v}=T_{uv}$ for all $u,v\in\mathcal{O}_K^{\times}$.} and let $A \in \mathcal{B}$ with $\mu(A)>0$. Let $\ell,\ell' \in \mathcal{O}_K^{\times}$ and $\alpha, \beta \in \OK$. Suppose that $q: \mathcal{O}_K^3 \to \C$ admits an $\OK$-factorization in $(x,y)$. Then, there exists a set of positive lower density\footnote{Recall that a set $E \subset (\mathcal{O}_K^{\times})^2$ is said to have positive lower density if $\liminf_{N \to \infty} \frac{|E \cap B_N|}{|B_N|}>0$, where $B_N=\{ u \in \mathcal{O}_K^{\times} : \NN(u)\leq N\}$} of pairs $(m,n)\in (\mathcal{O}_K^{\times})^2$ so that the elements $\ell(m+\alpha n)(m+\beta n)$ and $\ell'mn$ are distinct, and
%\begin{itemize}
%\item[(i)] The elements $\ell(m^2-n^2)$ and $\ell'mn$ are distinct, and
\begin{equation}\label{measurerec1plus}
\mu(T_{\ell(m+\alpha n)(m+\beta n)}^{-1}A \cap T_{\ell'mn}^{-1}A)>0.
\end{equation}
%By ``many'', as was mentioned above, we mean that the set of pairs $(m, n)$ for which \eqref{measurerec1} holds has positive lower density.
%\FM{Change this accordingly.}
\end{theorem}


Our next step is to use the Herglotz-Bochner theorem to reduce Theorem~\ref{firstreductionplus} to a property on the spectrum of the system. Recalling the definition of completely multiplicative function on $\mathcal{O}_K^{\times}$ we introduced before, we denote, through the rest of the paper
%We need to following definition.

%Proving Theorem~\ref{firstreduction} (or its reduced version Theorem~\ref{firstreductionplus}) and deducing Theorem~\ref{thmain} from it, allows us to employ techniques from ergodic theory. First, we note that, since the positivity of measures we wish to show only involves two sets, we can use the Herglotz-Bochner theorem, which will allow us to reduce proving Theorem~\ref{firstreduction} (Theorem~\ref{firstreductionplus} respectively) to positivity results for multiplicative functions on $\mathcal{O}_K^{\times}$ (taking values on $\U$, the complex unit disk), which we now introduce more precisely. 
%\begin{definition}
%Let $f:\mathcal{O}_K^{\times} \to \U$. We say that $f$ is \emph{completely multiplicative} if
%\[
%f(mn)=f(m)\cdot f(n), \quad \text{ for all } m,n \in \OK.
%\] \AK{We defined this on Page 7}
%We say $f$ is \emph{completely multiplicative} if the above holds for all $m, n \in \Z[i]^{\times}$
%\end{definition}
%Throughout the paper, we denote  
\[
\mathcal{M}:= \{ f: \mathcal{O}_K^{\times} \to \S^1 : f \text{ is completely multiplicative}\}.
\]
We give $\mathcal{M}$ the product topology, which makes it into a compact metric space (given that $\S^1$ is). There is a natural identification between the Pontryagin dual of $(K^{\times},\times)$ and $\mathcal{M}$, which we shall use to more easily apply the Herglotz-Bochner theorem. (Indeed, this follows by taking an integral basis and taking common factor, so any term of the form $\frac{p}{q}x_1+\frac{r}{s}x_2$ is of the form $\frac{u}{v}$ for some $u, v \in \mathcal{O}_K^{\times}.$) 

Consider the map $\varphi: K^{\times} \to [0,1]$ given by $\varphi\left(\frac{u}{v}\right):=\mu(T_u^{-1}A \cap T_v^{-1}A)$, for $u, v \in \mathcal{O}_K^{\times}$. It is easy to check that $\varphi$ is well defined and positive definite. Thus, by the Bochner-Herglotz theorem, there exists a finite positive Borel measure $\sigma$ on $\mathcal{M}$ such that $\sigma(\{1\})\geq \mu(A)^2$ (as a consequence of the $L^2$-mean ergodic theorem and the spectral theorem) such that for all $u, v \in \mathcal{O}_K^{\times}$ we have
\[
\int_{\mathcal{M}} f(u) \cdot \overline{f(v)} \ d\sigma(f) = \mu(T_u^{-1}A \cap T_v^{-1}A). 
\]
In particular, provided that none of the linear forms vanish on $m,n$ (which is a set of additive density $0$ in $(m,n)$), we can write
\[
\mu(T_{\ell (m+\alpha n)(m+\beta n)}^{-1}A \cap T_{\ell' mn}^{-1}A) = \int_{\mathcal{M}} f(\ell (m+\alpha n)(m+\beta n)) \cdot \overline{f(\ell' mn)} \ d\sigma(f).
\]
Thus, writing $$S_q:=\left\{(m,n) : mn(m+\alpha n)(m+\beta n) \neq 0 \text{ and } \ell(m+\alpha n)(m+\beta n) \neq \ell'mn\right\}$$ and using the parametrization  discussed above for $x,y$, 
Theorem~\ref{firstreductionplus} follows from the following result.
\begin{theorem}\label{secondreduction}
Let $\sigma$ be a finite positive Borel measure on $\mathcal{M}$ such that 
\begin{equation}\label{positivitymeasures}
    \text{$\sigma(\{1\})>0$ and } \int_{\mathcal{M}} f(u)\cdot \overline{f(v)} \ d\sigma(f)\geq 0 \quad \text{for every } u, v \in \mathcal{O}_K^{\times}.
\end{equation}
Then, for every $\ell, \ell' \in \mathcal{O}_K^{\times}$ and $\alpha, \beta \in \OK$ we have
\begin{equation}\label{multav1}
    \lim_{N \to \infty} \E_{\NN(n),\NN(m) \leq N, \ (m,n) \in S_q}\int_{\mathcal{M}}  f(\ell (m+\alpha n)(m+\beta n)) \cdot \overline{f(\ell' mn)}\ d\sigma(f)>0.
\end{equation}
%\begin{equation}\label{multav2}
%    \lim_{N \to \infty} \E_{\NN(n),\NN(m) \leq N, \ \NN(m)\neq\NN(n)}\int_{\mathcal{M}} f(\ell(m^2+ n^2)) \cdot \overline{f(\ell'mn)} \ d\sigma(f)>0.
%\end{equation}
\end{theorem}
\begin{remark}
The limits in \eqref{multav1} exist thanks to \cite[Theorem~1.12]{S23}, which differs from the situation in the $\N$ case, because of the irreducibility of one of the quadratic forms in $m$ and $n$.
\end{remark}


\subsection{A further break down for Theorem \ref{secondreduction}}\label{s23} 

\iffalse

\WS{

In the rest of Section \ref{s23}, we assume that $K=\Q(\sqrt{-d})$ with $d\in\N$ squarefree.
In order to ease the notation for the proof of Theorem \ref{thmain}, we apply some standard reductions on the linear forms. First, by a change of variable, we may assume without loss of generality that $a_{4,2} \neq 0$.  The linear substitution $(m,n) \mapsto (a_{4,2}m,n-a_{4,1}m)$ shows that we may assume that $a_{4,1}=0$ without loss of generality, and without changing the assumptions on our linear forms. 

Next, it follows that $a_{1,1},a_{2,1},a_{3,2}$ are all non-zero so if we now make the change $n \mapsto a_{1,1}a_{2,1}a_{3,1} n$ and factor each $a_{i,1}$ from $L_i$ for $i=1,2,3$, we can further assume that $a_{1,1}=a_{2,1}=a_{3,1}=1$ modulo multiples by elements of $\mathcal{O}_K^{\times}$. Lastly, the change $m \mapsto m-a_{3,2}n$ allows us to reduce Theorem \ref{thmain} to the special case when
$$L_{1}(m,n)L_{2}(m,n)=\ell(m+\alpha n)(m+\beta n) \text{ and } L_{3}(m,n)L_{4}(m,n)=\ell'mn$$
for some $\ell, \ell' \in \mathcal{O}_K^{\times}$ and $\alpha,\beta \in \mathcal{O}_K$.

The key of Theorem \ref{thmain} is to study recurrence properties for multiplicative functions.
\begin{definition}
Let $f:\mathcal{O}_K^{\times} \to \U$. We say that $f$ is \emph{completely multiplicative} if
\[
f(mn)=f(m)\cdot f(n), \quad \text{ for all } m,n \in \OK.
\]
%We say $f$ is \emph{completely multiplicative} if the above holds for all $m, n \in \Z[i]^{\times}$
\end{definition}
Throughout the paper, we denote  
\[
\mathcal{M}:= \{ f: \mathcal{O}_K^{\times} \to \S^1 : f \text{ is completely multiplicative}\}.
\]
We give $\mathcal{M}$ the product topology, which makes it into a compact metric space (given that $\S^1$ is).

To ensure  $L_{1}(m,n)L_{2}(m,n)$ and $L_{3}(m,n)L_{4}(m,n)$ take different nonzero values, we need to restrict our choices of the parameter $(m,n)$ from the set
\begin{equation}\label{ssqq}
 S_q:=\left\{(m,n)\in K^{2}: mn(m+\alpha n)(m+\beta n) \neq 0 \text{ and } \ell(m+\alpha n)(m+\beta n) \neq \ell'mn\right\}.   
\end{equation}
Following the same discussion as in Section 10.3 of \cite{S23} as well as the simplification discussed above, to prove Theorem \ref{thmain}, it suffices to show the following:

\begin{theorem}\label{secondreduction}
Let $\sigma$ be a finite positive Borel measure on $\mathcal{M}$ such that 
\begin{equation}\label{positivitymeasures}
    \text{$\sigma(\{1\})>0$ and } \int_{\mathcal{M}} f(u)\cdot \overline{f(v)} \ d\sigma(f)\geq 0 \quad \text{for every } u, v \in \mathcal{O}_K^{\times}.  
\end{equation}
Then, for every $\ell, \ell' \in \mathcal{O}_K^{\times}$ and $\alpha, \beta \in \OK$ we have
\begin{equation}\label{multav1}
    \lim_{N \to \infty} \E_{\NN(n),\NN(m) \leq N}\mathds{1}_{S_q}(m,n)\cdot\int_{\mathcal{M}}  f(\ell (m+\alpha n)(m+\beta n)) \cdot \overline{f(\ell' mn)}\ d\sigma(f)>0.
\end{equation}
%\begin{equation}\label{multav2}
%    \lim_{N \to \infty} \E_{\NN(n),\NN(m) \leq N, \ \NN(m)\neq\NN(n)}\int_{\mathcal{M}} f(\ell(m^2+ n^2)) \cdot \overline{f(\ell'mn)} \ d\sigma(f)>0.
%\end{equation}
\end{theorem}
\begin{remark}
The limits in \eqref{multav1} exist thanks to \cite[Theorem~1.12]{S23}, which differs from the situation in the $\N$ case, because of the irreducibility of one of the quadratic forms in $m$ and $n$.
\end{remark}

From this point on we stop following the  path taken in \cite{S23} and apply the schemes of \cite{FKM2,FKM1} instead.
 

}
\fi

%\AK{I don't like ``plan''. ``Strategy of the paper'', ``Strategy of the proof'' would be better (IMO).}

The next part of our proof strategy differs from the path taken in \cite{S23}. We do not attempt to use a decomposition result that works for all elements of $\mathcal{M}$ simultaneously (which, as we discussed, would not cover the case of Pythagorean pairs), but instead divide $\mathcal{M}$ into its aperiodic and pretentious parts. In order to ensure that the splitting is a truly disjoint union, we use the analog of Hal\'asz's theorem for number fields we mentioned: Theorem~\ref{C22intro}. 

We can then use linear concentration estimates that we develop in Section~\ref{SECconcest} to deal with the pretentious part, and results from \cite{S23} to deal with the aperiodic part (which will vanish).

In order to prove Theorem~\ref{secondreduction}, we will take the averages over the grid
\[
\{ (Qm+1, Qn) : m, n \in \mathcal{O}_K^{\times}\},
\]
for some $Q \in \OK$ that will be appropriately chosen later, depending only on the measure $\sigma$. Since we are only concerned with positivity, going along this grid is enough to establish \eqref{multav1}. Now, let $\ell,\ell' \in \mathcal{O}_K^{\times}$ and $\alpha,\beta \in \OK$ be fixed. For $\delta>0$, $f \in \mathcal{M}$ and $Q, m, n \in \mathcal{O}_K^{\times}$, we put
\begin{equation}
A_{\delta}(f,Q; m,n):= w_{\delta}(m,n)\cdot  f\Bigl(\frac{\ell(Qm+1+Q\alpha n)(Qm+1+Q\beta n)}{\ell' (Qm+1)Qn}\Bigr),
\end{equation}
%and
%\begin{equation}
%B_{\delta}(f,Q; m,n):= \tilde{w}_{\delta}(m,n)\cdot f(\ell((Qm+1)^2+(Qn)^2)) \cdot \overline{f(\ell'(Qm+1)Qn)},
%\end{equation}
where $w_{\delta}: (\mathcal{O}_K^{\times})^2 \to [0,1]$ is the weight defined in Lemma~\ref{weights} below which is supported on $S_{q}$. Since $0 \leq w_{\delta}\leq 1 $, and we also have the positivity property \eqref{positivitymeasures}, Theorem~\ref{secondreduction} will follow from the following.

\begin{theorem}\label{thirdreduction}
Let $\sigma$ be a Borel probability measure on $\mathcal{M}$ such that $\sigma(\{1\})>0$. Then, there exist $\delta>0$ and $Q \in \mathcal{O}_K^{\times}$ (depending only on $\sigma$) such that
\begin{equation}\label{adeltapos}
\lim_{N \to \infty} \E_{\NN(n),\NN(m) \leq N} \int_{\mathcal{M}} A_{\delta}(f,Q; m,n) \ d\sigma(f)>0.
%\text{ and}
\end{equation}
%\begin{equation}\label{bdeltapos}
%  \lim_{N \to \infty} \E_{\NN(n),\NN(m) \leq N} \int_{\mathcal{M}} B_{\delta_0}(f,Q; m,n) \ d\sigma(f)>0.
%\end{equation}
\end{theorem}
In order to analyse the limit in \eqref{adeltapos}, we begin by looking into the case where $f$ is aperiodic (the precise definition will be given later in Section~\ref{SEC: numthybackg}). %Using Proposition~\ref{P: appendix} we have:

\begin{proposition}\label{aperiodiclimit0}
Let $f: \mathcal{O}_K^{\times} \to \U$ be an aperiodic completely multiplicative function. Then, for every $\delta>0$ and $Q \in \mathcal{O}_K^{\times}$ we have
\begin{equation}
\lim_{N \to \infty} \E_{\NN(n),\NN(m) \leq N} A_{\delta}(f,Q; m,n)=0.
%\text{ and }
\end{equation}
%\begin{equation}
%  \lim_{N \to \infty} \E_{\NN(n),\NN(m) \leq N} B_{\delta}(f,Q; m,n)=0.  
%\end{equation}
\end{proposition}
Proposition \ref{aperiodiclimit0} is an analog of \cite[Proposition~4.1]{FKM1}.
The proof of Proposition \ref{aperiodiclimit0} is based on Proposition~\ref{P: appendix}, a variation of which was essentially proved in \cite{S23}. We postpone the details until Section \ref{sec: proofmain}.

We now turn our attention to the complement of the aperiodic completely multiplicative functions. As we will deduce from Theorem~\ref{C22intro}, these are exactly the pretentious completely multiplicative functions, so we introduce the following notation:
\begin{equation}
\mathcal{M}_p:=\{ f: \mathcal{O}_K^{\times} \to \S^1 : f \text{ is a pretentious completely multiplicative function}\}.
\end{equation}
Lemma~\ref{measurability} establishes that $\mathcal{M}_p$ is a Borel measurable subset of $\mathcal{M}$. Thus, by Proposition~\ref{aperiodiclimit0} and the dominated convergence theorem, we see that Theorem~\ref{thirdreduction} will follow if we find $\delta>0$ and $Q \in \mathcal{O}_K^{\times}$ such that the analog of \eqref{adeltapos} holds, replacing $\mathcal{M}$ with $\mathcal{M}_p$, i.e.,
\begin{equation}\label{adeltaposPret}
\lim_{N \to \infty} \E_{\NN(n),\NN(m) \leq N} \int_{\mathcal{M}_p} A_{\delta}(f,Q;m,n) \ d\sigma(f)>0.
%\text{ and}
\end{equation}
%\begin{equation}\label{bdeltaposPret}
%  \lim_{N \to \infty} \E_{\NN(n),\NN(m) \leq N} \int_{\mathcal{M}_p} B_{\delta_0}(f,Q;m,n) \ d\sigma(f)>0.
%\end{equation}

When $f$ is pretentious, it exhibits periodicity, which can be exploited with a suitable choice of $Q$. This simplifies matters considerably for the integrand that appears in \eqref{adeltaposPret}. In order to take advantage of this periodicity, we develop, in Section~\ref{SECconcest} a concentration estimate in Proposition~\ref{p25}. We do not restate it here, as it is fairly long and it requires introducing notation and technical terms that we will see later.
%\begin{proposition}\label{concEstMain}
%     Let $f\colon \Z[i]^{\times}\to\U$ be a multiplicative function with $f\sim \chi\cdot n^{i\tau}$ for some Dirichlet character $\chi$ with period $I$ and some $\tau\in\R$. Let 
%    \begin{equation}\label{folnerMult}
%    \Phi_{K}:=\left\{\prod_{\NN(p)\leq K}p^{a_{p}}\colon K<a_{p}\leq 2K\right\}
%    \end{equation}
%    and suppose that $K$ is large enough so that $I$ divides all elements of $\Phi_{K}$. Then, letting $q:=\prod_{\NN(p)\leq K} p$, we have 
%\begin{equation}
%    \begin{split}
%        \limsup_{N\to\infty}\max_{Q\in \Phi_{K}}\E_{\NN(u)\leq N}\vert f(Qu+a)-\chi(a)\NN(Qu)^{i\tau}\exp(F_{N}(f,K))\vert\ll \D(f,\chi\cdot \NN(u)^{i\tau};K,\infty)+K^{-1/2}
%    \end{split}
%\end{equation}
%for any $a\in\Z[i]^{\times}$ with $\NN(a)\leq \NN(Q)$ and $a \mod I\in (\Z[i]/I)^{\times}$,
%where the implicit constant is absolute and
%$$F_{N}(f,K):=\sum_{\NN(p)\leq N, p \nmid Q}\frac{1}{\NN(p)}(f(p)\cdot\overline{\chi}(p)\cdot \NN(p)^{-i\tau}-1).$$
%\end{proposition}
We mention in passing that we will use it mostly for the value $a=1$, and that it will be relevant for us that the implicit constant in the statement is independent of $K$; as well as having independence of the function $F_N(f,k)$ from $Q$ (all these terms will be introduced later).
%Note that we will mostly use this result with $a=1$.
%\begin{remark}
%It is important that the implicit constant is independent of $K$ and that the quantity $F_N(f,K)$ does not depend on $Q$ as long as $Q \in \Phi_K$ and $q \mid Q$.
%\end{remark}

To establish \eqref{adeltaposPret}, we further split the integral into two parts. On the one hand, we consider the multiplicative functions that are not purely Archimedean characters $(\NN(u)^{i\tau})_{u \in \mathcal{O}_K^{\times}}$, $\tau \in \R$, where the concentration estimate in Proposition~\ref{p25} allows us to show that their contribution is essentially non-negative if $Q$ is highly divisible. The other part is supported on Archimedean characters $\mathcal{A}$, which we define in \eqref{archimedean} below. Using the fact that $\sigma(\{1\})>0$ and for some $\delta>0$ small enough, the weight $w_{\delta}$ nullifies the effect of non-trivial Archimedean characters. 

To carry this out we will make good use of the properties of the multiplicative F\o lner sequence defined in \eqref{folnerMult} below. On a first reading, one can think of it as a suitable analog of the multiplicative F\o lner sequence
\[
\Psi_K:=\{ p_1^{a_1}\dots p_k^{a_k} : K \leq a_i \leq 2K, i=1,\dots,K\}
\]
in the integers.
Let
\begin{equation}\label{archimedean}
\mathcal{A}:= \left\{ (\NN(u)^{it})_{u \in \mathcal{O}_K^{\times}} : t \in \R\right\}.
\end{equation}
We will later show (it will follow from the proof of Proposition~\ref{fffss} given in Section~\ref{sec: proofmain}) that $\mathcal{A}$ is a Borel measurable subset of $\mathcal{M}$, but assuming that this is the case for now, the next step is to use the concentration estimate in Proposition~\ref{p25} to obtain the following result.

 
\begin{proposition}\label{zeroaveragesadeltabdelta}
Let $f \in \mathcal{M}_p \setminus \mathcal{A}$, $\delta>0$, $\ell,\ell' \in \mathcal{O}_K^{\times}$, and $\Phi_M$ as in \eqref{folnerMult}. Then,
\[
\lim_{M \to \infty} \E_{Q \in \Phi_M} \lim_{N \to \infty} \E_{\NN(n),\NN(m) \leq N} A_{\delta}(f,Q; m,n)=0.
%, \text{ and }
\]
%\[
%\lim_{K \to \infty} \E_{Q \in \Phi_K} \lim_{N \to \infty} \E_{\NN(n),\NN(m) \leq N} B_{\delta}(f,Q; m,n)=0.
%\]
\end{proposition}
Proposition~\ref{zeroaveragesadeltabdelta} can be viewed as a variation of \cite[Lemma~2.6]{FKM1} 
and its detail is postponed to Section \ref{sec: proofmain}. The essential idea is to use the concentration estimate in Proposition~\ref{p25} to replace it by an expression of the form $C_{\ell,\ell'}\cdot g(Q)\cdot \NN(Q)^{i\tau}$, for some $C_{\ell,\ell'} \in \U$ and $\tau \in \R$. Moreover, we can have $g \notin \mathcal{A}$, so the remaining outer limit will give us convergence to $0$ as $M \to \infty$ using Lemma~\ref{multav0}, whose proof is also deferred to Section~\ref{sec: proofmain}.

With all this taken into consideration we can deduce, using the dominated convergence theorem twice (as the two relevant limits exist), we have the following.

\begin{corollary}
Let $(\Phi_M)$ and $\mathcal{A}$ be given by \eqref{folnerMult} and \eqref{archimedean} respectively. Let $\sigma$ be a Borel probability measure on $\mathcal{M}_p$. Then, for every $\delta>0$ we have
\[
\lim_{M \to \infty} \E_{Q \in \Phi_M} \lim_{N \to \infty} \E_{\NN(n),\NN(m) \leq N} \int_{\mathcal{M}_p \setminus \mathcal{A}} A_{\delta}(f,Q; m,n) \ d\sigma(f)=0. 
%\text{ and }
\]
%\[
%\lim_{K \to \infty} \E_{Q \in \Phi_K} \lim_{N \to \infty} \E_{\NN(n),\NN(m) \leq N} \int_{\mathcal{M}_p \setminus \mathcal{A}} B_{\delta}(f,Q; m,n) \ d\sigma(f)=0.
%\]
\end{corollary}
Finally, only the functions from $\mathcal{A}$ remain. Here is where the weight $w_{\delta}$ helps with the positivity.

 \begin{proposition}\label{fffss}
Let $\sigma$ be a Borel probability measure on $\mathcal{M}$ such that $\sigma(\{1\})>0$ and $\mathcal{A}$ be as in \eqref{archimedean}. Then, there exist $\delta$ and $\rho$, depending only on $\sigma$, such that
\begin{equation}
\liminf_{N \to \infty} \inf_{Q \in \mathcal{O}_K^{\times}} \text{Re} \left( \E_{\NN(n),\NN(m) \leq N} \int_{\mathcal{A}} A_{\delta}(f,Q;m,n)\right) \geq \rho.
\end{equation}
\end{proposition}
%Finally, we sketch the proof of Theorem~\ref{thirdreduction}, assuming the previous results (with the proofs deferred). 
%Indeed, we can find $\delta_0,\rho_0>0$ so that
%\[
%\liminf_{M \to \infty}\E_{Q \in \Phi_M} \lim_{N \to \infty} \E_{\NN(n),\NN(m) \leq N} \int_{\mathcal{M}_p \setminus \mathcal{A}} A_{\delta}(f,Q; m,n) \ d\sigma(f) \geq \rho_0,%\text{ and}
%\]
%\[
%\liminf_{K \to \infty}\E_{Q \in \Phi_K} \lim_{N \to \infty} \E_{\NN(n),\NN(m) \leq N} \int_{\mathcal{M}_p \setminus \mathcal{A}} B_{\delta}(f,Q; m,n) \ d\sigma(f) \geq \rho_0.
%\]
%from where Theorem~\ref{thirdreduction} is easily deduced.

Proposition \ref{fffss} can be viewed as a variation of 
\cite[Lemma 2.8]{FKM1} and
again we defer its proof  to Section \ref{sec: proofmain}. It is here where we make essential use of the weight function $w_{\delta}$ to obtain positivity for the averages whose limit does not necessarily exist because of the nature of the archimedean characters.

%\AK{Last paragraph:
Assuming the previous results, Theorem~\ref{thirdreduction} easily follows from the fact that we can find $\delta_0,\rho_0>0$ so that
\[
\liminf_{M \to \infty}\E_{Q \in \Phi_M} \lim_{N \to \infty} \E_{\NN(n),\NN(m) \leq N} \int_{\mathcal{M}_p \setminus \mathcal{A}} A_{\delta_0}(f,Q; m,n) \ d\sigma(f) \geq \rho_0.%\text{ and}
\]

\subsection{Averaging schemes}\label{s:as}
In this subsection we wish to highlight the fact that the choice of the averaging scheme is very important to us for a number of reasons. 
In order to make the discussion easier to follow, we will focus our attention on the Gaussian integers, but the same points we shall discuss equally apply to other quadratic imaginary fields. 
%We remark that in order to make the above mentioned method work, the choice of the averaging scheme is very important. For simplicity we focus  on the case when $\OK=\Z[i]$. 
To study the averages of a multiplicative function of the form
$$\lim_{N\to\infty}\frac{1}{\vert\Phi_{N}\vert}\sum_{u\in\Phi_{N}}f(u),$$
there are at least 4 natural choices for the sequence $(\Phi_{N})_{N}$:
\begin{enumerate}
    \item   symmetric  balls $D_{N}:=\{a+bi\colon a,b\in\Z, a^{2}+b^{2}\leq N\}$;
    \item   symmetric boxes $B_{N}:=\{a+bi \colon -N\leq a,b\leq N\}$;
     \item  asymmetric  balls $D_{N}^{+}:=\{a+bi\colon a,b\in\N, a^{2}+b^{2}\leq N\}$;
    \item   asymmetric boxes $B_{N}^{+}:=\{a+bi \colon 1\leq a,b\leq N\}$.
\end{enumerate}


In \cite{S18,S23}, all the averages are taken over boxes. The  advantage  in doing so  is that the box average is well defined for any number field while the ball averages can only be defined for quadratic imaginary fields, as for a general number field $K$, the set $\{u\in\OK\colon \NN(u)\leq N\}$ can be infinite. However,  we are unable to prove an analog of Theorem \ref{characterization} if we define aperiodic functions using box averages as was done in \cite{S18,S23}, mainly because we do not know if an analog of Hal\'asz's Theorem (Theorem \ref{C22intro}) holds for box averages. Therefore, in this paper we favor ball averages over box averages.

Next we explain why it is more convenient to work with symmetric balls instead of the asymmetric ones. Consider the multiplicative function $f(u)=\frac{u}{\NN(u)^{1/2}}$. One can show that $f$ does not have finite distance to the product of Dirichlet characters and Archimedean characters either.  On the other hand, 
the average of $f$ over asymmetric boxes does not converge to 0 as $f$ takes values only in the first quadrant. So an analog of Hal\'asz's theorem fails to hold in this setting, unless we expand the definition of pretentious functions to include functions $f$ of this form (however we do not have this issue for symmetric ball averages as the symmetric ball average of $f$ does converge to 0, as can be checked with a Riemann sums argument).
Because of the above reasons, in this paper we choose to work with (i): averages over symmetric balls.



%\WS{(1) I think we should put two things together, the first is that why we average over full balls instead of 1/4 balls (as is explained below), the second is why we take ball averages instead of box averages. Also we could just mention in the introduction that the averging scheme is very important for us, and then leave the details to later sections.} \AK{Ditto}
\iffalse

To begin, one may try to consider averages on the first quadrant of the Gaussian integers $\Z[i]_+:=\{ a+bi : a, b \geq 0\}$. Given the geometry in $\C$, this seems reasonable. However, in this case, the multiplicative function $g(u)=\frac{u}{|u|}$ will not average to $0$, for
\[
\lim_{N \to \infty}  \frac{1}{N^2}\sum_{n,m=1}^N g(n+mi)
\]
is a non-zero element of the first quadrant with norm bigger than $0$. This creates an issue because $g$ does not have finite distance to one of the pretentious functions we will define later. Therefore, if we do not discard it and use a different averaging scheme, the class of functions with non-zero average would contain too many elements, and an analog of Hal\'asz's theorem would fail to hold. On the other hand, it is clear that this issue is not solved by averaging over balls (given by the usual norm) intersected with the first quadrant, as the same function $g(u)=\frac{u}{|u|}$ will fail to average to zero.
%\AK{I want more details here}

This naturally leads one to attempt to  average over full balls or rectangles in $\Z[i]$, as the problem we just mentioned is circumvented this problem. Indeed,
\[
\lim_{N \to \infty} \frac{1}{(2N+1)^2}\sum_{u \in [\pm N]^2} g(u) = \lim_{N \to \infty} \frac{1}{N}\sum_{\NN(u) \leq N} g(u) = 0.
\]
The explanation here is quite simple: both these sets enjoy a F\o lner-type property of absorbing the units of $\Z[i]$. More concretely, they are invariant under multiplication by $i$. However, $g(i) \neq 1$, so these averages necessarily have $0$ limit.

One may still worry that functions such as $g_m(u):=\left(\frac{u}{|u|}\right)^{4m}$, for $m \in \Z \setminus \{0\}$ may have non-zero averages, since the multiplication by $i$ trick no longer works. However, the following Riemann sums argument, which relies on the fact that $g_m(u)=g_m\left(\frac{u}{N}\right)$ for all $N \in \N$ still shows that they average to $0$ in the limit:
\[
\lim_{N \to \infty} \frac{1}{N} \sum_{\NN(u) \leq N} g_m(u) = \lim_{N \to \infty} \frac{1}{N} \sum_{\NN(u) \leq N} g_m\left(\frac{u}{N}\right) = \int_0^1\int_0^{2\pi} r \cdot f_m(r,\theta) \ drd\theta,\]
where we use the Riemann sums on the ball $\{ z \in\C : |z| \leq 1\}$, and the function $f_m(z)=\left(\frac{z}{|z|}\right)^{4m}$ interepreted with polar coordinates. This last integral is equal to (by the definition of $f_m(z)$)
\[ \int_0^{2\pi} \int_0^1 r^2 e^{i4m\theta} \ drd\theta = 0,
\]
since $m \neq 0$. This ensures that the pretentious functions that will be properly defined in Section~\ref{SEChalasz} will be exactly the only ones that fail to have zero average in the limit.

Lastly, one may wonder why work with balls instead of boxes. The answer here is not so straightforward: for technical reasons that will become more apparent in Section~\ref{SEChalasz}, working with balls instead of squares is more suitable for a set of the arguments we will use.

\fi
\section{Number theoretical background}\label{SEC: numthybackg}
In this section we will review some basic number theoretical notions and notation, and also introduce some basic estimates we will need to make use of in the sequel.
%\subsection{Basic facts on number fields and their rings of integers}\label{s31}

An \emph{(algebraic) number field} $K$ is a finite degree (and hence algebraic) field extension of the field of rational numbers $\Q$. The \emph{ring of integers} $\OK$ of a number field $K$ is the ring of all \emph{integral elements} in $K$ (i.e., roots of polynomials with integer coefficients and leading coefficient 1). Let $D=[K:\Q]$ denote the degree of the extension. It is classical that there exists an \emph{integral basis} $\mathcal{B}=\{b_{1},\dots,b_{D}\}$ of $\OK$, i.e., a basis of the $\Q$-vector space $K$ such that each element $x\in\OK$ can be uniquely represented as $x=\sum_{i=1}^{D}c_{i}b_{i}$ for some $c_{i}\in\Z$.   

Let $\iota\colon\Z^{D}\to\OK$ be the map given by $\iota(n_{1},\dots,n_{D})=n_{1}b_{1}+\dots+n_{D}b_{D}$. For $x\in K$, let $A_{x}$ be the unique $d\times d$ matrix such that $\begin{bmatrix}
    xb_{1} \\ \vdots\\ xb_{D}
\end{bmatrix}=A_{x}\begin{bmatrix}
    b_{1} \\ \vdots\\ b_{D}
\end{bmatrix}$. The $K$-\emph{norm} of $x\in K$ is defined to be $\NN_{K}(x):=\det(A_{x})$. When there is no risk of confusion regarding the underlying field $K$, we simply write $\NN$ instead of $\NN_{K}$. 



In this paper, our main focus is on quadratic fields, i.e., number fields of the form $K=\Q(\sqrt{-d})$ for some squarefree $d\in\Z$. 
%{\color{red} To do: throughout this paper, change all the $d$ to be square-free (positive) integers}
In this case, we have $\OK=\Z[\tau_{d}]$ and every $z\in\OK$ can be written as $m+n\tau_{d}$ for some $m,n\in\Z$ in a unique way. The element $\tau_d$ is given by $\tau_{d}=\sqrt{-d}$ if $d\equiv 1,2 \mod 4$ and $\tau_{d}=\frac{1+\sqrt{-d}}{2}$ if $d\equiv 3 \mod 4$. Moreover, we have that $\NN(m+n\tau_{d})$ is equal to $m^{2}+dn^{2}$ if $d\equiv 1,2 \mod 4$ and equal to $m^{2}+mn+\frac{d+1}{4}n^{2}$ if $d\equiv 3 \mod 4$.



We will use the following classical result on properties of the norm $\NN$.

\begin{lemma}
    Let $K$ be a number field. Then $\NN(xy)=\NN(x)\NN(y)$ for all $x,y\in K$. Also, for any $x \in \OK$ we have $\NN(x)\in\Z$.
\end{lemma}

We say that $\e\in \OK$ is a unit if $\NN(\e)=\pm 1$. 
It follows from Dirichlet's unit theorem \cite[Theorem~8.1]{neukirch99} that $\OK$ has finitely many units if and only if $K=\Q$ or $\Q(\sqrt{-d})$ for some square-free $d\in\N$. Moreover, 
we have a complete description for the units in this case.

\begin{lemma}\label{l32}
    Let $K=\Q(\sqrt{-d})$ for some squarefree $d\in\Z$. 
    \begin{enumerate}
        \item If $d=1$, then the units of $\OK$ are $\pm 1, \pm i$.
        \item If $d=3$, then the units of $\OK$ are $e^{k\pi i/3}, 0\leq k\leq 5$.
        \item If $d>0, d\neq 1,3$, then the units of $\OK$ are $\pm 1$.
    \end{enumerate}
\end{lemma}
This follows easily from computing the elements in $\OK$ whose norm is equal to $\pm 1$ and solving the resulting diophantine equations, which is straightforward, as in the imaginary quadratic case the norm is non-negative. 
%The lemma follows from https://math.stackexchange.com/questions/2732683/units-in-complex-quadratic-fields-mathbb-q-sqrt-d-with-d-equiv-1-pmod-4. 


 

Let $K$ be a number field. We use $\mathfrak{I}(K)$ to denote the set of all  ideals of $\OK$, $\mathfrak{Q}(K)$ to denote the set of all fractional ideals of $K$ (recall that a fractional ideal has the form $x^{-1}I$ for some $x \in \mathcal{O}_K^{\times}$ and non-trivial ideal $I \subseteq \OK)$. %{\color{red} define fractional ideal?}
For $I\in \mathfrak{I}(K)$, we say that its \emph{norm} is $\NN(I):=[\OK:I]$. In the case where $x^{-1}I\in \mathfrak{I}(K)$, the \emph{norm} is given by $\NN(x^{-1}I):=\NN(I)/\NN(x)$.

We say that two ideals $I, J \in \mf I(K)$ are \emph{coprime} if $I+J=(1)$. For Dedekind domains, this is equivalent to $I,J$ sharing no elements in their factorization into prime ideals.
 Let $a,b\in\OK^{\times}$, $\mf p\in \mf I(K)$ and $k\in\N$. We write $\mf p\vert a$ if $a\in\mf p$ (or equivalently $\mf p\vert (a)$). Write $\mf p^{k}|| a$ if $k$ is the largest integer for which $\mf p^{k}\vert a$. 
 



We recall a definition.
\begin{definition}
Let $K$ be a number field. %and write $\OK$ for its ring of integers. 
We say that the quotient $\mf Q(K)/\mf I(K)$ is the \emph{ideal class group of} $K$.
\end{definition}
Given a number field $K$, its ideal class group $G$ is always finite (e.g., see \cite[Theorem~6.3]{neukirch99}), and we say that the order of $G$ is the \emph{class number} of $K$.

The ideal class group of a number field $K$ depends on the algebraic properties of the number field under consideration. For example, if $K=\Q(i)$, its ideal class group is trivial (this is a consequence of $\Z[i]$ being a PID), but for $\Q(\sqrt{-5})$ it is isomorphic to $\Z/2\Z$.
%{\color{red} Define ideal class group, and ideal class numbers. And then give two examples: one is $\Q(i)$ whose ideal class group is trivial. one is $\Q(\sqrt{-5})$ whose ideal class group is $\Z/2\Z$.}

The following lemma will be used in the sequel.
\begin{lemma}\label{cl2}
    Let $K$ be a number field and $D=[K:\Q]$. For any $I\in\mathfrak{I}(K)$, there exists a finite union of $D$-dimensional infinite arithmetic progressions $P$ such that for $u,v\in\OK^{\times}$
$$\frac{(u)}{(v)}I\in\mathfrak{I}(K)\Leftrightarrow \frac{(u)}{(v)}=\frac{(u')}{(\NN(I))} \text{ for some } u'\in P.$$
In fact, $P$ is the set of $u'$ such that $\NN(I)^{D-1}\vert\NN(u')$.
\end{lemma} 
\begin{proof}
By multiplying with conjugates of $v$ if necessary, we may assume without loss of generality that $v\in\Z$. Then $\frac{(u)}{(v)}I\in\mathfrak{I}(K)\Leftrightarrow v^{D}\vert\NN(u)\NN(I)$.
Write $\NN(I)=u_{1}^{t_{1}}\dots u_{n}^{t_{n}}$ for some disticnt prime numbers $u_{i}$ and some powers $t_{i}\in\N$. 

We now proceed by cases. If a prime number $t$ is such that $t\vert v$ but $t\nmid \NN(I)$, then $t^{D}\vert \NN(u)$ and so $u/t\in\mathcal{O}_{K}$. So $\frac{(u)}{(v)}=\frac{(u/t)}{(v/t)}$. If $u_{i}^{t_{i}+1}\vert q$ for some $i$, then $p^{D}\vert\NN(u)\NN(I)$ implies that $u_{i}^{D}\vert \NN(p)$ and again $\frac{(u)}{(v)}=\frac{(u/u_{i})}{(v/u_{i})}$.
Thus, we may freely assume that $v=\NN(I)$.

We have that $v^{D}\vert\NN(u)\NN(I)\Leftrightarrow \NN(I)^{D-1}\vert\NN(u)$. The set of such $u$ is clearly a finite union of $D$-dimensional infinite arithmetic progressions. This completes the proof. 
\end{proof}

We conclude this section with a counting property on the number of 
 ideals in a given ideal class. These results are given in the following lemmas.

 \begin{lemma}[Theorem~2, \cite{MVO07}]\label{mv7}
     Let $K$ be a number field of degree $D$ and $C$ be an ideal class of it. For $x>0$, let $N(x,C)$ denote the number of ideals in the class $C$ whose norms are at most $x$. Then there exist constants $C_1, C_2>0$ depending on the number field $K$ only, such that 
     $$C_{1}(x^{\frac{1}{D}}-C_{2})^{D}\leq N(x,C)\leq C_{1}(x^{\frac{1}{D}}+C_{2})^{D}$$
      for all $x>1$ for some $C_{1}>0$ and $0<C_{2}<1$ depending only on $K$.
 \end{lemma}

 

 As a consequence of Lemma \ref{mv7}, we have (see  also \cite[Chapter 6, Theorem 39]{marcusbook} for a reference):
\begin{corollary}\label{lem: bounds for number of ideals of norm at most x}
Let $K$ be a number field of degree $D$. Then, there exists a universal constant $\gamma_K>0$ and some $0 \leq \eta<1$ %and some $0<\delta<1$ 
such that for any ideal class $C$ of $K$, we have that %the following holds:
\[
\frac{1}{x}\sum_{\mf u\in C\colon \NN(\mf u) \leq x}1 = \gamma_K+o(x^{\eta}).
\]
\end{corollary}
\begin{remark}
In particular, with the terminology from \cite{lucht2001}, Corollary~\ref{lem: bounds for number of ideals of norm at most x} above implies that $\mathfrak I(K)$ is an Axiom A arithmetic group using the standard norm of an ideal for $A=\gamma_{K}$. %It is also worth pointing out that with the notation as in Theorem~\ref{C22intro}, the constant $\gamma_K$ coincides with the residue $R_1$.
\end{remark}
%\begin{proof}
%This is a classical result. It follows from Lemma~\ref{mv7} see also \cite[Chapter 6, Theorem 39]{marcusbook} for a reference.
%\end{proof}



 \begin{lemma}\label{mv8}
     For any number field $K$, there exists a constant $C$ such that for any  $x>0$, the number of ideals of $\OK$ which is a power of a prime ideal and whose norm is $x$ is at most $[K\colon\Q]$.
 \end{lemma}
 \begin{proof}
 Let $D:=[K\colon\Q]$.
     By \cite[Lemma~2.9]{S23}, the norm of every prime ideal is a power of prime in $\N$. So we may assume that $x=p^{r}$ for some $r\in\N$ and prime $p\in\N$. Suppose that $\NN(\mf p^{k})=p^{r}$, then we must have that $r=kt$ for some $t\in\N$ and $\NN(\mf p)=p^{t}$. Again by \cite[Lemma~2.9]{S23}, we must have that $t\leq D$ and there are at most $D$ such $\mf p$. Now for each such $\mf p$, there is at most one choice of $k\in\N$ such that $\NN(\mf p^{k})=p^{r}$. So the number of ideals of the form $\mf p^{k}$ with norm equal to $p^{r}$ is at most $D$.
 \end{proof}

%\subsection{Sum of reciprocals of norms of primes and the Dedekind zeta function}

 


\begin{convention}
Throughout the paper, $\mf p$ and $\mf q$ always denote prime ideals of $\mathcal{O}_K$. Whenever we sum over $\mf p$ or $\mf q$, this sum is assumed to be taken along prime ideals.
\end{convention}


\iffalse

\begin{lemma}\label{ggee}
Let $K$ be a number field. We have
\[
\sum_{\mf p\in\mf I(K)}\frac{1}{\NN(\mf p)}=\infty \text{ and } \sum_{\mf p\in\mf I(K)}\frac{1}{\NN(\mf p)^{\sigma}}<\infty
\]
for all $\sigma>1$.
\end{lemma}
\begin{proof}
    We say that $p\in \OK$ is a \emph{prime element} if whenever $p=ab$ for some $a,b\in\OK$, one of $a,b$ must be a unit.
Then \cite[Theorem~2.18]{S23} says that
\begin{equation}\nonumber
    \sum_{ p\colon \text{ prime element in } \mathcal{O}_K}\frac{1}{\NN(p)}=\infty %\text{ and } \sum_{ p\colon \text{ prime in } \mathcal{O}_K}\frac{1}{\NN(p)^{\sigma}}<\infty
\end{equation}
    %for all $\sigma>1$ 
    (for a proof, see \cite[pp. 148--149]{murtyesmondebook}). This immediately implies that 
    $$\sum_{\mf p\in\mf I(K)}\frac{1}{\NN(\mf p)}\geq \sum_{ p\colon \text{ prime element in } \mathcal{O}_K}\frac{1}{\NN((p))}=\sum_{ p\colon \text{ prime element in } \mathcal{O}_K}\frac{1}{\NN(p)}=\infty.$$
For $\sigma>1$, it follows from Lemma \ref{mv7} that  there exists some constant $C$ depending only on $K$ such that
$$\sum_{\mf p\in\mf I(K)}\frac{1}{\NN(\mf p)}\leq C\sum_{\text{$p:$ prime element in $\N$}}\frac{1}{p^{\sigma}}<\infty.$$
\end{proof}


% where we used the bound for $|h(\mf p)|$ we gave at the beginning of the proof, and $T_0$ is some constant that depends on the number field (Lemma \ref{mv7}).

% \begin{remark}
%It is important to note that the sum of reciprocals of norms of primes diverges:
%\[
%\sum_{\mf p\colon \text{ prime in } \mathcal{O}_K}\frac{1}{\NN(\mf p)}=\infty,
%\]
%cf. \cite[Theorem~2.18]{S23}, and for a proof, see \cite[pp. 148--149]{murtyesmondebook}.
%\end{remark}

  

  The \emph{Dedekind zeta function} of a number field $K$ is the map $\zeta_K$ given by
 \begin{equation}\label{d1}
    \zeta_K(s):=\sum_{\mf u\in \mf I (K)}\NN(\mf u)^{-s}.
 \end{equation}
For simplicity, in this paper we write $\zeta$ instead of $\zeta_{K}$ since the underlying number field will always be clear from the context.
  As in the integer case, we have (see \cite[Theorem~5.11]{neukirch99}, for example) the following.
  
 \begin{lemma}\label{zp}
     The function $\zeta$ only has a simple pole at $s=1$ and admits an analytic extension in the rest of the complex plane.
 \end{lemma}
A few comments regarding this result: note that, just as is in the case of the standard Riemann-zeta function, we easily notice that $\zeta(s)$ is defined for $\text{Re}(s)>1$, and that it has a simple pole at $s=1$. We then extend it meromorphically on the whole complex plane using the analog of the functional equation that the Riemann-zeta function satisfies, and abusing notation, we keep the same symbol $\zeta$ for it. It is important to highlight that the functional equation is slightly different because our context is that of a number field $K$.

\fi

 %\WS{this deifnition is not rigorous. $\zeta$ should be first defined on a set where the sum converges, and then extend to the entire palne.}
 %The following estimate for the Dedekind zeta function holds.
%\FM{Double check this proof}
\iffalse
\begin{lemma}\label{l1}
Let $M>0$ and $\sigma,\tau\in\R$. We have
\[\zeta(\sigma+i\tau)=O\left(\frac{1}{\sigma-1}\right)\]
as $\sigma\to 1^+$ uniformly for $-M\leq \tau\leq M$.
\end{lemma}
\begin{proof}
%{\color{red} The first method doesn't give us what we want, it can be made into a remark or deleted altogether.}
%We have the following inequality:
%\[
%\left|\sum_{u \in \Z[i]\setminus{\{0\}}} \NN(u)^{-s}\right| \leq \sum_{n,m \in \Z \setminus \{0\}} \frac{1}{(n^2+m^2)^{\sigma}} \leq \sum_{n,m \in \Z \setminus \{0\}} \frac{1}{(2\sqrt{n^2m^2})^{\sigma}}=\frac{4}{2^{\sigma}}\left(\sum_{n=1}^{\infty}\frac{1}{n^{\sigma}}\right)^2,
%\]
%where we used the AM-GM inequality in the third step. Using the integral test on the last term of the inequality above, we get an upper bound of
%\[
%\frac{4}{2^{\sigma}}\frac{1}{(1-\sigma)^2}.
%\]
%Note: this argument could be generalized to get an upper bound of the form $O((\sigma-1)^d)$, where $d=[K:\Q],$ the degree of the extension of the number field, whose ring of integers $\mathcal{O}_K$ we would be interested in.

%(Question: is this inequality good enough for our purposes?)

%Using Lemma \ref{mv7}, which allows us to find a constant $K$ depending only on the number field used, bounding the number of elements of a certain norm:
Let $s=\sigma+i\tau$. Then, recall that we can write
\[
\left|\sum_{\mf u\in\mf I(K)} \NN(\mf u)^{-s}\right| =\lim_{x \to \infty} \sum_{1 \leq \NN(\mf u) \leq x} \NN(u)^{-s}.
\]
We now apply the Riemann-Stieltjes summation by parts trick, and we see that
\[
\frac{\zeta(s)}{s} = \int_1^{\infty} \frac{1}{x^{s+1}} H(x) \ dx,
\]
where $H(x)=\sum_{1 \leq \NN(\mf u) \leq x} 1$. Corollary~\ref{lem: bounds for number of ideals of norm at most x} implies that $H(x)=O(x)$, so the resulting integral is of order
\[
\int_1^{\infty} \frac{1}{x^{\sigma}} \ dx =\frac{1}{\sigma-1} .
\]
Notice then that for $s=1+i\tau$, $|s| \leq M+1$ is a uniform bound, so that 
\[
\zeta(\sigma+i\tau)=O\left(\frac{1}{\sigma-1}\right)
\]
uniformly in $-M \leq \tau \leq M$, as $\sigma \to 1^+$.
%\WS{Why this is bounded by a constant? It seems to me that the bound is of the form $C\sqrt{x}$. This issue appears many times in this paper. However, maybe we don't need a proof of this result because it follows from Theorem 4.1 by setting $g\equiv 1, L=1, \alpha=0$, $C=\lim_{N\to\infty} \frac{1}{x}\sum_{\NN(\mf u)\leq x} 1$ (the limit exists and is positive by Lemma \ref{mv7}) and by noticing that $\vert\frac{\sigma-1}{s-1}\vert\leq 1$}
%Lastly, by the integral test, we conclude that the last sum is of order $O\left( \frac{1}{\sigma-1}\right)$.
\end{proof}
\fi
%We also need the following estimates on certain tails of the Dedekind zeta function.

We conclude this section by some estimates related to the distributions of prime ideals which will be used in later sections.

%\WS{The old Lemma 3.11 is now part (iv) of the new Lemma 3.11. I commented the old proof and write a new one. Please double check. I also included some other properties that we will use later, whose proofs are all very similar.}
%{\color{red} \cite[Corollary 2.4]{lucht2001} implies the first statement of Lemma 4.8
%}
%I am not sure but \cite[Lemma 9.2]{lucht2001} could imply (v) with an appropriate choice for $h$}
%\WS{Only (i), (ii) and (vi) are used. I removed (iii), (iv), (v)}
\begin{lemma}\label{l2}
%Let $\varepsilon>0$. Then 
As $N\to\infty$, we have
\begin{enumerate}
    \item $\sum_{\NN(\mf p^{k})\leq N} 1=O\left(\frac{N}{\log N}\right);$
    \item $\sum_{\NN(\mf p^{k})\leq N} \NN(\mf p^{k})^{-1}=O(\log\log N);$
 %   \item $\sum_{N^{\varepsilon}\leq \NN(\mf p)\leq N}\NN(\mf p)^{-1}
% =O_{\e}(1);$
% \item $\sum_{\NN(\mf p)>N^{\varepsilon}}\NN(\mf p)^{-(1+(\log N)^{-1})}=O_{\e}(1);$
% \item  $\frac{1}{\log N}\sum_{\NN(\mf p)<N^{\varepsilon}}\frac{\log\NN(\mf p)}{\NN(\mf p)} =O(\e);$
 \item $\sum_{\NN(\mf p^{k}\mf q^{\ell})\leq N, \mf p\neq \mf q} 1=O\left(\frac{N\log\log N}{\log N}\right)$.
\end{enumerate}
%
%
%\[
%\sum_{\NN(\mf p)>x^{\varepsilon}}\NN(\mf p)^{-(1+(\log x)^{-1})}=O(1)
%\]
%and 
%\[\sum_{x^{\varepsilon}\leq \NN(\mf p)\leq x}\NN(\mf p)^{-1}=O(1).\]
\end{lemma}
\begin{proof}
For $i\in\N$, let  $\Pi_{i}(N)$ denote the set of integers  $x\in\{1,\dots,n\}$ with at most $i$ distinct factors. Let $\pi_{i}(N):=\vert\Pi_{i}(N)\vert$. 
Throughout the proof, we make use of the following facts without explicitly stating them:
\begin{itemize}
    \item It is known that (see \cite[Chapter II.6]{tenebaum_2015}, for example)
$$\pi_{i}(N)\sim C_{i}\frac{N(\log\log N)^{i-1}}{\log N};$$
\item By the Mertens' second theorem (see for example \cite[Chapter I.1, Theorem~9]{tenebaum_2015}), 
$$\sum_{p\leq N, \text{ $p$ is a prime}} p^{-1}=\log\log N+O(1);$$
\item By \cite[Lemma~2.9]{S23}, the norm of every prime ideal is of the form $p^{r}$ for some prime $p\in\N$ and some $1\leq i\leq [K:\Q]$;
\item By Lemma \ref{mv8}, we have 
$\sum_{\NN(\mf p^{k})=a} 1\ll 1$ for all $a\in\N$;
\item For any $I\subset\N\backslash\{1\}$ and $s\geq 1$, we have $$\sum_{a\in I, \text{ $a$ is a power of prime}}a^{-s}\leq \sum_{a\in I, \text{ $a$ is a prime}}a^{-s}+O(1),$$
since $\sum_{k=2}^{\infty}\sum_{n=2}^{\infty}n^{-sk}\leq \sum_{k=2}^{\infty}\sum_{n=2}^{\infty}n^{-k}=\sum_{n=2}^{\infty}\frac{1}{n(n-1)}=1.$
\end{itemize}
Let us show each of the asserted statements. 
First, for part (i), we have
$$\sum_{\NN(\mf p^{k})\leq N} 1=\sum_{a\in \Pi_{1}(N)}\sum_{\NN(\mf p^{k})=a} 1\ll \sum_{a\in \Pi_{1}(N)} 1=\pi_{1}(N)=O\left(\frac{N}{\log N}\right).$$
Next, part (ii) follows from the fact that
\begin{eqnarray*}
    \sum_{\NN(\mf p^{k})\leq N} \NN(\mf p^{k})^{-1} & = & \sum_{a\in \Pi_{1}(N)}\sum_{\NN(\mf p^{k})=a} a^{-1}
      \\&\ll & \sum_{a\in \Pi_{1}(N)} a^{-1}\ll \sum_{p\leq N, p \text{ is a prime}} p^{-1}=O(\log\log N).
\end{eqnarray*}
%$$\sum_{\NN(\mf p^{k})\leq N} \NN(\mf p^{k})^{-1}=\sum_{a\in \Pi_{1}(N)}\sum_{\NN(\mf p^{k})=a} a^{-1}\ll \sum_{a\in \Pi_{1}(N)} a^{-1}\ll \sum_{p\leq N, p \text{ is a prime}} p^{-1}=O(\log\log N).$$
\iffalse
As for part (iii), it is 
\begin{eqnarray*}
    \sum_{N^{\varepsilon}\leq \NN(\mf p)\leq N}\NN(\mf p)^{-1}
 & = & \sum_{N^{\varepsilon}\leq a\leq N, \text{ $a$ is a power of a prime}} \sum_{\NN(\mf p)=a} a^{-1}
 \\ &\ll & \sum_{N^{\varepsilon}\leq a\leq N, \text{ $a$ is a power of a prime}}a^{-1}
 \ll \sum_{N^{\varepsilon}\leq a\leq N, \text{ $a$ is a prime}}a^{-1}
 \\& = &\log\log N-\log\log N^{\varepsilon}+O(1)=O_{\varepsilon}(1).
\end{eqnarray*}    
We continue with part (iv). Note that
\begin{eqnarray*}
    \sum_{\NN(\mf p)>N^{\varepsilon}}\NN(\mf p)^{-(1+(\log N)^{-1})}
   & = &\sum_{a>N^{\varepsilon}, \text{ $a$ is a power of a prime}}\sum_{\NN(\mf p)=a}a^{-(1+(\log N)^{-1})}
   \\&\ll & \sum_{a>N^{\varepsilon}, \text{ $a$ is a power of a prime}}a^{-(1+(\log N)^{-1})}\\
 &  \ll & \sum_{a>N^{\varepsilon}, \text{ $a$ is a  prime}}a^{-(1+(\log N)^{-1})}+O(1),
\end{eqnarray*}  
where the right hand side is $O_{\varepsilon}(1)$ by following the same argument as in the bottom of \cite[p. 246]{E79book}.

Regarding part (v), we have 
\begin{eqnarray*}
    \frac{1}{\log N}\sum_{\NN(\mf p)<N^{\varepsilon}}\frac{\log\NN(\mf p)}{\NN(\mf p)} & = & \frac{1}{\log N}\sum_{a<N^{\varepsilon}, \text{ $a$ power of a prime}}\sum_{\NN(\mf p)=a}\frac{\log a}{a}
   \\&\ll & \frac{1}{\log N}\sum_{a<N^{\varepsilon}, \text{ $a$ power of a prime}}\frac{\log a}{a}\\
  & = & \frac{1}{\log N}\sum_{a<N^{\varepsilon}, \text{ $a$ prime}}\frac{\log a}{a}+\Theta
   =O(\varepsilon)+\Theta,
\end{eqnarray*}
where the last inequality follows from the same estimate as in \cite[p. 247]{E79book}, and 
\begin{equation}\nonumber
    \begin{split}
     \vert\Theta\vert\leq \frac{1}{\log N}\sum_{k=2}^{\infty}\sum_{n=2}^{\infty}\frac{\log n^{k}}{n^{k}}
      = \frac{1}{\log N}\sum_{n=2}^{\infty}\frac{(2n-1)\log n}{(n-1)^2n}=O\left(\frac{1}{\log N}\right),
    \end{split}
\end{equation}
after swapping the other of the sums, and summing the inner sum using differentiation of a geometric series. As for the last claim, simply note that the resulting series is clearly convergent to some absolute constant.
\fi
%Lastly, for part (vi),  
%\begin{equation}\label{mmvv6}
%    \begin{split}
%      &\qquad \sum_{\NN(\mf p^{k}\mf q^{\ell})\leq N, \mf p\neq \mf q} 1=\sum_{a,b\in \Pi_{1}(N),k,\ell\in\N,a^{k}b^{\ell}\leq N}\sum_{\NN(\mf p)=a, \NN(\mf q)=b}1
%       \ll \sum_{a,b\in \Pi_{1}(N),k,\ell\in\N,a^{k}b^{\ell}\leq N} 1.
%    \end{split}
%\end{equation}
%Fix any $a,b\in \Pi_{1}(N),k,\ell\in\N$ with $a^{k}b^{\ell}\leq N$. If $a,b$ are powers of different primes, then $a^{k}b^{\ell}$ belongs to $\pi_{2}(N)$, and for each $x\in \pi_{2}(N)$, $x$ can be written in the form $a^{k}b^{\ell}$ in exactly two ways \WS{not true. since $a$ is not necessarily a prime. but see the blue part below for a fix}. If $a,b$ are powers of the same prime, then $a^{k}b^{\ell}$ belongs to $\pi_{1}(N)$. In this case, for each $x=p^{r}\in \pi_{1}(N)$ with $p$ being a prime, the number of ways $x$ can be written in the form $a^{k}b^{\ell}$ is equal to the number of tuples $(u,v,k,\ell)\in\N^{4}$ such that $uk+v\ell=r$, which is at most $r^{3}$. So  the right hand side of (\ref{mmvv6}) can be bounded by
 %$$2\pi_{2}(N)+\sum_{p^{r}\in \pi_{1}(N)} r^{3}
  %     \leq O\left(\frac{N\log\log N}{\log N}\right)+\sum_{r=1}^{\lceil\frac{\log N}{\log 2}\rceil}r^{3}\pi_{1}({N}^{1/r}).$$
%The conclusion follows from the estimate
%\begin{equation}\nonumber
 %   \begin{split}
  %    &\qquad \sum_{r=1}^{\lceil\frac{\log N}{\log 2}\rceil}r^{3}\pi_{1}({N}^{1/r})
   %   \ll \sum_{r=1}^{\lceil\frac{\log N}{\log 2}\rceil}\frac{r^{4}{N}^{1/r}}{\log N}
    %  \leq \frac{N}{\log N}+\sum_{r=2}^{\lceil\frac{\log N}{\log 2}\rceil}\frac{r^{4}\sqrt{N}}{\log N}
    %  \\&\leq \frac{N}{\log N}+(\log N)^{4}\sqrt{N}
    %  \leq O\left(\frac{N\log\log N}{\log N}\right). 
    %\end{split}
%\end{equation}
Lastly, for part (iii), let $\Pi_{i,T}(N)$ denote the set of natural numbers no larger than $N$ which are of the form $p^{k}q^{\ell}$ for some primes $p,q$ and some $0\leq k, \ell \leq T$. Then,
\begin{equation}\label{mmvv6}
    \begin{split}
      &\qquad \sum_{\NN(\mf p^{k}\mf q^{\ell})\leq N, \mf p\neq \mf q} 1=\sum_{a,b\in \Pi_{1,D}(N),k,\ell\in\N,a^{k}b^{\ell}\leq N}\sum_{\NN(\mf p)=a, \NN(\mf q)=b}1
       \ll \sum_{a,b\in \Pi_{1,D}(N),k,\ell\in\N,a^{k}b^{\ell}\leq N} 1,
    \end{split}
\end{equation}
where $D:=[K:\Q]$.
Fix any $a,b\in \Pi_{1, D}(N),k,\ell\in\N$ with $a^{k}b^{\ell}\leq N$. If $a,b$ are powers of different primes, then $a^{k}b^{\ell}$ belongs to $\pi_{2}(N)$, and for each $x\in \pi_{2}(N)$, $x$ can be written in the form $a^{k}b^{\ell}$ for some  $a,b\in \Pi_{1,D}(N),k,\ell\in\N$ in at most $2D^{2}$ ways. If $a,b$ are powers of the same prime, then $a^{k}b^{\ell}$ belongs to $\pi_{1}(N)$. In this case, for each $x=p^{r}\in \pi_{1}(N)$ with $p$ being a prime, the number of ways $x$ can be written in the form $a^{k}b^{\ell}$ for some $a,b\in \Pi_{1,D}(N),k,\ell\in\N$ is  at most $D^{2}$. So  the right hand side of (\ref{mmvv6}) can be bounded by
 $$2D^{2}\pi_{2}(N)+D^{2}\sum_{p^{r}\in \pi_{1}(N)} 1
       \leq O\left(\frac{N\log\log N}{\log N}\right)+\sum_{r=1}^{\lceil\frac{\log N}{\log 2}\rceil}\pi_{1}({N}^{1/r}).$$
The estimate
\begin{eqnarray*}
    \sum_{r=1}^{\lceil\frac{\log N}{\log 2}\rceil}\pi_{1}({N}^{1/r})
      & \ll & \sum_{r=1}^{\lceil\frac{\log N}{\log 2}\rceil}\frac{r{N}^{1/r}}{\log N}
      \leq \frac{N}{\log N}+\sum_{r=2}^{\lceil\frac{\log N}{\log 2}\rceil}\frac{r\sqrt{N}}{\log N}
      \\& \leq & \frac{N}{\log N}+\log N\sqrt{N}
      \leq O\left(\frac{N\log\log N}{\log N}\right),
\end{eqnarray*} 
gives the conclusion.
\end{proof}

\iffalse
 %An upper bound for $\sum_{\NN(p)>x^{\e}}\NN(p)^{-(1+(\log x)^{-1})}=O(1)$ (this is the last equation of page 246 of the book, where they used integration by part. However, I think this part is fine for us as the range $\NN(p)>x^{\e}$ is contained in the union of two rectangles and so we can still do the 2-dim integration by part).
 Let $D$ be the degree of $K$.
It is known that prime ideals in $\mathcal{O}_K$ have norm equal to one of $\{p,\dots,p^{D}\}$. Therefore, since $\sigma:=1+(\log x)^{-1}>1$, if we argue as in Lemma~\ref{l2}, then it follows from Lemma~\ref{mv8} that \AK{I recomment to not use ``i'' except for the imaginary unit}
%(1) above (that is, we use the linear estimate bound on the number of elements with a given norm, see Lemma \ref{mv7}, we have the upper bound
\[
\sum_{\NN(\mf p)>x^{\varepsilon}}\NN(\mf p)^{-(1+(\log x)^{-1})} \leq \sum_{p>x^{\varepsilon}}\left( \sum_{j=1}^{D}\frac{1}{p^{j\sigma}}\cdot\sum_{\mf u \in \mf I(K) : \NN(\mf u) \in \{p,p^2.\dots,p^{D}\}} 1\right)\leq  
%[K:\Q]\cdot K_1 
D\sum_{p>x^{\varepsilon/D}}\frac{1}{p^{\sigma}}.
\]
%where $K_1>0$ is a universal constant depending on $K$. 
%for some $C>0$ depending only on $K$. \WS{This does not follow from Lemma \ref{mv7}, but from Lemma \ref{mv8}, and this $C$ can be replaced by $D$}
Finally, using the same argument as in the bottom of \cite[p. 246]{E79book}, %since we are only working over $\N$ now, gives
we have
\[
\sum_{p>x^{\varepsilon/D}}\frac{1}{p^{\sigma}}=O(1),
\]
which implies that
\[
\sum_{\NN(\mf p)>x^{\varepsilon}}\NN(\mf p)^{-(1+(\log x)^{-1})}=O(1),
\]
as desired.
\fi


\section{A classification result for completely multiplicative functions}\label{ss5}

%\WS{This section wa originally the last two subsection of the previous section and the last subsection of the next section. I put them together so that we can sell it as an number theory outcome of the paper (like what I did in the new introduction)}
The purpose of this section is to obtain a useful classification result for multiplicative functions, which will be the cornerstone of our analysis of the set $\mathcal{M}$ which we must split into different pieces according to the asymptotic behavior of the multiplicative function $f$. We begin with extensions of multiplicative functions.

\subsection{Extensions of multiplicative functions}
In this section, we prove Theorem \ref{characterization}.
Before that, we need to classify all the extensions of a completely multiplicative function $f$. 
Let $D=\mathcal{O}_K^{\times}$ or $\mathfrak{I}(K)$.
We say that $f\colon D\to\U$ is \emph{completely multiplicative} (on $D$) if $f(mn)=f(m)f(n)$ for all $m,n\in D$. 
%Defining $f(m/n)$ as $f(m)f(n)^{-1}$ for $m,n\in D$ when $D=\mathcal{O}_K^{\times}$ or $\mathfrak{I}(K)$, we see that every completely multiplicative function on $\mathcal{O}_K^{\times}$ and $\mathfrak{I}(K)$ extends uniquely to a multiplicative function on $K^{\times}$ and $\mathfrak{Q}(K)$ respectively. {\color{red} this is only true if $f$ is $\S$-valued, not $\U$ valued. Think which notation do we want to use}
%
%
For any completely multiplicative function $f\colon \mathcal{O}_K^{\times} \to\U$
with $f(\e)=1$ for all units $\e$,\footnote{It is clear that if $f$ admits an extension, then   $f(\e)=\tilde{f}((\e))=\tilde{f}((1))=f(1)=1$ for all units $\e$, so naturally we only define extensions of multiplicative functions $f$ if they satisfy this additional condition.} we say that a completely multiplicative function $\tilde{f}\colon \mathfrak{Q}(K)\to\U$ is an \emph{extension} of $f$ if  $\tilde{f}((x))=f(x)$ for all $x\in \mathcal{O}_K^{\times}$. One of the key ideas in this paper is the study of averages of completely multiplicative functions on a non-uniquely factorizable domain which is done by passing to extended versions of the original multiplicative functions.

Let $G$ be the ideal class group of $K$. By \cite[Theorem~6.3]{neukirch99}, $G$ is a finite group, and it is also clearly abelian by construction. Therefore, by the classification of finite abelian groups, we can find $k\in\N$, 
    and generators $g_{1},\dots,g_{k}$ of $G$ of orders $d_{1},\dots,d_{k}\in\N$ such that $G=\langle g_1,\dots,g_k\rangle$.
   
    For each $1\leq i\leq k$, let $I_{i}\in \mathfrak{Q}(K)$ be a representative of $g_{i}$. We may assume without loss of generality that $I_{i}\in\mathfrak{I}(K)$. If so, then we say that $(I_{1},\dots,I_{k})$ is an \emph{ideal class group representation} for $\OK$. 
    It follows that $I_{i}^{d_{i}}=(x_{i})$ for some $x_{i}\in \mathcal{O}_{K}^{\times}$.
%Let $f$ be a function, 
%we say that $(I_{1},\dots,I_{k})$ is \emph{compatible with} $f$ if
%  $f(x_{i})\neq 0$ for all $i$. 
It is clear that if $\tilde{f}$ is an extension of $f$, then we must have that $\tilde{f}(I_{i})^{d_{i}}=f(x_{i})$ for all $1\leq i\leq k$. 
Conversely, we have the following.





\begin{proposition}\label{prop: extensions}
    Let $f\colon \mathcal{O}_K^{\times}\to\S$ be a completely multiplicative function  with $f(\e)=1$ for all units $\e$. Suppose that %$f(x_{i})\neq 0$ 
    $(I_{1},\dots,I_{k})$ is an ideal class representation of $\OK$ %which is compatible with $f$
 %   {\color{red} new assumption}. Suppose also that 
 with
 $I_{i}^{d_{i}}=(x_{i})$ for some  $x_{i}\in \mathcal{O}_{K}^{\times}$ for  $1\leq i\leq k$. Then for any $a_{1},\dots,a_{k}\in\C$ with $a_{i}^{d_{i}}=f(x_{i}), 1\leq i\leq k$, there exists a unique extension $\tilde{f}$ of $f$ such that $\tilde{f}(I_{i})=a_{i}$ for all $1\leq i\leq k$. 
\end{proposition}
\begin{proof} %{\color{red} check the new proof}
The uniqueness part is obvious and so we now prove the existence of such an extension. Note that every ideal can be written as $\frac{(x)}{(y)}I$, where $I=I_{1}^{b_{1}}\dots I_{k}^{b_{k}}$, and for some $x,y\in \mathcal{O}_K^{\times}$ and $0\leq b_{i}\leq d_{i}-1$. 
By   Lemma \ref{cl2}, we may assume that $y=\NN(I)$. For such an ideal,
define
    $$\tilde{f}\left(\frac{(x)}{(\NN(I))}I\right):=f(x)f(\NN(I))^{-1}a_{1}^{b_{1}}\dots a_{k}^{b_{k}}$$
for all $x,y\in \mathcal{O}_K^{\times}$ and $0\leq b_{i}\leq d_{i}-1$, which is well defined since $x$ is uniquely determined up to a unit and $f(\e)=1$ for all units $\e$. %, and $f(\NN(I))$ is nonzero since $f(x_{i})\neq 0$ and thus $a_{i}\neq 0$, since $(I_{1},\dots,I_{k})$ is compatible with $f$. {\color{red} here we crucially used to assumption that $f(x_{i})\neq 0$} %This defines a function on $\mathfrak{Q}(K)$, since by the structure of the ideal class group, every element of $\mathfrak{Q}(K)$ can be written in the form $II_{1}^{b_{1}}\dots I_{k}^{b_{k}}$ for some $I$ of the form $\frac{(x)}{(y)}$ for some $x,y\in (\mathcal{O}_K)^{\times}$ and for some $0\leq b_{i}\leq d_{i}-1$ in a unique way (the representation $I=\frac{(x)}{(y)}$ is not unique, but the value $\tilde{f}(I):=f(x/y)$ is unique since $f(\e)=1$ for all units $\e$). 
Since clearly $a_{1},\dots,a_{k}\in\S$, it follows that $\tilde{f}$ takes values in $\S$.

It is clear that $\tilde{f}((x))=f(x)$ for all $x\in \mathcal{O}_K^{\times}$ (in which case we must have that $b_{1}=\dots=b_{k}=0$). 
We now show that $\tilde{f}$ is multiplicative.

Let $x,x'\in \mathcal{O}_K^{\times}$ and $0\leq b_{i},b'_{i}\leq d_{i}-1$ be such that $J:=\frac{(x)}{(\NN(I))}I$ and $J':=\frac{(x')}{(\NN(I'))}I'$ are ideals, where $I=I_{1}^{b_{1}}\dots I_{k}^{b_{k}}$ and $I'=I_{1}^{b'_{1}}\dots I_{k}^{b'_{k}}$.
Put $c_{i}=b_{i}+b'_{i}$ if $c_{i}\leq d_{i}-1$ and $c_{i}=b_{i}+b'_{i}-d_i$ otherwise. Denote $r_{i}=b_{i}+b'_{i}-c_{i}$.
Then 
$$JJ'=\frac{(xx'x_{1}^{r_{1}}\dots x_{k}^{r_{k}})}{(\NN(I_{1}^{b_{1}+b'_{1}}\dots I_{k}^{b_{k}+b'_{k}}))}I_{1}^{c_{1}}\dots I_{k}^{c_{k}}.$$
Since $JJ'$ is an ideal, by  Lemma \ref{cl2}, we have that 
$$JJ'=\frac{(y)}{(\NN(I_{1}^{c_{1}}\dots I_{k}^{c_{k}}))}I_{1}^{c_{1}}\dots I_{k}^{c_{k}}$$
for some $y\in\mathcal{O}_K^{\times}$.
Therefore, we have that 
$$y=xx'x_{1}^{r_{1}}\dots x_{k}^{r_{k}}/\NN(I_{1}^{r_{1}})\dots \NN(I_{k}^{r_{k}}).$$
By definition,
\begin{equation*}
    \begin{split}
        &\qquad \tilde{f}(JJ')
        =f(y)f(\NN(I_{1}^{c_{1}}\dots I_{k}^{c_{k}}))^{-1}a_{1}^{c_{1}}\dots a_{k}^{c_{k}}.
    \end{split}
\end{equation*}
On the other hand,
$$\tilde{f}(J)\tilde{f}(J')=f(x)f(\NN(I_{1}^{b_{1}}\dots I_{k}^{b_{k}}))^{-1}f(x')f(I_{1}^{b'_{1}}\dots I_{k}^{b'_{k}})a_{1}^{b_{1}+b'_{1}}\dots a_{k}^{b_{k}+b'_{k}}.$$
One can easily check that $\tilde{f}(JJ')=\tilde{f}(J)\tilde{f}(J')$, completing the proof.
%
\iffalse

To show that $\tilde{f}$ is multiplicative, it suffices to show the following claim.

 \textbf{Claim.} For all $x\in (\mathcal{O}_K)_{\times}$ and $b_{i}\in\N$, we have that $$\tilde{f}\left(\frac{(x)}{(y)}I_{1}^{b_{1}}\dots I_{k}^{b_k}\right)=\tilde{f}((x))\tilde{f}((y))^{-1}\tilde{f}(I_{1})^{b_{1}}\dots \tilde{f}(I_{k})^{b_{k}}.$$

Indeed, assume that $b_{i}=c_{i}d_{i}+r_{i}$ for some $c_{i}\in\N$ and $0\leq r_{i}\leq d_{i}-1$. Then 
$$\frac{(x)}{(y)}I_{1}^{b_{1}}\dots I_{k}^{b_{k}}=\frac{(xx_{1}^{c_{1}}\dots x_{k}^{c_{k}})}{(yy_{1}^{c_{1}}\dots y_{k}^{c_{k}})}I_{1}^{r_{1}}\dots I_{k}^{r_{k}}.$$
So 
\begin{equation*}
    \begin{split}
        &\qquad\tilde{f}\left(\frac{(x)}{(y)}I_{1}^{b_{1}}\dots I_{k}^{b_{k}}\right)=\tilde{f}\left(\frac{(xx_{1}^{c_{1}}\dots x_{k}^{c_{k}})}{(yy_{1}^{c_{1}}\dots y_{k}^{c_{k}})}I_{1}^{r_{1}}\dots I_{k}^{r_{k}}\right)
         =f\left(\frac{xx_{1}^{c_{1}}\dots x_{k}^{c_{k}}}{yy_{1}^{c_{1}}\dots y_{k}^{c_{k}}}\right)a_{1}^{r_{1}}\dots a_{k}^{r_{k}}
         \\&=f(x/y)a_{1}^{d_{1}c_{1}+r_{1}}\dots a_{k}^{d_{k}c_{k}+r_{k}}
         =\tilde{f}((x))\tilde{f}((y))^{-1}\tilde{f}(I_{1})^{b_{1}}\dots \tilde{f}(I_{k})^{b_{k}}.
    \end{split}
\end{equation*}
We are done.

\fi
\end{proof}

 \iffalse


\begin{lemma}\label{ccf}
%    Give a characterization for good $f$. {\color{red} I think it is not hard to give such a ahcaractrization, since there are a lot of freedoms. However, what we really need to use is: for any $\S$-valued completely multiplicative function $f$ and any Dirichlet character $\chi$, we have that $f\chi$ is good. This is very easy to prove.
 %   }
 For any $\S$-valued completely multiplicative function $f\colon\OK^{\times}\to\S$ and Dirichlet character $\chi$, $f\chi$ is compatible with some ideal class representation. As a consequence of Proposition \ref{prop: extensions}, $f\chi$ has exactly $\vert G\vert$ many extensions of $f$.
\end{lemma}
\begin{proof}
Assume that $\chi$ is of period $I_{0}$.
Let $(I_{1},\dots,I_{k})$ be an arbitrary ideal class representation with $I_{i}^{d_{i}}=(x_{i})$ for some $x_{i}\in \OK^{\times}$. It suffices to show that for each $1\leq i\leq k$, there exists an ideal $J_{i}$ which is in the same ideal class as $I_{i}$ such that $J_{i}^{d_{i}} \mod I_{0}$ is invertible in $(\mathfrak{I}(K)/I_{0})^{\times}$.

It follows from Lemma \ref{ccf} that $J_{i}=\frac{(p)}{(\NN(I_{i}))}I_{i}$ for some $p\in\OK^{\times}$ with $\NN(I_{i})\vert\NN(p)$.
On the other hand, $J_{i}^{d_{i}}=(p^{d_{i}}x_{i}/\NN(x_{i}))=(p^{d_{i}}/\overline{x_{i}})$.
So it suffices to find some $p\in\OK^{\times}$ with $\NN(I_{i})\vert\NN(p)$ such that $p^{d_{i}}/\overline{x_{i}} \mod I_{0} \in (\OK/I_{0})^{\times}$.  {\color{red} TBC.}


{\color{red} For example, if $I_{1}=I_{0}=\langle 2, 1-\sqrt{-5}\rangle$, then $d_{1}=x_{1}=2$, and we need $2\vert\NN(p)$ and $p^{2}/2 \mod I_{0}$ is invertible. We can take $p=1+\sqrt{-5}$, then $\NN(p)=6$ and $p^{2}/2=-2+\sqrt{-5}$ which is invertible mod $I_{0}$ (since $\NN(I_{0})=2$, an element mod $I_{0}$ is invertible as long as it is not in $I_{0}$)}
    
    {\color{red} TBC. I don't know how to do the general case now, but here is an example. Consider $K=\Q(\sqrt{-5})$. It has ideal class number two, with ideal class representation $I_{1}=\langle 2, 1-\sqrt{-5}\rangle$. Then $I_{1}^{2}=(2)$. If $\chi$ is not of perodic 2, then we are done since $f\chi(2)\neq 0$. If $\chi$ is of periodic 2 but (say) not of period 3, then consider $J_{1}=\langle 3,1-\sqrt{-5}\rangle$. It follows that $I_{1}$ is in the same ideal class as $I_{1}$, and $\frac{J_{1}}{I_{1}}=\frac{(1-\sqrt{-5})}{(2)}$, $J_{1}^{2}=(3)$. Since $f\chi(3)\neq 0$, $f$ is compatible with $J_{1}$ and we are done.
    }
\end{proof}

\fi

\subsection{Characterizations for aperiodic completely multiplicative functions}
%Our goal in this subsection is to provide a list of equivalent definitions for aperiodic functions, which is used to prove Theorem \ref{characterization}. 
 
We now begin to prove Theorem \ref{characterization} by providing a list of equivalent definitions of aperiodic multiplicative functions.
We begin with a short lemma that justifies an assumption we will make many times throughout the paper: that $f$ can be assumed to be trivial on units.
\begin{lemma}\label{tv}
    Let $d\in\N$ be squarefree and $f\colon \Z[\tau_{d}]^{\times}\to\C$ be a  completely multiplicative function. If $f(\e)\neq 1$ for some unit $\e$. Then $f$ is aperiodic.
 \end{lemma}
 \begin{proof}
 This follows from the F\o lner property that the balls $B_N:= \{ u \in \OK : \NN(u)\leq N\}$ enjoy. Indeed, notice that each $B_N$ has the property that for any unit $\e \in \Z[\tau_{d}]^{\times}$, it is $\e B_N=B_N$, $N \in \N$. %Therefore, any limit $L$ point of the averages that appear in \eqref{eq: aperiodicdef} is such that $L=f(\e) L$.
Let $L_{N}$ be the average in (\ref{eq: aperiodicdef}) without taking the limit over $N$. It is clear that $L_{N}=f(\e) L_{N}$. Thus, if $f(\e)\neq 1$ for some unit $\e$, then $L_{N}=0$ for all $N$. So $f$ is aperiodic. 
% Thus, if $f(\e)\neq 1$ for some unit $\e$, then $L=0$. By compactness, one such limit point must exist, from which the claim follows. 
 \end{proof}


The following proposition gives a list of equivalent definitions for aperiodic functions.
\begin{proposition}\label{prop: dircharequivalence}
 %   Suppose that $K$ has finitely many units and let 
 Let $d\in\N$ be squarefree and $f\colon \Z[\tau_{d}]^{\times}\to\S$ be a completely multiplicative function with  $f(\e)=1$ for all units $\e$. 
   % Let $b_{1},\dots,b_{D}$ be an integral basis of $\mathcal{O}_{K}$.
    The following are equivalent. 
    \begin{enumerate}
        \item $f$ is aperiodic.
        %, i.e, or all $a_{1},\dots,a_{D}\in\N$ and $b_{1},\dots,b_{D}\in\Z$ we have that $$\lim_{N\to\infty} \E_{m_{1},\dots,m_{D}\in\Z\colon \NN\left(\sum_{i=1}^{D}(a_{i}m_{i}+b_{i})e_{i}\right)\leq N}f\left(\sum_{i=1}^{D}(a_{i}m_{i}+b_{i})e_{i}\right)=0.$$
        \item The following limit vanishes.
        $$\lim_{N\to\infty}\sup_{P\in \mathcal{AP}[\tau_{d}]}\left|\E_{u\in\Z[\tau_{d}]^{\times}\colon \NN(u)\leq N}\mathds{1}_{P}(u)f(u)\right|=0.$$
        \item For any Dirichlet character $\chi\colon \Z[\tau_{d}]^{\times}\to\U$, we have that 
        \begin{equation}\nonumber%\label{dirproduct0}
        \lim_{N\to\infty}\E_{u\in\Z[\tau_{d}]^{\times}\colon \NN(u)\leq N}(f\chi)(u)=0.
        \end{equation}
         \item For any modified Dirichlet character $\chi'\colon \Z[\tau_{d}]^{\times}\to\S$, we have that $$\lim_{N\to\infty}\E_{u\in\Z[\tau_{d}]^{\times}\colon \NN(u)\leq N}(f\chi')(u)=0.$$
        \item For any modified Dirichlet character $\chi'\colon \Z[\tau_{d}]^{\times}\to\S$ %{\color{red} with $\chi(x_{i})\neq 0$ for $1\leq i\leq k$ (I think this should be the correect definition)} 
        and any extension $\widetilde{f\chi'}$ of $f\chi'$,%{\color{red} which exists by the previous proposition}, 
       \footnote{By Proposition \ref{prop: extensions}, $f\chi'$ admits $\vert G\vert$ extensions.}
        we have that 
        $$\lim_{N\to\infty}\E_{\mathfrak{v}\in \mathfrak{I}(\Q(\sqrt{-d}))\colon \NN(\mathfrak{v})\leq N}\widetilde{f\chi'}(\mathfrak{v})=0.$$
    \end{enumerate}
\end{proposition}

We remark that in (v), one can not replace modified Dirichlet characters by Dirichlet characters, since a $\U$-valued multiplicative function may not admit an extension.
\iffalse

{\color{blue}

Let $\chi(n)=1$ if $3\nmid n$ and $\chi(n)=0$ otherwise. Then this is a Dirichlet character. Now define $\chi'(3^{k}n):=\chi'(3)^{k}$ for $3\nmid n$, then this is a completely multiplicative function which coincides with $\chi$ outside $3\N$, and we don't need $\chi'(3)$ to be 1. {\color{red} This construction is not good because $\chi'$ is no longer periodic}


For convenience let's consider the $\Z$ case.
Let $P$ be an AP with period $q$.  Define $\chi_{+}(q^{k}n):=1$ and $\chi_{-}(q^{k}n):=(-1)^{k}$ for $q\nmid n$, then these are  completely multiplicative functions which coincides with Dirichlet characters outside $q\N$. .....


\

\

An example: consider $\E_{u}\bold{1}_{3\N+1}(u)h(u)$. There are two Dirichlet characters of period 3: $\chi_{1}(n)=0,1,1$ if $n\equiv 0,1,2\mod 3$, and  $\chi_{2}(n)=0,1,-1$ if $n\equiv 0,1,2\mod 3$. So the two modified Dirichlet characters are $\chi'_{1}(n)=1,1,1$ if $n\equiv 0,1,2\mod 3$, and  $\chi'_{2}(n)=1,1,-1$ if $n\equiv 0,1,2\mod 3$. 
Here $\bold{1}_{3\N+1}$ can not be written as a linear combination of $\chi'_{1}$ and $\chi'_{2}$. However, since $\chi'_{1}+\chi'_{2}=2(\bold{1}_{3\N}+\bold{1}_{3\N+1})$, we have that  $\E_{u}\bold{1}_{3\N+1}(u)h(u)$ is equal to
$$\E_{u}\frac{1}{2}(\chi'_{1}(u)+\chi'_{2}(u))h(u)-\E_{u}\bold{1}_{3\N}(u)h(u),$$
but $\E_{u}\bold{1}_{3\N}(u)h(u)=\frac{h(3)}{3}\E_{u}h(u)$. So if we know that the average of $h\chi$ is zero for all modified Dirichlet characters, then the average $\E_{u}\bold{1}_{3\N+1}(u)h(u)$ is zero. 


In general, if $f\chi$ has zero average for all Dirichlet characters, then in particular $f$ has zero average (taking $\chi\equiv 1$). So for any modified Dirichlet character $\chi'$, we have that 
$$\E_{n}f(n)\chi'(n)=\E_{n}f(n)\chi(n)+\E_{n}\bold{1}_{q\N}(n)f(n)=\E_{n}f(n)\chi(n)+\frac{f(q)}{q}\E_{n}f(n)=0.$$
So these two conditions should be equivalent.


}

\fi

%{\color{red} I have put together the proof of Propostion \ref{characterization}. But the details needs to be modified.}

 \iffalse
{\color{blue} Here is a potential way to fix the issue. We say that a multiplicative function $f$ is \emph{$I$-invertible} if for all $x\in\OK$ with $x\mod I\notin (\OK/I)^{\times}$, we have that $f(x)\in \S$, and otherwise $f(x)=0$. Let $\mathfrak{I}(K,I)$ denote the set of all ideals which are coprime to $I$ (meaning that $I+I'=\OK$?). We say that a completely multiplicative function $\tilde{f}\colon \mathfrak{I}(K,I)\to\S$ is an \emph{$I$-extention} of $f$ if $\tilde{f}((x))=f(x)$ for all $x\in\OK$ with $x\mod I\notin (\OK/I)^{\times}$. Then I think we can show that 

New Prop 4.3: every $I$-invertible completely multiplicative function admits an $I$-extention, which is also \emph{$I$-invertible}, meaning that $\tilde{f}(I')\in\S$ whenever $I'+I=\OK$.


Then for Proposition 4.7, condition (iv) becomes: for any $\chi$ with period $I$ and any $I$-exnteion $\tilde{f\chi}$ of $f\chi$ (note that $f\chi$ is $I$-invertible), we have....
}

\fi
 

We will split the proof of Proposition~\ref{prop: dircharequivalence} into a series of lemmas.
\begin{lemma}
In Proposition \ref{prop: dircharequivalence}, $(i)$ is equivalent to $(ii)$.
\end{lemma}
\begin{proof}
As was the case in \cite[Appendix A]{S23}, %one direction is trivial, and we will also prove the other one by contradiction.
the direction (ii)$\Rightarrow$(i) is trivial. We show that (i)$\Rightarrow$(ii).

Suppose that we can find an infinite sequence $(N_k)$ of natural numbers, a sequence of %infinite two-dimensional progressions 
$(P_k)$ in $\mathcal{AP}[\tau_{d}]$ and $\varepsilon>0$ such that
\[
\left| \E_{ \NN(u) \leq N_k} \mathds{1}_{P_k}(u) f(u)\right| \geq \varepsilon.
\]

%{\color{blue} I wrote a new proof.

 Assume that $P_{k}=
\{(a_{k,1}m+b_{k,1})+(a_{k,2}n+b_{k,2})\tau_{d}\in \Z[\tau_{d}]\colon m,n\in\Z\}$
        for some $a_{k,1},b_{k,1},a_{k,2},b_{k,2}\in\Z$ with $a_{k,1},a_{k,2}\in\N$. Then $\frac{\vert P_{k}\vert}{\vert \{u\colon \NN(u)\leq N_{k}\}\vert}\geq \varepsilon$. On the other hand, since $m$ and $n$ can take at most $\lceil\frac{10\sqrt{N_{k}}}{a_{k,1}}\rceil$ and $\lceil\frac{10\sqrt{N_{k}}}{a_{k,2}}\rceil$ consecutive integers in $P_{k}$, we have that
        $$\frac{\vert P_{k}\vert}{\vert \{u\colon \NN(u)\leq N_{k}\}\vert}\leq \frac{(\frac{10\sqrt{N}_{k}}{a_{k,1}}+1)(\frac{10\sqrt{N}_{k}}{a_{k,2}}+1)}{C_{d}N_{k}}\leq \frac{(\frac{10}{a_{k,1}}+\frac{1}{\sqrt{N_{k}}})(\frac{10}{a_{k,2}}+\frac{1}{\sqrt{N_{k}}})}{C_{d}}$$
        for some $C_{d}>0$ depending only on $d$. So 
        $(\varepsilon C_{d})^{1/2}\leq \frac{10}{\min\{a_{k,1},a_{k,2}\}}.$
If $k$ is sufficiently large, then $\min\{a_{k,1},a_{k,2}\}\ll_{d, \varepsilon} 1$.
By the pigeonhole principle, we may assume that all the $P_{k}$ are the same by passing to a subsequence if necessary. This contradicts to (i) and we are done.      
 % }      
  %\FM{Nice, I think we can comment out the old proof, then.}
%If we use the description provided in Section~\ref{SEChalasz}, we can write the inequality above (for quadratic imaginary fields) as
%\[
%\left| \E_{\NN(Q_{k}u+b_{k}) \leq N_k} h_j(Q_ku+b_k)\right| \geq \varepsilon
%\]
%{\color{red} do you mean $\NN(Q_{k}u+b_{k})\leq N_{k}$? -- Yes sorry}
%{\color{blue}
%\[
%\left| \E_{\NN(Q_{k}u+b_{k}) \leq N_k, u\in B_{k}} f(Q_ku+b_k)\right| \geq \varepsilon
%\]
%where $B_{k}$ is some 2 dimensional box. This implies that $\NN(Q_{k})\leq \varepsilon^{-1}$ and $\vert B_{k}\vert\gg \varepsilon N_{k}$. By the pigeonhole principle, we may assume that there is a single $Q_0, b_0 \in \mathcal{O}$ such that
%\[
%\left| \E_{\NN(Q_0u+b_0) \leq N_k, u\in B_{k}} f(Q_0u+b_0)\right| \geq \varepsilon,
%\]
%whence
%\[
%\lim_{k\to\infty}\left| \E_{\NN(Q_0u+b_0) \leq N_k} f(Q_0u+b_0)\right| \geq \varepsilon.
%\]
%One way to fix this is that when defining Type 2 aperiodic functions, we require the inner sup to be taken over infinite AP (instead of finite AP). This is because Proposition 2.1 is stil valid under this modifications (there are some places in the proof of Proposition 2.1 where the AP needs to be infinite. I have marked them in red.)
%}
%\[
%\left| \E_{\NN(Q_{0}u+b_{0}) \leq N_k} h_j(Q_0u+b_0)\right| \geq \varepsilon
%\]
%for all $k \in \N$. 
%This contradicts the assumption on $h_j,$ so we are done.
%{\color{red}
%I believe a similar approach would work if you think of APs as having the form
%\[
%P=\{ u \in \mathcal{O} : u=(a_1n_1+b_1)x_1+\dots+(a_kn_k+b_k)x_k\}
%\]
%where $x_i$ form an integral basis for $\mathcal{O}$ and $a_i$ are all non-zero. 
%Namely, we could also count how many of those elements are in the balls defined by norms, and obtain a bound of the form $\frac{1}{a_1\dots a_k+1} \leq \varepsilon$. This would also imply that there is a finite number of APs that "potentially give a contradiction". We can also use the same pigeonhole argument to extract only "one" AP that gives a contradiction, which immediately implies that both definitions are equivalent.
\end{proof}

 


Next we show the equivalence between (i) and (iii).
\begin{lemma}\label{lem: dircharequivalence}
%Let $\mathcal{O}_K$ be a number field, and let $g: \mathcal{O}_K^{\times} \to \C$ be a bounded multiplicative function. Further assume that $\mathcal{O}_K$ is a quadratic imaginary field. Then, $g$ is aperiodic if and only if, for every Dirichlet character $\chi$ on $\mathcal{O}_K$, it is
%\begin{equation}\label{dirichletproduct0}
%\lim_{N \to \infty} \E_{\NN(u) \leq N} \chi(u)g(u)=0.
%\end{equation}
In Proposition \ref{prop: dircharequivalence}, $(i)$ is equivalent to $(iii)$.
\end{lemma}

 

\begin{proof}

%{\color{blue} I rewrote the proof and commented the old one.}

We first show that (iii)$\Rightarrow$(i).
Let $P$ be the set given by (\ref{eq:adcd}). 
It suffices to show that
\[
\lim_{N\to\infty}\E_{\NN(n_1+n_2\tau_{d}) \leq N}\mathds{1}_{a_{1}\Z+b_{1}}(n_1)\mathds{1}_{a_{2}\Z+b_{2}}(n_2)f(n_1+n_2\tau_{d})=0.
\]
By taking common denominators, one can write $\mathds{1}_{a_{1}\Z+b_{1}}(n_1)\mathds{1}_{a_{2}\Z+b_{2}}(n_2)$ as a linear combination of characteristic functions of the form $\mathds{1}_{\alpha\Z+b}(n_1)\mathds{1}_{\alpha\Z+b'}(n_2)$ for some $\alpha\in\Z[\tau_{d}]^{\times}$ and $b, b'\in\Z$,\footnote{Here we can take $\alpha$ to be in $\Z^{\times}$; later we need to convert $\alpha$ to a number in $\Z[\tau_{d}]^{\times}.$} which means it is enough to show that
\[
\lim_{N\to\infty}\E_{\NN(n_1+n_2\tau_{d}) \leq N}\mathds{1}_{\alpha\Z+b}(n_1)\mathds{1}_{\alpha\Z+b'}(n_2)f(n_1+n_2\tau_{d})=0.
\]
Again, up to multiples, this is the same as showing that
\begin{equation}\label{progression0eq}
\lim_{N\to\infty}\E_{\NN(\alpha u+\rho) \leq N} f(\alpha u+\rho)=0,
\end{equation}
 where $\rho=b+b'\tau_{d}$.
 If $\rho=0$, then (\ref{progression0eq}) holds since $f(\alpha u)=f(\alpha)f(u)$ and since (iii) holds for $\chi\equiv 1$. Now we assume that $\rho\neq 0$.
 Because $d\in\N$, one can factor $\rho$ into irreducibles (although not necessarily in a unique way), so factoring out the common factors, it suffices to show that for every principal ideal $\langle \alpha\rangle \subseteq \Z[\tau_{d}]$, and every $\rho \in (\Z[\tau_{d}]/\langle \alpha \rangle)^{\times}$, the average in \eqref{progression0eq} is 0. We now simply notice that, $\mathds{1}_{\langle \alpha \rangle +\rho}(u)$ can be written as a finite linear combination of Dirichlet characters of period $\langle \alpha\rangle$, since $\rho$ is now necessarily invertible in $\Z[\tau_{d}]/\langle\alpha\rangle$ by construction. Thus we have that (iii)$\Rightarrow$(i).

Now we show that (i)$\Rightarrow$(iii).
 Let $\chi$ be a non-trivial Dirichlet character with period $I$. We observe that   there exists some non-unit $n \in \mathcal{O}_K^{\times}$   such that $nx \in I$ for every $x \in \OK$. Thus, we may write $\lim_{N\to\infty}\E_{u\in\Z[\tau_{d}]^{\times}\colon \NN(u)\leq N}(f\chi)(u)$ as a linear combination of averages of the form
\[
\lim_{N \to \infty} \E_{\NN(nu+\rho) \leq N} \chi(nu+\rho)f(nu+\rho)=\lim_{N \to \infty} \E_{\NN(nu+\rho) \leq N} \chi(\rho)f(nu+\rho).
\]
So as before, by using the basis $\{1,\tau_{d}\}$ we may easily convert this  into a linear combination of expressions of the form \eqref{eq: aperiodicdef}, completing the proof.
%
\iffalse



We wish to show that if all the averages of $\chi\cdot f$ are null, then $f$ is aperiodic. %Suppose that $\{x_1,x_2\}$ forms an integral basis for $\mathcal{O}_K$.
Let $a,b,c,d \in \Z$ with $a,b \neq 0$. Assume, without loss of generality that $a, b>0$. Then, we wish to show that
\begin{equation}\label{aperiodicity}
\lim_{N\to\infty}\E_{\NN((an+b)+(cm+d)\tau_{d}) \leq N}f((an+b)+(cm+d)\tau_{d})=0.
\end{equation}
This is equivalent to showing that (up to multiples)
\[
\lim_{N\to\infty}\E_{\NN(n_1+n_2\tau_{d}) \leq N}\mathds{1}_{a\N+b}(n_1)\mathds{1}_{c\N+d}(n_2)f(n_1+n_2\tau_{d})=0.
\]
By taking common denominators, one can write $\mathds{1}_{a\N+b}(n_1)\mathds{1}_{c\N+d}(n_2)$ as a linear combination of characteristic functions of the form $\mathds{1}_{\alpha\N+\beta}(n_1)\mathds{1}_{\alpha\N+\gamma}(n_2)$, which means it is enough to show that
\[
\lim_{N\to\infty}\E_{\NN(n_1+n_2\tau_{d}) \leq N}\mathds{1}_{\alpha\N+\beta}(n_1)\mathds{1}_{\alpha\N+\gamma}(n_2)f(n_1+n_2\tau_{d})=0.
\]
Again, up to multiples, this is the same as showing that
\begin{equation}\label{progression0eq}
\lim_{N\to\infty}\E_{\NN(u) \leq CN} f(\alpha u+\rho)=0,
\end{equation}
for some constant $C>0$, where $\rho \in \Z[\tau_{d}]^{\times}$. Because $d\in\N$, one can factor $\rho$ into irreducibles (although not necessarily in a unique way), so factoring out the common factors out of $g$, it is enough to be able to show that for every principal ideal $\langle \alpha\rangle \subseteq \Z[\tau_{d}]$, and every $\rho \in (\Z[\tau_{d}]/\langle \alpha \rangle)^{\times}$, the averages in \eqref{progression0eq} are 0. We now simply notice that, $\mathds{1}_{\langle \alpha \rangle +\rho}(u)$ can be written as a finite linear combination of Dirichlet characters, since $\rho$ is now necessarily invertible in $\Z[\tau_{d}]/\langle\alpha\rangle$ by construction. Thus, our hypothesis \eqref{dirproduct0} completes this direction.

On the other hand, if we assume that $f$ is aperiodic, let $\chi$ be a non-trivial Dirichlet character. In order to show that \eqref{dirproduct0} holds, we observe that if $I$ is the ideal such that $\chi(\cdot +I)=\chi(\cdot)$, there exists some $n \in \OK$ with $n\neq 0$, and $n$ not equal to a unit, such that for every $x \in \OK$ it is $nx \in I$. Thus, the average we want to show is zero is a linear combination of averages of the form
\[
\lim_{N \to \infty} \E_{\NN(u) \leq N} \chi(nu+\rho)f(nu+\rho)=\lim_{N \to \infty} \E_{\NN(u) \leq N} \chi(\rho)f(nu+\rho),
\]
so as before, by using an integral basis of $\mathcal{O}_K$ we easily convert this statement into a linear combination of expressions appearing in \eqref{aperiodicity}, completing the proof.
\fi
\end{proof}
\iffalse
\begin{lemma}\label{lc}
  (1) Let $\Lambda\subseteq \Z[i]$ be a two dimensional lattice and $u_{0}\in \Z[i]^{\times}$. There exist $K=O_{\Lambda}(1)$ such that 
    $$u_{0}+\Lambda=\sqcup_{j=1}^{K}v_{j}(u_{j}+I_{j})$$
    for some finite index ideal $I_{j}$, $v_{j}\in\Z[i]^{\times}$ and $u_{j}\in (\Z[i]/I)^{\times}$.
    
   (2) Let $I$ be a finite index ideal. Then for all $u\in (\Z[i]/I)^{\times}$, $\mathds{1}_{u+I}$ can be written as a linear combination of Dirichlet characters over $\Z[i]$ with period $I$.
\end{lemma}


\begin{proof}

%Need to ask Wenbo about this part: if I understand correctly, we are interested in multiplicative functions that are invariant under the action of a certain $\Z[i]$ lattice generated by $v_1\Z+v_2\Z$ for some $v_1, v_2 \in \Z[i]$ that are linearly independent over $\Q$. If this is the case, we can use the following trick to reduce to the case that was dealt with before.

%First note that if the function is invariant under a lattice $\Lambda,$ if we find a convenient sublattice $\Lambda'$, the function will still be invariant, and if we are only interested in being able to write it as linear combinations of Dirichlet characters, this will be enough (the idea here is, if a function is $4$-periodic, it is also $8$-periodic, but maybe the $8$-periodicity is useful to us in order to write it as a product of Dirichlet characters; I will try to argue that this is the case).

%Note: passing to a sublattice is not enough. We need to pass to an ideal included in the lattice $\Lambda$, not just a sublattice, since for now we only know how to defined Dirichlet characters for quotients of ideals. For example, $3\Z+4i\Z$ is a lattice but not an ideal, and the smallest ideal containing $3\Z+4i\Z$ is $\Z[i]$.



%First note that if $\Lambda = v_1\Z+v_2\Z$, then, we can find $q_1,q_2,q_3,q_4 \in \Q$ such that
%\[
%\begin{cases}
%    q_1v_1+q_2v_2=(1,0) \\
%    q_3v_1+q_4v_4=(0,1),
%\end{cases}
%\]
%where we use the standard identification of $\Z[i]$ with $\Z^2$. Now, clearing denominators, it is obvious that $\Lambda \supset \Lambda'=Me_1\Z+Me_2\Z$ for some big enough $M>0$ that will depend on the arithmetic properties of the vectors $v_1$ and $v_2$. We can further assume, without loss of generality (by taking squares, if needed), that $M$ is a perfect square. Therefore,


Part (1). First we claim that every two dimensional lattice $\Lambda\subseteq \Z[i]$ contains an ideal $I$. To see this, observe that $\Z[i]/\Lambda$ is a finite abelian group, and as such, there is some $n \in \N$ such that $nx=0$ for all $x \in \Z[i]/\Lambda$. It follows that the ideal generated by $n$, $I=\langle n \rangle \subseteq \Lambda$. Since $I$ is also a sublattice of $\Lambda$, we may write
$$u_{0}+\Lambda=\sqcup_{j=1}^{K}(w_{j}+I)$$
for some $K=O_{\Lambda}(1)$ and $w_{1},\dots,w_{K}\in \Z[i]$.
Since $u_{0}\neq 0$, we may also require that $w_{1},\dots,w_{K}\neq 0$.


Since $\Z[i]$ is a PID, there exists $a \in \Z[i]^{\times}$ such that $I=\langle a\rangle$ (note that this is the ideal generated by $a$, not the lattice generated by $a$) for some $a\in\Z[i]^{\times}$. Then for all $1\leq j\leq K$, we may write $w_{j}=u_{j}v_{j}$ and $a=u'_{j}v_{j}$ for some $u_{j},u'_{j},v_{j}\in \Z[i]^{\times}$ with $u_{j}$ being coprime to $u'_{j}$ (and therefore $u_{j}$ belongs to $(\Z[i]/\langle u'_{j}\rangle)^{\times}$). Then $w_{j}+I=v_{j}(u_{j}+\langle u'_{j}\rangle)$ and we are done.

{\color{red} for PID this proof works. For non-PID this method does not work.}

\


Part (2) of the conjecture is true as a consequence of Fourier analysis. Indeed, observe that the function $\mathds{1}_{u+I}(x)$ can be seen as a a composition of functions $\mathds{1}_{u} : (\Z[i]/I)^{\times} \to \C$ with the natural projection $\pi: \Z[i] \to \Z[i]/I$. Consider the dual group of $(\Z[i]/I)^{\times}$ given by $G:=\{ f : (\Z[i]/I)^{\times} \to \S^1 : f([a][b])=f([a])f([b]), \text{ for all } [a], [b] \in (\Z[i]/I)^{\times}\}$. Then, the function $\mathds{1}_{u}(x)$ admits a representation
\[
\mathds{1}_u(x)=\sum_{\chi \in G} a_u \chi,
\]
where the coefficients $a_u$ are the Fourier coefficients of the function $\mathds{1}_u$ for each of the characters in $G$. These characters can be viewed as multiplicative functions on $\Z[i]/I$, so they can extend to Dirichlet characters on $\Z[i]$ by using the natural projection $\pi$ considered above. Therefore, (2) holds, as claimed.
\end{proof}

By Lemma \ref{lc}, for any multiplicative function $\chi$, we have that
\begin{equation*}
    \begin{split}
        &\qquad\mathbb{E}_{u}\mathds{1}_{u_{0}+\Lambda}(u)\chi(u)
        =\sum_{j=1}^{K}\mathbb{E}_{u}\mathds{1}_{v_{j}(u_{j}+I_{j})}(u)\chi(u)
        =\sum_{j=1}^{K}\frac{\chi(v_{j})}{\NN(v_{j})}\mathbb{E}_{u}\mathds{1}_{u_{j}+I_{j}}(u)\chi(u),
    \end{split}
\end{equation*}
which is a linear combination of averages of functions of the form $\chi'\chi$ for some $\chi'\in\mathcal{B}$. We may now apply the previous result to the multiplicative function $\chi'\chi$ to study the average of $\chi$ along any 2-dimensional AP.
\fi









\begin{lemma}\label{dircharequivalence}
%Let $\mathcal{O}_K$ be a number field, and let $g: \mathcal{O}_K^{\times} \to \C$ be a bounded multiplicative function. Further assume that $\mathcal{O}_K$ is a quadratic imaginary field. Then, $g$ is aperiodic if and only if, for every Dirichlet character $\chi$ on $\mathcal{O}_K$, it is
%\begin{equation}\label{dirichletproduct0}
%\lim_{N \to \infty} \E_{\NN(u) \leq N} \chi(u)g(u)=0.
%\end{equation}
In Proposition \ref{prop: dircharequivalence}, $(i)$ is equivalent to $(iv)$.
\end{lemma}

\begin{proof}
We first show that (i)$\Rightarrow$(iv).
Let $\chi'$ be a modified Dirichlet character of period $I$. Then we may write $\chi'(u)=\chi(u)+\mathds{1}_{B}(u)$ for some  Dirichlet character $\chi$ of period $I$, where $B$ is the set of $u$ with $u \mod I\notin (\OK/I)^{\times}$. Note that the set of $u\in\OK$ for which $I\subseteq (u)$ is the set of $u$ for which $(\NN(I),\NN(u))\neq 1$, which is clearly the union of finitely many elements in $\mathcal{AP}[\tau_{d}]$.
So we may write $\chi'=\chi+\sum_{i=1}^{k}\mathds{1}_{P_{i}}$ for some $k\in\N\cup\{0\}$ and $P_{1},\dots,P_{k}\in \mathcal{AP}[\tau_{d}]$. Since Condition (i) implies Conditions (ii) and (iii), we have that 
\begin{eqnarray*}
\lim_{N\to\infty}\E_{u\in\mathcal{O}_{K}\colon \NN(u)\leq N}(f\chi')(u) & = & \lim_{N\to\infty}\E_{u\in\mathcal{O}_{K}\colon \NN(u)\leq N}(f\chi)(u)\\
& + & \sum_{i=1}^{k}\lim_{N\to\infty}\E_{u\in\mathcal{O}_{K}\colon \NN(u)\leq N}\mathds{1}_{P_{i}}(u)f(u)=0.
\end{eqnarray*}
So (iv) holds.

We now show that (iv)$\Rightarrow$(i). 
Similar to the argument in Lemma~\ref{lem: dircharequivalence},
it suffices to show that for every principal ideal $\langle \alpha\rangle \subseteq \Z[\tau_{d}]$, and every $\rho \in (\Z[\tau_{d}]/\langle \alpha \rangle)^{\times}$, the average in \eqref{progression0eq} is 0. Clearly \eqref{progression0eq} holds  when $\alpha=1$. Now assume that \eqref{progression0eq} holds  when $\NN(\alpha)\leq k$ for some $k\in\N$. Take $\alpha$ with $\NN(\alpha)\geq k+1$ and the additional property (that can be assumed without loss of generality) that if $\alpha=t\alpha'$ for some non-unit $t$, then $\NN(\alpha')\leq k$. Similar to the argument in Lemma~\ref{lem: dircharequivalence}, to show that \eqref{progression0eq} holds for this $\alpha$, it suffices to show that 
 \begin{equation}\nonumber%\label{dirproduct0}
        \lim_{N\to\infty}\E_{u\in\Z[\tau_{d}]^{\times}\colon \NN(u)\leq N}(f\chi)(u)=0
        \end{equation}
for all Dirichlet characters $\chi$ of period $\langle\alpha\rangle$. We may rewrite  $\chi=\chi'-\mathds{1}_{B}$ for some modified Dirichlet character $\chi'$ of period $\langle \alpha\rangle$,  where $B$ is the set of $u$ with $u \mod \langle \alpha\rangle\notin (\OK/\langle \alpha\rangle)^{\times}$. It is clear that the set $B$ can be expressed as the disjoint union $B=\sqcup_{i=1}^{\ell}(u_{i}+\langle \alpha\rangle)$ for some $\ell\in\N$ and $u_{i}\in\OK$. Since (iv) holds, it suffices to show that 
 \begin{equation}\nonumber%\label{dirproduct0}
\lim_{N\to\infty}\E_{u\in\Z[\tau_{d}]^{\times}\colon \NN(u)\leq N}\mathds{1}_{u_{i}+\langle \alpha\rangle}(u)f(u)=0
        \end{equation}
for all $1\leq i\leq \ell$. Since $u_{i}\notin (\OK/\langle \alpha\rangle)^{\times}$, there exists a non-unit $t\in\OK$ such that $u_{i}=tu'_{i}$ and $\alpha=t\alpha'$ for some $u'_{i},\alpha'\in\OK^{\times}$. So it suffices to show that
\begin{equation}\nonumber%\label{dirproduct0}
\lim_{N\to\infty}\E_{u'\in\Z[\tau_{d}]^{\times}\colon \NN(u')\leq N}\mathds{1}_{u'_{i}+\langle \alpha'\rangle}(u')f(tu')=0.
        \end{equation} 
 Since $\NN(\alpha')<\NN(\alpha)$, the conclusion follows from the induction hypothesis and the fact that $f(tu')=f(t)f(u')$.
%\FM{I think the proof works, I just had to modify the induction a little bit, because elements with norm exactly $k+1$ may not exist, but we can assume that if we "cross out" some divisor of $\alpha$, then we fall into an $\alpha'$ that has norm at most $k$ and the inductive hypothesis can be applied. I commented the old proof.}
%We now show that (iv)$\Rightarrow$(i). 
%Let $\chi$ be a   Dirichlet character of period $I$. Then we may write $\chi'(u)=\chi(u)+\bold{1}_{I\subseteq (u)}$ for some modified  Dirichlet character $\chi'$ of period $I$. 
%It follows from Lemma~\ref{lem: dircharequivalence} above that we may assume without loss of generality that the Dirichlet characters under consideration have period a principal ideal $I$. Under this further restriction, and assuming that $I=(\alpha)$, for some non-trivial $\alpha \in \OK$ with $\NN(\alpha)=n$, we can write, for each $N\in \N$
%\[
%\E_{\NN(u) \leq N} \mathds{1}_I(u) f(u) = \frac{1}{n}\E_{\NN(u) \leq nN} \sum_{\beta} \mathds{1}_{I}(\alpha u+\beta) f(\alpha u+\beta),
%\]
%{\color{red} there might be an issue with the proof: I think the set $I\subseteq\langle\alpha\rangle$ needs to be replaced by the set $B$ defined in the previous page, since $I$ may not be prime} 
%\FM{Is the issue that we cannot write $\chi'=\chi+\mathds{1}_{I \subset U}$? And in this case, what is the correct equality and why?} {\color{red} I wrote a new proof for this part}
%where $\beta$ ranges over all possible residues, since $(\alpha)$ is an ideal of finite index, given that $\OK$ is a Dedekind domain. But by assumption, $\alpha u \in I$. Now, all terms with $\beta \neq 0 \pmod I$, must have $\mathds{1}_I(\alpha u+\beta)=0$, for otherwise $\beta \in I$, a contradiction with our choice of representatives $\beta$. Thus, the averages reduce to a multiple of
%\[
%\E_{\NN(u)\leq nN} f(u),
%\]
%which are known to converge to $0$ as $N \to\infty$ given that the constant function $1$ is itself a modified Dirichlet character. Combining this fact together with our assumption that (iv) holds, we simply rewrite the Dirichlet character $\chi=\chi'-\mathds{1}_{I \subseteq (u)}$ to finish the proof.
%{\color{red} We need to think how to do it.  
%An example: consider $\E_{u}\bold{1}_{3\N+1}(u)h(u)$. There are two Dirichlet characters of period 3: $\chi_{1}(n)=0,1,1$ if $n\equiv 0,1,2\mod 3$, and  $\chi_{2}(n)=0,1,-1$ if $n\equiv 0,1,2\mod 3$. So the two modified Dirichlet characters are $\chi'_{1}(n)=1,1,1$ if $n\equiv 0,1,2\mod 3$, and  $\chi'_{2}(n)=1,1,-1$ if $n\equiv 0,1,2\mod 3$. 
%Here $\bold{1}_{3\N+1}$ can not be written as a linear combination of $\chi'_{1}$ and $\chi'_{2}$. However, since $\chi'_{1}+\chi'_{2}=2(\bold{1}_{3\N}+\bold{1}_{3\N+1})$, we have that  $\E_{u}\bold{1}_{3\N+1}(u)h(u)$ is equal to
%$$\E_{u}\frac{1}{2}(\chi'_{1}(u)+\chi'_{2}(u))h(u)-\E_{u}\bold{1}_{3\N}(u)h(u),$$
%but $\E_{u}\bold{1}_{3\N}(u)h(u)=\frac{h(3)}{3}\E_{u}h(u)$. So if we know that the average of $h\chi$ is zero for all modified Dirichlet characters, then the average $\E_{u}\bold{1}_{3\N+1}(u)h(u)$ is zero. 
%}
\end{proof}

 We are now in position to complete the proof of Proposition \ref{characterization}. 
 As usual, we let $K=\Q(\sqrt{d})$. Assume that $f\colon\mathcal{O}_K^{\times}\to\S$ is a completely multiplicative function and that $(I_{1},\dots,I_{k})$ is an ideal class representation. Let $\tilde{f}$ be an extension of $f$ (which exists by Proposition~\ref{prop: extensions}).
 Assume that $I_{i}^{d_{i}}=(x_{i})$ for some $x_{i}\in\mathcal{O}_K^{\times}$ and set $A=[d_{1}]\times \dots\times [d_{k}]$. For $\vec{i}=(i_{1},\dots,i_{k})\in A$, let $I_{\vec{i}}:=I_{1}^{i_{1}}\dots I_{k}^{i_{k}}$ and let $\mathfrak{I}_{\vec{i}}(K)$ denote the set of integer ideals in the same ideal class as $I_{\vec{i}}$. 
    
    
   Since $f(x_{i})\neq 0$, we have that $f(\NN(I_{i}))\neq 0$.
%{\color{blue}
Denote $$w_{\vec{i},N}:=\frac{\vert\{\mathfrak{v}\in \mathfrak{I}_{\vec{i}}(K)\colon \NN(\mathfrak{v})\leq N\}\vert}{\vert\{\mathfrak{v}\in \mathfrak{I}(K)\colon \NN(\mathfrak{v})\leq N\}\vert}.$$
%For convenient write $x\sim y$ if $x/y$ is a unit for $x,y\in\OK^{\times}$.
%Let $W\subseteq \OK^{\times}$ be a set where each equivalent class of $\sim$ has exactly one element in $W$.
 Then by  Lemma \ref{cl2},
    \begin{equation}\label{fweopf0}
    \begin{split}
        &\qquad \E_{\mathfrak{v}\in \mathfrak{I}(K), \NN(\mathfrak{v})\leq N}\tilde{f}(\mathfrak{v})
        \\&=\sum_{\vec{i}\in A} \E_{\mathfrak{v}\in \mathfrak{I}(K), \NN(\mathfrak{v})\leq N}\mathds{1}_{\mathfrak{v}\in \mathfrak{I}_{\vec{i}}(K)}\tilde{f}(\mathfrak{v})
        \\&=\sum_{\vec{i}\in A}w_{\vec{i},N} \E_{\frac{(u)}{(\NN(I_{\vec{i}}))}\in \mathfrak{Q}(K)\colon \frac{(u)}{(\NN(I_{\vec{i}}))}I_{\vec{i}}\in \mathfrak{I}(K), \NN(\frac{(u)}{(\NN(I_{\vec{i}}))}I_{\vec{i}})\leq N} \tilde{f}\left(\frac{(u)}{(\NN(I_{\vec{i}}))}I_{\vec{i}}\right)
        %\\&=\sum_{i_{1}\in[d_{1}],\dots,i_{k}\in[d_{k}]}\lim_{N\to\infty}       
        %\E_{\frac{(u)}{(q)}\in \mathfrak{Q}(K), \NN(u/q)\in \Z\cap [0,N/\NN(I_{1})^{i_{1}}\dots \NN(I_{k})^{i_{k}}]}\tilde{f}(\frac{(u)}{(q)}I_{1}^{i_{1}}\dots I_{k}^{i_{k}})
        \\&=\sum_{\vec{i}\in A}w_{\vec{i},N}\tilde{f}(I_{\vec{i}})\cdot  \frac{1}{T}\E_{u\in \mathcal{O}_K^{\times}\colon \frac{(u)}{(\NN(I_{\vec{i}}))}I_{\vec{i}}\in \mathfrak{I}(K), \NN(\frac{(u)}{(\NN(I_{\vec{i}}))}I_{\vec{i}})\leq N} f(u)f(\NN(I_{\vec{i}}))^{-1}
        \\&=\sum_{\vec{i}\in A}w_{\vec{i},N}C_{\vec{i},N}(f)\tilde{f}(I_{1})^{i_{1}}\dots \tilde{f}(I_{k})^{i_{k}}, 
    \end{split}
\end{equation}
where  $T$ is the number of units of $\Q(\sqrt{-d})$ which is finite by Lemma \ref{l32}.
$$C_{\vec{i},N}(f):=\frac{1}{T}\E_{u\in \mathcal{O}_K^{\times}\colon \frac{(u)}{(\NN(I_{\vec{i}}))}I_{\vec{i}}\in \mathfrak{I}(K), \NN(\frac{(u)}{(\NN(I_{\vec{i}}))}I_{\vec{i}})\leq N} f(u)f(\NN(I_{\vec{i}}))^{-1},$$
which is independent of the choices of the extension $\tilde{f}$.
%\FM{it looks like $p$ might be prime but I don't think it is, maybe we can use the letter $u$ or something more neutral}

Suppose first that (i) holds.  Now let $\chi'$ be a modified Dirichlet character. %By Lemma \ref{ccf}, $f\chi$   is complatible with some ideal class representation $(I_{1},\dots,I_{d})$. 
We first claim that the average of $f\chi'$ along every arithmetic progression in $\mathcal{AP}[\tau_{d}]$ is 0. Since the indicator function of every set in $\mathcal{AP}[\tau_{d}]$ is a linear combination Dirichlet characters, it suffices to show that the average of $f\chi'\chi$ is 0 for all Dirichlet characters $\chi$. However, since $\chi'\chi$ is a Dirichlet character, the claim follows from the fact that (i)$\Rightarrow$(iii).



Since the set of $u$ for which  $\frac{(u)}{(\NN(I_{\vec{i}}))}I_{\vec{i}}\in \mathfrak{I}(K)$ is the disjoint union of finitely many elements in $\mathcal{AP}[\tau_{d}]$, it follows from (iv) (which is equivalent to (i)) that   $\lim_{N\to\infty}C_{\vec{i},N}(f\chi')=0$ for all $\vec{i}$. 

So for any extension $\widetilde{f\chi'}$ of $f\chi'$, it follows from \eqref{fweopf0} that 
$$\lim_{N\to\infty} \E_{\mathfrak{v}\in \mathfrak{I}(K), \NN(\mathfrak{v})\leq N}\widetilde{f\chi'}(\mathfrak{v})=0.$$
This implies (v).

Conversely, assume that (v) holds. We show that this implies (iv). Let $\chi'$ be a modified Dirichlet character. %By Lemma \ref{ccf}, $f\chi$  is complatible with some ideal class representation $(I_{1},\dots,I_{d})$.
Let $\mathcal{F}$ denote the set of all extensions of $f\chi'$.
Then it follows from (v) that $$\lim_{N\to\infty} \sum_{g\in\mathcal{F}}\E_{\mathfrak{v}\in \mathfrak{I}(K), \NN(\mathfrak{v})\leq N}g(\mathfrak{v})=0.$$
On the other hand, by \eqref{fweopf0}, we have that 
   \begin{eqnarray*}
 \sum_{g\in\mathcal{F}}\E_{\mathfrak{v}\in \mathfrak{I}(K), \NN(\mathfrak{v})\leq N}g(\mathfrak{v}) 
    & = & \sum_{g\in\mathcal{F}}\sum_{\vec{i}\in A}w_{\vec{i},N}C_{\vec{i},N}(f\chi')g(I_{1})^{i_{1}}\dots g(I_{k})^{i_{k}}
        \\& = &\sum_{a_{j}^{d_{j}}=(f\chi')(x_{j}), 1\leq j\leq k}\sum_{\vec{i}\in A}w_{\vec{i},N}C_{\vec{i},N}(f\chi')a_{1}^{i_{1}}\dots a_{k}^{i_{k}}
        \\& = &\sum_{\vec{i}\in A}w_{\vec{i},N}C_{\vec{i},N}(f\chi')\cdot\sum_{a_{j}^{d_{j}}=(f\chi')(x_{j}), 1\leq j\leq k}a_{1}^{i_{1}}\dots a_{k}^{i_{k}}
        \\& = & w_{\vec{0},N}C_{\vec{0},N}(f\chi').
 \end{eqnarray*}
So,
$$\lim_{N\to\infty}w_{\vec{0},N}C_{\vec{0},N}(f\chi')
        =\lim_{N\to\infty} \sum_{g\in\mathcal{F}}\E_{\mathfrak{v}\in \mathfrak{I}(K), \NN(\mathfrak{v})\leq N}g(\mathfrak{v})=0.$$
        On the other hand, it follows from \cite[Theorem~11.1.5]{murtyesmondebook} that
        $\lim_{N\to\infty}w_{\vec{0},N}=\frac{1}{|A|}$. So
        $\lim_{N\to\infty}C_{\vec{0},N}(f\chi')=0$ by taking their quotients. So (iv) holds for such $\chi'$ and 
we are done. 

%So
%   \begin{equation}\nonumber
%    \begin{split}
%        &\qquad \lim_{N\to\infty}C_{\vec{0},N}(f\chi')= \lim_{N\to\infty}w_{\vec{0},N}C_{\vec{0},N}(f\chi')
%        =\lim_{N\to\infty} \sum_{g\in\mathcal{F}}\E_{\mathfrak{v}\in \mathfrak{I}(K), \NN(\mathfrak{v})\leq N}g(\mathfrak{v})=0,
%    \end{split}
%\end{equation}
%where the first equality follows from \cite[Theorem~11.1.5]{murtyesmondebook}. 
%\FM{The last line is a bit hard to follow for me}So (iv) holds for such $\chi'$ and 
%we are done. 
%}



   
  %  Then by  Lemma \ref{cl2},
  %  \begin{equation}\label{fweopf}
  %  \begin{split}
  %      &\qquad\lim_{N\to\infty} \E_{\mathfrak{v}\in \mathfrak{I}(K), \NN(\mathfrak{v})\leq N}\tilde{f}(\mathfrak{v})
  %      \\&=\sum_{\vec{i}\in A}\lim_{N\to\infty} \E_{\mathfrak{v}\in \mathfrak{I}(K), \NN(\mathfrak{v})\leq N}\mathds{1}_{\mathfrak{v}\in \mathfrak{I}_{\vec{i}}(K)}\tilde{f}(\mathfrak{v})
  %      \\&=\sum_{\vec{i}\in A}\lim_{N\to\infty}\frac{\vert\{\mathfrak{v}\in \mathfrak{I}_{\vec{i}}(K)\colon \NN(\mathfrak{v})\leq N\}\vert}{\vert\{\mathfrak{v}\in \mathfrak{I}(K)\colon \NN(\mathfrak{v})\leq N\}\vert} \E_{\frac{(p)}{(\NN(I_{\vec{i}}))}\in \mathfrak{Q}(K)\colon \frac{(p)}{(\NN(I_{\vec{i}}))}I_{\vec{i}}\in \mathfrak{I}(K), \NN(\frac{(p)}{(\NN(I_{\vec{i}}))}I_{\vec{i}})\leq N} \tilde{f}\left(\frac{(p)}{(\NN(I_{\vec{i}}))}I_{\vec{i}}\right)
        %\\&=\sum_{i_{1}\in[d_{1}],\dots,i_{k}\in[d_{k}]}\lim_{N\to\infty}       
        %\E_{\frac{(p)}{(q)}\in \mathfrak{Q}(K), \NN(p/q)\in \Z\cap [0,N/\NN(I_{1})^{i_{1}}\dots \NN(I_{k})^{i_{k}}]}\tilde{f}(\frac{(p)}{(q)}I_{1}^{i_{1}}\dots I_{k}^{i_{k}})
  %      \\&=\E_{\vec{i}\in A}\lim_{N\to\infty}  \E_{\frac{(p)}{(\NN(I_{\vec{i}}))}\in \mathfrak{Q}(K)\colon \frac{(p)}{(\NN(I_{\vec{i}}))}I_{\vec{i}}\in \mathfrak{I}(K), \NN(\frac{(p)}{(\NN(I_{\vec{i}}))}I_{\vec{i}})\leq N} \tilde{f}\left(\frac{(p)}{(\NN(I_{\vec{i}}))}I_{\vec{i}}\right)  \text{ (\cite[Theorem~11.1.5]{murtyesmondebook})}
  %      \\&=\E_{\vec{i}\in A}\tilde{f}(I_{\vec{i}})\cdot\lim_{N\to\infty}  \E_{\frac{(p)}{(\NN(I_{\vec{i}}))}\in \mathfrak{Q}(K)\colon \frac{(p)}{(\NN(I_{\vec{i}}))}I_{\vec{i}}\in \mathfrak{I}(K), \NN(\frac{(p)}{(\NN(I_{\vec{i}}))}I_{\vec{i}})\leq N} f(p)f(\NN(I_{\vec{i}}))^{-1}
  %      \\&=\E_{\vec{i}\in A}C_{\vec{i}}(f)\tilde{f}(I_{1})^{i_{1}}\dots \tilde{f}(I_{k})^{i_{k}} 
  %  \end{split}
%\end{equation}
%for  
%$$C_{\vec{i}}(f):=\lim_{N\to\infty}  \E_{\frac{(p)}{(\NN(I_{\vec{i}}))}\in \mathfrak{Q}(K)\colon \frac{(p)}{(\NN(I_{\vec{i}}))}I_{\vec{i}}\in \mathfrak{I}(K), \NN(\frac{(p)}{(\NN(I_{\vec{i}}))}I_{\vec{i}})\leq N} f(p)f(\NN(I_{\vec{i}}))^{-1}$$
%which is independent of the choices of the extention $\tilde{f}$, provided that all th elimits  $C_{\vec{i}}(f)$ exist. 


%Suppose first that (i) holds.  Now let $\chi'$ be a modified Dirichlet chracter. %By Lemma \ref{ccf}, $f\chi$   is complatible with some ideal class representation $(I_{1},\dots,I_{d})$. 
%Then (i) implies (iv) which implies that $C_{\vec{i}}(f\chi')=0$ for all $\vec{i}$ (since the set of $p$ for which  $\frac{(p)}{(\NN(I_{\vec{i}}))}I_{\vec{i}}\in \mathfrak{I}(K)$ is the disjoint union of finitely many elements in $\mathcal{AP}[\tau_{d}]$) {\color{blue} another mistake here: for this to be true, we need the average of $f\chi'$ along AP to be zero, which reuqires some proof}.
%So for any extension $\widetilde{f\chi'}$ of $f\chi'$, it follows from \eqref{fweopf} that 
%$$\lim_{N\to\infty} \E_{\mathfrak{v}\in \mathfrak{I}(K), \NN(\mathfrak{v})\leq N}\widetilde{f\chi'}(\mathfrak{v})=0.$$
%This implies (v).

%Conversely, assume that (v) holds. We show that this implies (iv). Let $\chi'$ be a modified Dirichlet character. %By Lemma \ref{ccf}, $f\chi$  is complatible with some ideal class representation $(I_{1},\dots,I_{d})$.
%Then it follows from (v) that $$\lim_{N\to\infty} \E_{\mathfrak{v}\in \mathfrak{I}(K), \NN(\mathfrak{v})\leq N}\widetilde{f\chi'}(\mathfrak{v})=0$$
%for any extension $\widetilde{f\chi'}$ of $f\chi'$. 
%Similar to \eqref{fweopf}, we have that 
%$$P(z_{1},\dots,z_{k}):=\E_{\vec{i}\in A}C'_{\vec{i}}(f\chi')z_{1}^{i_{1}}\dots z_{k}^{i_{k}}=0$$
%for all $z_{1},\dots,z_{k}$ with $z_{i}^{d_{i}}=f(x_{i})$,
%where $I_{i}^{d_{i}}=(x_{i})$ and
% $$C'_{\vec{i}}(f):=\liminf_{N\to\infty}  \E_{\frac{(p)}{(\NN(I_{\vec{i}}))}\in \mathfrak{Q}(K)\colon \frac{(p)}{(\NN(I_{\vec{i}}))}I_{\vec{i}}\in \mathfrak{I}(K), \NN(\frac{(p)}{(\NN(I_{\vec{i}}))}I_{\vec{i}})\leq N} f(p)f(\NN(I_{\vec{i}}))^{-1}.$$
% Fix any such $z_{1},\dots,z_{k-1}$, since $f(x_{i})\neq 0$, we have that $P(z_{1},\dots,z_{k-1},\cdot)$ is a polynomial of degree at most $d_{k}-1$ with $d_{k}$ different roots. %{\color{red} this is where we need to use the assumption that $f(x_{i})\neq 0$, since $f$ is compatible with $(I_{1},\dots,I_{k})$}. 
%Thus, $P$ is independent of the the last variable. Similarly, the polynomial $P$ is also independent of all the variables, and thus $P\equiv 0$. So $C'_{\vec{i}}(f\chi)=0$ for all $\vec{i}$. In particular, $C'_{\vec{0}}(f\chi')=0$ and so (iv) holds for such $\chi'$.
%We are done. 




\iffalse

%{\color{red} there is a problem with this equality. In order for $f(p/q)$ to be well defined (or to extend to domain of $f$ from $\OK$ to $K$, we need $f$ to be $\S$-valued instead of $\U$-valued). Note also that in Proposition 4.7 we have to work with $\U$-valued functions, because even if $f$ is $\S$-valued, $f\chi$ is always $\U$-valued (or at least $\S\cup\{0\}$ valued).}

\textbf{Claim 1.} We have $$\lim_{N\to\infty}\E_{\mathfrak{v}\in \mathfrak{I}(K)\colon \NN(\mathfrak{v})\leq N}\tilde{f}(\mathfrak{v})=0$$
for all the extensions $\tilde{f}$ of $f$
if and only if $C_{\vec{i}}(f)=0$ for all $\vec{i}$.

For convenience we assume that all the limits under consideration exists {\color{red} this assumption is fine when proving the claim.}
If $C_{\vec{i}}(f)=0$ for all $\vec{i}$, then clearly (iii) implies that $$\lim_{N\to\infty}\E_{\mathfrak{v}\in \mathfrak{I}(K)\colon \NN(\mathfrak{v})\leq N}\tilde{f}(\mathfrak{v})=0$$
for all the extensions $\tilde{f}$ of $f$. Conversely, we have that 
$$P(z_{1},\dots,z_{k}):=\E_{\vec{i}\in A}C_{\vec{i}}(f)z_{1}^{i_{1}}\dots z_{k}^{i_{k}}=0$$
for all $z_{i}^{d_{i}}=f(p_{i}/q_{i})$.
Fix any $z_{1},\dots,z_{k-1}$ with $z_{i}^{d_{i}}=f(x_{i})$, we have that $P(z_{1},\dots,z_{k-1},\cdot)$ is a polynomial of degree at most $d_{k}-1$ with $d_{k}$ different roots {\color{red} this is where we need to use the assumption that $f(x_{i})\neq 0$}. So $P$ is independent of the the last variable. Similarly $P$ is independent of all the variables, and thus $P\equiv 0$. So $C_{\vec{i}}(f)=0$ for all $\vec{i}$. This proves the claim.

\

Now suppose that (iv) holds, then Claim 1 implies (iii) (whose LHS is the term $C_{\vec{0}}(f)$).


Now we assume that (i) holds. Clearly (iv) holds by Claim 1 if we can prove the following:

\textbf{Claim 2.} For any $I\in\mathfrak{I}(K)$, there exist $r\in\Z$ and a union of finitely many $D$-dimensional infinite arithematic progressions $P$ such that 
$$\frac{(p)}{(q)}I\in\mathfrak{I}(K)\Leftrightarrow \frac{(p)}{(q)}=\frac{(p')}{(r)} \text{ for some } p'\in P.$$

By multiplying with soem conjugates of $q$, we may assume WLOG that $q\in\Z$ {\color{red} check}. Then $\frac{(p)}{(q)}I\in\mathfrak{I}(K)\Leftrightarrow q^{D}\vert\NN(p)\NN(I)$.
Assume that $\NN(I)=u_{1}^{t_{1}}\dots u_{n}^{t_{n}}$ for some disticnt prime numbers $u_{i}$ and $t_{i}\in\N$. 
If $t\vert q$ but $t\nmid \NN(I)$ for some prime number $t$, then $t^{D}\vert \NN(p)$ and so $p/t\in\mathcal{O}_{K}$. So $\frac{(p)}{(q)}=\frac{(p/t)}{(q/t)}$. If $u_{i}^{t_{i}+1}\vert q$ for some $i$, then $p^{D}\vert\NN(p)\NN(I)$ implies that $u_{i}^{D}\vert \NN(p)$ and again $\frac{(p)}{(q)}=\frac{(p/u_{i})}{(q/u_{i})}$.
So we may assume WLOG that $q=\NN(I)$.

Then $q^{D}\vert\NN(p)\NN(I)\Leftrightarrow \NN(I)^{D-1}\vert\NN(p)$. The set of such $p$ is clearly a finite union of $D$-dimensional infinite arithematic progressions. This completes the proof of  Propostion \ref{characterization}. 
 

\fi



%\subsection{Measurability of aperiodic and pretentious functions}

%In this subsection we will culminate our previous work which will allow us to obtain the measurability result in Lemma~\ref{measurability}, crucial to do the precise analysis we wish to perform on $\mathcal{M}$, by splitting it into ``random'' and ``structured'' parts.

 
 
 
%In order to show that aperiodicity or pretentiousness are distinct, but more importantly, complementary notions, we need the following lemma:


We are now ready to prove Theorem \ref{characterization}:

%\begin{proof}[Proof of Theorem \ref{characterization}]
%Suppose that $f$ is a completely multiplicative function which is not aperiodic.  
%    It follows from Proposition \ref{dircharequivalence} that there exists a Dirichlet character $\chi\colon \OK\to\C$ and some extension $\tilde{f\chi}$ of $f\chi$, we have that 
%        $$\lim_{N\to\infty}\E_{\mathfrak{v}\in \mathfrak{I}(K)\colon \NN(\mathfrak{v})\leq N}\tilde{f\chi}(\mathfrak{v})=0.$$
%Since $g$ is completely multiplicative, case (iii) of Conjecture \ref{c27k} cannot happen, and thus we have that $\D(\tilde{f\chi},\NN(\mathfrak{v})^{i\tau})<\infty$ for some $\tau\in\R$. %Restricting to the sum over principle prime ideals in $\D(\tilde{f\chi},\NN(\mathfrak{v})^{i\tau})$, this implies that
%%  $\D_{0}(f\chi,\NN(u)^{i\tau})\leq \D(\tilde{f\chi},\NN(\mathfrak{v})^{i\tau})<\infty$.    {\color{red} Here I am not using the full strength of the fact that $\D(\tilde{f\chi},\NN(\mathfrak{v})^{i\tau})<\infty$, but I hope this is good enough for our purposes.}
%\end{proof}
 
%\begin{theorem}\label{characterization}
%{\color{red} $\S$-valued} completely multiplicative functions over %$\mathcal{O}_K$ are either aperiodic or pretentious.
%\end{theorem}
\begin{proof}[Proof of Theorem \ref{characterization}]  
Suppose that $f$ is not aperiodic. By  Lemma \ref{tv}, we have that $f(\e)=1$ for all units $\e$. By Proposition \ref{prop: dircharequivalence}, 
there exists a modified Dirichlet character $\chi\colon \Z[\tau_{d}]^{\times}\to\S$ and some 
         extension $\widetilde{f\chi}$ of $f\chi$ 
        such that the following fails to be true
        $$\lim_{N\to\infty}\E_{\mathfrak{v}\in \mathfrak{I}(\Q(\sqrt{-d}))\colon \NN(\mathfrak{v})\leq N}\widetilde{f\chi}(\mathfrak{v})=0.$$
 By Theorem \ref{C22intro}, we must have that $\D(\widetilde{f\chi},\NN(\mf u)^{i\tau})<\infty$ for some $\tau\in\R$.
 It is not hard to see that this implies that $f$ is  pretentious.
 %
%This is a consequence of Theorem~\ref{C22} and Lemma~\ref{lem: dircharequivalence}.  Namely, if a function is not aperiodic, it must mean that the averages of $\widetilde{g\chi}$ are not $0$, for some modified Dirichlet character $\chi$ (here we allow the trivial one) and some extension $\widetilde{g\chi}$ of $g\chi$  by Proposition~\ref{prop: dircharequivalence}. 
%Now, the averages of the multiplicative function $\widetilde{g\chi}$ fail to be zero exactly when $\widetilde{g\chi} \sim \NN(u)^{i\tau}$ for some $\tau \in \R$, completing the proof.
%    {\color{red} TBC. It should follow from interpreting 0-averages along any progression as the same as 0-averages when multiplied by any Dirichlet character. From this it follows that one of these is not 0 if and only if $g\cdot \chi \sim \NN^{i\tau}$ for some $\tau$. In fact one can then show that both $\tau$ and $\chi$ are unique, but I don't think we need that.}
\end{proof}
A consequence of Theorem~\ref{characterization} above is that it allows us to easily obtain the following measurability result in $\mathcal{M}$:
\begin{lemma}\label{measurability}
The set of pretentious completely multiplicative functions $\mathcal{M}_p$ is Borel.
\end{lemma}
\begin{proof}
By Theorem~\ref{characterization} and the definition of aperiodic functions, we may write 
\[
\mathcal{M}_p = \bigcup_{P\in\mathcal{AP}[\tau_{d}]} \left\{ f \in \mathcal{M} : \limsup_{N \to \infty} \left| \E_{u \in P, \NN(u) \leq N } f(u)\right|>0\right\}.
\]
Since $\mathcal{AP}[\tau_{d}]$ is countable, the displayed union is countable. It now follows by standard methods that $\mathcal{M}_p$ is a countable union of subsets that are each Borel measurable, since they are achieved as the nullset of $\limsup$ of sequences of continuous functions on $\mathcal{M}$, and the set $\mathcal{M}$ comes equipped with the topology of pointwise convergence.
\iffalse
Observe that Proposition~\ref{prop: dircharequivalence} allows us to write $\mathcal{M}_p$ %(see also Lemma~\ref{lem: dircharequivalence}) 
as
\[
\mathcal{M}_p = \bigcup_{\alpha,\beta \in \OK} \left\{ g \in \mathcal{M} : \limsup_{N \to \infty} \left| \E_{u \in \mathcal{O}_K^{\times}, \NN(u) \leq N } g(\alpha u+\beta)\right|>0\right\}.
\]
Since $\OK$ is countable, the displayed union is countable. It now follows by standard methods that $\mathcal{M}_p$ is a countable union of subsets that are each Borel measurable, since they are achieved as the nullset of $\limsup$ of sequences of continuous functions on $\mathcal{M}$, and this comes equipped with the pointwise convergence topology.
\fi
\end{proof}
\iffalse
\section{Hal\'asz's theorem and applications}\label{SEChalasz}
The goal of this section is to prove an analog of Hal\'asz's theorem (Theorem \ref{C22intro}) as well as a classification result for multiplicative functions (Theorem \ref{characterization}) over imaginary quadratic number fields. We are interested in obtaining Hal\'asz's theorem  because it will give us a convenient characterization of pretentious multiplicative functions on the set of ideals of a number field $\mathcal{O}_K$, which we denote by $\mathfrak{I}(K)$. This, in turn, will allow us to show a measurability result we need, in order to carry out the strategy of the proof, as delineated in Section~\ref{sec: proofstrat}. 

We will adapt the ideas from \cite[Chapter~6]{E79book} to the context of quadratic imaginary rings of integers. It is worth pointing out that many of the proofs in this section are much simpler in the few cases where the quadratic imaginary field is also a UFD: ideal-domained functions in that case do not need to be invoked, and one can work with standard multiplicative functions on the ring of integers.

Throughout this section, $g(\mathfrak{u})$ will denote an arbitrary multiplicative function on $\mathfrak{I}(K)$.% We denote by $\NN(I)$ the norm of the ideal $I$. \WS{move the definition of $\NN(I)$ to earlier sections} 
When using the complex variable $s \in \C$, we will often write $s=\sigma+i\tau$.  Also throughout, the function $G(s)$  will be given by:
$$G(s):=\sum_{\mf{u}\in \mathfrak{I}(K)}g(\mf u)\NN(\mf u)^{-s} = \lim_{N \to \infty} \sum_{1\leq \NN(\mf u)\leq N } g(\mf u)\NN(\mf u)^{-s}.$$
 

 
%    For convenience write $a:=a_{1}+ia_{2}$ for any letter $a$.



\subsection{Asymptotic average formula for multiplicative functions}
The purpose of this section is to prove the following analog of \cite[Theorem~6.1]{E79book}, which will be the main tool we use to prove Theorems \ref{characterization} and \ref{C22intro}.

 \begin{theorem}\label{C21}
 Let $g\colon\mf I(K)\to\C$ be a multiplicative function with $\vert g\vert\leq 1$. Then there exist $C>0,$ $\alpha\in\R$ and a slowly oscillating function\footnote{We say that $L:\R \to \C$ is slowly oscillating if for every $u>0$ it is $\lim_{x \to \infty} \frac{L(ux)}{L(x)}=1$.} $L\colon \R\to \C$ with $\vert L\vert=1$ so that 
 \begin{equation}\label{1}
     \sum_{1\leq\NN(\mf u)\leq x}g(\mf u)=\frac{x^{1+i\alpha}}{1+i\alpha}CL(\log  x)+o(x), \text{ as }x\to\infty
 \end{equation}
if and only if, in the same notation, we have that  
  \begin{equation}\label{2}
     G(s)=\frac{C}{s-(1+i\alpha)}L\left(\frac{1}{\sigma-1}\right)+o\left(\frac{1}{\sigma-1}\right), \text{ as }\sigma\to 1
 \end{equation}
  uniformly on each bounded interval $-M\leq \tau\leq M$.
  \end{theorem}

Throughout the proof, we denote
$$H(x):=\sum_{1\leq \NN(\mf u)\leq x}g(\mf u).$$



We first show that \eqref{1} implies \eqref{2}. We remark that this implication is not needed in this paper, but we provide its proof for completeness. Our approach is to apply a multi-dimensional integration by parts.
  
  \begin{proof}[Proof of \eqref{1} implies \eqref{2}] 
 % Let
 % $$H(x):=\sum_{1\leq \NN(\mf u)\leq x}g(\mf u).$$
 Note that $H(1-)=0$ and 
    $\lim_{N\to\infty}\left| \frac{H(x)}{sx^{s}}\right|\leq \lim_{N\to\infty}\frac{O(1)}{sx^{\sigma-1}}=0$ if $\sigma>1$ by Corollary~\ref{lem: bounds for number of ideals of norm at most x}.    
    Using the properties of the Riemann-Stieltjes integral,
   %(write $u=u_{1}+iu_{2}$)
    \begin{equation}\label{early}
    \begin{split}
    &\qquad \frac{G(s)}{s}=\frac{1}{s}\lim_{R\to\infty}\int_{1_-}^R\frac{1}{x^{s}}\,dH(x)=-\lim_{R\to\infty}\int_{1}^R H(x)\,d\left(\frac{1}{sx^{s}}\right)
    \\&\qquad\;\;\;\quad\quad=\lim_{R\to\infty}\int_1^R\frac{1}{x^{s+1}}H(x)\,dx = \int_1^{\infty} \frac{1}{x^{s+1}}H(x)\,dx.
    \end{split}
 \end{equation}
 %  Indeed, it is clear that $H(1-)=0$ by construction, and $\left| \frac{H(x)}{su^{s}}\right|\leq C_K\frac{R}{sR^{\sigma}}$, for some constant $C_K$ depending on the number field $\mathcal{O}_K$   provided that $x\geq R$ which, in turn, goes to $0$ as $R \to \infty$ with our choice of $s$.
   
Let $\varepsilon>0$. By (\ref{1}), there exists $x_{0}\geq 1$ such that if $x \geq x_{0}$, then we have that 
 $$\vert H(x)-\lambda x^{1+i\alpha}L(\log x)\vert<\varepsilon x,$$
 where $\lambda=C/(1+i\alpha)$. Then
     \begin{equation}\nonumber
    \begin{split}
    &\qquad \left| \int_{x \geq x_{0}}\frac{1}{x^{s+1}}H(x)\,dx-\int_{x\geq x_{0}}\frac{1}{x^{s+1}}\lambda x^{1+i\alpha}L(\log x) \ dx\right|
   \leq %\int_{\max\{u_{1},u_{2}\}\geq x_{0}} \e \NN(x)^{-\sigma}\,d(\NN(x))
    \varepsilon\int_{x\geq 1} x^{-\sigma}\,dx = \frac{\varepsilon}{\sigma-1}.
    %=\varepsilon\frac{1}{\sigma-1}\left(-\int_{\NN(u)\geq 1} \,d(\NN(u)^{1-\sigma})\right)
    %\leq \frac{\varepsilon C_0}{\sigma-1},
    \end{split}
 \end{equation} 
%for some universal constant $C_{0}$. 
  Therefore, we have that 
  \begin{equation}\nonumber
    \begin{split}
    \frac{G(s)}{s}=\lambda \int_{x\geq 1} x^{-s+i\alpha}L(\log  x)\,dx+o\left(\frac{1}{\sigma-1}\right).
    \end{split}
 \end{equation}
 Fix $\Delta>1$. Let $y_{1}=\exp(\Delta^{-1}(\sigma-1)^{-1})$ and  $y_{2}=\exp(\Delta(\sigma-1)^{-1})$. Then since $L$ is slowly oscillating, we have that
 $$L(\log  x)=L\left(\frac{1}{\sigma-1}\right)+o(1),$$
 uniformly for all $y_{1}\leq  x\leq y_{2}$. Therefore, 
   \begin{equation}\nonumber
    \begin{split}
   &\qquad \int_{y_{1}\leq x\leq y_{2}}x^{-s+i\alpha}L(\log x)\,dx
   =\int_{ y_{1}\leq x\leq y_{2}}x^{-s+i\alpha}L\left(\frac{1}{\sigma-1}\right)\,dx+o\left(\frac{1}{\sigma-1}\right).
  % \\&=\int_{ y_{1}\leq \NN(x)\leq y_{2}}\frac{1}{\NN(x)^{s+1}}\omega(x)^{1+i\alpha}L\left(\frac{1}{\sigma-1}\right)\,d(\NN(x))+,
    \end{split}
 \end{equation}
Therefore, given that $|L|=1$, using crude upper bounds, as well as the computation above, we have
  \begin{eqnarray*}
      \left|G(s)-\frac{s\lambda}{s-(1+i\alpha)} L\left(\frac{1}{\sigma-1}\right) \right|
   & = & \vert s\vert\cdot\left|\frac{G(s)}{s}-\lambda L\left(\frac{1}{\sigma-1}\right) \int_{x\geq 1}x^{-s+i\alpha}\,dx\right|
   \\ &\leq & 2\vert s\lambda\vert\int_{x<y_{1} \text{ or } x>y_{2}} x^{-\sigma}\,dx+o\left(\frac{1}{\sigma-1}\right)\\
 & \leq & \frac{1}{\sigma-1}(2\vert s\lambda\vert(\Delta^{-1}+e^{-\Delta})+o(1)).
  \end{eqnarray*}
Since $\vert L\vert=1$ and 
\[
 \frac{s\lambda}{s-(1+i\alpha)} = \lambda+\frac{C}{s-(1+i\alpha)},
 \]
 we have that 
  \begin{equation}\nonumber
    \begin{split}
   \left|G(s)-\frac{C}{s-(1+i\alpha)} L\left(\frac{1}{\sigma-1}\right) \right|
  \leq \frac{1}{\sigma-1}(2\vert s\lambda\vert(\Delta^{-1}+e^{-\Delta})+o(1)).
       \end{split}
 \end{equation}
 
   %\\&=\frac{1}{\sigma-1}(2\vert\lambda\vert((1-e^{-\Delta^{-1}})^{2}+(2e^{-0.5\Delta}))+o(1))
 %  \leq \frac{1}{\sigma-1}(2\vert\lambda\vert(\Delta^{-2}+(2e^{-0.5\Delta}))+o(1)).
% {\color{red} Here we view $\int \NN(u)^{-\sigma}\,d(\NN(u))$ as $-\frac{1}{\sigma-1}\int d(\NN(u)^{1-\sigma})$ and do integration by parts. For the part $\NN(u)<y_{1}$, we pick the end points as $1+i$, $1+i\sqrt{y_{1}-1}$, $\sqrt{y_{1}-1}+i$ and $\sqrt{y_{1}-1}+i\sqrt{y_{1}-1}$. For the part $\NN(u)>y_{2}$, we pick the end points of the two infinite boxes as $\sqrt{y_{2}}/2$ and  $\sqrt{y_{2}}i/2$.
% }
% {\color{blue} Careful here with the norm of $1+i$ being equal to $\sqrt{2}$. I think this changes some of the bounds by constants in the third and fourth lines, but after all, since we are going to send $\Delta$ to infinity, this will not matter. I think the result we care about holds modulo constants.}
 Let $\sigma$ approach $1$ from the right first, and then let $\Delta$ grow very large, we get (\ref{2}).
 %
 \iffalse
 
 . It follows from the computation above that
 \[
  G(s)=\frac{\lambda}{s-(1+i\alpha)} L\left(\frac{1}{\sigma-1}\right) +o\left(\frac{1}{\sigma-1}\right)
 \]
 so we complete this part of the proof by using the following two facts. First, that $|L|=1$, and second, that
 \[
 \frac{s\lambda}{s-(1+i\alpha)} = \lambda+\frac{C}{s-(1+i\alpha)}
 \]
  % \begin{equation}\nonumber
  %  \begin{split}
  % &\qquad\int_{u_{1},u_{2}\geq 1}\frac{s}{\NN(u)^{s+1}}\omega(u)^{1+i\alpha}\,d(\NN(u))+o\left(\frac{1}{\sigma-1}\right)
  % \\&=\frac{s\lambda}{s-(1+i\alpha)} L\left(\frac{1}{\sigma-1}\right) \int_{u_{1},u_{2}\geq 1}\,d(-\NN(u)^{(1+i\alpha)-s})+o\left(\frac{1}{\sigma-1}\right)
     % \\&=-\lambda L(\frac{1}{\sigma-1}) \int_{u_{1},u_{2}\geq 1}\omega(u)^{1+i\alpha}\,d(\NN(u)^{-s})+o(\frac{1}{\sigma-1})
   %    \\&=\frac{s\lambda}{s-(1+i\alpha)}K(s-(1+i\alpha))  L\left(\frac{1}{\sigma-1}\right)+o\left(\frac{1}{\sigma-1}\right)
    %    \\&=(\lambda+\frac{C}{s-(1+i\alpha)})K(s-(1+i\alpha)) L\left(\frac{1}{\sigma-1}\right)+o\left(\frac{1}{\sigma-1}\right)
     %   \\&=\frac{C}{s-(1+i\alpha)}K(s-(1+i\alpha)) L\left(\frac{1}{\sigma-1}\right)+o\left(\frac{1}{\sigma-1}\right).
    %\end{split}
 %\end{equation}
 %{\color{red} take care of this constant $K(s-(1+i\alpha))=2^{s-(1+i\alpha)}$ in the statement of the theorem}
\fi
  \end{proof}
For the implication \eqref{2} implies \eqref{1}, we begin noticing the following property of $G(s)$:
%  {\color{blue} If you have time, verify some of the red parts in the outline. Most of the red parts should be very similar to the proof in Elliott79, but need to be checked carefully. Most red parts are independent, and working on them does not require to understand the full picture of the proof}
 $$G'(s)=-\sum_{\mf u\in\mf I (K)}g(\mf u)\log(\NN(\mf u))\NN(\mf u)^{-s},$$
when $\sigma>1$.

As shall become apparent, we will first consider the case where $\alpha=0$, and then employ an integral trick to deduce the general case for arbitrary values of $\alpha$. The following equality will be very convenient for us:
  $$\frac{1}{2\pi i}\int_{\sigma-i\infty}^{\sigma+i\infty}\frac{x^{s}}{s^{2}}\,ds=\mathds{1}_{x>1}\log x,$$
for all $x,\sigma>0$. By the dominated convergence theorem, %, which allows us to exchange the sum and the integral,
we have that  (if $\sigma>1$)
 \begin{equation}\label{1234}
    \begin{split}
       -\sum_{1 \leq \NN(\mf u)\leq x}g(\mf u)\log(\NN(\mf u))\log\left(\frac{x}{\NN(\mf u)}\right)=\frac{1}{2\pi i}\int_{\sigma-i\infty}^{\sigma+i\infty}\frac{x^{s}}{s^{2}}G'(s)\,ds.
    \end{split}
 \end{equation}
% {\color{red} I do not know if something similar holds if we replace the range $\NN(u)<x$ by a box}
% {\color{blue} I suspect we might be able to obtain a similar result by defining a $G(s)$ with two variables, like $G(s,t):=\sum_{u \in \Z[i]^+} g(u)u_1^{-s} u_2^{-t}$, for complex variables, $s, t$, and then using a two-variable version of the integral from $\sigma-i\infty$ to $\sigma+i\infty$. Then, we should be getting a rectangle, instead of a box in the formula above. It's unclear how this change in the $G$ would affect the rest of the arguments.}
 
 Let $M_{1},M_{2}\geq 2$ be real numbers that we fix temporarily. The range of integration of the integral above will be divided into three pieces $I_1, I_2, I_3$. We will obtain bounds on the contributions of each of these terms using slightly different methods.

 \subsubsection{Estimate on $I_{1}$.}
 We begin by considering the range $I_{1}:= \{ \vert\tau\vert>M_{2}\}$. Then (let $1<\sigma\leq 2$) the Cauchy-Schwarz inequality yields
  \begin{equation}\label{a1}
    \begin{split}
       \left|\frac{1}{2\pi i}\int_{I_{1}}\frac{x^{s}}{s^{2}}G'(s)\,ds\right|^{2}
       \leq \left(\frac{x^{\sigma}}{2\pi}\int_{I_{1}}\left|\frac{G'(s)}{s^{2}}\right|\,d\tau \right)^{2}
       \leq  \frac{x^{2\sigma}}{4\pi^{2}}\int_{I_{1}}\left|\frac{G'(s)}{s G(s)}\right|^{2}\,d\tau\cdot \int_{I_{1}}\left|\frac{G(s)}{s}\right|^{2}\,d\tau.
    \end{split}
 \end{equation}
%We will find it helpful to establish the following inequality 
 % \begin{equation}\label{123first}
 %   \begin{split}
 %      \int_{\vert\tau-\rho\vert\leq 1}\Bigl\vert G(s)\Bigr\vert^{2}\,d\tau\leq \frac{O(1)}{\sigma-1}
 %   \end{split}
 %\end{equation}
%for all $\rho\in\R$.

%We will formulate this inequality in the lemma below; to prove it, we will make use of an application of Parseval's formula that, as we will see, will keep reappearing in slightly different forms.

We will make use of the following estimate:

\begin{lemma}\label{firstplancherel}
Let $g, G$ and $H$ be as above. Then, for all $\rho \in \R$, it is
      \begin{equation}\nonumber%\label{123second}
    \begin{split}
       \int_{\vert\tau-\rho\vert\leq 1}\Bigl\vert G(s)\Bigr\vert^{2}\,d\tau\leq \frac{O(1)}{\sigma-1}.
    \end{split}
 \end{equation}
\end{lemma}
\begin{proof}
Recall from (\ref{early}) that
\[
\frac{G(s)}{s}=\int_{1}^{\infty} \frac{1}{u^{s+1}} H(u) \ du.
\]
Using the change of variables $u=e^w$ in this integral allows us to rewrite it as:
\[
\frac{G(s)}{s}=\int_{0}^{\infty} \frac{H(e^{w})}{e^{w(s+1)}} \ e^w \ dw.
\]
In other words, it is
\[
\frac{G(s)}{s}= \int_{0}^{\infty} H(e^{w}) e^{-w\sigma}e^{-wi\tau}\ dw.
\]
This exhibits the function of interest $\frac{G(s)}{s}$ as a Fourier transform of $H(e^w)e^{-w\sigma}$, which is viewed as a function of $w$.

Moreover, since, as a function of $\tau$, the function $\frac{G(s)}{s}$ is $O(|\tau|^{-1})$, it follows that it is in $L^2(\R)$. Thus, Parseval's theorem implies that
\[
\int_{-\infty}^{\infty} \left| \frac{G(s)}{s} \right|^2 \ d\tau = 2\pi \int_{0}^{\infty} \left| H(e^{w}) e^{-w\sigma}\right|^2 \ dw.
\]
%Now, if we take the trivial bound $|g(\mf u)|\leq 1$, the innermost integral can be bounded above by a constant multiple (depending on the number field that controls the sizes of the balls determined by the associated norm) of $e^w$, giving us an upper bound not bigger than
Since $|g(\mf u)|\leq 1$, we have that $\vert H(e^{w})\vert\leq C_{1}e^{w}$ by Lemma \ref{mv7}, and so
\[
2\pi \int_{0}^{\infty} \left| H(e^{w}) e^{-w\sigma}\right|^2 \ dw\leq 2\pi C_{1}\int_{0}^{\infty} e^{-2w(\sigma-1)} \ dw \leq \frac{C_2 }{\sigma-1},
\]
where $C_1$ and $C_{2}$ are absolute constants that depend only on the number field. Thus, whenever $1<\sigma \leq 3$, we have
\[
\int_{-1}^1 |G(s)|^2 \ d\tau \leq 10\int_{-\infty}^{\infty} \left| \frac{G(s)}{s}\right|^2 \ d\tau \leq \frac{10 C_2}{\sigma-1}.
\]
Replacing $g(\mf u)$ by $g(\mf u)\NN(\mf u)^{-i\rho}$ in the above argument yields that the bound (because we used the crude bound on $g$) 
\[
\int_{|\tau-\rho| \leq 1} |G(s)|^2 \ d\tau \leq \frac{10  C_2}{\sigma-1}
\]
is uniform over all real $\rho$.
\end{proof}


%{\color{red} check this part. This is done on page 228 of Elliot79 using Fourier transform. I don't know if Fourier transform still apply in our case}

  It follows from Lemma  \ref{firstplancherel} that
 \begin{equation}\label{a2}
    \begin{split}
      & \int_{I_{1}}\left|\frac{G(s)}{s}\right|^{2}\,d\tau
       \leq \sum_{\vert m\vert>M_{2}}\int_{\vert\tau-m\vert\leq 1} \left|\frac{G(s)}{s}\right|^{2}\,d\tau \\& \quad\quad\quad\quad\quad\quad
       \leq 4\sum_{\vert m\vert>M_{2}} m^{-2}\int_{\vert\tau-m\vert\leq 1} \left|G(s)\right|^{2}\,d\tau
       <\frac{O(1)}{M_{2}(\sigma-1)}.
    \end{split}
 \end{equation}
 
 
 %  \begin{equation}\nonumber
%    \begin{split}
 %      &\qquad\int_{-1}^{1}\Bigl\vert G(s)\Bigr\vert^{2}\,d\tau
%       =\int_{-1}^{1}\sum_{u,u'\in\Z[i]_{+}}g(u)\overline{g}(u')\log(\NN(u))\log(\NN(u'))\NN(uu')^{-\sigma}\NN(u/u')^{-i\tau}\,d\tau
 %      \\&\leq 2\sum_{u,u'\in\Z[i]_{+}}\log(\NN(u))\log(\NN(u'))\NN(uu')^{-\sigma}
 %   \end{split}
% \end{equation}


%{\color{red}
%There is an issue that needs discussion here: the proof would mostly go through without many changes, but it makes strong use of Euler products. There are analogs of Euler products for number fields in general, but they rely on IDEALS instead of primes.

%This is an important issue and not easy to get around: there are primes of the form $4n+1$ that will not be primes in $\Z[i]$, because, say $5=(2+i)(2-i)$, and notice that the decomposition into now primes of $\Z[i]$ will always be a prime times its conjugate, which is bad, because $2-i \notin \Z[i]_+$. We could still perhaps restrict to primes of certain forms, like $4n+3,$ which will stay in our grid of interest
%}
%In the following lemma, we will find it useful to consider multiplicative functions that map all the units to $1$. It is okay to assume this, for the balls $\NN(u) \leq x$ are asymptotically invariant under multiplication by units (even if we consider dilations and shifts of them). This can be seen as a F\o lner property of these balls: they absorb units.

%Indeed, if $f(i) \neq 1$, then $\E_{1\leq \NN(u) \leq x} g(u) = \E_{1 \leq \NN(u)\leq x} g(iu) = g(i) \E_{1\leq \NN(u) \leq x} g(u)$, which implies that $\lim_{x \to \infty} \E_{1 \leq \NN(u) \leq x} g(u)=0$.

%Therefore, such functions would not give any meaningful contribution to the averages we want to control, and we can safely take this extra assumption without loss of generality.

Now we need to estimate the term
$\int_{I_{1}}\left|\frac{G'(s)}{s G(s)}\right|^{2}\,d\tau$
 on the right hand side of  (\ref{a1}).   
We need the following analog of  \cite[Lemma~6.6]{E79book}.

\begin{lemma}\label{eulerproduct1}
Let $h(\mf p):=h_{s}(\mf p):=\sum_{k=1}^{\infty}g(\mf p^{k})\NN(\mf p^{k})^{-s}$, where $g: \mf I(K) \to \U$ is a multiplicative function. Then,
$$G(s)=\left(\prod_{\mf p \in \mf I(K), \NN(\mf p)=2}(1+h(\mf p))\right)\exp\left(\sum_{\mf p \in \mf I(K), \NN(\mf p)\geq 3}g(\mf p)\NN(\mf p)^{-s}\right)G_{1}(s),$$
for $\sigma>1$ for some $G_{1}$ that is analytic when $\sigma>1/2$ and satisfies the bounds
$$C_0^{-1}\leq \vert G_{1}(s)\vert\leq C_0, \vert G'_{1}(s)\vert\leq C_1 \quad (\text{for }\sigma \geq 1),$$
where $C_0, C_1$ are some absolute constants.
%\AK{are the bounds ONLY for $\sigma>1$?} \FM{I believe these bounds work for $\sigma \geq 1$ because $G_1$ itself is analytic for $\sigma>1/2$}
%{\color{red} Check if 100 is optimal. For example, can we replace 100 by 3?

%It seems that we can replace 100 by 3. The reason is that in our case we have $\vert h_{s}(p)\vert\leq 2/3$ when $\NN(p)\geq 3$. So the same argument in Elliot79 applies to our case.

%}

\end{lemma}
\begin{proof}
Observe that $h$ is well-defined for $\sigma>0$ by construction. Indeed, observe that for $\mf p \in \mf I(K)$ prime,  
\[
|h(\mf p)| \leq \sum_{k=1}^{\infty} \NN(\mf p^k)^{-\sigma}=\NN(\mf p)^{-\sigma}(1-\NN(\mf p)^{-\sigma})^{-1} \leq
\begin{cases}
    4\NN(\mf p)^{-\sigma}, \quad \text{if } \sigma \geq 1/2, \\
    2\NN(\mf p)^{-1}, \quad \text{if } \sigma \geq 1.
\end{cases}
\]

Since $|g|\leq 1$,   the series
$\sum_{\mf u \in \mf I(K)} |g(\mf u)| \NN(\mf u)^{-\sigma}$ converges when $\sigma>1$ by Lemma \ref{zp}.

By the fundamental theorem of arithmetic for ideals in Dedekind domains, ideals in $\mf I(K)$ can be written as products of prime ideals in a unique way (up to permutations). Therefore, we may apply the same argument as \cite[Lemma~6.6]{E79book} to conclude  that 
%Our goal is to express $G(s)$ with an Euler product. In order to do that, we need several things. 
%Firstly, the series
%\[
%\sum_{\mf u \in \mf I(K)} |g(\mf u)| \NN(\mf u)^{-\sigma}
%\]
%needs to converge, which it does for $\sigma>1$ given that $|g(\mf u)|\leq 1$ for all $\mf u \in \mf I(K)$. Next, we need to use the analog of the fundamental theorem of arithmetic for ideals in Dedekind domains. This allows us to write any ideal $\mf u \in \mf I(K)$ uniquely as a product of prime ideals in $\mf I(K)$. Thus,
\[
G(s)=\prod_{\mf p \in \mf I(K), \NN(\mf p)\geq 2} (1+h(\mf p)).
\]
Now, let $$G_1(s):= \prod_{\mf p \in \mf I(K), \NN(\mf p)\geq 3} (1+h(\mf p)) \exp(-g(\mf p)\NN(\mf p)^{-s}).$$ We claim that this gives us the desired representation of $G$. Indeed, it only remains to see that $G_1$ satisfies the asserted bounds, because the equality holds trivially.

We first consider the case where $\sigma>\frac{1}{2}$. Note that    $|h(\mf p)| \leq \frac{2}{3}$ when $\NN(\mf p)$ is large enough. Using the estimate
\[
|\log(1+w)-w| \leq 2|w|^2,
\]
it follows as in \cite[Lemma~6.6]{E79book} that
\begin{eqnarray*}
    |\log(1+h(\mf p))-g(\mf p)\NN(\mf p)^{-s}| & \leq & |\log (1+h(\mf p))-h(\mf p)|+|h(\mf p)-g(\mf p) \NN(\mf p)^{-s}| \\ & \leq &
2|h(\mf p)|^2+\sum_{k=2}^{\infty} |g(\mf p^k)| \NN(\mf p)^{-k\sigma}\\ & \leq & 36\NN(\mf p)^{-2\sigma}.
\end{eqnarray*}
So $G_{1}$ is defined by a product which converges uniformly absolutely in any half plane $\sigma\geq\sigma_{0}$ with $\sigma_{0}>1/2$. 
So $G_{1}$ is analytic when $\sigma>1/2$.

Next, if $\sigma \geq 1$, we will have $|h(\mf p)|\leq \frac{2}{3}$ again, but now for all $\mf p$ with $\NN(\mf p)>2$, i.e. $\NN(\mf p) \ge 3$. Thus,
\[
\log |G_1(s)| \leq \sum_{\mf p \in \mf I(K), \NN(\mf p) \geq 3} \left( 2|h(\mf p)|^2+\sum_{k \geq 2} |g(\mf p^k)| \NN(\mf p)^{-k}\right) \leq 10\sum_{\mf p \in \mf I(K), \NN(\mf p) \geq 3} \NN(\mf p)^{-2},\]
%\[
%\leq 10T_0\sum_{p \geq 3} p^{-2}<5T_0 
%\]

%where we used the bound for $|h(\mf p)|$ we gave at the beginning of the proof, and $T_0$ is some constant that depends on the number field (Lemma \ref{mv7}). 
which is bounded by a universal constant, say $C_{00}$ by Lemma \ref{ggee}. Note now, that if $\NN(p)\geq 11$ and $\sigma \geq 3/4$, then $|h(\mf p)| \leq \frac{4}{\NN(\mf p)^{\sigma}}<2/3$, so we have the upper bound
\[
|G_1(s)| \leq e^{C_{00}} \prod_{\NN(p) \leq 7} (1+|h(\mf p)|)\leq C_0.
\]
Changing the signs and standard Euler product representation facts, one can obtain the upper bound for $G_1(s)$. See for example \cite[Lemma~2.13]{E79book} and the discussion right below it.

Now, the inequality above for $G_1(s)$ holds uniformly in the half-plane $\sigma \geq 3/4$. Thus, by Cauchy's theorem, we see that for $\sigma \geq 1$ it is
\[
|G_1'(s)| = \left| \frac{1}{2\pi i} \int_{|s-z|=1/4} \frac{G_1(z)}{(z-s)^2} \ dz \right|\leq C_1,
\]
for some other universal constant $C_1$.
\end{proof}
%{\color{red} give a proof. it should be correct}

We can now proceed by taking the derivative of $\log G(s)$, we have, using the representation in Lemma~\ref{eulerproduct1}, that
\begin{equation}\label{gergw2}
    \begin{split}
    \frac{G'(s)}{G(s)}=\sum_{\NN(\mf p)<3}\frac{h'(\mf p)}{1+h(\mf p)}-\sum_{\NN(\mf p)\geq 3}g(\mf p)\NN(\mf p)^{-s}\log(\NN(\mf p))+\frac{G'_{1}(s)}{G_{1}(s)}.
    \end{split}
 \end{equation}
Note that
  \begin{equation}\label{gergw}
    \begin{split}
       \int_{-\infty}^{\infty}\Bigl\vert\frac{G'_{1}(s)}{s G_{1}(s)}\Bigr\vert^{2}\,d\tau
       \leq O(1)\int_{-\infty}^{\infty}\vert s\vert^{-2}\,d\tau=O(1).
    \end{split}
 \end{equation}
 
 Next, for $\NN(\mf p)<3$, (i.e. $\NN(\mf p)=2$) consider (given $\sigma>1$)
 $$A_{\mf p}(s):=\frac{1}{1+h(\mf p)}=1+\sum_{m=1}^{\infty}(-h(\mf p))^{m}=1+\sum_{m=1}^{\infty}\left(-\sum_{k=1}^{\infty}g(\mf p^{k})2^{-ks}\right)^{m}=\sum_{\ell=1}^{\infty} a_{\ell}2^{-\ell s}.$$
If we write
 $$B_{2}(s):=1+\sum_{m=1}^{\infty}\left(\sum_{k=1}^{\infty}2^{-ks}\right)^{m}=\sum_{\ell=1}^{\infty} b_{\ell}2^{-\ell s} = \frac{2^s-1}{2^s-2},$$
 then clearly $\vert a_{\ell}\vert\leq b_{\ell}$.
Since
 \[
 |h'(\mf p)| = \left| \sum_{k=1}^{\infty}-g(\mf p^k)2^{-ks} \log 2^k \right| \leq \sum_{k=1}^{\infty} 2^{-k}k\log 2 <2\quad  (\sigma>1),
 \]
  whenever $1<\sigma<2$.
 we may use \cite[Lemma~6.5]{E79book} (due to Montgomery) to conclude 
 that 
  \begin{equation}\label{sdvv}
    \begin{split}
      &\qquad \int_{-1}^{1}\left|\frac{h'(\mf p)}{1+h(\mf p)}\right|^{2}\,d\tau
       \leq O(1)\int_{-1}^{1}\left|\frac{1}{1+h(\mf p)}\right|^{2}\,d\tau
       \leq O(1)\int_{-2}^{2}\left| B_{2}(s)\right|^{2}\,d\tau
       \\&=O(1)\int_{-2}^{2}\left| \frac{2^{s}-1}{2^{s}-2}\right|^{2}\,d\tau
       \leq O(1)\int_{-2}^{2}\frac{1}{\left| 2^{s}-2\right|^{2}}\,d\tau \leq O(1) \int_{-2}^2 |s|^{-2} \left| \sum_{l=1}^{\infty} 2^l\cdot 2^{-ls}\right|^2 \ d\tau \\& \leq O(1)\int_{-\infty}^{\infty} |s|^{-2} \left| \sum_{l=1}^{\infty} 2^l \cdot 2^{-ls} \right|^2 \ d\tau = O(1) \int_{0}^{\infty} \left| \sum_{2^l\leq e^w} 2^le^{-w\sigma}\right|^2 \ dw \\&\leq O(1) \int_{0}^{\infty} e^{-2w\sigma} \left(\sum_{2^l \leq e^w} 2^l\right)^2 \ dw  \leq O(1) \int_{0}^{\infty} e^{-2w(\sigma-1)} \leq \frac{O(1)}{\sigma-1}.
    \end{split}
 \end{equation}
%Using the
%\WS{  I think here we can just say that $O(1) \int_{0}^{\infty} \left| \sum_{2^l\leq e^w} 2^le^{-w\sigma}\right|^2 \ dw\leq  \frac{O(1)}{\sigma-1}$ by the first inequality on page 236 of the book. This will save about half page since we can also keep the next paragraph to ourself
%}
 
   Here, we used the same Fourier trick as in Lemma~\ref{firstplancherel}, but adapted for the function of interest. %Namely, we define $H(s):=\sum_{l \geq 1} 2^l \cdot 2^{-ls}$, and we use that
  %\[
  %s^{-1}H(s)=s^{-1}\int_{1^-}^{\infty} 2^{-us} \ d\tilde{H}(u) = \int_1^{\infty} 2^{-us} \log (2) \tilde{H}(u) \ du,
  %\]
  %where $\tilde{H}(u) := \sum_{l \leq u} 2^l$. Now, use the change of variables $2^u=e^w$ (so that $u \cdot \log (2)=w$ on the integral above to see that it is equal to
  %\[
   %\int_{\log 2}^{\infty} e^{-ws} H\left( \frac{w}{\log 2}\right) \ dw= \int_{\log 2}^{\infty} e^{-w\tau} H\left( \frac{w}{\log 2}\right) e^{-iw\tau }\ dw,
  %\]
  %which means that $e^{-ws}H\left( \frac{w}{\log 2}\right)$ is the Fourier transform of $s^{-1}H(s)$. Using Plancherel allows us to justify the last two lines in the computation above.
   On the other hand, we can apply this same argument by shifting $s$ to $s+im$, for $|m| \geq 2$, obtaining
    \begin{equation}\nonumber
    \begin{split}
      &\qquad \int_{-\infty}^{\infty}\left|\frac{h'(\mf p)}{s(1+h(\mf p))}\right|^{2}\,d\tau
       \leq \frac{O(1)}{\sigma-1}+\sum_{\vert m\vert\geq 2}\frac{4}{m^{2}}\int_{\vert\tau-m\vert\leq 1}\left|\frac{h'(\mf p)}{1+h(\mf p)}\right|^{2}\,d\tau=\frac{O(1)}{\sigma-1}.
    \end{split}
 \end{equation} 
 
 Finally, we again adapt this same trick to obtain the following upper bound:
  \begin{equation}\label{Jplanchereltrick}
    \begin{split}
       \int_{-\infty}^{\infty}\left| s^{-1}\sum_{\NN(\mf p)\geq 3}g(\mf p)\NN(\mf p)^{-s}\log(\NN(\mf p))\right|^{2}\,d\tau
       \leq \frac{O(1)}{\sigma-1}.
    \end{split}
 \end{equation} 
 Indeed, consider the function $J(s):=\sum_{\NN(\mf p) \geq 3} g(\mf p)\NN(\mf p)^{-s} \log \NN(\mf p)$ and its associated function
 \[
 \tilde{J}(u):=\sum_{\NN(\mf p) \leq u} g(\mf p)\log \NN( \mf p).
 \]
Since the estimates of elements of a given norm in a number field grows at most linearly by Lemma \ref{mv8}, since a prime ideal has norm a power of a prime $p$ that is at most $p^{[K: \Q]}$, and since $\sum_{p \leq u} \log p \leq 4u,$ for $u>0$ (this follows from Tchebycheff's inequality; see \cite[p. 235]{E79book}), we have that $\vert\tilde{J}(u)\vert\leq O(u)$.
So we may write
 \begin{equation}\label{early2}
s^{-1}J(s)=s^{-1}\int_{1-}^{\infty} u^{-s} \ d\tilde{J}(u)=\int_{1}^{\infty} u^{-s-1} \tilde{J}(u) \ du.
 \end{equation}
 
 Next, we effect the same change of variables as in the proof of Lemma~\ref{firstplancherel} to see that the left hand side of \eqref{Jplanchereltrick} is equal to
 
 \[
 2\pi \int_0^{\infty} \left|e^{-w\sigma} \tilde{J}(e^w)\right|^2  \ dw,
 \]
 %Thus, the integral in \eqref{Jplanchereltrick} is at most
 which is bounded by
 \[
 O(1) \int_0^{\infty} e^{-2w(\sigma-1)} \ dw \leq \frac{O(1)}{\sigma-1},
 \]
 as desired.
% {\color{red} check if this is correct. the book claim that this part is similar to (\ref{123})}
 Combining (\ref{gergw2}), (\ref{gergw}), (\ref{sdvv}) and (\ref{Jplanchereltrick}), we have that 
 \begin{equation}\label{a3}
    \begin{split}
       \int_{-\infty}^{\infty}\left|\frac{G'(s)}{s G(s)}\right|^{2}\,d\tau
       \leq \frac{O(1)}{\sigma-1}.
    \end{split}
 \end{equation}
 Combining  (\ref{a1}), (\ref{a2}) and (\ref{a3}), we have 
  \begin{equation}\label{e1}
    \begin{split}
        \left|\frac{1}{2\pi i}\int_{\vert\tau\vert>M}\frac{x^{s}}{s^{2}}G'(s)\,ds\right|\leq \frac{O(x^{\sigma})}{M_{2}^{1/2}(\sigma-1)}.
    \end{split}
 \end{equation}
%{\color{red}
%Note that this expression still holds if we replace $I_{1}$ by $\vert\tau-\alpha\vert>M$. This is because (\ref{a1}), (\ref{a2}) and (\ref{a3}) still hold ((\ref{a2}) relies on (\ref{123}) which still holds).
 % \begin{equation}\nonumber
  %  \begin{split}
 %      &\qquad \Bigl\vert\frac{1}{2\pi i}\int_{\vert\tau-\alpha\vert>M}\frac{x^{s}}{s^{2}}G'(s)\,ds\Bigr\vert^{2}\leq \Bigl(\int_{-\infty}^{\infty}\Bigl\vert\frac{G'(s)}{s G(s)}\Bigr\vert^{2}\,d\tau\Bigr)\Bigl\vert\frac{1}{2\pi i}\int_{\vert\tau-\alpha\vert>M}\frac{x^{s}G'(s)}{s^{2}}ds\Bigr\vert
 %      \\&\qquad \leq \frac{O(1)}{\sigma-1}\Bigl\vert\frac{1}{2\pi i}\int_{\vert\tau-\alpha\vert>M}\frac{x^{s}G'(s)}{s^{2}}ds\Bigr\vert
  %     \\&\leq \frac{O(1)}{\sigma-1} \frac{x^{2\sigma}}{4\pi^{2}}\int_{I_{1}}\Bigl\vert\frac{G'(s)}{s G(s)}\Bigr\vert^{2}\,d\tau\cdot \int_{I_{1}}\Bigl\vert\frac{G(s)}{s}\Bigr\vert^{2}\,d\tau
  %  \end{split}
 %\end{equation}
%}

  \subsubsection{Estimate on $I_{2}$.}
We next consider the range $I_{2}:=\{ \tau : M_{1}(\sigma-1)<\vert\tau\vert\leq M_{2}\}$. We will follow a lot of the same arguments as for the range $I_1$. First, given $1<\sigma<2$, we still have
  \begin{equation}\label{bdnd1}
    \begin{split}
       \left|\frac{1}{2\pi i}\int_{I_{2}}\frac{x^{s}}{s^{2}}G'(s)\,ds\right|^{2}
       \leq \left(\frac{x^{\sigma}}{2\pi}\int_{I_{2}}\left|\frac{G'(s)}{s^{2}}\right|\,d\tau\right)^{2}
       \leq  \frac{x^{2\sigma}}{4\pi^{2}}\int_{I_{2}}\left|\frac{G'(s)}{s G(s)}\right|^{2}\,d\tau\cdot \int_{I_{2}}\left|\frac{G(s)}{s}\right|^{2}\,d\tau.
    \end{split}
 \end{equation}
We already have a bound for the first integral thanks to (\ref{a3}),
%that appears as a result of the Cauchy-Schwarz inequality from our previous work on the range $I_1$, since we were able to show that
%\[
%\int_{-\infty}^{\infty} \left| \frac{G'(s)}{sG(s)}\right|^2 \ d\tau \leq \frac{O(1)}{\sigma-1}.
%\]
but we need to treat the second term differently.
%Now we look at the aforementioned second term, but we will deal with it differently.
To do so, we make use of the trivial estimate
\begin{equation}\label{bdnd2}
   \int_{I_2} \left| \frac{G(s)}{s} \right|^2 \ d\tau \leq \max_{s\in I_2} |G(s)|^{1/2} \int_{I_2} |s|^{-2} |G(s)|^{3/2} \ d\tau.
\end{equation}
The assumption \eqref{2}  implies that
\begin{equation}\label{bdnd3}
    \max_{s\in I_2} |G(s)|^{1/2} \leq \left(\frac{O(1)}{M_{1}(\sigma-1)}\right)^{1/2}.
\end{equation}
%{\color{red} this is still the case when $I_{2}$ is centered at $\alpha$, since 
%$\vert s-(1+\alpha i)\vert\geq \vert\sigma-1\vert$.}
It remains to estiamte $\int_{I_2} |s|^{-2} |G(s)|^{3/2} \ d\tau$.
By Lemma~\ref{eulerproduct1}, we can bound
\[
|G(s)| \leq O(1) \exp\left(\text{Re} \sum_{ \NN(\mf p)\geq 3} g(\mf p)\NN(\mf p)^{-s}\right)
\]
which implies that
\begin{eqnarray*}
    |G(s)|^{3/4} & \leq & O(1)  \exp\left(\sum_{\NN(\mf p)\geq 3} \text{Re} \ \frac{3}{4}g(\mf p)\NN(\mf p)^{-s}\right) \\ 
    & = & O(1) \exp\left(\text{Re} \ \sum_{\NN(\mf p)\geq 3} \left(\frac{3}{4}\right)^{\Omega(\mf p)}g(\mf p)\NN(\mf p)^{-s}\right),
\end{eqnarray*}
where $\Omega(\mf u)$ is the number of prime ideals that divide $\mf u$ counted with multiplicity. Applying Lemma~\ref{eulerproduct1} again (replacing $g(\mf u)$ by $g(\mf u)\left(\frac{3}{4}\right)^{\Omega(\mf u)}$ and using the fact that $|e^{z}|=e^{\text{Re }z}$ as well as using the bound $|1+h(\mf p)| \geq 1-\frac{2}{3}$ for $\NN(\mf p)=2,$ which appears in the denominator), we obtain (once more, we use the unique factorization property that ideals of a Dedekind domain enjoy) the upper bound
\begin{equation}\label{wgwgdg}
    |G(s)|^{3/4} \leq O(1)\left| \sum_{\mf u \in \mf I_{\geq 3}(K)} \left(\frac{3}{4}\right)^{\Omega(\mf u)} g(\mf u)\NN(\mf u )^{-s}\right|,
\end{equation}
where $\mf I_{\geq 3}(K)$ denotes the set of ideals whose prime ideal factors are all of norm at least 3.
Next, we need a variant of \cite[Lemma~6.5]{E79book}, which we show next. It crucially relies on the fact that if $y \neq 0$ (or trivially in the case $y=0$) we have
\[
\int_{-1}^1 (1-|t|)e^{ity} \ dt = \left(\frac{\sin \left(\frac{y}{2}\right)}{\frac{y}{2}}\right)^2 \geq 0.
\]
\begin{lemma}\label{montgomeryadapted}
Let
\[
A(s):=\sum_{\mf u \in \mf I(K)} a_{\mf u} \NN(\mf u)^{-s} \quad B(s)= \sum_{\mf u \in \mf I(K)} b_{\mf u} \NN(\mf u)^{-s}
\]
be series that converge absolutely for $\text{Re } s=\sigma$. Suppose that $|a_{\mf u}|\leq b_{\mf u}$ for all $\mf u \in \mf I(K)$. Then, for any $T \geq 0$ we have
\[
\int_{-T}^T |A(s)|^2 \ d\tau \leq 2 \int_{-2T}^{2T} |B(s)|^2 \ d\tau.
\]
\end{lemma}
\begin{proof}
Notice that, by splitting the integral on the left hand side above and adding more positive terms we have
\begin{eqnarray*}
    \int_{-T}^T |A(s)|^2 \ d\tau & \leq & 2 \int_{-2T}^{2T} \left(1-\frac{|\tau|}{2T}\right) |A(s)|^2 \ d\tau \\ & = & 2\sum_{\mf u \in \mf I(K)}\sum_{\mf v \in \mf I(K)} a_{\mf u} \overline{a_{\mf v}} \NN(\mf u \mf v)^{-\sigma} \int_{-2T}^{2T} \left(1-\frac{|\tau|}{2T}\right) \NN(\mf u\mf v^{-1})^{-i\tau} \ d\tau.
\end{eqnarray*}
Since we know that each of these integrals is positive, we apply the trivial bounds on the coefficients to see that the last double sum is at most
\[
2\sum_{\mf u \in \mf I(K)}\sum_{\mf v \in \mf I(K)} b_{\mf u} \overline{b_{\mf v}} \NN(\mf u \mf v)^{-\sigma}  \int_{-2T}^{2T} \left(1-\frac{|\tau|}{2T}\right) \NN(\mf u\mf v^{-1})^{-i\tau} \ d\tau = \]
\[
2\int_{-2T}^{2T} \left(1-\frac{|\tau|}{2T}\right) |B(s)|^2 \ d\tau \leq 2 \int_{-2T}^{2T} |B(s)|^2 \ d\tau,
\]
and we are done.
\end{proof}
 Since $\sigma>1$, by (\ref{wgwgdg}) and Lemma~\ref{montgomeryadapted},
\begin{eqnarray*}
   \int_{-1}^1 |G(s)|^{3/2} \ d\tau  & \leq & O(1) \int_{-1}^1 \left| \sum_{\mf u \in  \mf I_{\geq 3}(K)} \left(\frac{3}{4}\right)^{\Omega(\mf u)}g(\mf u)\NN(\mf u)^{-s}\right|^2 \ d\tau \\
& \leq & O(1) \int_{-1}^1 \left| \sum_{\mf u \in \mf I_{\geq 3}(K)} \left(\frac{3}{4}\right)^{\Omega(\mf u)}\NN(\mf u)^{-s}\right|^2 \ d\tau.
\end{eqnarray*}
Another application of Lemma~\ref{eulerproduct1} ensures the following upper bound of the previous integrand:
\[
O(1) \left|\exp\left(\sum_{\mf p \in \mf I(K), \NN(\mf p) \geq 3} \frac{3}{4} \NN(\mf p)^{-s}\right)\right|^2 \leq O(1) |\zeta(s)|^{3/2},
\]
where we recall that $\zeta$ is the Dedekind zeta function defined in (\ref{d1}). %As in the integer case, it only has a simple pole at $s=1$ and is analytic in the rest of the complex plane (see \cite[Theorem~5.11]{neukirch99}, for example).  
%This implies the existence of 
It follows from Lemma \ref{zp} that there exist absolute constants $c_2, c_3$ such that
\[
\int_{-1}^1 |G(s)|^{3/2} \ d\tau \leq c_2\int_{-2}^2 \left(\frac{1}{|s-1|^{3/2}}+1\right) \ d\tau<c_3 (\sigma-1)^{-1/2}.
\]
Similarly, we have
\[
\int_{|\tau-m|\leq 1} |G(s)|^{3/2} \ d\tau<c_3(\sigma-1)^{-1/2}, \text{ for } m \in \Z \setminus \{0\}. 
\]
This, in turn, implies that   
\begin{equation}\label{fffp}
    \begin{split}
   &\qquad \int_{-\infty}^{\infty} \frac{|G(s)|^{3/2}}{|s|^{2}} \ d\tau  \leq  \int_{-1}^1 |G(s)|^{3/2}\ d\tau +4 \sum_{|m| \geq 2} m^{-2} \int_{|\tau-m|\leq 1} |G(s)|^{3/2} \ d\tau\\& \qquad\qquad\qquad\qquad\;\;\;\quad < O(1)(\sigma-1)^{-1/2}.
   \end{split}
\end{equation}
Combining (\ref{a3}), (\ref{bdnd1}), (\ref{bdnd2}), (\ref{bdnd3}) and (\ref{fffp}),
 for $x$ large enough in terms of $K$ and $M$, we have that  
   \begin{equation}\label{e2}
    \begin{split}
        \left|\frac{1}{2\pi i}\int_{I_{2}}\frac{x^{s}}{s^{2}}G'(s)\,ds\right|\leq \frac{O(x^{\sigma})}{M_{1}^{1/4}(\sigma-1)}.
    \end{split}
 \end{equation}
 %{\color{red} the conclusion holds for the case when $I_{2}$ is centered at $\alpha$}

  \subsubsection{Estimate on $I_{3}$ and completion of the proof.}
 We finally consider the range $I_{3}: \vert\tau\vert\leq M_{1}(\sigma-1)$.
 In the rest of the proof, we set $\sigma=\sigma_0=1+(\log x)^{-1}$.
 By using the same argument as in \cite[p. 239]{E79book}, %which is not influenced at all by changing to a different ring of integers, 
 we have that
    \begin{equation}\label{e3}
    \begin{split}
       \frac{1}{2\pi i}\int_{I_{3}}\frac{x^{s}}{s^{2}}G'(s)\,ds=-Cx\log x\cdot L(\log x)+O(K^{-1}x\log x).
    \end{split}
 \end{equation}

%{\color{blue}We will use brute force on the contour which is the line segment $I_3$ for $\sigma=\sigma_0$ together with the semicircle, call it $C,$ with center at $(\sigma_0,0)$ and radius $K(\sigma_0-1)=K\log^{-1}x.$ Since $x$ can be taken arbitrarily large, we can assume that $K<\log x+1$ (so the contour doesn't contain $0$). $C$ has parametrization
%\[r(t)=(\sigma_0+K\log^{-1}x\cos t)+i(K\log^{-1}x\sin t),\;\;t\in [\pi/2,3\pi/2].\]
%We will estimate 
%\[\int_C\frac{x^s}{s^2(s-1)}\;ds=\int_{\frac{\pi}{2}}^{\frac{3\pi}{2}}\frac{x^{r(t)}r'(t)}{r(t)^2(r(t)-1)}\;dt.\]
%{\color{red} what about the case 
%\[r(t)=(\sigma_0+K\log^{-1}x\cos t)+i(\alpha+K\log^{-1}x\sin t),\;\;t\in [\pi/2,3\pi/2]?\]
%}
%Notice that $x^{\log^{-1}x}=e,$ so
%\[x^{r(t)}\;\;\text{is bounded by}\;\;x^{\sigma_0}\sim x.\]
%On the other hand
%\[|r'(t)|=K\log^{-1}x,\]
%\[|r(t)-1|=\frac{\sqrt{1+2K\cos t+K^2}}{\log x}\geq \frac{K-1}{\log x}\sim \frac{K}{\log x},\]
%{\color{red} (not sure if this $\sim$ here is cheating)}
%and 
%\begin{eqnarray*}
%|r(t)|^2& = &\frac{(\log x+1)^2+K^2+2(\log x+1)K\cos t}{\log^2x}\\
%&\geq & \frac{(\log x+1)^2+K^2-2(\log x+1)K}{\log^2 x}\\
%& = &\frac{(\log x-K)^2+2\log x+1-2K}{\log^2x}.
%\end{eqnarray*}

%{\color{red} the lower bound for $\vert r(t)\vert$ is the distance between $0$ and $\sigma_{0}+i\alpha$ (the center of the circle) minues $K\log^{-1}x$ (the radius). So
%$$\vert r(t)\vert\geq \sqrt{(1+\log^{-1}x)^{2}+\alpha^{2}}-K\log^{-1}x\geq 1-(K-1)\log^{-1}x>1/2.$$

%So putting everything together should give us an upper bound of the form $O(x).$ 
%}


%Taking $x$ large enough so that $2K\leq \log x+1,$ we get 
%\[|r(t)|^2\geq \frac{\log x}{\log^2x}=\frac{1}{\log x},\]
%so, putting everything together should give us an upper bound of the form $O(K^{-1}x\log x).$ This, while a terrible upper bound, seems to work.
%}  
Combining all the three parts, for $\sigma=1+(\log x)^{-1}$, it follows from (\ref{1234}), (\ref{e1}), (\ref{e2}) and (\ref{e3}) that
     \begin{equation}\nonumber
    \begin{split}
       \left| \sum_{ \NN(\mf u)\leq x}g(u)\log(\NN(\mf u))\log\left(\frac{x}{\NN(\mf u)}\right)-Cx\log x\cdot L(\log x)\right|\leq O(1)(M_{2}^{-1/2}+M_{1}^{-1/4})x\log x
    \end{split}
 \end{equation}
 if $x$ is sufficiently large depending on $M_{1},M_{2}$.
 %\WS{I thought we discussed about it before but I forgot the answer: shouldn't the RHS be $x^{1+(\log x)^{-1}}\log x$?}
 %\FM{$x^{(\log x)^{-1}}=e$, so we might as well ignore it, since we already have an $O(1)$ term.}
 Dividing by $x\log x$, then letting $x\to\infty$ and finally letting $M_{2}\to\infty$, and finally $M_{1} \to \infty$ we have that 
  \begin{equation}\label{ggg}
    \begin{split}
      \sum_{\NN(\mf u)\leq x}g(\mf u)\log(\NN(\mf u))\log\left(\frac{x}{\NN(\mf u)}\right)=Cx\log x \cdot L(\log x)+o(x\log x).
    \end{split}
 \end{equation}
 %\WS{Sometimes we sum over $\NN(\mf u)<x$ and sometimes we sum over $\NN(\mf u)\leq x$. Maybe we can make them consistent}
 Next, we make use of the argument on \cite[p. 240]{E79book}. 
We give a sketch of how the computations work out:
first, we use (\ref{ggg}) for $x$ and $x(1+\varepsilon)$ for some $0<\varepsilon<1/2$ and take their differences to conclude that
\[
\sum_{\NN(\mf u) \leq x} g(\mf u)\log(\NN(\mf u))\log(1+\varepsilon)+\sum_{x \leq \NN(\mf u) \leq x(1+\varepsilon)} g(\mf u) \log(\NN(\mf u)) \log\left(\frac{x}{\NN(\mf u)}\right) = \]
\[
C\varepsilon x\log x \cdot L(\log x)+o(x\log x)+O(\varepsilon^2x\log x).
\]
Here, part of the sums from the left hand side are absorbed by the big $O$ on the right hand side, where Corollary~\ref{lem: bounds for number of ideals of norm at most x} gets used.

Using the expansion $\log(1+\varepsilon)=\varepsilon+O(\varepsilon^2)$, we can divide by $\varepsilon$ to conclude
\[
\limsup_{x \to \infty} \left| (x\log x)^{-1}\sum_{\NN(\mf u) \leq x} g(\mf u)\log(\NN(\mf u))-CL(\log x)\right| \leq O(\varepsilon),
\]
This inequality is based, first on the fact that $\frac{1}{x\log x}\sum_{\NN(\mf u) \leq x} g(\mf u)\log \NN(\mf u) \cdot \varepsilon^2=O(\varepsilon^2)$ again by Corollary~\ref{lem: bounds for number of ideals of norm at most x}, together with the estimate
\[
\frac{1}{x\log x}\sum_{x \leq \NN(\mf u) \leq x(1+\varepsilon)} g(\mf u) \log \NN(\mf u) \log \left(\frac{x}{\NN(\mf u)}\right) \leq \]
\[
\frac{\log (x(1+\varepsilon))}{x\log x}\sum_{x \leq \NN(\mf u) \leq x(1+\varepsilon)} \varepsilon  = O(\varepsilon^2),
\]
but $\varepsilon$ was arbitrary, which means that
\[
\sum_{\NN(\mf u) \leq x} g(\mf u)\log(\NN(\mf u)) =Cx\log xL(\log x)+o(x\log x) \quad (x \to \infty).
\]
But we also have the following two estimates:
\begin{eqnarray*}
    \sum_{x/\log x \leq \NN(\mf u) \leq x} g(\mf u)\log(\NN(\mf u)) & = & \sum_{x/\log x \leq \NN(\mf u) \leq x} g(\mf u)(\log x+ O(\log \log x)) \\ & = &
\log x \sum_{\NN(\mf u) \leq x}g(\mf u) + O(x(\log x)^{1/2}),
\end{eqnarray*}
from crude bounds and enlarging the sum if necessary, together with the fact that $\log\log x$ is dominated by $\sqrt{\log x}$, and finally combined with the following inequality which is obtained using crude bounds, also using Corollary~\ref{lem: bounds for number of ideals of norm at most x}
\[
\sum_{\NN(\mf u) \leq x/ \log x}g(\mf u) \log(\NN(\mf u)) = O(x).
\]
Combining all these yields
\[
\sum_{\NN(\mf u)\leq x} g(\mf u) = CxL(\log x)+o(x)
\]
when $x \to \infty$, which proves Theorem \ref{C21} when $\alpha=0$.

\

Finally, we analyse the case when $\alpha \neq 0$. Let
\[
D(x):=\sum_{\NN(\mf u)\leq x} g(\mf u)\NN(\mf u)^{-i\alpha}.
\]
Using the same analysis as above, but replacing $g(\mf u)$ by $g(\mf u)\NN(\mf u)^{-i\alpha}$ we see that
\[D(x)=CxL(\log x)+o(x).\]
Notice that integration by parts yields
\[
\sum_{\NN(\mf u)\leq x} g(\mf u)=\int_{1-}^x y^{i\alpha} \ dD(y) = x^{i\alpha}D(x)-i\alpha\int_1^x y^{i\alpha-1}D(y) \ dy.
\]
 Let $\varepsilon \in (0,1/2)$. In the range $\varepsilon x <y\leq x$, we have
 \[
 D(y)=CyL(\log x)+o(x),
 \]
 which implies that
 \begin{eqnarray*}
      i\alpha \int_{\varepsilon x}^x y^{i\alpha-1}D(y) \ dy & = & i\alpha CL(\log x)\int_{\varepsilon x}^x y^{i\alpha} \ dy+o(x)\\ & = & \frac{i\alpha}{1+i\alpha}Cx^{1+i\alpha}L(\log x)+O(\varepsilon x),
 \end{eqnarray*}
 and using the trivial bounds on the range $1\leq y <\varepsilon x$, we see that, in fact,
 \[
 \sum_{\NN(\mf u)\leq x} g(\mf u) = \frac{C}{1+i\alpha}x^{1+i\alpha}L(\log x)+O(\varepsilon x).
 \]
 Sending $\varepsilon \to 0$ completes the proof of Theorem \ref{C21}.
% We then use the argument on page 240 of Elliot79 to show that
% {\color{red} check the details}
% \begin{equation}\nonumber
 %   \begin{split}
  %    \sum_{u\in\Z[i]_{+}, \NN(u)<x}g(u)=Cx L(\log x)+o(x).
  %  \end{split}
% \end{equation}
 
 %{\color{red} The next step is to apply the previous equation for $g(u)\omega(u)$ for some multiplicative function $\omega$ and then use integration by parts. However, I do not know how to do  integration by parts in our case, as the range is a ball not a rectangle.

%{\color{blue} We could attempt the following: first transform the sums into integrals over balls. Then, use a change of coordinates into suitable polar coordinates so that the balls become rectangles. Then use integration by parts over there for the appropriate Riemann-Stieltjes integral. Finally convert it back to balls for the sums.}
 
 %On the other hand, at least we have the following partial conclusion: if (\ref{2}) holds for $C=0$, then $\sum_{u\in\Z[i]_{+}, \NN(u)<x}g(u)=o(x)$. Combining with the techinque in the proof of Theorem 6.2, we should be able to show that if 
%$$\sum_{p}\NN(p)^{-1}(1-\text{Re}(g(p)\NN(p)^{-i\tau}))<\infty$$
%for all real number $\tau$,
%then (\ref{2}) holds with $C=0$ and thus $\sum_{\NN(u)\leq x}g(u)=o(x).$

%Again the problem is: here the average $\NN(u)\leq x$ is in a ball not a rectangle. 
 
% }
 

%Strategy of the proofs:

%Step 1: show the difficult direction of Conjecture \ref{C21} for ball averages, namely if 
%$$G(s)=\frac{C}{s-(1+i\alpha)}L\left(\frac{1}{\sigma-1}\right)+o\left(\frac{1}{\sigma-1}\right), \sigma\to 1$$
% hold uniformly on each bounded interval $-M\leq \tau\leq M$, then 
%$$ \sum_{\NN(u)<x}g(u)=\frac{x^{1+i\alpha}}{1+i\alpha}CL(\log \vert x\vert)+o(\vert x\vert^{2}), \vert x\vert\to\infty.$$

%I checked the proof and it seems that the easy direction is not needed. 
\subsection{Hal\'asz's theorem: the pretentious approach}
\iffalse
We will define the following notion of distance between multiplicative functions on $\mf I(K)$:
\begin{definition}
Let $f,g\colon\mf  I(K)\to\C$ be multiplicative functions, and denote by $\mf p$ a prime ideal of $\mathcal{O}_K$. The \emph{distance} function between $f$ and $g$, denoted by $\D(f,g)$ is given by
$$\D(f,g):=\sum_{\mf p\colon \text{prime in } \mathcal{O}_K}\frac{1}{\NN(\mf p)}(1-\text{Re}(f(\mf p)\overline{g}(\mf p))).$$
Similarly, the truncated distance between $f$ and $g$ is given by:
$$\D(f,g;K,M):=\sum_{K<\NN(\mf p)\leq M}\frac{1}{\NN(\mf p)}(1-\text{Re}(f(\mf p)\overline{g}(\mf p))).$$
\end{definition}

 

 


We now move to prove the following version of \cite[Theorem~6.2]{E79book}:
\begin{theorem}\label{C22}
     Let $g\colon\mf I(K)\to\C$ be a multiplicative function with $\vert g\vert\leq 1$. Then, 
     \begin{enumerate}
         \item If $\inf_{\tau\in\R}\D(g,\NN(\mf u)^{i\tau})=\infty$, then $\lim_{x \to \infty} \E_{\NN(\mf u)\leq x}g(\mf u)=0$.
         \item If, on the other hand, $\D(g,\NN(\mf u)^{i\tau})<\infty$, for some $\tau\in\R$, then
         $$\sum_{1 \leq \NN(\mf u)\leq x}g(u) = \frac{x^{1+i\tau}}{1+i\tau}\cdot R_1\prod_{\mf p\colon \text{prime in } \mathcal{O}_K,\NN(\mf p)\leq x}(1-1/\NN(\mf p))(1+h_{1+i\tau}(\mf p))+o(x),$$
         where $h_{s}(\mf p):=\sum_{k=1}^{\infty}g(\mf p^{k})\NN(\mf p^{k})^{-s}$, and $R_1$ is the residue of the Dedekind zeta function of our number field at its simple pole at $s=1$.
     \end{enumerate}
\end{theorem}

\fi

%In this Section, we prove Theorem \ref{C22intro}.
%We need the following lemma from \cite{E79book}:
%\begin{lemma}[Lemma 6.7 in \cite{E79book}]\label{uniformityeliott}
%Let $l_n(w)$, $n \in \N$ be a sequence of real functions, defined and continuous for $-M \leq w \leq M$. For each $w \in [-M,M]$, suppose that $l_n(w)$ increases (decreases) monotonically to the value of a continuous function $l(w)$. Then, this convergence is uniform on $[-M,M]$. 
%\end{lemma}
%\begin{remark}
%Note that if $l_n(w) \to \infty$ (as $n \to \infty$) the result also holds by considering the sequence $\min(1, 1/l_n(w)).$ \WS{is this limit taken in $n$ or $w$? do we allow $l(w)$ to be infinity?} \FM{Yes, I think it is allowed but with the change using the minimum it becomes finite again.}
%\end{remark}
In this Section, we prove Theorem \ref{C22intro}.
We need the following lemma:
\begin{lemma}\label{uniformityeliott}
Let $l_n(w)$, $n \in \N$ be a sequence of real functions, defined and continuous for $-M \leq w \leq M$. For each $w \in [-M,M]$, suppose that $l_n(w)$ increases monotonically to $\infty$. Then, this divergence is uniform for  $w\in [-M,M]$. 
\end{lemma}
\begin{proof}
  The conclusion follows from \cite[Lemma~6.7]{E79book} by observing that the sequence $l'_{n}(w):=\min(1, 1/l_n(w))$ decreases monotonically to 0.
\end{proof}
 







\begin{proof}[Proof of Theorem \ref{C22intro}]
   We first assume that  $\inf_{\tau\in\R}\D(g,\NN(\mf u)^{i\tau})=\infty$. 
   By Lemma~\ref{uniformityeliott}, we have (taking the infinite sum as a limit in $N$ of the partial sum of balls determined by the norm) that 
    $$\sum_{\mf p\colon \text{prime ideal in } \mathcal{O}_K}\frac{1}{\NN(\mf p)^{\sigma}}(1-\text{Re}(f(\mf p)\overline{g}(\mf p)))\to \infty \quad (\sigma\to 1+)$$
   uniformly for all $\tau$ belonging to a given bounded interval $\vert\tau\vert\leq M$.
   %Let $\zeta(s)=\sum_{\mf u\in\mf I(K)}\NN(\mf u)^{-s}$.
   By Lemma~\ref{eulerproduct1}, there is some absolute constant $C$ such that
   \begin{eqnarray*}
       \vert G(s)/\zeta(s)\vert & \leq & C\left|\exp\left(\text{Re}\sum_{\mf p \in \mf I(K), \NN(\mf p)\geq 3}(g(\mf p)-1)\NN(\mf p)^{-s}\right)\right|
  \\& = & C\left|\exp\left(-\sum_{\mf p \in \mf I(K), \NN(\mf p)\geq 3}\NN(\mf p)^{-\sigma}(1-\text{Re}(g(\mf p)\NN(\mf p)^{-i\tau}))\right)\right|\to 0,
   \end{eqnarray*} 
 as $\sigma\to 1+$ with the same uniformity for all  $\vert\tau\vert\leq M$. Therefore, by Lemma \ref{l1}, we have that 
   \begin{equation}\nonumber
   \begin{split}
    \vert G(s)\vert\to o\left(\frac{1}{\sigma-1}\right)  
   \end{split}
   \end{equation}
 as $\sigma\to 1+$ with the same uniformity for all  $\vert\tau\vert\leq M$.
 It then follows from Theorem~\ref{C21} that  $\lim_{N \to \infty} \E_{\NN(\mf u) \leq N} g(\mf u)=0$ (setting $C=\alpha=0$ and $L\equiv 1$).
 
   
   \
   
   We now assume that $\D(g,\NN(\mf u)^{i\alpha})<\infty$ for some $\alpha\in\R$. As we saw in the very last steps of the proof of Theorem~\ref{C21}, it is enough to establish the asymptotic estimate in the case where $\alpha=0$, so assume that this is the case, namely
   \begin{equation}\label{aiso}
       \sum_{\mf p}\NN(\mf p)^{-1}(1-\text{Re}(g(\mf p)))<\infty.
   \end{equation}
   %{\color{red} we assume that $\alpha=0$ (here similar to the last step of Conjecture \ref{C21}, we need to think how to pass to the general case).

   %The book claims that the  comments in the part of the proof where $D(x)$ is introduced (for Theorem 6.1 in [Elliott79]) are enough to allow us to assume $\alpha=0$. In our case, we have the same argument working out, so it should be okay to also assume $\alpha=0$ WLOG.
   %}
   We caution the reader that in the rest of the proof, the constant term $O(1)$ is allowed to depend implicitly on the function $g$.
   Let $K\geq 1$. For $\sigma>1$ and $\vert\tau\vert \leq K(\sigma-1)$, by Lemma~\ref{eulerproduct1}, we may write 
   \begin{equation}\label{e41}
  G(s)/\zeta(s)=Q(s)\exp\left(-\sum_{\mf p \in \mf I(K), \NN(\mf p)\geq 3}(1-g(\mf p))\NN(\mf p)^{-s}\right)   \end{equation}
   for some function $Q$ which is continuous in every rectangle $1\leq \sigma\leq 2, \vert\tau\vert\leq M$. 
   
   
   
   
%   {\color{red}




%If $\alpha\neq 0$, then
%$$\zeta(s-i\alpha)=\sum_{u}\NN(u)^{i\alpha}\NN(u)^{-s}$$
%and
%\begin{equation}
%  G(s)/\zeta(s-i\alpha)=Q(s)\exp\left(-\sum_{p \in \Z[i]_+, \NN(p)\geq 100}(\NN(p)^{-i\alpha}-g(p))\NN(p)^{-s}\right)   
%  \end{equation}
%
%  and
%\begin{equation}
%  G(s-i\alpha)/\zeta(s)=Q(s)\exp\left(-\sum_{p \in \Z[i]_+, \NN(p)\geq 100}(1-\NN(p)^{i\alpha}g(p))\NN(p)^{-s}\right)   
%  \end{equation}
%}

   
   
   
   Denote
   $$L\left(\frac{1}{\sigma-1}\right):=\exp\left(i\text{Im} \sum_{\mf p \in \mf I(K), \NN(\mf p)\geq 3}\NN(\mf p)^{-\sigma}g(\mf p)\right).$$
%   and $$C:=Q(1)\exp\left(-\sum_{\mf p \in \mf I(K), \NN(\mf p)\geq 3}\NN(\mf p)^{-1}(1-\text{Re}(g(\mf p)))\right).$$
To show that $L(w)$ is a slowly varying function, we need the following Lemma.


  \begin{lemma}\label{lem68elliott}
   Let $M_{1}\geq 1$. If \eqref{aiso} holds,
   then
   $$\sum_{\mf p}\vert\NN(\mf p)^{-\sigma'}-\NN(\mf p)^{-s}\vert\cdot\vert 1-g(\mf p)\vert\to 0,$$
   as $\sigma'\to 1+$, uniformly for $\vert\tau\vert\leq M_{1}(\sigma'-1), (\sigma'-1)/2\leq \sigma-1\leq \sigma'-1$. 
   \end{lemma}
   \begin{proof}
The proof is similar to \cite[Lemma~6.8]{E79book}.
       Note that since $|1-z|^2 \leq 2(1-\text{Re } z)$ for $\vert z\vert\leq 1$. We have
       \begin{eqnarray*}
           \left| \NN(\mf p)^{-\sigma'}-\NN(\mf p)^{-s}\right|^{2}
               & = & \NN(\mf p)^{-2\sigma'}\left| 1-\NN(\mf p)^{-(\sigma'-s)}\right|^{2}
               \\&\leq & \NN(\mf p)^{-2\sigma} 2(1-\text{Re}(\exp(-(\sigma'-s)\log\NN(\mf p))))
               \\&\leq & 2\NN(\mf p)^{-2\sigma}\vert\sigma'-s\vert \log \NN(\mf p)\\
              & \leq & 2(1+2K)\NN(\mf p)^{-2\sigma}(\sigma-1)\log \NN(\mf p),
       \end{eqnarray*} 
       where we used that $1-\text{Re }(e^{-(\sigma'-s)\log \NN(\mf p)}) = \text{Re } \left((\sigma'-s)\int_0^1 e^{-t(\sigma'-s)\log \NN(\mf p)} \ dt\right)$, which is bounded above (using trivial bounds, since $\text{Re }(\sigma'-s)\geq 0$) by $|\sigma'-s|\log \NN(\mf p)$. Thus, there exists some universal constant $C>0$ such that
 \begin{eqnarray*}
     \sum_{\mf p}\NN(\mf p)\vert \NN(\mf p)^{-\sigma'}-\NN(\mf p)^{-s}\vert^{2}
               & \leq & 6K(\sigma-1)\sum_{\mf p}\NN(\mf p)^{-(2\sigma-1)}\log \NN(\mf p)
               \\&\leq & 6K(\sigma-1)\left(-\frac{\zeta'(2\sigma-1)}{\zeta(2\sigma-1)}+O(1)\right)\leq CK,
 \end{eqnarray*} 
  %     for some universal consant $K$. {\color{red} check the last step}
where in the second to last step we applied Lemma~\ref{eulerproduct1} once more to the derivative of $\log \zeta(s)$. Let $0<\varepsilon<1$. Again, since $\vert 1-z\vert^{2}\leq 2(1-\text{Re}(z))$, by (\ref{aiso}), there exists $N$  such that 
$$\sum_{\NN(\mf p)>N}\frac{\vert 1-g(\mf p)\vert^{2}}{\NN(\mf p)}<\varepsilon^{2}.$$
Therefore, by the Cauchy-Schwartz inequality, we have that 
 \begin{equation}\nonumber
           \begin{split}
               &\qquad \sum_{\NN(\mf p)>N}\left| \NN(\mf p)^{-\sigma'}-\NN(\mf p)^{-s}\right| \cdot \vert 1-g(\mf p)\vert\leq (CK)^{1/2}\varepsilon.              
           \end{split}
       \end{equation}
 On the other hand, for all $\NN(\mf p)<N$,
       \begin{equation}\nonumber
           \begin{split}
               \vert \NN(\mf p)^{-\sigma'}-\NN(\mf p)^{-s}\vert \cdot \vert 1-g(\mf p)\vert\to 0              
           \end{split}
       \end{equation}
       as $\sigma'\to 1+$ uniformly. So the sum we want is at most $o(1)+O(K^{1/2}\varepsilon)$ as $\sigma'\to 1+$. Since $\varepsilon$ is arbitrary, we are done.
   \end{proof}

For $(\sigma-1)/2\leq \sigma'-1\leq \sigma-1$,  by Lemma~\ref{lem68elliott},
    \begin{equation}\nonumber
   \begin{split}
    &\qquad L\left(\frac{1}{\sigma'-1}\right)/L\left(\frac{1}{\sigma-1}\right)
    \\&=\exp\left(i\text{Im}\left(\sum_{\mf p \in \mf I(K), \NN(\mf p)\geq 3}(\NN(\mf p)^{-\sigma'}-\NN(\mf p)^{-\sigma})(1-g(\mf p))\right)\right)=\exp(o(1))\to 1
   \end{split}
   \end{equation}
 %  Since
%    \begin{equation}\nonumber
 %  \begin{split}
%    &\qquad \Bigl\vert\sum_{p \in \Z[i]_+, \NN(p)\geq 100}\{\NN(p)^{-\sigma'}-\NN(p)^{-\sigma}\}\{1-g(p)\}\Bigr\vert
%    \\&\leq 2\sum_{p \in \Z[i]_+, \NN(p)\geq 100}\NN(p)^{-(\sigma+1)/2}(1-\NN(p)^{-(\sigma-1)/2})=o(1) 
  % \end{split}
   %\end{equation}
   as $\sigma\to 1+$ (note that we used that $\text{Im  }(1-g(\mf p))=-\text{Im }g(\mf p)$).
  Thus, $L$ is a slowing varying function. 
  
  Now let
  $$C:=Q(1)\exp\left(-\sum_{\mf p \in \mf I(K), \NN(\mf p)\geq 3}\NN(\mf p)^{-1}(1-\text{Re}(g(\mf p)))\right).$$
  By Lemma~\ref{lem68elliott} again,  we have
    \begin{equation}\label{ffew1}
   \begin{split}
    &\qquad \frac{G(s)}{L\left(\frac{1}{\sigma-1}\right)\zeta(s)}
    =Q(s)\exp\left(-\sum_{\mf p \in \mf I(K), \NN(\mf p)\geq 3} \left(\NN(\mf p)^{-s}(1-g( \mf p))+i\text{Im}\NN(\mf p)^{-\sigma}g(\mf p)\right)\right)
    \\&=Q(s)\exp\left(-\sum_{\mf p \in \mf I(K), \NN(\mf p)\geq 3}((\NN(\mf p)^{-s}-\NN(\mf p)^{-\sigma})(1-g(\mf p))+\NN(\mf p)^{-\sigma}(1-\text{Re}(g(\mf p))))\right)
  %  \\&=Q(s)\exp\left(-\sum_{\mf p \in \mf I(K), \NN(\mf p)\geq 3}(\{\NN(\mf p)^{-s}-\NN(\mf p)^{-\sigma}\}\{1-g(\mf p)\}+\NN(\mf p)^{-\sigma}\{1-\text{Re}(g(\mf p))\})\right)
    \\&=(C+o(1))\exp(o(1))\to C
   \end{split}
   \end{equation}
   as $\sigma\to 1+$ uniformly for $\vert\tau\vert\leq M_{1}(\sigma-1)$. Let $R_1$ be the residue of the Dedekind $\zeta$ function at the pole $s=1$ (see for example \cite[Theorem~VII.7.6]{neukirch99}). 
   Then
 \begin{equation}\label{ffew2}
   \begin{split}
   &\qquad \left| G(s)-\frac{C\cdot R_1\cdot L\left(\frac{1}{\sigma-1}\right)}{s-1} \right| = \left| \frac{1}{s-1}\right|\left|\frac{G(s)}{L\left(\frac{1}{\sigma-1}\right)\zeta(s)} \cdot \zeta(s)(s-1) -CR_1\right| 
    \\&\qquad\qquad\qquad\qquad\qquad\qquad\;\;\quad\leq 
   \frac{1}{\sigma-1}\left|\frac{G(s)}{L\left(\frac{1}{\sigma-1}\right)\zeta(s)} \cdot \zeta(s)(s-1) -CR_1\right|.
   \end{split}
   \end{equation}
   Combining (\ref{ffew1}) with (\ref{ffew2}), we have that 
   \begin{equation}\label{ffew3}
   \begin{split}
  G(s)-\frac{CR_1L\left(\frac{1}{\sigma-1}\right)}{s-1}=O\left(\frac{1}{K^{1/2}(\sigma-1)}\right)
   \end{split}
   \end{equation}
   uniformly for $\vert\tau\vert\leq M_{1}(\sigma-1)$ if $\sigma$ is sufficiently close to 1. 

   Next we need to establish (\ref{ffew3}) for $M_{1}(\sigma-1)<\vert\tau\vert\leq M_{2}$ for some  $M_{2}>M_{1}(\sigma-1)$. We need the following estimate:

   \begin{lemma}\label{lem69elliott}
    Let $M_{2}>M_{1}(\sigma-1)$. If \eqref{aiso} holds, then
   \[
   |G(s)| \leq O(M_{1}^{-1/2}(\sigma-1)^{-1})
   \]
   %where the implicit absolute constant depends only on $g$ and the field $K$ and the inequality is valid 
   uniformly for $M_{1}(\sigma-1) \leq |\tau| \leq M_{2}$, for all $\sigma$ sufficiently near to $1, \sigma>1$.  
   \end{lemma}
   \begin{proof}
   For each prime ideal $\mf p$ and $\tau \in \R$, since %it is (as a simple application of the triangle inequality)
   \begin{equation*}
       \begin{split}
         &\qquad  2(1-\text{Re } \NN(\mf p)^{i\tau})=|1-\NN(\mf p)^{i\tau}|^2 \leq 2|1-g(\mf p)|^2+2|g(\mf p)-\NN(\mf p)^{i\tau}|^2 
         \\&\leq 2|1-g(\mf p)|^2+4(1-\text{Re }\NN(\mf p)^{-i\tau}g(\mf p)),
       \end{split}
   \end{equation*}
  we have
   \[
   2\sum_{\mf p} \NN(\mf p)^{-\sigma}(1-\text{Re } \NN(\mf p)^{i\tau}) \leq 2\sum_{\mf p} \NN(\mf p)^{-\sigma}|1-g(\mf p)|^2+4\sum_{\mf p} \NN(\mf p)^{-\sigma} (1-\text{Re }\NN(\mf p)^{-i\tau}g(\mf p)).
   \]
   By Lemma~\ref{eulerproduct1}, we have
   \[
   \exp\left( \text{Re } \sum_{\mf p} \NN(\mf p)^{-\sigma} (1-\NN(\mf p)^{i\tau})\right) \geq C_1 \frac{\zeta(\sigma)}{|\zeta(\sigma-i\tau)|},
   \]
   for some universal constant $C_1>0$, and also
   \[
    \exp\left( \text{Re } \sum_{\mf p} \NN(\mf p)^{-\sigma} (1-\NN(\mf p)^{-i\tau}g(\mf p))\right) \leq C_2 \frac{\zeta(\sigma)}{|G(\sigma+i\tau)|},
   \]
    for some universal constant $C_2>0$. On the other hand, we have that
   \[
   \sum_{\mf p} \NN(\mf p)^{-\sigma} |1-g(\mf p)|^2 \leq \sum_{\mf p} \NN(\mf p)^{-1} 2(1-\text{Re }g(\mf p)) \leq C_3
   \]
   for some constant $C_3>0$. Thus,
 \begin{equation*}
       \begin{split}
         &\qquad |G(s)| \leq \frac{C_{2}\zeta(\sigma)}{\exp\left(\text{Re } \sum_{\mf p} \NN(\mf p)^{-\sigma}(1-\NN(\mf p)^{-i\tau}g(\mf p))\right)}
         \\&\leq \frac{C_{2}\zeta(\sigma)\exp\left( \frac{1}{2} \text{Re } \sum_{\mf p}\NN(\mf p)^{-\sigma}|1-g(\mf p)|^2\right)}{\exp\left(\frac{1}{2}\text{Re } \sum_{\mf p}\NN(\mf p)^{-\sigma} (1-\text{Re } \NN(\mf p)^{i\tau})\right)} \leq  C_{1}^{-1/2}C_{2}\exp(C_{3}/2)(\zeta(\sigma)|\zeta(\sigma-i\tau)|)^{1/2}.
       \end{split}
   \end{equation*}


Finally, since
$$\left|\frac{1}{\sigma-i\tau-1}\right| = \frac{1}{\sqrt{(\sigma-1)^2+M_{1}^2(\sigma-1)^2}} \leq \frac{1}{M_{1}(\sigma-1)},$$
the  desired bound follows from the  fact that $\zeta(s)$ has a simple pole at $s=1$, but is otherwise analytic in the disc $|s-1|<2M_{2}$.
   %
%   \[
%   |G(s)| \leq \frac{\zeta(\sigma)}{\exp\left(\text{Re } \sum_{\mf p} \NN(\mf p)^{-\sigma}(1-\NN(\mf p)^{-i\tau}g(\mf p))\right)}\leq \frac{\zeta(\sigma)\exp\left( \frac{1}{2} \text{Re } \sum_{\mf p}\NN(\mf p)^{-\sigma}|1-g(\mf p)|^2\right)}{\exp\left(\frac{1}{2}\text{Re } \sum_{\mf p}\NN(\mf p)^{-\sigma} (1-\text{Re } \NN(\mf p)^{i\tau})\right)} \leq\]
%   (where we used the second inequality with $\exp$), and this is no bigger than
%   \[
%   C_4(\zeta(\sigma)|\zeta(\sigma-i\tau)|)^{1/2},
%   \]
%   where we used the second $\exp$ inequality, together with the bound with $C_4$ we found above.
%
\iffalse
   To obtain the desired bound, we finally used the fact that $\zeta(s)$ has a simple pole at $s=1$, but is otherwise analytic in the disc $|s-1|<2M$, completing the proof. Note also that $\left|\frac{1}{\sigma-i\tau-1}\right| = \frac{1}{\sqrt{(\sigma-1)^2+K^2(\sigma-1)^2}} \leq \frac{1}{K(\sigma-1)}$, whence the claim (putting together a similar estimate, since $\zeta(\sigma) \leq \frac{1}{\sigma-1}$).
   \fi
   \end{proof}

 Note that when $M_{1}(\sigma-1)\leq \vert\tau\vert\leq M_{2}$, we have $\frac{1}{s-1}=O\left(\frac{1}{M_{1}(\sigma-1)}\right)$. So  Lemma~\ref{lem69elliott} implies that  (\ref{ffew3}) also holds
   uniformly for $M_{1}(\sigma-1)<\vert\tau\vert\leq M_{2}$ if $\sigma$ is sufficiently close to 1.
   %
 %  It follows from the discussions above that conveniently split the range of $\tau$ that:
So we have
   $$\limsup_{\sigma\to 1+}\sup_{\vert\tau\vert\leq M_{2}}(\sigma-1)\left| G(s)-\frac{C\cdot R_1\cdot L\left(\frac{1}{\sigma-1}\right)}{s-1}\right|\leq O(M_{1}^{-1/2}).$$
 %  for some absolute constant $C_{1}$ depedning only on $g$ and $K$.
 Setting $M_{2}=2M_{1}$ and then letting $M_{1}\to\infty$,
 it  follows from Theorem~\ref{C21}  that $\E_{\NN(\mf u)\leq x}g(\mf u)$ converges to $CL(\log x)$. %\WS{why can we apply Theorem 4.1 here? (how can you let $K$ go to infinity?)} 
% \FM{$M$ is a number we have freedom to fix and will send to infinity, so essentially we can choose $K$ as large as necessary.}
 Since we  assume that $\alpha=0$, it suffices to show that
   $$CL(\log x)=\prod_{\NN(\mf p)\leq x}(1-1/\NN(\mf p))(1+h_{1}(\mf p))+o_{x\to\infty}(1).$$
  % as $x\to\infty$.

 From now on we set  
   $\sigma:=1+(\log x)^{-1}$. By (\ref{ffew1}) and (\ref{e41}),
 \begin{equation}\label{CL}
   \begin{split}
  CL(\log x)=\frac{G(\sigma)}{\zeta(\sigma)}+o_{x\to\infty}(1)=Q(\sigma)\exp\left(-\sum_{\NN(\mf p)\geq 3}\NN(\mf p)^{-\sigma}(1-g(\mf p))\right)+o_{x\to\infty}(1).
   \end{split}
   \end{equation}
   Let $0<\varepsilon\leq 1$ to be chosen later. By Cauchy-Schwarz, Part (iv) of Lemma \ref{l2} and the assumption that  $\D(g,1)<\infty$, we have
   \begin{equation}\nonumber
   \begin{split}
    & \Bigl\vert\sum_{\NN(\mf p)\geq x^{\varepsilon}}\NN(\mf p)^{-\sigma}(1-g(\mf p))\Bigr\vert
     \leq \left(\sum_{\NN(\mf p)\geq x^{\varepsilon}}\NN(\mf p)^{-\sigma}\right)^{1/2}\left(\sum_{\NN(\mf p)\geq x^{\epsilon}}\NN(\mf p)^{-\sigma}\left| 1-g(\mf p)\right|^{2}\right)^{1/2}
    \\&\ll \Bigl(\sum_{\NN(\mf p)\geq x^{\varepsilon}}\NN(\mf p)^{-\sigma}\vert 1-g(\mf p)\vert^{2}\Bigr)^{1/2}  
     \leq \Bigl(2\sum_{\NN(\mf p)\geq x^{\varepsilon}}\NN(\mf p)^{-1}(1-\text{Re}(g(\mf p)))\Bigr)^{1/2}=o_{x\to\infty}(1).
   \end{split}
   \end{equation}
  % Therefore, 
 %  $$\sum_{\NN(\mf p)\geq x^{\varepsilon}}\NN(\mf p)^{-\sigma}(1-g(\mf p))=o(1).$$
 %  Similarly but more simply, we have that 
  %  \begin{equation}\nonumber
 %  \begin{split}
 %  \sum_{x^{\varepsilon}\leq \NN(\mf p)\leq x}\NN(\mf p)^{-1}(1-g(\mf p))=o(1).
 %  \end{split}
 %  \end{equation}
 %   Indeed, since $|1-z|^2\leq 2(1-\text{Re}(z))$ for $|z|\leq 1,$ we have 
 Similarly, by Part (iii) of Lemma \ref{l2}, %\WS{I commented the old proof and give a new proof in Part (iii) of Lemma \ref{l2}}
    \begin{equation}\nonumber
   \begin{split}
    & \Bigl\vert\sum_{x^{\varepsilon}\leq \NN(\mf p)\leq x}\NN(\mf p)^{-1}(1-g(\mf p))\Bigr\vert
    \leq \left(\sum_{x^{\varepsilon}\leq \NN(\mf p)\leq x}\NN(\mf p)^{-1}\right)^{1/2}\left(\sum_{x^{\epsilon}\leq \NN(\mf p)\leq x}\NN(\mf p)^{-1}\vert 1-g(\mf p)\vert^{2}\right)^{1/2}
    \\&\ll \Bigl(\sum_{x^{\varepsilon}\leq \NN(\mf p)\leq x}\NN(\mf p)^{-1}\vert 1-g(\mf p)\vert^{2} \Bigr)^{1/2}
      \leq \Bigl(2\sum_{\NN(\mf p)\geq x^{\varepsilon}}\NN(\mf p)^{-1}(1-\text{Re}(g(\mf p)))\Bigr)^{1/2}=o_{x\to\infty}(1).
   \end{split}
   \end{equation}
   \iffalse
   where the second inequality is because
   we can find an absolute constant $K_1$ such that
    \[\sum_{x^{\varepsilon}\leq \NN(\mf p)\leq x}\NN(\mf p)^{-1}
    \leq K_1 \sum_{x^{\varepsilon} \leq p \leq x} \frac{1}{p} = K_1 \int_{x^{\varepsilon}}^x y^{-1} d\pi(y)=K_1[y^{-1}\pi(y)]_{x^{\varepsilon}}^x+\int_{x^{\varepsilon}}^x y^{-2}\pi(y),
    \]
    where $\pi(y)=\sum_{p \leq y} 1$. We now use the fact that $\pi(y)=O\left(\frac{y}{\log y}\right)$ to obtain the upper bound
    \[
    K_1\cdot C_0 \left( \frac{1}{\log x}+\frac{1}{\varepsilon \log x} \right)+C_0 \int_{x^{\varepsilon}}^x \frac{1}{y\log y} \ dy = K_1C_0'\frac{1}{\log x}+C_0(-\log \varepsilon)=O(1).
    \]
    \fi
   Therefore, we have that
  \begin{equation}\label{cd1}
   \begin{split}
    \Bigl\vert\sum_{\NN(\mf p)\leq x}\NN(\mf p)^{-1}(1-g(\mf p))\Bigr\vert=o_{x\to\infty}(1).
   \end{split}
   \end{equation}
   
   Finally, by  the mean value inequality and Part (v) of Lemma \ref{l2},
    \begin{equation}\label{cd2}
   \begin{split}
   &\qquad\sum_{\NN(\mf p)\leq x^{\varepsilon}}(\NN(\mf p)^{-1}-\NN(\mf p)^{-\sigma})\vert 1-g(\mf p)\vert
   \\&\leq 2\sum_{\NN(\mf p)\leq x^{\varepsilon}}\NN(\mf p)^{-1}(1-\exp(-\log \NN(\mf p)/\log x))
  \leq \frac{2}{\log x}\sum_{\NN(\mf p)\leq x^{\varepsilon}}\frac{\log\NN(\mf p)}{\NN(\mf p)} %=\frac{2}{\log x}K_1 \sum_{p<x^{\varepsilon}} \frac{\log p}{p}
  % \leq C_{11}\varepsilon
  =O(\varepsilon).
   \end{split}
   \end{equation}
  % where $K_1$ is a constant that only depends on the degree of the number field, so using the same estimate as in \cite[p. 247]{E79book} we obtain the last inequality with an absolute constant $C_{11}$ for all sufficiently large $x$. To pass from the second to the third line, we also used the mean value inequality.
 % \WS{I merged your argument to Lemma \ref{l2}, and fixed a gap where you forgot the consider the case when $\NN(\mf p)$ is a power of a prime instead of a prime}
%   Putting everything together and using the triangle inequality
Combining (\ref{cd1}) and (\ref{cd2}), we have
\begin{equation*}
    \begin{split}
        \qquad & \left| \sum_{\mf p}\NN(\mf p)^{-\sigma}(1-g(\mf p))-\sum_{\NN(\mf p)  \leq  x}\NN(\mf p)^{-1}(1-g(\mf p)) \right| 
        \\&\leq \sum_{\NN(\mf p)\leq x^{\varepsilon}}(\NN(\mf p)^{-1}-\NN(\mf p)^{-\sigma})\vert 1-g(\mf p)\vert
          +\left| \sum_{x^{\varepsilon}\leq \NN(\mf p) \leq x} \NN(\mf p)^{-1}(1-g(\mf p)) \right| 
        \leq O(\varepsilon).
    \end{split}
\end{equation*}
%   we have that 
%   $$\limsup_{x\to\infty}\left|\sum_{\mf p}\NN(\mf p)^{-\sigma}(1-g(\mf p))-\sum_{\NN(\mf p)\leq x}\NN(\mf p)^{-1}(1-g\mf (p))\right|\leq O(\varepsilon).$$
%  So we get the claim when we take $\limsup_{x \to \infty}$ on both sides.
   Since $\varepsilon$ is arbitrary, we have 
   $$\sum_{\mf p}\NN(\mf p)^{-\sigma}(1-g(\mf p))=\sum_{\NN(\mf p)\leq x}\NN(\mf p)^{-1}(1-g(\mf p))+o_{x\to\infty}(1).$$
   Combining this with (\ref{CL}), together with the fact that $Q(1)$ can be written as a quotient of two converging infinite products, we are done.
\end{proof}
 


\iffalse

We need to check the proof in the book, but this is likely to be true. Although we do need some inputs for the Riemann Zeta function for Gaussian integers. We need:

(1) $\sum_{u\in\Z[i]_{+}}\NN(u)^{-s}=O(\frac{1}{s-1})$ as $\sigma\to 1^+$ (uniformly for $-M\leq \tau\leq M$)
{\color{red}
We have the following inequality:
\[
\left|\sum_{u \in \Z[i]\setminus{\{0\}}} \NN(u)^{-s}\right| \leq \sum_{n,m \in \Z \setminus \{0\}} \frac{1}{(n^2+m^2)^{\sigma}} \leq \sum_{n,m \in \Z \setminus \{0\}} \frac{1}{(2\sqrt{n^2m^2})^{\sigma}}=\frac{4}{2^{\sigma}}\left(\sum_{n=1}^{\infty}\frac{1}{n^{\sigma}}\right)^2,
\]
where we used the AM-GM inequality in the third step. Using the integral test on the last term of the inequality above, we get an upper bound of
\[
\frac{4}{2^{\sigma}}\frac{1}{(1-\sigma)^2}.
\]
%Note: this argument could be generalized to get an upper bound of the form $O((\sigma-1)^d)$, where $d=[K:\Q],$ the degree of the extension of the number field, whose ring of integers $\mathcal{O}_K$ we would be interested in.

%(Question: is this inequality good enough for our purposes?)

Improvement using [Murty, van Order: Counting integral ideals in a number field]:
\[
\left|\sum_{u \in \Z[i]\setminus{\{0\}}} \NN(u)^{-s}\right| \leq
\sum_{n=1}^{\infty} \sum_{u \in \Z[i]_+ : \NN(u)=n} |N(u)|^{-s} = \sum_{n=1}^{\infty} |n^{-s}| \sum_{u \in \Z[i]_+ : \NN(u) = n} 1 =
\]
\[\sum_{n=1}^{\infty} \sum_{u \in \Z[i]_+ : \NN(u)=n} |N(u)|^{-s} \leq \sum_{n=1}^{\infty} |n^{-s}| \left(\sum_{u \in \Z[i]_+ : \NN(u) \leq n+1} 1-\sum_{u \in \Z[i]_+ : \NN(u) \leq n-1} 1  \right) \leq \sum_{n=1}^{\infty} n^{-\sigma} K,
\]
where $K>0$ is an absolute constant that depends on the degree of the extension. Now, by the integral test, the last sum is of order $O\left( \frac{1}{\sigma-1}\right)$.


}
(2) An upper bound for $\sum_{\NN(p)>x^{\e}}\NN(p)^{-(1+(\log x)^{-1})}=O(1)$ (this is the last equation of page 246 of the book, where they used integration by part. However, I think this part is fine for us as the range $\NN(p)>x^{\e}$ is contained in the union of two rectangles and so we can still do the 2-dim integration by part).
{\color{red}
It is known that primes in $\Z[i]$ have norm equal to either $p$ or $p^2$. Therefore, since $\sigma=1+(\log x)^{-1}>1$, if we argue as in (1) above (that is, we use the linear estimate bound on the number of elements with a given norm, see Lemma \ref{mv7}, we have the upper bound
\[
\sum_{\NN(p)>x^{\e}}\NN(p)^{-(1+(\log x)^{-1})} \leq \sum_{p>x^{\varepsilon}}\left( \frac{1}{p^{\sigma}}+\frac{1}{p^{2\sigma}}\right)\sum_{u \in \Z[i]_+ : \NN(u) \in \{p,p^2\}} 1\leq 2K \sum_{p>x^{\varepsilon}}\frac{1}{p^{\sigma}},
\]
where $K>0$ is a universal constant that gives us the number of elements in $\Z[i]$ whose norm is either $p$ or $p^2$. Finally, using the same argument as in the bottom of \cite[p. 246]{E79book}, we obtain
\[
\sum_{\NN(p)>x^{\e}}\NN(p)^{-(1+(\log x)^{-1})}=O(1),
\]
as desired.

}

\fi

\

%Step 3: Show the following version of Theorem 6.3 of Elloit79: 

%\begin{conjecture}\label{c27}
 %    Let $g\colon\Z[i]\to\C$ be a multiplicative function with $\vert g\vert\leq 1$. 
  %   \begin{enumerate}
   %      \item If $\D(g,1)<\infty$ and if for all prime $p$ with $\NN(p)=2$, we have $g(p^{k})\neq -1$ for some $k\geq 1$, {\color{red} i.e. $1+h_{1+i\tau}(p)\neq 0$ for all $p$}  then $\E_{u\in\Z[i]_{+}}g(u)$ exists and but is not equal to 0.
    %     \item If $\D(g,\NN(u)^{i\tau})<\infty$ for some $\tau\neq 0$ and if for all prime $p$ with $\NN(p)=2$, {\color{red} i.e. $1+h_{1+i\tau}(p)\neq 0$ for all $p$} we have $g(p^{k})\neq -2^{ik\tau}$ for some $k\geq 1$, then $\E_{u\in\Z[i]_{+}}g(u)$ does not exist.
     %    \item If $\D(g,\NN(u)^{i\tau})<\infty$ for some $\tau\in\R$ and if for some prime $p$ with $\NN(p)=2$, we have $g(p^{k})=-2^{ik\tau}$ for all $k\geq 1$, {\color{red} i.e. $1+h_{1+i\tau}(p)=0$ for some $p$} then $\E_{u\in\Z[i]_{+}}g(u)$ exists and equals to 0.
      %   \item If $\inf_{\tau\in\R}\D(g,\NN(u)^{i\tau})=\infty$, then $\E_{u\in\Z[i]_{+}}g(u)$ exists and equals to 0.
     %\end{enumerate}
%\end{conjecture}

%We need to check the proof in the book, but this part is easy and it follows from the previous result. Similar to Lemma 6.6 from the book, we should be able to get something like
%$$G(s)=\left(\prod_{\NN(p)<100}(1+h_{s}(p))\right)\exp\left(\sum_{\NN(p)\geq 100}g(p)\NN(p)^{-s}\right)G_{1}(s).$$
%In the integer case, the term $1+h_{s}(p)$ could be zero only if $p=2$. In our case, the term $1+h_{s}(p)$  could be zero only if $\NN(p)=2$.  


% {\color{blue} now everything between the two blue sentences can be deleted}

% Once we finish these steps, we can move the averages along AP.


%   I did not find the proof of the classification for averages along AP, but the following is probably an approach (for the integer case):

% Let $A(\chi):=\E_{u\in\Z} g(u)\chi(u)$ and $\mathcal{B}_{q}$ denote the set of all Dirichlet charater of period $q$. Then {\color{red} need to verify. The function $\mathds{1}_{n \equiv r \mod q}(n)$ can be identified with the function $f:(\Z/q\Z)^{\times} \to \C$ given by $f(x)=\mathds{1}_{\{r\}}(x)$ by extending it to $0$ on the $n$ that are not coprime with $q$. The Dirichlet characters are the characters of $(\Z/q\Z)^{\times}$ extended in the same fashion to all of $\Z$, so by Fourier analysis on finite abelian groups, $f$ is a finite linear combinations of Dirichlet characters. Since they are extended in the same way when $n$ and $q$ are not coprime, the claim follows.} the function $\mathds{1}_{n\equiv r \mod q}(n)$ is a linear combination of elements in $\mathcal{B}_{q}$ for all $(r,q)=1$. Conversely, obviously every element of $\mathcal{B}_{q}$ is a linear combination of $\mathds{1}_{n\equiv r \mod q}(n), (r,q)=1$. So $A(\mathds{1}_{n\equiv r \mod q})=0$ for all $(r,q)=1$ if and only if $A(\chi)=0$ for all $\chi\in \mathcal{B}_{q}$. By the previous result, $g$ is aperiodic if and only if for all $\chi$, either $\inf_{\tau\in\R}\D(g\chi,u^{i\tau})=\inf_{\tau\in\R}\D(g,\overline{\chi}u^{i\tau})=\infty$ or there exists $\tau_{\chi}\in\R$ with $\D(g\chi,u^{i\tau_{\chi}})<\infty$ and $g(2^{k})\chi(2)^{k}=-2^{i k\tau_{\chi}}$ for all $k\geq 1$. {\color{red} If we can exclude the second case (think how to do it)}, then we can deduce that $g$ is aperiodic if and only if $\D(g,\overline{\chi}u^{i\tau})=\infty$ for all Dirichlet character $\chi$ and $\tau\in\R$.

% For Gaussian integers, we want something like this. Let $\mathcal{B}_{a,b}$ be the set of all 
% completely multiplicative function $\chi\colon\Z[i]\to\C$ with $\chi((m+ka)+(n+\ell b)i)=\chi(m+ni)$ for all $k,\ell,m,n\in\Z$ and with $\chi(m+ni)\neq 0 \Leftrightarrow (m,a)=(n,b)=1$ {\color{red} maybe we need to consider a larger class of functions where $\chi(m+ni)\neq 0 \Leftrightarrow \min\{(m,a),(n,b)\}=1$}. The above argument works if we can show the following: {\color{red} for all $m,n$ with $(m,a)=(n,b)=1$, the function $\mathds{1}_{x\equiv m \mod a, y\equiv n \mod b}(x+yi)$ is a linear combination of functions in $\mathcal{B}_{a,b}$.}

% %Let $\chi: \Z[i] \to \C$ be a character. Then we say, that $\chi$ is periodic if $\chi(\vec{n}+\vec{a})=\chi(\vec{n})$ for all $\vec{n} \in \Z[i]$ and for a set of $\vec{a} \in \Z[i]$ that forms an additive subgroup $J \subset \Z[i]$ such that $|\Z[i]:J|<\infty$



%  %The case for $\Z[i]$ seems to be more tricky. For example, for all $m,n$, can we write $\mathds{1}_{x\equiv m \mod 4, y\equiv n \mod 4}(x+yi)$ as a linear combination of multiplicative functions of period 4? One of the many issues is that there is no completely multiplicative function of period 4 with $f(1+i)\neq 0$. This is because $(1+i)^{4}\in 4\Z[i]$.

% %An even simpler example: can we write $\mathds{1}_{x\equiv m \mod 2, y\equiv n \mod 2}(x+yi)$ as a linear combination of multiplicative functions of period 2? One of the many issues is that there is no completely multiplicative function of period 4 with $f(1+i)\neq 0$. This is because $(1+i)^{2}\in 2\Z[i]$.

% %Your subgroup is $\langle (4,4)\rangle \subset \Z^2$

% %I think the following should be true:  $f(n,m)=\sum_i a_i\chi_i(n) \sum_j b_j \chi_j(m)$ by lifting $(\Z/4\Z)^{\times} \times (\Z/4\Z)^{\times} \to \Z[i]$




% For $v_{1},v_{2}\in\Z[i]$ with $(v_{1}\Z+v_{2}\Z)(v_{1}\Z+v_{2}\Z)\subseteq v_{1}\Z+v_{2}\Z$, let $\Z[i]/(v_{1}\Z+v_{2}\Z)$ be the set of all elements in $\Z[i]$ modulo $v_{1}\Z+v_{2}\Z$. Let $(\Z[i]/(v_{1}\Z+v_{2}\Z))^{\times}$ be the set of all $u\in \Z[i]$ (modulo $v_{1}\Z+v_{2}\Z$) such that $uu'-1\in v_{1}\Z+v_{2}\Z$ for some $u'\in\Z[i]$. Let $\mathcal{B}_{v_{1},v_{2}}$ denote the set of all completely multiplicative function $\chi\colon \Z[i]\to\C$ with $\chi(\cdot+v_{1}\Z+v_{2}\Z)=\chi$ and with $\chi(x)=0$ if and only if $x\mod (v_{1}\Z+v_{2}\Z)\in (\Z[i]/(v_{1}\Z+v_{2}\Z))^{\times}$. %Now can we show that the indicator of any 2d arithmetic progression can be written as linear combinations of $\cup_{v_{1},v_{2}}\mathcal{B}_{v_{1},v_{2}}$.


% %Consider the following example, we want to write $\mathds{1}_{(1+i)+4\Z+4i\Z}$ as a linear combination of periodic multiplicative functions. The issue is that $1+i$ is not an invertible element in $4\Z+4i\Z$. However, for any completely multiplicative function, we have that 
% %$$\mathbb{E}_{m,n}\chi((1+i)+(4m+4ni))=\mathbb{E}_{m,n}\chi(1+i)\chi(1+(2-2i)m+(2+2i)n).$$
% %But 1 is an invertible element in $(2-2i)\Z+(2+2i)\Z$. Also note that $((2-2i)\Z+(2+2i)\Z)^{2}\subseteq (2-2i)\Z+(2+2i)\Z$.

% %One can compute that the set of all invertible elements in $(2-2i)\Z+(2+2i)\Z$ are  



% %In this case, $\mathds{1}_{1+(2-2i)\Z+(2+2i)\Z}$ is a multiplicative function.
% %So 
% %$$\mathbb{E}_{m,n}\chi((1+i)+(4m+4ni))$$


%  I reformulate the issue and solution as follows: suppose we wish to study 
% $$\mathbb{E}_{m,n}\chi((1+i)+(4m+4ni))=\frac{1}{16}\mathbb{E}_{u}\mathds{1}_{(1+i)+4\Z+4i\Z}(u)\chi(u).$$
% Note that $1+i$ does not belong to $(\Z[i]/(4\Z+4i\Z))^{\times}$ since $(1+i)^{4}\equiv 0 \mod 4\Z+4i\Z$. So we cannot write $\mathds{1}_{(1+i)+4\Z+4i\Z}$ as linear combintions of elements in $\mathcal{B}_{4,4i}$. However, note that 
% $$\mathbb{E}_{m,n}\chi((1+i)+(4m+4ni))=\mathbb{E}_{m,n}\chi(1+i)\chi(1+(2-2i)m+(2+2i)n)=\frac{1}{C}\chi(1+i)\mathbb{E}_{u}\mathds{1}_{1+(2-2i)\Z+(2+2i)\Z}(u)\chi(u).$$
% One can compute that $(\Z[i]/(2-2i)\Z+(2+2i)\Z))^{\times}$ is a ciclic group or order 4 (isomorphic to $(\Z/5\Z)^{\times}$), with $a=2+i$ being its generator ($a=2+i$, $a^{2}=3$, $a^{3}=2-i$ and $a^{4}=1$ are the only invertible elements). So $\mathds{1}_{1+(2-2i)\Z+(2+2i)\Z}(u)$ for sure can be written as a linear combination of elements in $\mathcal{B}_{2-2i,2+2i}$ (since they can be viewed as Dirichlet characters of order 5 on integers).




% {\color{blue} now everything between the two blue sentences can be deleted}


\iffalse
%\section{Concentration estimates}\label{SECconcest}
In this section we will obtain some linear concentration estimates, which will allow us to easily treat the pretentious part of the multiplicative functions. In particular, we are going to obtain an analog of Proposition 2.5 of [FKM]. To do so, we are going to prove a suitable form of Tur\'an-Kubilius for our purposes. 

We begin with a short counting lemma:
\begin{lemma}\label{countingCRT}
Let $I$ be an ideal in $\Z[i]$, let $M \in \N$ be large, and suppose that $M$ is large enough so that $I$ divides all elements of $\Phi_{M}$. Consider natural numbers $k, l \in \N$, let $Q:=\prod_{\NN(p)\leq M} p$, and $p, q$ be distinct primes such that $p, q \nmid Q$. We have
\begin{equation}\label{countingCRTeq1}
\lim_{N\to \infty} \frac{| \{u \in Q[\pm N]^2+a : Q \mid u-a, q^l \mid u \text{ and } p^k \mid u\}}{\NN(Q)(2N)^2} = \frac{1}{\NN(Qp^k) \NN(q^l)}.
\end{equation}
\end{lemma}
\begin{proof}
Let $p, q, l, k$ and $a, Q$ be as in the statement. We begin by rewriting the left hand side of \eqref{countingCRTeq1} as 

\[
\frac{1}{\NN(Q)(2N)^2}\sum_{u \in Q[\pm N]^2+a} \mathds{1}_{\{0\}}(u \pmod{q^l}) \mathds{1}_{\{0\}
}(u \pmod{ p^k}) \mathds{1}_{\{a\}}(u \pmod Q)\]
This is the same as counting all the $u$ in the shifted box $Q[\pm N]^2+a$ such that they have residue class $(a,0,0)$ in $\Z[i]/(Q) \times \Z[i]/(p^k) \times \Z[i]/(q^l)$ (where we take the ideal generated by each element).

By the Chinese Remainder Theorem, the product ring above is isomorphic to $\Z[i]/(Qp^kq^l)$, since $\Z[i]$ is a UFD and the elements are relatively prime by assumption. Moreover, it is known that the size of this latter quotient ring is exactly $\NN(Qp^kq^l)$. Now the proof can easily be finished as a trivial application of the mean ergodic theorem, since the $\Z[i]$-action by addition on the quotient ring is ergodic, the limit displayed above will converge to the integral of the characteristic function of the point in the aforementioned quotient ring $\Z[i]/(Qp^kq^l)$. This trivially has size $\frac{1}{\NN(Qp^kq^l)}$, and we are done.
\end{proof}
With the help of Lemma~\ref{countingCRT}, we can now prove the main result of this section; the following concentration estimate:
\begin{proposition}\label{p25}
    Let $f\colon \Z[i]\to\U$ be a multiplicative function with $f\sim \chi\cdot n^{i\tau}$ for some Dirichlet character $\chi$ with period $I$ {\color{red} which is an ideal} and some $\tau\in\R$. Let 
    $$\Phi_M:=\left\{\prod_{\NN(p)\leq M}p^{a_{p}}\colon M<a_{p}\leq 2M\right\}$$
    and suppose that $K$ is large enough so that $I$ divides all elements of $\Phi_{M}$. Then, letting $q:=\prod_{\NN(p)\leq M} p$, we have 
\begin{equation}
    \begin{split}
        \limsup_{N\to\infty}\max_{Q\in \Phi_{M}}\E_{u\in[\pm N]^{2}}\vert f(Qu+a)-\chi(a)\NN(Qu)^{i\tau}\exp(F_{N}(f,K))\vert\ll \D(f,\chi\cdot \NN(u)^{i\tau};K,\infty)+K^{-1/2}
    \end{split}
\end{equation}
for any $a\in\Z[i]$ with $\NN(a)\leq \NN(Q)$ and $a \mod I\in (\Z[i]/I)^{\times}$,
where  the implicit constant is absolute and
$$F_{N}(f,K):=\sum_{K<\NN(p)\leq 2N^{2}, p \nmid Q}\frac{1}{\NN(p)}(f(p)\cdot\overline{\chi}(p)\cdot \NN(p)^{-i\tau}-1).$$
 {\color{red} Note here $Q$ are integers in $\Z[i]$ not $\Z$, and $2N^{2}$ is the norm of $N+iN$.}
 \end{proposition}
\begin{proof}

    {\color{red} I copied below the proof of Lemma 2.5 of [KMPT]. please check the details.}

    Write $f(u)=f'(u)\chi(u)\NN(u)^{i\tau}$. Then $f'$ is a multiplicaitve function taking values in $\U$. Note that if $I$ divides $Q$, then {\color{blue} I think this is independent of the choice of $I$ and $Q$: we are using $|e^{ix}-e^{iy}| \leq |x-y|$ and then getting a term of the form $\tau\cdot \ln \frac{\NN(Qu+a)}{\NN(Qu)}$ which will go to $0$ as $\NN(u) \to \infty$.}
    $$f(Qu+a)=f'(Qu+a)\chi(a)\NN(Qu+a)^{i\tau}=f'(Qu+a)\chi(a)\NN(Qu)^{i\tau}+O_{\tau}(\NN(u)^{-1}).$$
    Since $\lim_{N\to\infty}\E_{u\in[\pm N]^{2}}\NN(u)^{-1}=0$ {\color{red} check} {\color{blue} Yes, $\NN(u) \to 0$ as $u \to \infty$ in the sense of leaving compact sets, and convergence implies Ces\`aro convergence}, it suffices to show that {\color{blue} Here we are estimating sums of differences of $F_N(f,Q)-F_N(f',Q)$ which can be bounded by the sum of the reciprocals of norms of primes which grow like $\log \log x$, so when taking an exponential, we get something that grows like $\log x$. This dies when dividing by $x$ in the expectation, so this step is correct.}
\begin{equation}
    \begin{split}
        \limsup_{N\to\infty}\max_{Q\in \Phi_{K}}\E_{u\in[\pm N]^{2}}\vert f'(Qu+a)-\exp(F_{N}(f',K))\vert\ll \D(f',1;K,2N^{2})+K^{-1/2},
    \end{split}
\end{equation}
where $$F_{N}(f',K)=\sum_{K<\NN(p)\leq 2N^{2}, p\nmid Q}\frac{1}{\NN(p)}(f'(p)-1).$$
Let $h\colon\Z[i]\to\C$ be the additive function given by $h(p^{k}):=f'(p^{k})-1$ for prime powers. Since $z=e^{z-1}+O(\vert z-1\vert^{2})$, we have
$$f'(Qu+a)=\prod_{p^{k}\Vert Qu+a, \NN(p)>K} f'(p^{k})=\prod_{p^{k}\Vert Qu+a, \NN(p)>K}\exp(h(p^{k}))+O(\vert h(p^{k})\vert^{2}).$$
By the inequality $\left|\prod_{i=1}^{k}z_{i}-\prod_{i=1}^{k}w_{i}\right|\leq\sum_{i=1}^{k}\vert z_{i}-w_{i}\vert$ which can be proved by induction, we have that
$$f'(Qu+a)=\exp(h(Qu+a))+O\left(\sum_{p^{k}\Vert Qu+a, \NN(p)>K}\vert h(p^{k})\vert^{2}\right).$$
Also
note that
\begin{equation}
    \begin{split}
        &\quad F_{N}(f',K)=\sum_{K<\NN(p^{k})\leq 2N^{2}, p \nmid Q}h(p^{k})\NN(p)^{-k}(1-\NN(p)^{-1})+O(1/K):=\mu_{h,N}+O(1/K).
    \end{split}
\end{equation}
 {\color{red} check this step}
 {\color{blue} Idea: split into $k=1$ and $k >1$. The $k=1$ part we compare to the actual definition of $F_N(f',K)$, and it gives us an error bounded above by $\sum_{K<\NN(p)<2N^2} \NN(p)^{-2}\ll 1/K$. On the other hand, the second term is made of primes to the power at least $2$, so simply bounding trivially $h$ and $(1-\NN(p)^{-1})$ by absolute constants, we get a similar error term to the one claimed.}
Since $\vert e^{w}-e^{z}\vert\ll \vert w-z\vert$, we have that
\begin{equation}
    \begin{split}
        &\quad\E_{u\in[\pm N]^{2}}\vert f'(Qu+a)-\exp(F_{N}(f',K))\vert
        \\&\ll \E_{u\in[\pm N]^{2}}\vert \exp(h(Qu+a))-\exp(F_{N}(f',K))\vert+\E_{u\in[\pm N]^{2}}\sum_{p^{k}\Vert Qu+a, \NN(p)>K}\vert h(p^{k})\vert^{2}
        \\&\ll \E_{u\in[0,N]^{2}}\vert h(Qu+a)-F_{N}(f',K)\vert+\E_{u\in[\pm N]^{2}}\sum_{p^{k}\Vert Qu+a, \NN(p)>K}\vert h(p^{k})\vert^{2}
        \\&=\E_{u\in[\pm N]^{2}}\vert h(Qu+a)-\mu_{h,N}\vert+\E_{u\in[\pm N]^{2}}\sum_{p^{k}\Vert Qu+a, \NN(p)>K}\vert h(p^{k})\vert^{2}+O(1/K).
    \end{split}
\end{equation}

For the second term, we have
\begin{equation}
    \begin{split}
        &\quad \E_{u\in[\pm N]^{2}}\sum_{p^{k}\Vert Qu+a, \NN(p)>K}\vert h(p^{k})\vert^{2}
        =\sum_{\NN(p)>K, k\in\N}\vert h(p^{k})\vert^{2}\frac{1}{N^{2}}\sum_{u\in [\pm N]^{2}}\mathds{1}_{p^{k}\Vert Qu+a}
        \\&\ll \sum_{\NN(p)>K, k\in\N}\vert h(p^{k})\vert^{2}\NN(p^{k})^{-1} {\color{red} check}
    \end{split}
\end{equation}
{\color{blue} This is essentially count of the size of an AP with a certain step. The bound is fairly trivial over $\N$. If we move to $\Z[i]$ we simply make use of Lemma \ref{mv7} which estimates the number of elements of certain norms to deduce that the same will hold up to a certain absolute constant. This is not a problem because we are already writing $\ll$.}

For the first term,
we have 
\begin{equation}\label{20}
    \begin{split}
        &\quad \E_{u\in[ \pm N]^{2}}\vert h(Qu+a)-\mu_{h,N}\vert 
        \\&=\E_{u\in[\pm N]^{2}}\vert h(Qu+a)-\sum_{\NN(p^{k})\leq 2N^{2}, p\nmid Q}h(p^{k})\NN(p)^{-k}(1-\NN(p)^{-1})\vert
        \\&\leq \left(\E_{u\in[\pm N]^{2}}\vert h(Qu+a)-\sum_{\NN(p^{k})\leq 2N^{2},p\nmid Q}h(p^{k})\NN(p)^{-k}(1-\NN(p)^{-1})\vert^{2}\right)^{1/2}
        \\&\ll (1+o_{N\to\infty}(1)) \left(\sum_{\NN(p)>K, k\in\N}\vert h(p^{k})\vert^{2}\NN(p^{k})^{-1}\right)^{1/2}
    \end{split}
\end{equation}
%{\color{red} (1): to get the last step, we need a Turan-Kubilius inequality for $\Z[i]$. (2): also I don't see why in [FKM] this step follows from Turan-Kubilius (the one on wikipedia). I see that the last inequality holds if the term $h(Qu+a)$ were $h(u)$, but how did they pass to $Qu+a$? One thing to look at from (5.2) to (5.8) of [FKM2], this is the place where they replace $Qu+a$ by $u$.}




%{\color{red} reference for the theorem: Introduction To Analytic And Probabilistic Number Theory(Tenenbaum), page 319 of the pdf file.}



{\color{red}

In order to fully justify the last inequality of (\ref{20}), we need to show that
\begin{equation}
    \begin{split}
        &\quad \E_{u\in [\pm N]^2}\vert h(Qu+a)-\sum_{\NN(p^{k})\leq N,p\nmid Q}h(p^{k})\NN(p)^{-k}(1-\NN(p)^{-1})\vert^{2}
        \\&\ll (1+o_{N\to\infty}(1)) \sum_{\NN(p)>K, k\in\N}\vert h(p^{k})\vert^{2}\NN(p^{k})^{-1}.
    \end{split}
\end{equation}
In order to do that, as in the proof of the Tur\'an-Kubilius theorem, we let $c_{N,K,Q}:=\sum_{ \NN(p^{k})\leq N,p\nmid Q}h(p^{k})\NN(p)^{-k}(1-\NN(p)^{-1})$. It suffices to show that 
$\E_{u\in [0,N]^2}h(Qu+a)\approx c_{N,K,Q}$ and $\E_{u\in [\pm N]^2}\vert h(Qu+a)\vert^{2}\approx c_{N,K,Q}^{2}$.
First observe that

\begin{equation}
    \begin{split}
        &\quad \E_{u\in [\pm N]^2}h(Qu+a)
        =\frac{1}{(2N)^2}\sum_{u\in Q[\pm N]^2+a, u\equiv a \mod Q}h(u)
        \\&=\frac{1}{(2N)^2}\sum_{u\in Q[\pm N]^2+a, u\equiv a \mod Q}\sum_{p^{k}|| u, p \in \Z[i]_+} h(p^{k})
        \\&=\frac{1}{(2N)^2}\sum_{u\in Q[\pm N]^2+a, u\equiv a \mod Q}\sum_{p^{k}|| u, p\nmid Q, p \in \Z[i]_+} h(p^{k})
        \\&=\frac{1}{(2N)^2}\sum_{p^{k}\in Q[\pm N]^2+a, p\nmid Q, p \in \Z[i]_+}h(p^{k})\sum_{u\in Q[\pm N]^2+a, p^{k}||u, u\equiv a \mod Q} 1
     %   \\&=\frac{1}{N^2}\sum_{p^{k}\in  [0,QN+a]^2, p\nmid Q}h(p^{k})(\sum_{u\in [0,QN+a]^2, p^{k}|u, u\equiv a \mod Q} 1-\sum_{u\in [0,QN+a]^2, p^{k+1}|u, u\equiv a \mod Q} 1)
  %  \\&\approx \frac{1}{N^2}\sum_{K<p^{k}\leq QN+a, p\nmid Q}h(p^{k})(\frac{QN+a}{p^{k}Q}-\frac{QN+a}{p^{k+1}Q})
    \\&\approx \sum_{p^{k}\in Q[\pm N]^2+a, p\nmid Q, p \in \Z[i]_+}h(p^{k})\NN(p)^{-k}(1-\NN(p)^{-1})
    \end{split}
\end{equation}
There are two important remarks to make that justify the previous calculations. First, for the last line, we used Lemma~\ref{countingCRT}. Second, we are using that the factorization into primes of $u$ can be taken so that they all lie in the first quadrant. This is, as we have seen before, because we can take $f$ to be such that $f(i)=1$ if we average on the grid $[\pm N]^2$. This, in turn, implies that $h(\pm i)=h(\pm 1)=0$, which allows us to ensure that all the prime elements in the prime decomposition of each $u$ lie in the first quadrant. 

For the second term, we have

\begin{equation}
    \begin{split}
        &\quad \E_{u\in [\pm N]^2}\vert h(Qu+a)\vert^{2}
        =\frac{1}{(2N)^2}\sum_{u\in  Q[\pm N]^2+a, u\equiv a \mod Q}\vert h(u)\vert^{2}
        \\&=\frac{1}{(2N)^2}\sum_{u\in  Q[\pm N]^2+a, u\equiv a \mod Q}\left| \sum_{p^{k}|| u, p \in \Z[i]_+} h(p^{k})\right|^{2}
        \\&=\frac{1}{(2N)^2}\sum_{u\in  Q[\pm N]^2+a, u\equiv a \mod Q}\left(\sum_{p^{k}|| u, p \in \Z[i]_+} \vert h(p^{k})\vert^{2}+\sum_{p^{k},q^{\ell}|| u, p\neq q, p, q \in \Z[i]_+} h(p^{k})\overline{h}(q^{\ell})\right). 
    \end{split}
\end{equation}
Now, the first term can be dealt with in the same way as we did before (essentially, exchanging the sums and using Lemma~\ref{countingCRT}). This yields an error term $\Theta_1$ of the form
\[
\Theta_1=\sum_{p \in Q[\pm N]^2+a, p \nmid Q, p \in \Z[i]_+} |h(p^k)|^2 \left( \frac{1}{\NN(p^k)}-\frac{1}{\NN(p^{k+1})}\right)
\]
%Note that $\frac{1}{N}\sum_{u\leq QN+a, u\equiv a \mod Q}\sum_{p^{k}|| u} \vert h(p^{k})\vert^{2}$ can be bounded by
As for the remaining main term, we also swap the order of the sums and rewrite it as
%\begin{equation}
 %   \begin{split}
  %      &\quad \frac{1}{(2N)^2}\sum_{p^{k}, q^l \in Q[\pm N]^2+a, p \neq q, p, q \in \Z[i]_+}\vert h(p^{k})\vert^{2}\sum_{u\leq QN+a, u\equiv a \mod Q, p^{k}\Vert u} 1 
   %     \\&\approx \frac{1}{N}\sum_{p^{k}\leq QN+a}\vert h(p^{k})\vert^{2}p^{-k}(1-p^{-1}) 
    %    \\&\leq \frac{1}{N}\sum_{p^{k}\leq QN+a}\vert h(p^{k})\vert^{2}p^{-k}
   % \end{split}
%\end{equation}
%(the RHS of (\cite{20}) is similar to this.)
% On the other hand, we have
\begin{equation}
    \begin{split}
        &\quad \frac{1}{(2N)^2}\sum_{p^{k},q^{\ell}\in Q[\pm N]^2+a, p\neq q, p,q\in \Z[i]_+} h(p^{k})\overline{h}(q^{\ell})\sum_{u\in Q[\pm N]+a, u\equiv a \mod Q, p^{k}||u, q^{\ell}||u}1
        \\&\approx \frac{1}{(2N)^2}\sum_{p^{k},q^{\ell}\in Q[\pm N]^2+a, p\neq q, p,q\in \Z[i]_+} h(p^{k})\overline{h}(q^{\ell})\NN(p^k)\NN(q^l)\left(1-\frac{1}{\NN(p)}\right) \left(1-\frac{1}{\NN(q)}\right)
        \\&= \left|\sum_{p \in Q[\pm N]^2+a, p \in \Z[i]_+} h(p^{k})\left( \frac{1}{\NN(p^k)}-\frac{1}{\NN(p^{k+1})}\right)\right|^2+\Theta_2,
      %  \\&=\vert\sum_{K<p^{k}\leq QN+a, p\nmid Q}h(p^{k})p^{-k}(1-p^{-1})\vert^{2}+error,
    \end{split}
\end{equation}
where in the second step we, once again, used Lemma~\ref{countingCRT}, and $\Theta_2$ is another error term of the form
\[
\Theta_2=\sum_{p \in Q[\pm N]^2+a, p \nmid Q, p \in \Z[i]_+} |h(p^k)|^2 \left( \frac{1}{\NN(p^k)}-\frac{1}{\NN(p^{k+1})}\right)^2
\]
}
 
 So now it remains to control $\sum_{\NN(p)>K, k\in\N}\vert h(p^{k})\vert^{2}\NN(p^{k})^{-1}$. Clearly, 
 $$\sum_{\NN(p)>K, k\geq 2}\vert h(p^{k})\vert^{2}\NN(p^{k})^{-1}\ll K^{-1}.$$
 {\color{red} check} {\color{blue} Yes, since $k\geq 2$, and again using Lemma \ref{mv7}, as well as a trivial constant bound for $|h(p^k)|,$ we see that it is enough to bound
\[
 \sum_{p>K, k \geq 2} \frac{\text{const}}{p^k}=\sum_{p>K} \frac{1}{p^2} \frac{p}{p-1} \ll \sum_{p>K} \frac{1}{p^2} \ll K^{-1},
\]
 as desired.
 }. On the other hand,
 $$\sum_{\NN(p)>K}\vert h(p)\vert^{2}\NN(p)^{-1}=\sum_{\NN(p)>K}2(1-\text{Re}f'(p))\NN(p)^{-1}=2\D(f',1;K,\infty),$$
we are done.
\end{proof}
%\section{Pythagorean pairs for square coefficients}

%{\color{blue} I realized that in FKM, Pythagorean pairs of type 1 do not require concentration estimates. It might be a good idea to try to see if we can prove that part with our current knowledge in $\Z[i]$. This would correspond to Section 4 in [FKM] (not [FKM2]).}
%{\color{red} Yes the next step is to check if Section 4 in [FKM] works for the following sample case (say, take $a=b=1$ and $c=-1$):
In this section, we will obtain the following positivity results. As we saw in Section~\ref{sec: proofstrat}, they will imply the partition regularity results promised in the Introduction.
\begin{equation}\label{e14} 
    \liminf_{N \to \infty} \frac{1}{N^4}\sum_{n, m \in \Z[i] \cap [-N,N]^2} \int_{\mathcal{M}} f(\ell(m^2-n^2))\overline{f(\ell'mn)} \ d\sigma(f)>0.
\end{equation}
\begin{equation}\label{e15} 
    \liminf_{N \to \infty} \frac{1}{N^4}\sum_{n, m \in \Z[i] \cap [-N,N]^2} \int_{\mathcal{M}} f(\ell(m^2+n^2))\overline{f(\ell'mn)} \ d\sigma(f)>0.
\end{equation}
{\color{red} here please feel free to replace $[-N,N]^2$ by any convex set which is convenient.
}

%}
We will make use of the next two simple lemmas:
\begin{lemma}\label{multav0}
Let $f: \Z[i] \to \U$ be a non-trivial multiplicative function (that is, $f \neq 1$). Let $\Phi_M$ be a multiplicative F\o lner sequence in $\Z[i]$. Then,
\[
\lim_{M \to \infty} \E_{u \in \Phi_M} f(u) = 0.
\]
\end{lemma}
\begin{proof}
Since $\U$ is compact, we consider limit points of the sequence $(\E_{u \in \Phi_M}f(u))_{M \in \N}$. Abusing notation, suppose that the limit of this sequence exists. 

Since $f\neq 1$, let $v \in \Z[i]$ be such that $f(v)\neq 1$. Since $\Phi_M$ is a F\o lner sequence, we can write
\[
\lim_{M \to \infty} \E_{u \in \Phi_M} f(u) = \lim_{M \to \infty} \E_{u \in v\Phi_M} f(u) = \lim_{M \to \infty} \E_{u \in \Phi_M} f(vu) = f(v) \lim_{M \to \infty} \E_{u \in \Phi_M} f(u).
\]
Since $f(v) \neq 1$, this implies that the limit in question must be $0$. Therefore, $0$ is the only possible accumulation point for this sequence, so compactness of $\U$ completes the proof.
\end{proof}
\begin{lemma}\label{sequenceavfact}
Let $a: \Z[i] \to \U$ be a sequence. {\color{red} I made a mistake here. In this lemma all the averages should be along boxes not balls, as this is what we need in the next lemma. Does the same argument works for boxes? (If not we need to change the average scheme from the beginning)} {\color{blue} This argument also works for balls, if needed} Let $l_1, l_2 \in \Z[i]$ not both $0$. Suppose that for some $\varepsilon>0$ and some sequence $L_N: \N \to \U$ we have
\[
\limsup_{N \to \infty} \E_{u \in [-N,N]^2} |a(u)-L_N| \leq \varepsilon.
\]
Then,
\[
\limsup_{N \to \infty} \E_{n, m \in [-N,N]^2} |a(l_1m+l_2n)-L_{lN}| \ll \varepsilon,
\]
where $l:=\NN(l_1)+\NN(l_2)+2$ {\color{blue} I added the +2 just in case, but it may not be needed. It's not important anyway}.
\end{lemma}
\begin{proof}
We can estimate
\begin{equation}\label{sequenceavfactE1}
\E_{n,m \in [-N,N]^2} |a(l_1m+l_2n)-L_{lN}| \ll \frac{1}{N^4}\sum_{u \in [-lN,lN]^2} w_N(u) |a(u)-L_{lN}|,
\end{equation}
where for $u \in \Z[i]$ we put $w_N(u):=|\{(m,n) \in [-N,N]^2 : l_1m+l_2n=u\}|$. By the estimate Lemma \ref{mv7} on the number of elements of a certain norm, we see that $w_N(u) \ll N^2$, whence we see that the right hand side of \ref{sequenceavfactE1} is bounded above (disregarding constants) by
\[
\E_{u \in [-lN, lN]^2} |a(u)-L_{lN}|.
\]
\end{proof}
Before moving onto the main proof, we need a couple of preliminary results. Consider the sets on $\S^1$ given by
\[
I(\delta):=\{\exp(2\pi i\phi) : \phi \in (-\delta,\delta)\}.
\]
\begin{lemma}\label{weights}
Let $0<\delta<\frac{1}{2}$ and consider the trapezoidal function $F_{\delta}: \S^1 \to [0,1]$ which is equal to $1$ on $I_{\delta/2}$ and $0$ outside of $I_{\delta}$. Let
\[
w_{\delta}(m,n):= F_{\delta}((\NN(\ell(m^2-n^2))^i\cdot (\NN(\ell'mn))^{-i}), \quad m, n \in \Z[i], 
\]
and
\begin{equation}\label{weight1}
\tilde{w}_{\delta}(m,n):= F_{\delta}((\NN(\ell(m^2+n^2))^i\cdot (\NN(\ell'mn))^{-i}), \quad m, n \in \Z[i]
\end{equation}
Then,
\begin{equation}\label{weight2}
\lim_{N \to \infty} \E_{n, m \in [-N,N]^2 } w_{\delta}(m,n)>0, \text{ and } \lim_{N \to \infty} \E_{n, m \in [-N,N]^2 } \tilde{w}_{\delta}(m,n)>0,.
\end{equation}
\end{lemma}
\begin{proof}
We begin by showing that the limit in \eqref{weight1} is positive. Indeed, said limit can be rewritten as
\[
\lim_{N \to \infty} \E_{n, m \in [-N,N]^2 }F_{\delta}((\NN(\ell((m/N)^2-(n/N)^2))^i\cdot (\NN(\ell'\cdot m/N\cdot n/N))^{-i}).
\]
Let $\NN: \C \to [0,\infty)$ be given by the usual norm in $\C$, i.e., $\NN(z):=|z|^2$. Then, since the function $G_{\delta}(w,z):=F_{\delta}( (\NN(\ell(w^2-z^2)))^i \cdot (\NN(\ell'\cdot w \cdot z))^{-i} \mathds{1}_{([-1,1] \times [-i,i])^2, wz \neq 0}(w,z)$ is continuous except for a set of measure zero with respect to the Lebesgue measure on $S \times S:=[-1,1] \times [-i,i]$. Reinterpreting the averages above as Riemann sums shows that the limit we are interested in, is in fact equal to
\[
\int_{S}\int_{S} G_{\delta}(w,z) \ dwdz,
\]
where $dwdz$ is the appropriately normalized Lebesgue measure on the square $S \times S$. It remains to show that the integral of $G_{\delta}$ over $S \times S$ is positive. Indeed, as long as we show that $G_{\delta}$ is not identically $0$ outside the set of points where either $w$ or $z$ equals $0$, since it is non-negative and continuous almost everywhere, we are done. 

In order to see that this is the case, observe that if we set $w=az$ for $a \in \R$, the property $\NN(\ell(w^2-z^2))=\NN(\ell'wz)$ will be satisfied, provided that we can solve the equation 
\[
\frac{\NN(\ell)}{\NN(\ell')}(a^2-1)=a.
\]
Setting $\lambda:=\frac{\NN(\ell)}{\NN(\ell')}$, we see that it has the solution $a = \frac{\sqrt{4\NN(l)^2+\NN(l)^2}+\NN(l)}{2\NN(l')}$, which ensures that indeed, we can solve the equation and have $a>1 $. Thus, on all points of the line $w=az$ inside the square $S \times S$, we have that $G_{\delta} \neq 0$, which in particular means that the integral under consideration is positive, since we are obtaining non-zero values for $w, z$ which is important given the way the function $G_{\delta}$ is defined.

For the second weight in \eqref{weight2} the argument we use is essentially the same, so we only give a brief outline.

First, we consider the continuous function $\tilde{G}_{\delta}(w,z):=F_{\delta}( (\NN(\ell(w^2+z^2)))^i \cdot (\NN(\ell'\cdot w \cdot z))^{-i}\mathds{1}_{S \times S, wz \neq 0}(w,z)$. A similar argument used for the first weight, shows that the limit in \eqref{weight2} is, in fact, equal to
\[
\int_{S}\int_{S} \tilde{G}_{\delta}(w,z) \ dwdz,
\]
where $dwdz$ is the same measure as for the first weight. Thus, in order to complete the proof, it only remains to show that $\tilde{G}_{\delta}$ is not identically zero. In order to show this, we take $b:=\NN(\ell'/\ell)e^{2k\pi}$ for some $k$ large enough so that $b>2$. Then, setting $a:=\frac{b+\sqrt{b^2-4}}{2}$ gives us an element such that
\[
\frac{\NN(\ell)}{\NN(\ell')}(a^2+1) = e^{2k\pi} a,
\]
where $a>1$. This being the case ensures that the argument of $\tilde{G}_{\delta}$ is equal to $1$ for non-zero values of $w, z$. This implies that said function is non-zero off the points where it is not defined, so its integral is positive, as we wished to show.
\end{proof}

Next we prove an analogue of \cite[Proposition~4.1]{FKM1}:
\begin{proposition}\label{aperiodic0}
Let $f: \Z[i] \to \U$ be an aperiodic completely multiplicative function, let $\ell,\ell', Q \in \Z[i]$ and $\delta>0$. Then, with the weights $w_{\delta}$ and $\tilde{w}_{\delta}$ defined in Lemma~\ref{weights} we have
\[
\lim_{N \to \infty} \E_{m, n \in [-N,N]^2} w_{\delta}(m,n) \cdot f(\ell(Qm+1)^2-(Qn)^2)\cdot \overline{f(\ell'(Qm+1)(Qn))} = 0,
\]
and
\[
\lim_{N \to \infty} \E_{m, n \in [-N,N]^2} \tilde{w}_{\delta}(m,n) \cdot f(\ell(Qm+1)^2+(Qn)^2)\cdot \overline{f(\ell'(Qm+1)(Qn))} = 0.
\]
\end{proposition}
\begin{proof}
By the definition of $w_{\delta}(m,n)$, we can use the Stone-Weierstrass theorem on $\S^1$, so that together with linearity of the limit, the continuous function $F_{\delta}$ appearing in the weights can be replaced by a power $z^k$. This will simplify matters considerably, transforming the averages we need to show converge to $0$ into
\[
\lim_{N \to \infty} \E_{m, n \in [-N,N]^2} (\NN(\ell(m^2-n^2)))^{ki} \cdot (\NN(\ell'mn))^{-ki} \cdot f(\ell(Qm+1)^2-(Qn)^2)\cdot \overline{f(\ell'(Qm+1)(Qn))},
\]
for $k \in \Z$. Note that, with the same trick we used in Lemma~\ref{weights} we may as well replace $n, m$ by $Qn$ and $Qm$ respectively inside the norms, because the changes cancel out with the complex conjugate.

We next argue that, in fact, we can change $Qm$ by $Qm+1$ without affecting the averages. Indeed, for given $n \in \N$, $\lim_{m \to \infty} \log \frac{\NN((\ell(Qm)^2-(Qn)^2))}{\NN(\ell((Qm+1)^2-(Qn)^2))}=0$, and similarly, $\lim_{m \to \infty} \log \frac{\NN(\ell(Qmn))}{\NN(\ell(Qm+1)Qn)}=0$, so this change is valid. 

Therefore, it is enough to establish the convergence to $0$ of the averages
\[
\E_{m, n \in [-N,N]^2} f_k(\ell(Qm+1)^2-(Qn)^2)\cdot \overline{f_k(\ell'(Qm+1)(Qn))},
\]
where $f_k(n):=f(n) \cdot (\NN(n))^{ik}$.
{\color{red} Double check this, it should be correct}. Since $f$ is aperiodic, so is $f_k$ and $\overline{f_k}$, so using the analog of Frantzikinakis and Host's theorem for $\Z[i]$, that is, Proposition~\ref{P: appendix} we are done.

The argument for the second expression is exactly the same, we simply replace the minus with the plus. Note that we can still apply Proposition~\ref{P: appendix} because over $\Z[i]$ we have the factorization $(Qm+1)^2+(Qn)^2=(Qm+1+i(Qn))(Qm+1-i(Qn))$, which is needed for this application to be valid.
\end{proof}
Using the weights above, we now use a variation of \cite[Lemma~2.6]{FKM1}:

\begin{lemma}
    Let $f\colon \Z[i]\to\S^{1}$ be a completely multiplicative function which is pretentious but not super pretentious. Let $\delta>0$. Then
    \begin{equation} \label{e116}
    \lim_{K\to\infty}\E_{Q\in \Phi_{K}}\lim_{N \to \infty} \E_{m,n\in \Z[i] \cap [-N,N]^2} w_{\delta}(m,n)f(\ell(Qm+1)^2-(Qn)^2)\overline{f(\ell'(Qm+1)Qn)}=0,
\end{equation}
and
\begin{equation}\label{weight2withf}
    \lim_{K\to\infty}\E_{Q\in \Phi_{K}}\lim_{N \to \infty} \E_{m,n\in \Z[i] \cap [-N,N]^2} \tilde{w}_{\delta}(m,n)f(\ell(Qm+1)^2+(Qn)^2)\overline{f(\ell'(Qm+1)Qn)}=0
\end{equation}
where $w_{\delta}$ and $\tilde{w}_{\delta}$ are as defined in Lemma~\ref{weights}, and the limit exists because of Proposition~\ref{aperiodic0}.
\end{lemma}
\begin{proof}
We begin by showing \eqref{e116}. As we will point out later, the proof for \eqref{weight2withf} is very similar and only requires some small modifications.

By hypothesis on $f$, there exist $\tau \in \R$ and a non-trivial Dirichlet character $\chi$ such that $f(u)=g(u)\cdot \NN(u)^{i\tau}$, where $g\sim \chi, g\neq 1$. Let
    $$\tilde{L}_{\delta}(f,Q):=f(Q)\cdot \NN(Q)^{-i\tau}\cdot\lim_{N \to \infty} \E_{m,n\in \Z[i] \cap [-N,N]^2} w_{\delta}(m,n)f(\ell(Qm+1)^2-(Qn)^2)\overline{f(\ell'(Qm+1)Qn)}.$$

Fix $\varepsilon>0$ and take $K_{0}=K_{0}(\varepsilon,f)$ so that 
$$\sum_{\NN(p)>K_{0}}\frac{1}{p}(1-\text{Re}(f(p)\cdot \overline{\chi}(p)\cdot \NN(p)^{-i\tau}))+K_{0}^{-1/2}\leq \varepsilon.$$
By Proposition \ref{p25}, we have that
$$\E_{u\in [-N,N]^{2}}\vert f(Qu+1)-\NN(Qu)^{i\tau}\exp(F_{N}(f,K))\vert\ll\varepsilon$$
for every $N>K>K_{0}$ and $Q\in\Phi_{K}$. By Lemma~\ref{sequenceavfact}, this implies (for any choice of plus and minus, that
\[
\limsup_{N\to\infty}\E_{m,n\in [-N,N]^{2}}\left| f(Q(m\pm n))+1)-\NN(Q(m\pm n))^{i\tau}\exp(F_{10N}(f,K))\right|\ll\varepsilon.
\]
{\color{red} Note: in [FKM], for the $m-n$ case, the average was taken over those $m>n$ to ensure $m^{2}-n^{2}>0$.}
Combining the above 3 estimates together with the triangle inequality (since all the involved functions have absolute value less than one), we have that
 \begin{equation} 
    \begin{split}
       &\limsup_{N\to\infty}\E_{m,n\in [-N,N]^{2}}\Bigl\vert f((Qm+1)^2-(Qn)^2)\overline{f(Qm+1)}
       \\&-\NN\Bigl(\frac{(Qm)^2-(Qn)^2}{Qm}\Bigr)^{i\tau}\exp(2F_{10N}(f,K)-F_{N}(f,K))\Bigr\vert\ll\varepsilon. 
    \end{split}
\end{equation}
We now multiply by $f(\ell)\overline{f(\ell')}\cdot w_{\delta} \cdot \overline{f(Qn)}\NN(Q)^{i\tau}$, which is bounded in absolute value by $1$ to obtain
 \begin{equation} 
    \begin{split}
       &\limsup_{N\to\infty}\E_{m,n\in [-N,N]^{2}}\left| w_{\delta}(m,n)f(\ell((Qm+1)^2-(Qn)^2))\overline{f(\ell'(Qm+1)Qn)}(f(Q)\NN(Q)^{-i\tau})\right.
       \\&-w_{\delta}(m,n)f(\ell)\overline{f(\ell')}\left.\NN\left(\frac{m^2-n^2}{mn}\right)^{i\tau}\overline{f(n)}\exp(2F_{10N}(f,K)-F_{N}(f,K))\right|\ll\varepsilon. 
    \end{split}
\end{equation}
Since the second term is independent of $Q$, if $K>K_0$, it is $|\tilde{L}_{\delta}(f,Q)-\tilde{L}_{\delta}(f,Q')|\ll \varepsilon$ (for $Q, Q' \in \Phi_M$). Therefore, letting $K \to \infty$ we obtain
$$\limsup_{M\to\infty}\max_{Q,Q'\in\Phi_{M}}\vert\tilde{L}_{\delta}(f,Q)-\tilde{L}_{\delta}(f,Q')\vert=0.$$
For each $M\in\N$, pick some $Q_{M}\in\Phi_{M}$. Then
the left hand side of (\ref{e116}) is equal to
$$\limsup_{M\to\infty}\tilde{L}_{\delta}(f,Q_{M})\E_{Q\in\Phi_{M}}\overline{f}(Q)\NN(Q)^{i\tau}.$$
Since $Q\mapsto \overline{f}(Q)Q^{i\tau}$ is a nontrivial multiplicative function, it follows from Lemma~\ref{multav0} that $\lim_{M \to \infty} \E_{Q\in\Phi_{M}}\overline{f}(Q)\NN(Q)^{i\tau}=0$ and we are done.

The changes needed for \eqref{weight2withf} are very minor. Namely, all the expressions of the form $(Qm+1)^2-(Qn)^2$ get replaced by $(Qm+1)^2+(Qn)^2$. Then notice that Lemma~\ref{sequenceavfact} also implies that 
\[
\limsup_{N\to\infty}\E_{m,n\in [-N,N]^{2}}\left| f(Q(m\pm in))+1)-\NN(Q(m\pm in))^{i\tau}\exp(F_{10N}(f,K))\right|\ll\varepsilon.
\]
We replace all the $m^2-n^2$ for $m^2+n^2$ in the other parts of the proof to conclude in the same manner.
\end{proof}

We also need the following variation of \cite[Lemma~2.8]{FKM1}. 
\begin{lemma}
    Let $\sigma$ be a Borel probability measure on $\mathcal{M}_{p}$ (the set of pretentious functions) with $\sigma(\{1\})>0$ and let $\mathcal{A}=\{(\NN(u)^{i\tau})_{u\in\Z[i]}\colon \tau\in\R\}$.  There exist $\delta_{0},\rho_{0}>0$ dedpending only on $\sigma$ such that
    \begin{equation}\label{measurepositivity1}
    \liminf_{N\to\infty}\inf_{Q\in\Z[i]_{+}}\text{Re}\left(\E_{m,n\in [\pm N]^{2}}\int_{\mathcal{A}}w_{\delta}(m,n)f(\ell((Qm+1)^2-(Qn)^2))\overline{f(\ell'(Qm+1)Qn)}\,d\sigma(f)\right)\geq \rho_{0},
    \end{equation}
    and
    \begin{equation}\label{measurepositivity2}
    \liminf_{N\to\infty}\inf_{Q\in\Z[i]_{+}}\text{Re}\left(\E_{m,n\in  [\pm N]^{2}}\int_{\mathcal{A}}\tilde{w}_{\delta}(m,n)f(\ell((Qm+1)^2+(Qn)^2))\overline{f(\ell'(Qm+1)Qn)}\,d\sigma(f)\right)\geq \rho_{0}
    \end{equation}
\end{lemma}
\begin{proof}
We will only prove inequality \eqref{measurepositivity1}, since inequality \eqref{measurepositivity2} follows with the exact same methods, simply changing the minus for the plus throughout. In particular, once the two inequalities are established, we can make sure the $\delta_0,\rho_0>0$ work for both without loss of generality.

Let $a:=\sigma(\{1\})>0$ and for $\delta>0$ set
    $$\mu_{\delta}:=\lim_{N\to\infty}\E_{m,n\in [\pm N]^{2}}w_{\delta}(m,n).$$
By Lemma~\ref{weights}, we have that the limit exists and $\mu_{\delta}>0$. For $T>0$, denote
$$\mathcal{A}_{T}:=\{(\NN(u)^{i\tau})_{u\in\Z[i]}\colon \tau\in [-T,T]\}$$
which we claim is a closed set. Indeed, we can write
\[
\mathcal{A}_T= \bigcap_{u \in \Z[i]} \pi_u^{-1}\left(\{ \NN(u)^{i\tau} : \tau \in [-T,T] \}\right),
\]
where $\pi_u$ is the projection onto the $u$-th coordinate, with $u \in \Z[i]$. Now, given $\tau \in \R$, the set $\{\NN(u)^{i\tau} : \tau \in [-T,T] \}$ is the image of the compact set $[-T,T]$ under the continuous map $\tau \mapsto\NN(u)^{-i\tau}$, which means it is compact, and since the target space is Hausdorff, it must be closed. The topology taken on $\mathcal{M}$ makes each of the $\pi_u$ continuous, and an aribtrary intersection of closed sets is closed. Thus, $\mathcal{A}_T$ is a Borel set.

By the monotone convergence theorem for sets, there exists $T_{0}=T_{0}(\sigma)>0$ such that $\sigma(\mathcal{A}\backslash\mathcal{A}_{T_{0}})\leq a/2$.
Since $\lim_{N\to\infty}\sup_{Q\in \Z[i]_{+}}\vert\log\NN(Qu+1)-\log\NN(Qu)\vert=0$,
we have that 
\begin{equation}
    \begin{split}
        &\lim_{N\to\infty}\sup_{f\in \mathcal{A}_{T_{0}},Q\in\Z[i]_{+}}\E_{m,n\in \Z[i]\cap [\pm N]^{2}}\vert f(\ell((Qm+1)^2-(Qn)^2))\overline{f(\ell'(Qm+1)Qn)}
        \\&-f(\ell(m^{2}-n^{2}))\overline{f(\ell'mn)}\vert=0.
    \end{split}
\end{equation}
On the other hand, by the definition of $w_{\delta}$, {\color{blue} above, the limit in $N$ after the supremum will only have contributions close to the value of $w_{\delta}$ when the value of $f(\ell(m^2-n^2))\overline{f(\ell'mn)}$ is very close to $1$, but in the arc $I_{\delta}$. This will be independent of the choice of multiplicative function $f$ provided it is in the set $\mathcal{A}_{T_0}$, so its difference with the average without the $f$'s is $O(\delta)$, i.e.
\[\lim_{N\to\infty}\sup_{f\in \mathcal{A}_{T_{0}}} \vert \E_{m,n\in \Z[i]\cap [\pm N]^{2}}w_{\delta}(m,n)f(\ell(m^{2}-n^{2}))\overline{f(\ell'mn)}-\E_{m,n\in \Z[i]\cap [\pm N]^{2}}w_{\delta}(m,n)\vert=O(\delta)\]

Thus, taking the last limit in $\delta$ yields the equality that was claimed.

I think we only may need to change the average to be on $[-N,N]^2$ so that all the other results can be applied.}
    \begin{equation}
    \begin{split}
        \lim_{\delta\to 0^{+}}\lim_{N\to\infty}\sup_{f\in \mathcal{A}_{T_{0}}} \left| \E_{m,n\in  [\pm N]^{2}}w_{\delta}(m,n)f(\ell(m^{2}-n^{2}))\overline{f(\ell'mn)}-\E_{m,n\in  [\pm N]^{2}}w_{\delta}(m,n)\right|=0.
    \end{split}
\end{equation}
So if $\delta_{0}$ is sufficiently small depending only on $T_{0}$, we have that
$$\liminf_{N\to\infty}\inf_{Q\in\Z[i]_{+}}\text{Re}\left(\E_{m,n\in [\pm N]^{2}}\int_{\mathcal{A}_{T_{0}}}w_{\delta}(m,n)f(\ell((Qm+1)^2-(Qn)^2))\overline{f(\ell'(Qm+1)Qn)}\,d\sigma(f)\right)$$
is at least $\sigma(\mathcal{A}_{T_{0}})\mu_{\delta_{0}}\geq \sigma(\{1\})\mu_{\delta_{0}}=a\mu_{\delta_{0}}$.
On the other hand, 
$$\liminf_{N\to\infty}\sup_{Q\in\Z[i]_{+}}\left|\E_{m,n\in [\pm N]^{2}}\int_{\mathcal{A}\backslash\mathcal{A}_{T_{0}}}w_{\delta}(m,n)f(\ell((Qm+1)^2-(Qn)^2))\overline{f(\ell'(Qm+1)Qn)}\,d\sigma(f)\right|$$
is bounded above by $a\mu_{\delta_{0}}/2$. We are done by setting $\rho_{0}:=a\mu_{\delta_{0}}/2$.
\end{proof}

{\color{red}
Combining everything from this section and writing some of the analogs as in Section 2 of [FKM] ensures that we have the following result:
\begin{theorem}
Let $a, b, c \in \N$ be squares. Then, the equation
\[
ax^2\pm by^2 = \pm cz^2
\]
is partition regular in $x, y$ and $x, z$. 
{\color{blue} Theorem 1.2 of [FKM] shows that $ax^{2}+by^{2}+cz^{2}=0$ is p.r. over $\mathbb{Z}$ for any two of $x,y,z$ if one or two of $a,b,c$ are squares and if the rests are negative squares. So the only case that was not covered in [FKM] is when all of $a,b,c$ are positive squares. In this case p.r. over $\mathbb{Z}$ obviously fails, but our method (for type 1 tuples) shows that p.r. over $\mathbb{Z}[i]$ holds. (Similar to the comment in the beginning of Section 6, we don't need type 2 estimates for this equation)
}

More generally, let $a,b,c \in \Z[i]$. Then, the equation
\[
a^2x^2+b^2y^2=c^2z^2
\]
is partition regular with respect to $x, y$ and $x, z$.
\end{theorem}
Note that we could also use this method to deal with more general quadratic forms that 
}

{\color{blue}{Things that we should check: If we have other UFDs (notion that coinsides to the one of PIDs for rings of integers of number fields), can we get the corresponding result? Here, one potential problem, e.g., in $\mathbb{Z}[\sqrt{2}],$ is that we may have infinitely many units. We crucially used that $\mathbb{Z}[i]$ has 4 units. What about finitely many? (Check the concentration estimates for the quadratic term in the parametrization that cannot be split as a product of linear terms. Here, we have 
\[x^2+2y^2=z^2,\] with parametrization, $x=k(m^2-2n^2),$ $y=k2mn,$ $z=k(m^2+2n^2),$ so we expect $(x,y)$ and if the other ``details'' work, $(y,z).$) Also, we may get negative values in the norm.

Can we get something more, say for $x^2+2y^2=z^2$ in $\mathbb{Z}[i\sqrt{2},\sqrt{2}]=\langle 1,i\sqrt{2},\sqrt{2}\rangle$? Here everything factors linearly. $R:=\mathcal{O}(\Q(\sqrt{2},i))$
{\color{red} This equation is known to be p.r. over $R$ by [Sun2] (w.r.t. any two of $x,y,z$). So what remains open is: (1) p.r. over $\mathbb{Z}[\sqrt{2}]$ for $x,y$; (2) p.r. over $\mathbb{Z}[\sqrt{-2}]$ for $y,z$; (3) p.r. over $\mathbb{Z}[\sqrt{-2}]$ for $x,z$; (4) p.r. over $\mathbb{Z}[\sqrt{2}]$ for $x,z$.
In terms of difficulty, I think it is Q2<Q3<Q1<Q4.
We have Q2<Q3 and Q1<Q4 because Q1 and Q2 are type 1 pairs whereas Q3 and Q2 are type 2 pairs. We also have that Q2,Q3<Q1,Q4 because $\mathbb{Z}[\sqrt{-2}]$ has infinitely many units.


I have realized we may be able to use to our advantage that, at least, for $\Z[\sqrt{2}]$, the infinitely many units are all of the form $(1+\sqrt{2})^k, k \in \Z$, so as soon as $f(1+\sqrt{2})$ is determined, so is the value of any of the other units. As for $\Z[\sqrt{-2}]$ here we should only have $\pm 1$ as units, since $\NN(a+\sqrt{-2}b)=a^2+2b^2$, and the norm of a unit must be equal to one, so $b=0$, whence $a=\pm 1$, as claimed.
}


Maybe here we can also get conclusions that with the known results (even in FKM2) are not known. E.g., in FKM2, they say that they don't know whether $x^2+17y^2=34z^2$ is pr in $x,y.$ But, working in $\mathbb{Z}[\sqrt{17}]$ we have the parametrization ???
{\color{red}
$x=k34(m^{2}-34mn-17n^{2}), y=k34(m^{2}+2mn-17n^{2}), z=6k\sqrt{17}(m^{2}+17n^{2})$, where $x$ is factorizable over $\Z[\sqrt{34}]$, $y$ is factorizable over $\Z[\sqrt{2}]$ and $z$ is factorizable over $\Z[\sqrt{17}]$. Another factorization is 
$x=k\sqrt{17}(m^{2}-2\cdot 17^{2}n^{2}), y=-k(m^{2}+4\cdot 17mn+2\cdot 17^{2}n^{2}), z=-k(m^{2}+34mn+2\cdot 17^{2}n^{2})$.
}

Note that not all quadratic equations over $\Z[i]$ are partition regular in pairs. Let $p$ be a Gaussian prime and consider the equation $x^2+y^2=pz^2$. For a Gaussian integer $n$, associate it with the color that gives the parity of the largest power of $p$ that divides $n$. Then no solution of the same color can exist. How about two variables? Maybe same idea, with $p=1+i$. 

}}
{\color{red} IMPORTANT NOTE: Since we are working on $\Z[i]$, the parametrization of the type 2 pairs, actually reduces to a type 1 pair: i.e. a product of linear forms, and it appears we will not need the concentration estimates (i.e. this whole section 5, which is only to deal with the quadratic forms that are irreducible) to obtain the partition regularity that we sought. 
{\color{blue} Therefore, a related question is: is there any p.r. result we can deduce assuming that our method works for type 2 tuples? For example, if we can consider the p.r. problem over $x^{2}+y^{2}+z^{2}=0$, then by symmetry there is no type 2 tuple at all.
}

It is still interesting to see if such concentration estimates can be achieved in $\Z[i]$, though I believe a different approach may be required when dealing with Lemma 5.3 of [FKM] because writing out the argument in a bit more detail, it seems that $N^2\log N$ is a lower bound.

My feeling is that the issue is that $x^2+y^2$ is not irreducible over $\Z[i],$ which leads me to believe that, perhaps, if we were to replace it for a different IRREDUCIBLE quadratic form, then we could achieve the estimates, and extend this method to deal with fairly general quadratic forms, as in [FKM].
}
We will attempt to see first what can be recovered for $\Z[i]$ without taking into account the balls vs rectangles issue.

First, let $f, g:\Z[i] \to \U$ be multiplicative functions, $\chi$ a Dirichlet character in the sense defined above and $t\in \R$. For $K_0 \in \N$, let

{\color{red}  May need to change the definitions to boxes; double check later; it should hopefully not make a big difference.}
\[
G_N(f,K_0):= 2 \sum_{K_0<\NN(p)\leq N, \NN(p) \equiv 1 \pmod{4}} \frac{1}{\NN(p)}(f(p)\cdot \overline{\chi(p)}\cdot \NN(p)^{-it}-1),
\]
and
\[
\D_1(f,g;x,y)^2:=\sum_{x<\NN(p)\leq y; \NN(p) \equiv 1 \pmod{4}} \frac{1}{\NN(p)}(1-\text{Re}(f(p)\cdot\overline{g(p)}))
\]
We want to show the following result:
\begin{proposition}\label{concentrationestimate}
    Let $K_0, N \in \N$ and $f: \Z[i] \to \U$ be a multiplicative function. Let $t \in \R$ and $\chi$ be a Dirichlet character with index $q$. Put $Q:=\prod_{\NN(p) \leq K_0} p^{a_p}$ for some natural numbers $a_p$ and suppose that $q \mid \NN(Q)$. If $N$ is large enough, depending ONLY on $Q$, then for all $a, b \in \Z[i]$ with $\NN(a), \NN(b)\leq \NN(Q)$ and $(a^2+b^2, Q)=1$ we have
    \[
    \E_{m, n \in [\pm N]^2} |f((Qm+a)^2+(Qn+b)^2)-\chi(a^2+b^2)\cdot \NN(Q)^{2it}(\NN(m^2+n^2))^{it}\exp(G_N(f,K_0))| \ll
    \]
    \[
    (\D_1+\D_1^2)(f, \chi\cdot  \NN(n)^{it};K_0,\sqrt{N})+Q^2\cdot\D_1(f,\chi \cdot \NN(n)^{it};N,3Q^2N^2)+Q\cdot \D_1(f,\chi\cdot\NN(n)^{it};\sqrt{N},N)+K_0^{-1/2},
    \]
    with the implicit constant being absolute.
\end{proposition}
{\color{red} From what  I saw, the conditions I tried to imitate that $p, q \equiv 1 \pmod 4$ in the definitions of the distance functions and $G_N$ are there because of the discussion of number of solutions of quadratic equations modulo $p$ where the involved equation is a sum of squares. But for us, $-1=i^2$ is always true, so $-1$ is always a square for us and we don't need to distinguish such primes.

Theorem 2 from https://www.ams.org/journals/tran/1940-048-03/S0002-9947-1940-0003000-5/S0002-9947-1940-0003000-5.pdf should be helpful.

%Lemma 5.5 and 5.6 lift and improve the results from additive functions to multiplicative functions. It seems that here the restrictions on primes may also be needed because the proof uses when a prime divides a sum of two squares, it must divide one of the two summands if $p$ is congruent to $3$ mod $4$.
}
 
{\color{red} I am having some trouble with Lemma 5.3.}
In order to prove the concentration estimate given in Proposition~\ref{concentrationestimate} above, we begin with the following lemma:
\begin{lemma}
For $N\in \N$ and $Q \in \Z[i]$, $a, b \in \Z[i]$, and primes $p,q$ such that $(pq, Q)=1$, let
\[
w_{N,Q}(p,q):=\frac{1}{(2N)^4} \sum_{n, m \in [\pm N]^2; p, q|| (Qm+a)^2+(Qn+b)^2} 1.
\]
Then,
\begin{equation}\label{asymptoticpqequal}
w_{N, Q}(p,p) = \frac{2}{\NN(p)}\left(1-\frac{1}{\NN(p)} \right)^2+O\left(\frac{1}{N^2}\right),
\end{equation}
and if $p \neq q$, it is
\begin{equation}\label{asymptoticpqunequal}
w_{N,Q}(p,q) = \frac{4}{\NN(pq)}\left(1-\frac{1}{\NN(p)}\right)^2\left(1-\frac{1}{\NN(q)}\right)^2+O\left(\frac{1}{N^2}\right).
\end{equation}
\end{lemma}
\begin{proof}
As in [FKM], let $\epsilon \in \{0,1,2,3,4\}$. We begin by showing the asymptotic behavior in equation~\eqref{asymptoticpqequal}. Suppose that $p$ is a prime with $(p, Q)=1$. First note that if $p \mid Qn+b$, then, if $p \mid (Qm+a)^2+(Qn+b)^2$, it is also $p^2 \mid (Qm+a)^2+(Qn+b)^2$ since $\Z[i]$ is a UFD. In other words, such $p$ will not contribute anything to the sum defining $w_{N,Q}(p,p)$. Thus, we may as well suppose that $p \nmid (Qn+b)$. Since $-1=i^2$ and $p$ is a prime element of $\Z[i]$, the equation
\[
(Qm+a)^2+(Qn+b)^2 \equiv 0 \pmod p
\]
will only have two solutions in $m$, provided the $n$ variable is fixed. So if $n$ varies in the rectangle $[\pm N]^2$, there are $2\left[ \frac{(2N)^2}{\NN(p)}\right]+\epsilon$ solutions in $m \in [\pm N]^2$ to the quadratic equation above (Recall that the size of $\Z[i]/(p)$ is equal to $\NN(p)$). Now, there are $(2N)^2-\lfloor \frac{(2N)^2}{\NN(p)}\rfloor+\epsilon$ elements with $p \nmid Qn+b$ as we assumed that $(p, Q) =1$. This makes for a total of (after removing the epsilons and the floors and absorbing them into the error term)
\[
2\frac{(2N)^4}{\NN(p)}-2\frac{(2N)^4}{\NN(p)^2}+O(N^2)
\]
solutions to the quadratic equation in the two variables. This is almost what we need, but we need to remove the possible solutions to this congruence modulo $p^2$, since we are interested only in the $p$ such that $p \mid\mid (Qm+a)^2+(Qn+b)^2$. We can argue in the same way to now deduce the estimate
\[
2 \frac{(2N)^2}{\NN(p)^2}-2\frac{(2N)^2}{\NN(p)^3}+O(N^2)
\]

for the number of solutions to $(Qm+a)^2+(Qn+b)^2 \equiv 0 \pmod{p^2}$. It is worth mentioning that this case is slightly different in $\Z[i]$ when compared to $\Z$, but again, the equations look like $n^2-a^2\equiv 0 \pmod{p^2}$ and even though $\Z[i]/(p^2)$ is not a field, the two solutions $n=\pm a$ are unique, for yet another solution for which $n \equiv a \pmod p$ and $n \equiv -a \pmod p$ implies $2n \equiv 0 \pmod p$, which will yield a contradiction with $n^2 \equiv a^2 \pmod{p^2}$ provided that $2$ is invertible. This is not a problem for us, since our choice of $Q$ later will always be such that $p, q$ will be ``large''.

Finally, with these terms, an easy computation shows that
\[
w_{N,Q}(p,p) = \frac{2}{\NN(p)}\left( 1-\frac{1}{\NN(p)}\right)^2+O\left(\frac{1}{N^2}\right).
\]
We now proceed to establish \eqref{asymptoticpqunequal}. As with the case of \eqref{asymptoticpqequal}, we may assume that both $p, q$ do not divide $Qn+b$, for otherwise they do not contribute to the sum, so we will have $(pq, Qn+b)=1$. Given $r, s \in \Z[i]$, consider the sums
\[
A_{r,s}:=\frac{1}{(2N)^2}\sum_{m, n \in [\pm N]^2, r,s \mid (Qm+a)^2+(Qn+b)^2,  (rs,Qn+b)=1}1.
\]
These help us realize, by inclusion-exclusion that, in fact,
\[
w_{N}(p,q)=A_{p,q}-A_{p^2,q}-A_{p,q^2}+A_{p^2,q^2}.
\]
Now, using the Chinese Remainder Theorem, the congruence
\begin{equation}\label{congruencemodpq}
(Qm+a)^2+(Qn+b)^2 \equiv 0 \pmod{pq}
\end{equation}
will have $4$ solutions in $m$ for each fixed $n \in [\pm N]^2$ (recall again that $-1=i^2$), since we will have $2$ solutions $\pmod p$ and $2$ more $\pmod q$, since $p \nmid Qn+b$. This amounts to saying that given $n \in [\pm N]^2$ with $(pq, Qn+b)=1$, equation \eqref{congruencemodpq} will have $4\left[ \frac{(2N)^2}{\NN(pq)} \right]+\epsilon$ solutions in the variable $m \in [\pm N]^2$. 

Using again the fact that $\Z[i]$ is a UFD and inclusion-exclusion again, we see that the number of $n \in [\pm N]^2$ such that $(pq, Qn+b)=1$ is $(2N)^2-\left[\frac{(2N)^2}{\NN(p)}\right]-\left[\frac{(2N)^2}{\NN(q)}\right]+\left[\frac{(2N)^2}{\NN(pq)}\right]$. It follows that the number of solutions to \eqref{congruencemodpq} for $n, m \in [\pm N]^2$ with $(pq, Qn+b)=1$ is
\[
(2N)^4\frac{4}{\NN(pq)}\cdot \left( 1-\frac{1}{\NN(p)}-\frac{1}{\NN(q)}+\frac{1}{\NN(pq)}\right)+O(N^2), 
\]
so we obtain
\[
A_{p,q} = \frac{4}{\NN(pq)}\left(1-\frac{1}{\NN(p)}\right)\left(1-\frac{1}{\NN(q)}\right)+O\left(\frac{1}{N^2}\right).
\]
We can argue in the same way for $A_{p^2,q}, A_{p,q^2}$ and $A_{p^2,q^2}$, again using that $i^2=-1$, the Chinese Reminder Theorem, and the same argument as in the $p=q$ case for the congruences modulo $p^2$ and $q^2$ respectively. These yield the estimates
\[
A_{p^2,q} = \frac{4}{\NN(p^2q)}\left(1-\frac{1}{\NN(p)}\right)\left(1-\frac{1}{\NN(q)}\right)+O\left(\frac{1}{N^2}\right)
\]
and
\[
A_{p,q^2} = \frac{4}{\NN(pq^2)}\left(1-\frac{1}{\NN(p)}\right)\left(1-\frac{1}{\NN(q)}\right)+O\left(\frac{1}{N^2}\right)
\]
Lastly,
\[
A_{p^2,q^2} = \frac{4}{\NN(p^2q^2)}\left(1-\frac{1}{\NN(p)}\right)\left(1-\frac{1}{\NN(q)}\right)+O\left(\frac{1}{N^2}\right)
\]
Putting this all together with the expression for $w_N(p,q)$ in terms of the $A_{r,s}$ we derived above yields \eqref{asymptoticpqunequal}, completing the proof.
\end{proof}
We also need the following estimate for $w_{N,Q}(p,q)$ in the case that $p, q$ are not necessarily primes.
\begin{lemma}
    
\end{lemma}
We would like to have the following estimate:
\[
\sum_{n \in [\pm N]^2} r_2(n) \ll N^2?,
\]
where $r_2(n)$ counts the number of ways that $n \in \Z[i]$ can be written as $n=x^2+y^2$, where we can assume that $x, y \in [\pm CN]^2$ as well.
{\color{red} The equation $x^2+y^2=n$ can be written as $(x+iy)(x-iy)=n$. If we consider the $\C$-linear map $\varphi(x,y):=(x-iy,x+iy)$, and the region $R:=\{(w,z) : wz \in [-N,N]^2\}$, then the count we are interested in is the set $R':=\{(w,z) : \varphi(w,z) \in R\}$. Since we essentially want to estimate the area of $R'$, it will be enough to estimate the area of $R$ up to constants, since $\varphi$ is an invertible linear map with $|\det\varphi|=2$. Thus, the problem is transformed into estimating $|R|$. The area/volume of $R$ will contain that of the set $\{(z,w) : 1 \leq |zw| \leq N\} =\{(x,y,z,t) : 1 \leq (x^2+y^2)(z^2+t^2)\leq N^2\}$ thought of in $\R^4 \cong \C^2$. Using a change of variables $(x,y,z,t)=(r_1\cos\theta_1,r_1\sin\theta_1,r_2\cos\theta_2,r_2\sin\theta_2)$, this is at least, $\gg \int_1^N \int_1^{N/r_1} r_1r_2 \ dr_2dr_1 = N^2\log N$. I think this may not be good enough for us for Lemma 5.3 because the growth of $N^2\log N$ is a bit too fast.}
%{\color{blue}
%The sum of $r_{2}$ was only used to esitmate $w_{N,Q}(l)$. But
%do we really need an upper bound of the form $w_{N,Q}(l)\leq \frac{Q^{2}}{l}$? I think in [FMK], this estimate was only used when $l>\sqrt{N}$ (and some times even $l>N$, please check). In this case is the error term $\log N$ accpectable? Update: maybe it is not good enough...
%}


{\color{red} Next thing to check: try to see how to deal with the error term (5.22) from [FKM] in the Gaussian integer case. On the other hand, it seems that (5.21) follows from Andreu's lemma, but it also needs to be checked carefully.}

{\color{blue} check if Lemma 4.15 of [FKM2] is helpful for us, they deal with more general quadrtic expressions instead of sum of squares.}

\begin{lemma}\label{turankubsumsofsquares}
{\color{red} Here the averaging scheme may need to be balls instead of rectangles. I think Lemma 5.1 works just as well for balls without any significant changes, so it should be fine.}
Let $K_0, N \in \N$, $a, b \in\Z[i]$. Assume that $\NN(a), \NN(b) \leq \NN(Q)$.
Let $h: \Z[i] \to \U$ be an additive function such that
\begin{itemize}
    \item[(i)] $h(p)=0$ for all primes with $\NN(p) \leq K_0$ and $\NN(p) >N$
    \item[(ii)] {\color{red} Maybe we do not need this} $h(p)=0$ for all primes with $\NN(p) \equiv 3 \pmod 4$.
    \item[(iii)] $h(p^k)=0$ for all primes $p$ and all $k \geq 2$.
\end{itemize}
Let $Q=\prod_{ \NN(p) \leq K_0} p^{a_p}$ for some $a_p \in \N$. Then for all large enough $N$, depending only on $K_0$ we have
\[
\E_{m,n \in [\pm N]^2} |h((Qm+a)^2+(Qn+b)^2)-H_N(h,K_0)|^2 \ll \D^2(h; K_0, \sqrt{N})+Q^2\D^2(h; \sqrt{N},N)+K_0^{-1},
\]
where the implicit constant is absolute, and
\[
H_N(h,K_0):= 2\sum_{K_0 < \NN(p)\leq N} \frac{h(p)}{\NN(p)},
\]
and
\[
\D^2(h; K_0,N):=\sum_{K_0 <\NN(p) \leq N} \frac{|h(p)|^2}{\NN(p)}
\]
\end{lemma}
\begin{proof}

\end{proof}
\fi
\begin{corollary}\label{c27k}
     Let $g\colon\mathfrak{I}(K)\to\C$ be a multiplicative function with $\vert g\vert\leq 1$. 
     \begin{enumerate}
         \item If $\D(g,1)<\infty$ and if for all prime ideal $\mathfrak{p}$ with $\NN(\mathfrak{p})=2$, we have $g(\mathfrak{p}^{k})\neq -1$ for some $k\geq 1$,  then $\E_{\mathfrak{u}\in\mathfrak{I}(K)}g(\mathfrak{u})$ exists  but is not equal to 0.
         \item If $\D(g,\NN(\mathfrak{u})^{i\tau})<\infty$ for some $\tau\neq 0$ and if for all prime ideals $\mathfrak{p}$ with $\NN(\mathfrak{p})=2$, we have $g(\mathfrak{p}^{k})\neq -2^{ik\tau}$ for some $k\geq 1$, then $\E_{\mathfrak{u}\in\mathfrak{I}(K)}g(\mathfrak{u})$ does not exist.
         \item If $\D(g,\NN(\mathfrak{u})^{i\tau})<\infty$ for some $\tau\in\R$ and if for some prime ideal $\mathfrak{p}$ with $\NN(\mathfrak{p})=2$, we have $g(\mathfrak{p}^{k})=-2^{ik\tau}$ for all $k\geq 1$,  then $\E_{\mathfrak{u}\in\mathfrak{I}(K)}g(\mathfrak{u})$ exists and equals to 0.
         \item If $\inf_{\tau\in\R}\D(g,\NN(\mathfrak{u})^{i\tau})=\infty$, then $\E_{\mathfrak{u}\in\mathfrak{I}(K)}g(\mathfrak{u})$ exists and equals to 0.
     \end{enumerate}
\end{corollary}
\begin{proof}
Let the notations be the same as in Theorem \ref{C22intro} and Theorem~\ref{C21}. We begin by showing (i). If $\D(g,1)<\infty$, then by Theorem \ref{C22intro}, the limit of the averages $\E_{u \in \mf I(K)} g(\mf u)$ is equal to
\[
L_0:=\lim_{x \to \infty} \prod_{\mf p: \text{ prime ideals in }\OK, \NN(\mf p) \leq x}(1-1/\NN(\mf p))(1+h_1(\mf p)).
\]
% $h$ is defined in \ref{eulerproduct1}

By our previous work in Lemma~\ref{lem69elliott} above, the limit $L_0$ exists, and can only be zero (by general theory of infinite products) if one of the factors $(1+h_1(\mf p))$ is equal to zero. This can only happen, given the definition of $h_s(\mf p)$ as a geometric sum, if $\NN(\mf p)=2$, and simultaneously, $g(\mf p^k)=-1$ for all $k \in \N$.

Given the conditions of (i), it follows that the limit exists, but is non-zero.

A similar argument (but to show that the limit is now actually equal to $0$) works for (iii).

First, note that, the limit
\[
L_{\tau}:=\lim_{x \to \infty} \prod_{\mf p: \text{ prime ideals in }\OK, \NN(\mf p) \leq x}(1-1/\NN(\mf p))(1+h_{1+i\tau}(\mf p)),
\]
will exist for the $\tau \in \R$ appearing in assertion (iii) given the distance assumption on $g$.

Furthermore, the condition stated in (iii) forces one of the factors in the product defining $L_{\tau}$ to be $0$, and we get the result.

For Part (ii), assume for the sake of contradiction, that $L_\tau$ exists. Then, as we have argued above, it must be non-zero.
%\WS{Why it must be non-zero? each term is non-zero, but there products may have 0 limit}
Thus, it must be that
\[
\frac{1}{x}\E_{\NN(\mf u) \leq x} g(\mf u) \NN(\mf u)^{i\tau} \to L,
\]
which, by the summation by parts trick, tells us that, in fact,
\[
\frac{1}{x}\E_{\NN(\mf u) \leq x} g(\mf u)=Lx^{i\tau}(1+i\tau)^{-1}+o_{x\to \infty}(1).
\]
By taking complex conjugates we may assume that $\tau \geq 0$, so now suppose that $\tau>0$. We set $x=\exp(2\pi m\tau^{-1})$, with $m\in \N$, and this allows us to deduce that $L(1+i\tau)^{-1}$ is real, but now, if we let $x=\exp(2\pi(m+1/2)\tau^{-1})$, with $m\in \N$, we see that $iL(1+i\tau)^{-1}$ is also real. This is an impossibility, since $L \neq 0$, which means that the averages do not converge, as claimed.
%Note that for any ideal $\mf p$, 
% $$(1-1/\NN(\mf p))\cdot\vert 1+h_{1+i\tau}(\mf p)\vert\leq (1-1/\NN(\mf p))(1+\sum_{k=1}^{\infty}\NN(\mf p)^{-k})=1,$$
 %where the equality holds if and only if $(1-1/\NN(\mf p))(1+h_{1+i\tau}(\mf p))=1$.
%So 
%By the proof of Theorem~\ref{C22intro} (see especially the case where $\tau=0$ in Lemma~\ref{lem69elliott}), the limit  
%$$L:=\lim_{x\to\infty}\prod_{\mf p\colon \text{prime in } \mathcal{O}_K,\NN(\mf p)\leq x}(1-1/\NN(\mf p))(1+h_{1+i\tau}(\mf p))$$
% exists and is finite. Moreover, it is not zero unless one of the terms is zero {\color{red} the precise proof is on pages 248-250 of the book, also maybe the proof of Lemma \ref{eulerproduct1} maybe useful}. By Theorem \ref{C22intro} (ii), if $L=0$, then $\E_{\mathfrak{u}\in\mathfrak{I}(K)}g(\mathfrak{u})$ exists and is equal to 0; if $L\neq 0$ but $\tau=0$, then  $\E_{\mathfrak{u}\in\mathfrak{I}(K)}g(\mathfrak{u})$ exists and is not equal to 0; if $L,\tau\neq 0$, then $\E_{\mathfrak{u}\in\mathfrak{I}(K)}g(\mathfrak{u})$ does not exist.
% Therefore, to prove parts (i), (ii) and (iii), it suffices to show that $L$ is zero if and only if 
% \begin{equation}\label{fefe}
%    \text{$\NN(\mathfrak{p})=2$, we have $g(\mathfrak{p}^{k})=-2^{ik\tau}$ for all $k\geq 1$.}
%\end{equation}
% for some prime ideal $\mf p$ (with $\NN(\mf p)=2$).
 
 %Note that 
 %$$\vert 1+h_{1+i\tau}(\mf p)\vert\geq 1-\sum_{k=1}^{\infty}\NN(\mf p)^{-k}=1-\frac{1}{\NN(\mf p)-1}\geq 0,$$
%where the equality holds if and only if (\ref{fefe}) holds. So we are done.

Finally, Part (iv) follows from Theorem \ref{C22intro} (i).
\end{proof}


 \fi


\section{Proof of the Tur\'an-Kubilius inequality and applications}\label{SECconcest}
In this section we prove the  Tur\'an-Kubilius inequality Theorem \ref{TK} and then use it to obtain an analog concentration estimates of \cite[Proposition~2.5]{FKM1}.
%
%
%we obtain some linear concentration estimates, which allow us to easily treat the pretentious multiplicative functions. In particular, we are going to obtain an analog of \cite[Proposition~2.5]{FKM1}. To do so, we are going to prove a suitable form of a theorem of Tur\'an-Kubilius.  
%
Throughout this section, when $K=\Q(\sqrt{-d})$ for some squarefree $d\in\N$,
let $\gamma_{K}:=\lim_{N \to \infty} \frac{1}{N}\sum_{\NN(u) \leq N} 1$
%\footnote{The limit defining $\gamma_K$ exists as one can rewrite the double sums in questions as suitable Riemann sums. Thus, the actual value of $\gamma_K$ is given by the area $\iint_{\NN(x+\tau_d y) \leq 1} \ \mathds{1}_{[-1,1]^2}(x,y) \ dxdy$.}
%\WS{explain why the limit $\gamma_{K}$ exists}. 
%\footnote{The limit exists and is positive by Lemma \ref{mv7}.}
 be the constant given by Corollary~\ref{lem: bounds for number of ideals of norm at most x}.
 We will need a counting lemma, but in order to properly show it, we must first have a certain estimate for $C$-Lipschitz regions of $\R^2$, the definition of which we recall next.

%Let $K$ be a quadratic imaginary number field, and let $\mathcal{O}_K$ denote its ring of integers. Recall that given $x \in \mathcal{O}_K$ we put $\NN(x):=x\sigma(x) ,$ where $\sigma$ is the non-trivial Galois automorphism of $K.$
%We begin with a short counting lemma, where we will use the notation $B_N:=\{ u\in\mathcal{O}_{K} : \NN(u) \leq N\}$:
%\begin{lemma}\label{countingCRT}
%Let $I$ be an ideal in $\mathcal{O}$, let $K \in \N$, and suppose that $K$ is large enough so that $\Phi_K \subset I$. Consider natural numbers $k, l \in \N$, let $Q$ be an element of smallest possible norm that belongs to $\prod_{\NN(\mathfrak{p})\leq K} \mathfrak{p}$, and let $\mathfrak{p}, \mathfrak{q}$ be distinct prime ideals such that $\mathfrak{p}, \mathfrak{q} \nmid Q$. We have
%\begin{equation}\label{countingCRTeq1}
%\lim_{N\to \infty} \frac{| \{u \in QB_N+a : Q \mid u-a, \mathfrak{q}^l \mid u \text{ and } \mathfrak{p}^k \mid u\}}{\NN(Q)\gamma_{K} N} = \frac{1}{\NN(Q\mathfrak{p}^k) \NN(\mathfrak{q}^l)},
%\end{equation}
%where $\gamma_{K}:=\lim_{N \to \infty} \frac{1}{N}\sum_{\NN(u) \leq N} 1$.
%\end{lemma}

%{\color{red}

%Replace the previous lemma with the following:

%\WS{I commented the old lemma, and I think this is the lemma we need. The old lemma was incorrect because $\Bigl\vert \frac{\gamma_{K} N}{\NN(I)}-\vert(QB_N+a)\cap I|\Bigr\vert$ can not be bounded a constant, For example, the number of integer points in a ball of raduis $\sqrt{N}$ is approximately $\pi N+O(\sqrt{N})$. It is possible to lower the error term $\O(\sqrt{N})$ a little, but one can not reduce it to $O(1)$. So we can only expect the following weaker estimate to hold. But I think this is also good for our purposes.}

\subsection{Proof of the Tur\'an-Kubilius inequality} In this section we prove Theorem \ref{TK}.

First we need the definition of $C$-Lipschitz regions of $\R^2$.
\begin{definition}
Let $S\subseteq \mathbb{R}^{2}$ be the region given in polar coordinates by
     $$S:=\{(\rho\cos\theta,\rho\sin\theta)\colon \theta\in [0,2\pi), 0\leq \rho\leq r(\theta)\}$$
     for some continuous $2\pi$-periodic map $r\colon \R\to \R_{\geq 0}$, i.e., with $r(\cdot+2\pi)=r(\cdot)$. We say that $S$ is \emph{$C$-Lipschitz} if
     $$r(\theta)\leq C\text{ and } \vert r(\theta)-r(\theta')\vert\leq C\vert\theta-\theta'\vert$$
      for all $\theta,\theta'\in\R$.
\end{definition}
We can now give a counting lemma, based on \cite[Lemma~10.1]{kubiliusbook}.
 \begin{lemma}\label{dafp}
 Let $C\geq 1$.
 For any  $C$-Lipschitz region $S\subseteq \mathbb{R}^{2}$, any $v\in\R^{2}$ and any $N\in\R_{+}, N\geq 1$,  we have that
$$\left|\left| (S_{N}+v)\cap \Z^{2}\right|-N^{2}\cdot m(S)\right|\leq O(CN),$$
where $S_{N}+v:=\{x\in\R^{2}\colon \frac{1}{N}(x-v)\in S\}$ and $m(S)$ is the area of $S$.
 \end{lemma}
 \begin{proof}
 \iffalse
     Let $N'=\lfloor N\rfloor$. For $1\leq i\leq N'$, let 
     $$A_{i}=\{(\rho\cos\theta,\rho,\sin\theta)\colon \theta\in [2(i-1)\pi/N',2i\pi/N'), 0\leq \rho\leq r(\theta)\}$$
     and $A_{i,N}+v:=\{x\in\R^{2}\colon \frac{1}{N}(x-v)\in A_{i}\}$. Let $B_{i}$ be the union of unit boxes in $\R^{2}$ whose lower left corner belongs to $(A_{i,N}+v)\cap \Z^{2}$. Then $m(B_{i})=\vert(A_{i,N}+v)\cap \Z^{2}\vert$.
     Let $D_{i}$ be the set of all points whose distance to $A_{i,N}+v$ is at most 10. Then it is clear that $m(D_{i}\backslash(A_{i,N}+v))$ is at most 
     $$100 (\sup_{\theta\in [2(i-1)\pi/N',2i\pi/N')}r(\theta)N+\frac{2\pi}{N'}(\sup_{\theta\in [2(i-1)\pi/N',2i\pi/N')}(r(\theta)N)^{2}-\inf_{\theta\in [2(i-1)\pi/N',2i\pi/N')}(r(\theta)N)^{2})),$$
     which is bounded by
     $$100(CN+\frac{2\pi}{N'}\cdot 2C^{2}\cdot \frac{2\pi}{N'})$$
      \fi
      
     %  Let $B_{N}$ be the union of unit boxes in $\R^{2}$ whose lower left corner belongs to $(S_{N}+v)\cap \Z^{2}$. Then $m(B_{N})=\vert(S_{N}+v)\cap \Z^{2}\vert$.
   
     %Let $D_{N,+}$ be the set of all points whose distance to $S_{N}+v$ is at most 10, and $D_{N,-}$ be the set of all points whose distance to the complement of $S_{N}+v$ is at most 10. Then it is clear that  
     %$$m(D_{N,+}\setminus D_{N,-})\leq O(CN)$$
      %  So 
     %$m(B_{N}\Delta (S_{N}+v))\leq m(D_{N,+}\setminus D_{N,-})\leq O(CN)$. %Then 
     %$$\Bigl\vert\vert (S_{N}+v)\cap \Z^{2}\vert-N^{2}\cdot m(S)\Bigr\vert=\vert m(B_{N})-m(S_{N}+v)\vert\leq O(CN).$$
     %{\color{red}
     Since $\lceil C\rceil\leq 2C$, we may assume without loss of generality that $C\in\N$.
      Let $B_{N}$ be the union of unit boxes in $\R^{2}$ whose lower left corner belongs is a lattice point belonging to $(S_{N}+v)\cap \Z^{2}$. Then, by construction, we ensured that $m(B_{N})=\vert(S_{N}+v)\cap \Z^{2}\vert$. Fix $N$. Divide $[0,2\pi)$ into intervals $I_{i}=[\frac{2\pi(i-1)}{CN},\frac{2\pi i}{CN})$
for $1\leq i\leq CN$, let $r_{i,\inf}:=\inf_{\theta\in I_{i}}r(\theta)$ and $r_{i,\sup}:=\sup_{\theta\in I_{i}}r(\theta)$, and define 
$$D_{N,i}:=\{(\rho\cos \theta,\rho\sin\theta)\colon \theta\in I_{i}, \rho\in [r_{i,\inf}N,r_{i,\sup}N]\}$$
and $D'_{N,i}$ be the set of points whose distance to $D_{N,i}$ is at most 2.

We first claim that $m(D'_{N,i})=O(1)$.
Fix $N$ and $i$. Let $$E_{1}:=\{(\rho\cos \theta,\rho\sin\theta)\colon \theta\in I_{i}, \rho\in [r_{i,\inf}N-10,r_{i,\sup}N+10]\},$$
%\FM{This $-10$ should be with the radius, I imagine?}
$E_{2}$ be the rectangle adjacent to $E_{1}$ with one of its edge being $$\left\{(\rho\cos \theta,\rho\sin\theta)\colon \theta=\frac{2\pi (i-1)}{CN}, \rho\in [r_{i,\inf}N-10,r_{i,\sup}N+10]\right\}$$  and the other edge of length 10 pointing out of $E_{1}$, and $E_{3}$ be the rectangle adjacent to $E_{1}$ with one of its edge being $$\{(\rho\cos \theta,\rho\sin\theta)\colon \theta=\frac{2\pi i}{CN}, \rho\in [r_{i,\inf}N-10,r_{i,\sup}N+10]\}$$  and the other edge of length 10 pointing out of $E_{1}$. Then $D'_{N,i}\subseteq E_{1}\cup E_{2}\cup E_{3}$.

Since
\begin{equation*}
    \begin{split}
        \qquad &
        m(E_{1})=\frac{2\pi}{CN}((r_{i,\sup}+10)^{2}-(r_{i,\inf}-10)^{2})
        =\frac{2\pi}{CN}(r_{i,\sup}+r_{i,\inf})(r_{i,\sup}-r_{i,\inf}+20)
       \\& = \frac{2\pi}{CN}2CN(CN\cdot\frac{2\pi}{CN}+20)=O(1)
    \end{split}
\end{equation*}
and
$$m(E_{i})=10 (r_{i,\sup}-r_{i,\inf}+20)\leq 10 (CN\cdot\frac{2\pi}{CN}+20)=O(1)$$
for $i=2,3$. The claim follows.

\


Note that if $x\in B_{N}\Delta (S_{N}+v)$. Then there exists $1\leq i\leq CN$ such that the distance between $x$ and the boundary
$$\partial S_{N,i}:=\{(Nr(\theta)\cos\theta,Nr(\theta)\sin\theta)\colon \theta\in I_{i}\}$$ is at most 2. Since $\partial  S_{N,i}\subseteq D_{N,i}$, we have that the distance between $x$ and $D_{N,i}$ is at most 2 and thus $x\in\cup_{i=1}^{CN}D'_{N,i}$.
Thus, by the previous claim we obtain the desired bound
\begin{equation*}
    \begin{split}
        %\qquad &
        m(B_{N}\Delta (S_{N}+v))\leq \sum_{i=1}^{CN}m(D'_{N,i})=O(CN)%\sum_{i=1}^{CN}\frac{2\pi}{CN}((r_{i,\sup}+100)^{2}-(r_{i,\inf}-100)^{2})
        %\\&=2\pi \sum_{i=1}^{CN}(r_{i,\sup}+r_{i,\inf})(r_{i,\sup}-r_{i,\sup}+200)
       % \leq 2\pi \sum_{i=1}^{CN}2CN(CN\cdot\frac{2\pi}{CN}+200)=O(CN)
    \end{split}
\end{equation*}
and we are done. %So it remains to prove the claim.
   %  }
 \end{proof}
 
%\WS{Throughout the paper, sometimes we write $\Q(\sqrt{-d})$ and sometimes we write $\Q(\sqrt{-d})$}

\begin{lemma}\label{countingCRT}
Let $K=\Q(\sqrt{-d})$ for some squarefree $d\in\N$, $k, l \in \N, a,Q\in\mathcal{O}_K^{\times}$, and $I$ be an ideal of $\OK$ such that $(Q)$ is coprime to $I$. Let $B_N:=\{ u\in\mathcal{O}_{K} : \NN(u) \leq N\}$. There exists $C>0$ depending only on $d$   such that for any $N>0$, we have that 
\begin{equation}\nonumber%\label{newieq}
    \Bigl\vert \frac{\gamma_{K} N}{\NN(I)}-\vert(QB_N+a)\cap I|\Bigr\vert\leq C\sqrt{N/\NN(I)}.
\end{equation}
\end{lemma}

%\WS{
%Sketch of the proof: 
%We think of elements in $\Z[\tau_{d}]$ as points in $\Z^{2}$, which is embedded in $\R^{2}$.
%Let $W=Q^{-1}I$ viewed as a subset of $\R^{2}$. Then $W$ is a lattice in $\R^{2}$. There exists a linear transformation $T\colon\R^{2}\to\R^{2}$ such that $T(W)=\Z^{2}$. Let $B'=T(B_{1})$. Then $B'$ is an ellipse. It suffices to show that for any shift $x$ and dilation $\sqrt{N}$, the difference between the area of $B'_{N}:=\sqrt{N}\cdot B'$ and the number of integer points in $B'_{N}+x$ is bounded by $C\sqrt{N}$.

%}


\iffalse

\begin{lemma}\label{countingCRT}
Let $K=\Q(\sqrt{-d})$ for some squarefree $d\in\N$, $k, l \in \N, a,Q\in\mathcal{O}_K^{\times}$, and $\mathfrak{p}, \mathfrak{q}$ be distinct prime ideals of $\OK$ such that $Q\notin\mathfrak{p}$ and $Q\notin\mathfrak{q}$. Let $B_N:=\{ u\in\mathcal{O}_{K} : \NN(u) \leq N\}$. We have
%\begin{equation}\label{countingCRTeq1}
%\lim_{N\to \infty} \frac{| \{u \in QB_N+a : Q \mid u-a, u\in\mathfrak{q}^l \cap\mathfrak{p}^k\}|}{\NN(Q)\gamma_{K}_{K} N} = \frac{1}{\NN(Q\mathfrak{p}^k) \NN(\mathfrak{q}^l)},
%\end{equation}
\begin{equation}\label{countingCRTeq1}
\lim_{N\to \infty} \frac{| (QB_N+a)\cap\mathfrak{q}^l \cap\mathfrak{p}^k|}{\gamma_{K} N}=\frac{1}{\NN(\mathfrak{p}^k) \NN(\mathfrak{q}^l)}.
\end{equation}
\WS{In order to proof Turan-Kubilius as on page 303 of the book "Introduction to analytic and Probabilistic number theory", we need the following: there is a universal $C$ such that for all $N$ and ideal $I$,
$$0\leq \frac{\gamma_{K} N}{\NN(I)}-\vert(QB_N+a)\cap I\backslash\{0\}| \leq C.$$
(when $(a,Q)=1$, then the 0 does not matter. But when $a=0$, then the 0 matter)
 For example, in the integer case, $\vert B_{N}\cap p\Z\backslash\{0\}\vert=2[N/p]$ and $\gamma=2$. So
 $$0\leq 2N/p-[2N/p]\leq 1.$$


We think of elements in $\Z[\tau_{d}]$ as points in $\Z^{2}$, which is embedded in $\R^{2}$.
Let $W=Q^{-1}I$ viewed as a subset of $\R^{2}$. Then $W$ is a lattice in $\R^{2}$. There exists a linear transformation $T\colon\R^{2}\to\R^{2}$ such that $T(W)=\Z^{2}$. Let $B'=T(B_{1})$. Then $B'$ is a Jordan measurable set. It suffices to show that for any shift $x$ and dilation $N$, the areas of $B'_{N}:=N\cdot B'$ minus the number of integer points in $B'_{N}+x$ is between -1 and C, and the lowder bound can be upgraded to 0 if $x\notin\Z^{2}$.

}
\FM{We don't have to worry about the $a=0$ case because we take $a, Q \in \mathcal{O}_K^{\times}$.}
%where $\gamma_{K}_{K}:=\lim_{N \to \infty} \frac{1}{N}\sum_{\NN(u) \leq N} 1$.\WS{explain why the limit $\gamma_{K}_{K}$ exists}

\WS{At a second thought, I think this statement maybe incorrect. For example, the number of integer points in a ball of raduis $\sqrt{N}$ is approximately $\pi N+O(\sqrt{N})$. It is possible to lower the error term $\O(\sqrt{N})$ a little, but one can not reduce it to $O(1)$.} {\color{red} In the proof below, I will write in red on the part that is affacted by this issue (now we only assume that 
\begin{equation}\label{newieq}
    \vert \frac{\gamma_{K} N}{\NN(I)}-\vert(QB_N+a)\cap I\backslash\{0\}|\vert\leq C\sqrt{N/\NN(I)}
\end{equation}
 
)}
\end{lemma}

\fi

%}

%\begin{lemma}\label{countingCRT}
%Let $I$ be an ideal in $\mathcal{O}$, let $K \in \N$, and suppose that $K$ is large enough so that $\Phi_K \subset I$. Consider natural numbers $k, l \in \N$, let $Q:=\prod_{\NN(\mathfrak{p})\leq K} \mathfrak{p}$, and let $\mathfrak{p}, \mathfrak{q}$ be distinct prime ideals such that $\mathfrak{p}, \mathfrak{q} \nmid Q$. We have
%\begin{equation}\label{countingCRTeq1}
%\lim_{N\to \infty} \frac{| \{u \text{ with } \NN(u-a)\leq \NN(Q)N : Q \mid u-a, \mathfrak{q}^l \mid u \text{ and } \mathfrak{p}^k \mid u\}}{\NN(Q)A_K N} = \frac{1}{\NN(Q\mathfrak{p}^k) \NN(\mathfrak{q}^l)},
%\end{equation}
%where $A_K$ is the constant such that $\frac{1}{K_1N}\sum_{\NN(u) \leq N}1 \to 1$ as $N \to \infty$.
%\end{lemma}
\begin{proof}
%Since the balls $B_N:=\{ u \in \OK: \NN(u)\leq N\}$ are asymptotically translation invariant, and given the definition of $\gamma_{K}$ above, we have 
%$$\lim_{N\to\infty}\frac{|\{u \in \OK: \NN(\frac{u-a}{Q})\leq N\}|}{\gamma_{K} N}=\NN(Q)$$
%\WS{give a reference?}
%We may rewrite\footnote{We use the notation $a \equiv u \pmod Q$ to mean that $a-u \in (Q)$.} the left hand side of ?? as 
%  \[\NN(Q)\cdot\lim_{N\to\infty}\E_{u \in \OK: \NN(\frac{u-a}{Q})\leq N} \mathds{1}_{\{0\}}(u \pmod{\mf q^l}) \mathds{1}_{\{0\}
%}(u \pmod{ \mf p^k}) \mathds{1}_{\{a\}}(u \pmod Q).\]
%
% We begin by rewriting the left hand side of \eqref{countingCRTeq1} as 
%
%\[
%\frac{1}{\NN(Q)\gamma_{K} N}\sum_{u \in QB_N+a} \mathds{1}_{\{0\}}(u \pmod{\mf q^l}) \mathds{1}_{\{0\}
%}(u \pmod{ \mf p^k}) \mathds{1}_{\{a\}}(u \pmod Q)\]
%This is the same as counting all the $u$ in the shifted ball $\NN(\frac{u-a}{Q})\leq N$ such that they have residue class $(a,0,0)$ in $\mathcal{O}_K/(Q) \times \mathcal{O}_K/(\mf p^k) \times \mathcal{O}_K/(\mf q^l)$ (where we take the ideal generated by each element).

%\WS{here and throughout, we need to define the meaning of $a\equiv u \mod Q$. It means that $a-u\in Q\Z[\tau_{d}]$, or $a-u\in (Q)$.}

%\WS{The statement of the lemma has been updated. Perhaps Lemma 10.1 and 10.2 of "book K" can be used to prove this lemma}
%By the Chinese Remainder Theorem, and given that $\mathcal{O}_K$ is a ring of integers, the product ring above is isomorphic to $\mathcal{O}_K/\langle (Q)\mf p^k\cdot \mf q^l \rangle$, since the ideals are relatively prime by assumption. Moreover, it is known that the size of this latter quotient ring is exactly $\NN(Q) \NN(\mf p^k)\NN( \mf q^l)$. 
%Now the proof can easily be finished as a trivial application of the mean ergodic theorem, since the $\mathcal{O}_K$-action by addition on the quotient ring is ergodic, the limit displayed above will converge to the integral of the characteristic function of the point in the aforementioned quotient ring $\mathcal{O}_K/\langle(Q)\mf p^k \cdot \mf q^l\rangle$. This trivially has size $\frac{1}{\NN(Q) \NN(\mf p^k) \NN(\mf q^l)}$, and we are done.

%\WS{I wrote an outline for a new proof. Check the red part}

 Given that $I$ is an ideal in $K$, an imaginary extension or $\Q$, it follows that $I$ is a free $\Z$-module of rank 2. Thus, we may write $I$ as 
     $$I=\{m\alpha+n\beta\colon m,n\in\Z\}$$
     for some $\alpha=a_{1}+b_{1}\tau_{d},\beta=a_{2}+b_{2}\tau_{d}\in\Z[\tau_{d}]$. As a consequence of Minkowski's bound (see \cite[Theorems 6.6 and 7.4]{neukirch99}), we see that the fundamental parallelepiped that generates the ideal lattice $I$ has volume proportional to the norm of the ideal, where the constant only depends on the number field $K$. Thus, we may take generators $\alpha, \beta$ whose norm is $O(\NN(I))$. % \FM{See if we can derive this directly from \cite{S23} ?}
     This means that we may further assume that $\vert a_{i}\vert, \vert b_{i}\vert\leq C\sqrt{\NN(I)}$ for some universal constant $C$ depending only on $d$. %{\color{red} give a proof}. 
     Since $(Q)$ is coprime to $I$, there exist unique $(m_{0},n_{0})\in\{0,\dots,Q-1\}^{2}$ such that $Q\vert m\alpha+n\beta-a$ if and only if $m\equiv m_{0} \mod Q$ and $n\equiv n_{0} \mod Q$.
     Therefore, the set $(QB_{N}+a)\cap I$ consists of elements of the form $(Qm+m_{0})\alpha+(Qn+n_{0})\beta$ with $m,n\in\Z$ such that $\NN((Qm+m_{0})\alpha+(Qn+n_{0})\beta-a)\leq Q^{2}N$.
     
  %   Let $S$ be the set of $(m,n)\in\R^{2}$ such that $\NN(Qm\alpha+Qn\beta)\leq Q^{2}$. Then 
   %  $\vert (QB_{N}+a)\cap I\vert=\vert (S_{\sqrt{N}}+a-(m_{0}\alpha+n_{0}\beta))\cap\Z^{2}\vert.$ Let $T\colon \Q^{2}\to\Q^{2}$ be the bijective linear transformation given by 
    %  $$T(m,n):=\iota(Qm\alpha+Qn\beta),$$
     % where $\iota\colon\Q(\sqrt{-d})\to\Q^{2}$ is the natural map given by $\iota(m+n\tau_{d}):=(m,n)$.
     % Let $S':=\{(m,n)\in\R^{2}\colon m^{2}+\vert\tau_{d}^{2}\vert n^{2}\leq Q^{2}\}$. Clearly $S=T^{-1}S'$.
     % Thus 
     % $$m(S)=(Q^{2}\NN(I))^{-1}m(S')=\gamma_{K}/\NN(I).$$


      % Using polar coordinates for ellipses, it is not hard to see that the Lipschitz consant of an ellipse is bounded by its diameter. So $S'$ is $O_{d,Q}(1)$-Lipschitz. Since $\vert a_{i}\vert, \vert b_{i}\vert\leq O_{d}(\sqrt{\NN(I)})$, we have that $S$ is $O_{d,Q}(\sqrt{\NN(I)})$-Lipschitz. By Lemma \ref{dafp}, 
      %$$\Bigl\vert\vert (QB_{N}+a)\cap I\vert-\frac{\gamma_{K}N}{\NN(I)}\Bigr\vert=\Bigl\vert\vert (S_{\sqrt{N}}+a-(m_{0}\alpha+n_{0}\beta))\cap \Z^{2}\vert-N\cdot m(S)\Bigr\vert\leq O_{d,Q}\left(\sqrt{\frac{N}{\NN(I)}}\right).$$
      
      Let $S$ be the set of $(m,n)\in\R^{2}$ such that $\NN(m\alpha+n\beta)\leq 1$ (where $\NN$ is the canonical extension of the norm on $\Q(\sqrt{-d})$ to $\C$). Then 
     $$\vert (QB_{N}+a)\cap I\vert=\Bigl\vert (S_{\sqrt{N}}+\frac{1}{Q}(a-(m_{0}\alpha+n_{0}\beta)))\cap\Z^{2}\Bigr\vert.$$ Let $T\colon \Q^{2}\to\Q^{2}$ be the bijective linear transformation given by 
      $$T(m,n):=\iota(m\alpha+n\beta),$$
      where $\iota\colon\Q(\sqrt{-d})\to\Q^{2}$ is the natural map given by $\iota(m+n\tau_{d}):=(m,n)$.
      Let $S':=\{(m,n)\in\R^{2}\colon \NN(m+\tau_{d} n)\leq 1\}$. Clearly $S=T^{-1}S'$.
      Thus 
      $$m(S)=\NN(I)^{-1}m(S')=\gamma_{K}/\NN(I).$$
      
      
       Using polar coordinates for ellipses, it is not hard to see that the Lipschitz consant of an ellipse is bounded by its diameter. So $S'$ is $O_{d}(1)$-Lipschitz. Since $\vert a_{i}\vert, \vert b_{i}\vert\leq O_{d}(\sqrt{\NN(I)})$, we have that $S$ is $O_{d}(\sqrt{\NN(I)})$-Lipschitz. By Lemma~\ref{dafp}, we have 
      \[\left\vert\vert (QB_{N}+a)\cap I\vert-\frac{\gamma_{K}N}{\NN(I)}\right\vert=\Bigl\vert\vert (S_{\sqrt{N}}+\frac{1}{Q}(a-(m_{0}\alpha+n_{0}\beta)))\cap \Z^{2}\vert-N\cdot m(S)\Bigr\vert\leq O_{d}\left(\sqrt{\frac{N}{\NN(I)}}\right),\]
      as was to be shown.
\end{proof}
%{\color{red} This proof may need some modifications for arbitrary imaginary quadratic number fields, since we may not have a good equivalent of the CRT.}

 
\begin{proof}[Proof of Theorem \ref{TK}]
%%Let 
%$$c_{M,K,Q}:=\sum_{M<\NN(\mf p^{k})\leq N,\mf p\nmid Q}h(\mf p^{k})\NN(\mf p)^{-k}(1-\NN(\mf p)^{-1})$$ 
%It suffices to show that 
%$\E_{\NN(u) \leq N}h(Qu+a)\approx c_{M,K,Q}$ and $\E_{\NN(u) \leq N}\vert h(Qu+a)\vert^{2}\approx c_{M,K,Q}^{2}$.
% $B_N:=\{ u\in\OK : \NN(u) \leq N\}$.
%Then
%Throughout the proof we allow implicit constants to be dependent on $Q$. 
%\FM{Do we need to be careful about this when we want to use it in Section 6? Or is the comment that we wrote yesterday enough to guarantee independence from $Q$, i.e., should we say something in Section 6 like, the $O_Q()$ that appears in the statement of prop 5.4 is such that $Q$ can be forgotten?}\WS{I think this is an issue. 
%To overcome this difficulty, I changed the RHS of (49).}
We first assume that $h$ is non-negative.
 Note that
\begin{equation}\label{firstTK}
    \begin{split}
        &\quad \E_{\NN(u) \leq N}h((Qu+a))
        =\frac{1}{\gamma_{K} N}\sum_{\NN(\frac{u-a}{Q})\leq N, u\equiv a \mod Q}h((u))
        \\&\quad\quad\quad\quad\quad\quad\quad\quad\quad\quad\;\;=\frac{1}{\gamma_{K} N}\sum_{\NN(\frac{u-a}{Q})\leq N, u\equiv a \mod Q}\sum_{\mf p^{k}|| u} h(\mf p^{k})
        \\&\quad\quad\quad\quad\quad\quad\quad\quad\quad\quad\;\;=\frac{1}{\gamma_{K} N}\sum_{\NN(\frac{u-a}{Q})\leq N, u\equiv a \mod Q}\sum_{\mf p^{k}|| u, \mf p\nmid Q} h(\mf p^{k}).
  %      \\&=\frac{1}{\gamma_{K} N}\sum_{\NN(\mf p^{k})\leq N', \mf p\nmid Q}h(\mf p^{k})\sum_{\NN(\frac{u-a}{Q})\leq N, \mf p^{k}||u, u\equiv a \mod Q} 1,
     %   \\&=\frac{1}{N^2}\sum_{p^{k}\in  [0,QN+a]^2, p\nmid Q}h(p^{k})(\sum_{u\in [0,QN+a]^2, p^{k}|u, u\equiv a \mod Q} 1-\sum_{u\in [0,QN+a]^2, p^{k+1}|u, u\equiv a \mod Q} 1)
  %  \\&\approx \frac{1}{N^2}\sum_{K<p^{k}\leq QN+a, p\nmid Q}h(p^{k})(\frac{QN+a}{p^{k}Q}-\frac{QN+a}{p^{k+1}Q})
  %  \\& =\sum_{\NN(\mf p^{k})\leq \NN(Q)N, \mf p\nmid Q}h(\mf p^{k})\NN(\mf p)^{-k}(1-\NN(\mf p)^{-1})+\Theta_{1}+\Theta_{2},
    \end{split}
\end{equation}
%where the upper bound $2\NN(Q)(N+1)$ in the last line follows from the estimate that if $\mf p^{k}\vert u$ for some $u$ with $\NN(\frac{u-a}{Q})\leq N$, then
%$$\NN(\mf p^{k})\leq \NN(u)\leq 2(\NN(u-a)+\NN(a))\leq 2(\NN(Q)N+\NN(Q))=2\NN(Q)(N+1).$$
%
 %
 Note that if $\mf p^{k}\vert u$ for some $u$ with $\NN(\frac{u-a}{Q})\leq N$, then
$$\NN(\mf p^{k})\leq \NN(u)\leq 2(\NN(u-a)+\NN(a))\leq 2(\NN(Q)N+\NN(Q))=N'.$$ 
By Lemma~\ref{countingCRT},
\begin{equation}\label{firstTK2}
    \begin{split}
        &\quad 
        \frac{1}{\gamma_{K} N}\sum_{\NN(\frac{u-a}{Q})\leq N, u\equiv a \mod Q}\sum_{\mf p^{k}|| u, \mf p\nmid Q} h(\mf p^{k})
        \\&=\frac{1}{\gamma_{K} N}\sum_{\NN(\mf p^{k})\leq N', \mf p\nmid Q}h(\mf p^{k})\sum_{\NN(\frac{u-a}{Q})\leq N, \mf p^{k}||u, u\equiv a \mod Q} 1
    \\& =\frac{1}{\gamma_{K} N}\sum_{\NN(\mf p^{k})\leq N', \mf p\nmid Q}h(\mf p^{k})\left(\sum_{\NN(\frac{u-a}{Q})\leq N, \mf p^{k}|u, u\equiv a \mod Q} 1-\sum_{\NN(\frac{u-a}{Q})\leq N, \mf p^{k+1}|u, u\equiv a \mod Q} 1\right)
    \\&=\sum_{\NN(\mf p^{k})\leq N', \mf p\nmid Q}h(\mf p^{k})\left((\NN(\mf p)^{-k}-\NN(\mf p)^{-(k+1)})+O\left(\frac{1}{\sqrt{N\NN(\mf p^{k})}}\right)\right).
    \end{split}
\end{equation}
%Finally, 
%\begin{equation}\label{firstTK3}
%    \begin{split}
 %       &\quad \frac{1}{\gamma_{K} N}\sum_{N<\NN(\mf p^{k})\leq 2\NN(Q)(N+1), \mf p\nmid Q}h(\mf p^{k})\sum_{\NN(\frac{u-a}{Q})\leq N, \mf p^{k}||u, u\equiv a \mod Q} 1
%    \\& \leq \frac{1}{\gamma_{K} N}\sum_{N<\NN(\mf p^{k})\leq \NN(Q)N, \mf p\nmid Q}h(\mf p^{k})\cdot 6(2\NN(Q)+1),
%    \end{split}
%\end{equation}
%where the constant 6 comes from the fact that $K$ has at most 6 units by Lemma \ref{l32}.
%Combining (\ref{firstTK}) and (\ref{firstTK2}), %and (\ref{firstTK3}), 
%we have that 
%\begin{equation}\nonumber
%        \left|\E_{u\in\OK\colon\NN(u) \leq N}h((Qu+a))-A_{N'}\right|\leq \frac{C}{N}\cdot\sum_{\NN(\mf p^{k})\leq N', \mf p\nmid Q}\vert h(\mf p^{k})|
%\end{equation}
%\FM{Why do we have the bound
%\[
%\frac{1}{\sqrt{N}}\sum_{\NN(\mf p^k) \leq N', \mf p \nmid Q} \frac{|h(\mf p^k)|}{\sqrt{\NN(\mf p^k)}}  \leq \frac{1}{N} \sum_{\NN(\mf p^k) \leq N', \mf p \nmid Q} |h(\mf p^k)|?
%\]
%I believe the only bound we have for $\NN(\mf p^k)$ under these conditions is that it must be at least $M$, so it should be $\sqrt{NM}$ in the denominator, I believe.
%}
for some $C>0$ depending only on $K$ and $Q$. Therefore,
%\FM{The first inequality here seems to confirm what I said in the comment above, in which case I would just rewrite the inequality to just have the first term I wrote in purple.}
combining (\ref{firstTK}) and (\ref{firstTK2}), we have that
\begin{equation}\label{firstTK0}
    \begin{split}
       &\qquad \left| A_{Q,N}\cdot\left(\E_{u\in\OK\colon\NN(u) \leq N}h((Qu+a))-A_{N'}\right)\right|
       \\&\ll \frac{1}{\sqrt{N}}\sum_{\NN(\mf p^{k})\leq N', \mf p\nmid Q}\vert h(\mf p^{k})\vert \cdot \NN(\mf p^{k})^{-1/2}\cdot \sum_{\NN(\mf p^{k})\leq N', \mf p\nmid Q}\vert h(\mf p^{k})\vert\cdot \NN(\mf p^{k})^{-1}
        \\&\leq \frac{1}{\sqrt{N}}\sum_{\NN(\mf p^{k})\leq N', \NN(\mf p)>M}\vert h(\mf p^{k})\vert\cdot \NN(\mf p^{k})^{-1/2}\cdot \sum_{\NN(\mf p^{k})\leq N', \NN(\mf p)>M}\vert h(\mf p^{k})\vert\cdot \NN(\mf p^{k})^{-1}
          \\&\leq \frac{1}{\sqrt{N}}\left(B_{N'}\cdot\sum_{\NN(\mf p^{k})\leq N'} 1\right)^{1/2}\cdot\left(B_{N'}\cdot\sum_{\NN(\mf p^{k})\leq N'}\NN(\mf p^{k})^{-1}\right)^{1/2}
          \\&=B_{N'}\cdot \frac{1}{\sqrt{N}}\left(\sum_{\NN(\mf p^{k}),\NN(\mf q^{\ell})\leq N'}\NN(\mf q^{\ell})^{-1}\right)^{1/2}.
          \end{split}
\end{equation}




Now we estimate the last term of (\ref{ltk}). We have
\begin{equation}\label{secondTK}
    \begin{split}
        &\quad \E_{\NN(u) \leq N}\vert h((Qu+a))\vert^{2}
        =\frac{1}{\gamma_{K} N}\sum_{\NN(\frac{u-a}{Q})\leq N, u\equiv a \mod Q}\vert h((u))\vert^{2}
        \\&=\frac{1}{\gamma_{K} N}\sum_{\NN(\frac{u-a}{Q})\leq N, u\equiv a \mod Q}\left| \sum_{\mf p^{k}|| u} h(\mf p^{k})\right|^{2}
        \\&=\frac{1}{\gamma_{K} N}\sum_{\NN(\frac{u-a}{Q})\leq N, u\equiv a \mod Q}\left(\sum_{\mf p^{k}|| u} \vert h(\mf p^{k})\vert^{2}+\sum_{\mf p^{k},\mf q^{\ell}|| u, \mf p\neq \mf q} h(\mf p^{k})h(\mf q^{\ell})\right)
          \\&=\frac{1}{\gamma_{K} N}\sum_{\NN(\frac{u-a}{Q})\leq N, u\equiv a \mod Q}\left(\sum_{\mf p^{k}|| u, \mf p\nmid Q} \vert h(\mf p^{k})\vert^{2}+\sum_{\mf p^{k},\mf q^{\ell}|| u, \mf p\neq \mf q, \mf p,\mf q\nmid Q} h(\mf p^{k})h(\mf q^{\ell})\right).
    \end{split}
\end{equation}
The first term in the last line of \eqref{secondTK} can be dealt with in the same way as in  \eqref{firstTK2} %and \eqref{firstTK3} 
(replacing $h$ by $h^{2}$), namely 
\begin{equation}\label{secondTK2}
    \begin{split}
        &\quad 
        \frac{1}{\gamma_{K} N}\sum_{\NN(\frac{u-a}{Q})\leq N, u\equiv a \mod Q}\sum_{\mf p^{k}|| u, \mf p\nmid Q} h(\mf p^{k})^{2}
 %   \\&=\sum_{\NN(\mf p^{k})\leq N', \mf p\nmid Q}h(\mf p^{k})^{2}((\NN(\mf p)^{-k}-\NN(\mf p)^{-(k+1)})+O(\frac{1}{N}))
 %   \\&\leq B_{N'}+\frac{C}{N}\sum_{\NN(\mf p^{k})\leq N', \mf p\nmid Q}h(\mf p^{k})^{2}
  %  \\&\leq B_{N'}+\frac{C}{N}\sum_{\NN(\mf p^{k})\leq N', \NN(\mf p)>M}h(\mf p^{k})^{2}
     \\&=\frac{1}{\gamma_{K} N}\sum_{\NN(\mf p^{k})\leq N', \mf p\nmid Q}h(\mf p^{k})^{2}\sum_{\NN(\frac{u-a}{Q})\leq N, \mf p^{k}||u, u\equiv a \mod Q} 1
     \\&=\sum_{\NN(\mf p^{k})\leq N', \mf p\nmid Q}h(\mf p^{k})^{2}\NN(\mf p)^{-k}+O(1)\cdot\frac{1}{\sqrt{N}}\sum_{\NN(\mf p^{k})\leq N', \mf p\nmid Q}h(\mf p^{k})^{2}\cdot \NN(\mf p^{k})^{-1/2}
      \\&\leq\sum_{\NN(\mf p^{k})\leq N', \mf p\nmid Q}h(\mf p^{k})^{2}\NN(\mf p)^{-k}+O(C_{N}).%\Theta,
    \end{split}
\end{equation}
%\WS{The error term $\Theta$ is replaced by $C_{N}$. I do not know how to deal with $C_{N}$ in this proposition. But later I will deal with it since $h$ is bounded in our applications (see (\ref{208})).}
\iffalse
where 
$$\vert\Theta\vert\ll \frac{1}{\sqrt{N}}\sum_{\NN(\mf p^{k})\leq N', \mf p\nmid Q}h(\mf p^{k})^{2}\cdot \NN(\mf p^{k})^{-1/2}\leq O(\sqrt{Q}) \sum_{\NN(\mf p^{k})\leq N', \mf p\nmid Q}h(\mf p^{k})^{2}\cdot \NN(\mf p^{k})^{-1}\leq O(\sqrt{Q}B_{N'}).$$
\WS{fixed, but this bound is not good enough. we need the bound to be $\leq(2+o_{N\to\infty}(1),Q) B_{N'}$. This can be done easily when $h$ is bounded on powers of prime ideals (which is our case). But can we prove it for the general case?}
\FM{I don't understand this step because $\NN(\mf p^k)^{-1/2} \leq \NN(\mf p^k)^{-1}$ is if and only if $\NN(\mf p^k)^{1/2}\leq 1$. And if you are using some sort of Cauchy-Schwartz I think the squares etc, don't add up.}
\WS{this term probablity should be $o(B_{N'})$. If so then the RHS of (\ref{ltk})
     can be upgraded to $\leq(2+o_{N\to\infty}(1),Q) B_{N'}$, i.e. the main term is independent of $Q$ and the moreover the exact constant is 2. This is the case when $K=\Q$.} 
\fi


For the second term in \eqref{secondTK}, similarly to (\ref{firstTK}), we also swap the order of the sums and rewrite it as
%\begin{equation}
 %   \begin{split}
  %      &\quad \frac{1}{(2N)^2}\sum_{p^{k}, q^l \in Q[\pm N]^2+a, p \neq q, p, q \in \Z[i]_+}\vert h(p^{k})\vert^{2}\sum_{u\leq QN+a, u\equiv a \mod Q, p^{k}\Vert u} 1 
   %     \\&\approx \frac{1}{N}\sum_{p^{k}\leq QN+a}\vert h(p^{k})\vert^{2}p^{-k}(1-p^{-1}) 
    %    \\&\leq \frac{1}{N}\sum_{p^{k}\leq QN+a}\vert h(p^{k})\vert^{2}p^{-k}
   % \end{split}
%\end{equation}
%(the RHS of (\cite{20}) is similar to this.)
% On the other hand, we have
\begin{equation}\nonumber
    \begin{split}
        &\quad \frac{1}{\gamma_{K} N}\sum_{\NN(\mf p^{k}\mf q^{\ell}) \leq N', \mf p\neq \mf q,\mf p,\mf q\nmid Q} h(\mf p^{k})h(\mf q^{\ell})\sum_{\NN(\frac{u-a}{Q})\leq N, u\equiv a \mod Q, \mf p^{k}||u, \mf q^{\ell}||u}1.
    \end{split}
\end{equation}
As was the case in \eqref{firstTK2}, we have that
\begin{equation}\label{secondTK3}
    \begin{split}
        &\quad \frac{1}{\gamma_{K} N}\sum_{\NN(\mf p^{k}\mf q^{\ell}) \leq N', \mf p\neq \mf q,\mf p,\mf q\nmid Q} h(\mf p^{k})h(\mf q^{\ell})\sum_{\NN(\frac{u-a}{Q})\leq N, u\equiv a \mod Q, \mf p^{k}||u, \mf q^{\ell}||u}1
        \\&=\sum_{\NN(\mf p^{k}\mf q^{\ell})\leq N', \mf p\neq \mf q, \mf p,\mf q\nmid Q} h(\mf p^{k})h(\mf q^{\ell})\left(\NN(\mf p^k)^{-1}\NN(\mf q^l)^{-1}\left(1-\frac{1}{\NN(\mf p)}\right) \left(1-\frac{1}{\NN(\mf q)}\right)\right. \\&\quad\quad\quad\quad\quad\quad\quad\quad\quad\quad\quad\quad\quad\quad\quad\quad\quad\quad\quad\quad\quad\quad\quad\quad\quad\left.+O\left(\frac{1}{\sqrt{N\NN(\mf p^{k}\mf q^{\ell})}}\right)\right)
        \\&\leq A_{Q,N'}^{2}+\frac{C}{\sqrt{N}}\sum_{\NN(\mf p^{k}\mf q^{\ell})\leq N', \mf p\neq \mf q, \NN(\mf p),\NN(\mf q)>M} h(\mf p^{k})h(\mf q^{\ell})\NN(\mf p^{k})^{-1/2}\NN(\mf q^{\ell})^{-1/2}     %  \\&=\vert\sum_{K<p^{k}\leq QN+a, p\nmid Q}h(p^{k})p^{-k}(1-p^{-1})\vert^{2}+error,
       % \\&=A_{N'}^{2}+\frac{C}{\sqrt{N}}\sum_{\NN(\mf p^{k}\mf q^{\ell})\leq N', \mf p\neq \mf q, \NN(\mf p),\NN(\mf q)>M} \frac{h(\mf p^{k})}{\NN(\mf p^{k})^{1/2}}\cdot\frac{h(\mf q^{\ell})}{\NN(\mf q^{\ell})^{1/2}}\cdot \NN(\mf p^{k})^{1/2}\cdot \NN(\mf q^{\ell})^{1/2}
        \\&\leq A_{Q,N'}^{2}+B_{N'}\cdot\frac{C}{\sqrt{N}}\left(\sum_{\NN(\mf p^{k}\mf q^{\ell})\leq N', \mf p\neq \mf q} 1\right)^{1/2}.
            \end{split}
\end{equation}

Combining (\ref{secondTK}), (\ref{secondTK2}) and (\ref{secondTK3}), we have that
\begin{equation}\label{secondTK0}
    \begin{split}
        &\quad \E_{\NN(u) \leq N}\vert h((Qu+a))\vert^{2}
        \\&\leq %O(\sqrt{Q}B_{N'})
        B_{N'}+A_{Q,N'}^{2}+O(C_{N})+\frac{C}{N}\sum_{\NN(\mf p^{k})\leq N', \NN(\mf p)>M}h(\mf p^{k})^{2}
        \\&
        \quad\quad\quad\quad\quad\quad\quad\quad\quad+B_{N'}\cdot\frac{C}{N}\left(\sum_{\NN(\mf p^{k}\mf q^{\ell})\leq N', \mf p\neq \mf q} \NN(\mf p^{k})\cdot \NN(\mf q^{\ell})\right)^{1/2}.
    \end{split}
\end{equation}   

\iffalse
Finally, note that 
 \begin{equation}\label{snnn2}
    \begin{split}
        &\quad A_{N'}-A_{N}
        =\sum_{N<\NN(\mf p^{k})\leq N',\mf p\nmid Q}h(\mf p^{k})\NN(\mf p)^{-k}(1-\NN(\mf p)^{-1})
        \\&\leq \frac{2}{N}\sum_{N<\NN(\mf p^{k})\leq N',\mf p\nmid Q}h(\mf p^{k})
        \leq \frac{2}{N}\sum_{N<\NN(\mf p^{k})\leq N',\mf \NN(\mf p)>M}h(\mf p^{k})
        \\&\leq B_{N'}^{1/2}\cdot \frac{2}{N}\left(\sum_{N<\NN(\mf p^{k})\leq N'}\NN(\mf p^{k})\right)^{1/2}
        \leq B_{N'}^{1/2}\cdot \frac{2}{N}\left(N'\sum_{N<\NN(\mf p^{k})\leq N'}1\right)^{1/2}
        \ll B_{N'}^{1/2}, 
        \end{split}
\end{equation} 
where the last inequality follows from Lemma \ref{mv8}.  
 \fi

Finally, by Part (i) of Lemma \ref{l2}, 
\begin{equation}\label{snnn2}
    \begin{split}
        &\quad A_{Q,N'}-A_{Q,N}
        =\sum_{N<\NN(\mf p^{k})\leq N',\mf p\nmid Q}h(\mf p^{k})\NN(\mf p)^{-k}(1-\NN(\mf p)^{-1})
        \\&\leq \frac{2}{N}\sum_{N<\NN(\mf p^{k})\leq N',\mf p\nmid Q}h(\mf p^{k})
        \leq \frac{2}{N}\sum_{N<\NN(\mf p^{k})\leq N',\mf \NN(\mf p)>M}h(\mf p^{k})
        \\&\leq B_{N'}^{1/2}\cdot \frac{2}{N}\left(\sum_{N<\NN(\mf p^{k})\leq N'}\NN(\mf p^{k})\right)^{1/2}
        \leq B_{N'}^{1/2}\cdot \frac{2}{N}\left(N'\sum_{N<\NN(\mf p^{k})\leq N'}1\right)^{1/2}
        \\&\ll B_{N'}^{1/2}\cdot \frac{2}{N}\left(\frac{{N'}^{2}}{\log N'}\right)^{1/2}=o_{Q;N\to\infty}(B_{N'}^{1/2}).
        \end{split}
\end{equation} 
Combining %(\ref{ltk2}), 
(\ref{firstTK0}), (\ref{secondTK0}) and (\ref{snnn2}), we have
%the left hand side of (\ref{ltk}) can be bounded by
%
% \begin{equation}\label{snnn}
%    \begin{split}
%        &\quad \E_{\NN(u) \leq N}\left| h((Qu+a))-A_{N'}\right|^{2}  
%        \\&=-A_{N'}^{2}+2A_{N'}(\E_{\NN(u) \leq N}h((Qu+a))-A_{N'})+\E_{\NN(u) \leq N}\vert h((Qu+a))\vert^{2}
 %       \\&\leq B_{N'}(1+\frac{1}{N}\Bigl(\sum_{\NN(\mf p^{k}),\NN(\mf q^{\ell})\leq N'}\NN(\mf p^{k})\NN(\mf q^{\ell})^{-1}\Bigr)^{1/2}+???+\frac{C}{N}\Bigl(\sum_{\NN(\mf p^{k}\mf q^{\ell})\leq N', \mf p\neq \mf q} \NN(\mf p^{k})\cdot \NN(\mf q^{\ell})\Bigr)^{1/2}).
 %       \end{split}
%\end{equation}
%
 \begin{equation}\label{snnn}
    \begin{split}
        &\quad \E_{\NN(u) \leq N}\left| h((Qu+a))-A_{Q,N}\right|^{2}  
        \\&=(A_{Q,N}-A_{Q,N'})^{2}-2A_Q,{N}\left(\E_{\NN(u) \leq N}h((Qu+a))-A_{Q,N'}\right)\\&\quad\quad\quad\quad\quad\quad\quad\quad\quad\quad\quad\quad\quad\quad\quad\quad\quad+\E_{\NN(u) \leq N}\vert h((Qu+a))\vert^{2}-A_{Q,N'}^{2}
        \\&\leq B_{N'}\cdot\left(1+C\left(\frac{1}{N}\sum_{\NN(\mf p^{k}),\NN(\mf q^{\ell})\leq N'}\NN(\mf q^{\ell})^{-1}\right)^{1/2}\right.\\&\quad\quad\quad\quad\quad\quad\quad\quad\quad\left.+C\left(\frac{1}{N}\sum_{\NN(\mf p^{k}\mf q^{\ell})\leq N', \mf p\neq \mf q} 1\right)^{1/2}+o_{Q;N\to\infty}(1)\right)+O(C_{N})
        \\&=B_{N'}(%O(\sqrt{Q}) 
        1+o_{Q;N\to\infty}(1))+O(C_{N}),
        \end{split}
\end{equation}
where the last inequality follows from %Parts (i), (ii) and (iii) of 
Lemma \ref{l2}.

\iffalse
 \textbf{Claim.} The right hand side of (\ref{snnn}) is $B_{N'}(1+o_{N\to\infty}(1))$.

 \iffalse
give a proof. %This is the function $\epsilon(x)$ on page 302 of the book "Introduction to analytic and Probabilistic number theory"

Recall from ??? that 
$$\sum_{p\leq N} 1\ll \frac{N}{\log N} \text{ and } \sum_{p\leq N} \frac{1}{p}\ll \log\log N,$$
where $p$ is taken over primes in $\N$.

Note that $\NN(\mf p)$ is either a prime or a square of a prime. So it follows from Lemma \ref{mv7} that
$$\sum_{\NN(\mf p)\leq N'} 1=\sum_{\NN(\mf p)=p, p\leq N'} 1+\sum_{\NN(\mf p)=p^{2}, p^{2}\leq N'} 1\ll \sum_{p\leq N'} 1+\sum_{p\leq \sqrt{N'}} 1\ll \frac{N'}{\log N'}.$$
Therefore,
$$\sum_{\NN(\mf p^{k})\leq N'} 1=\sum_{k=1}^{\infty} \sum_{\NN(\mf p)\leq {N'}^{1/k}} 1\ll \sum_{k=1}^{\infty} \frac{k{N'}^{1/k}}{\log N'}\ll \frac{N'}{\log N'}.$$
Similarly, by  Lemma \ref{mv7}, 
$$\sum_{\NN(\mf p)\leq N'} \frac{1}{\NN(\mf q)}\leq \sum_{\NN(\mf p)=p, p\leq N'} \frac{1}{p}+\sum_{\NN(\mf p)=p^{2}, p^{2}\leq N'} \frac{1}{p^{2}}\ll\sum_{p\leq N'} \frac{1}{p}+\sum_{p\leq \sqrt{N'}} \frac{1}{p^{2}}\ll \log\log N'.$$
Therefore,
$$\sum_{\NN(\mf q^{\ell})\leq N'} \frac{1}{\NN(\mf q^{\ell})}\leq \sum_{\NN(\mf q^)\leq N'} \sum_{\ell=1}^{\infty}\frac{1}{\NN(\mf q^{\ell})}=\sum_{\NN(\mf q)\leq N'}  \frac{1}{\NN(\mf q)-1}=\sum_{\NN(\mf q)\leq N'}  \frac{1}{\NN(\mf q)}+O(1)\ll\log\log N'.$$
So we have that 
$$\frac{1}{N}\sum_{\NN(\mf p^{k}),\NN(\mf q^{\ell})\leq N'}\NN(\mf q^{\ell})^{-1}\ll \frac{\log\log N'}{\log N'}.$$

On the other hand, let $\pi_{2}(N)$ denote the number of integers  $x\in\{1,\dots,n\}$ with at most 2 disinct factors. By  Lemma \ref{mv7},
\begin{equation}\nonumber
    \begin{split}
       \sum_{\NN(\mf p^{k}\mf q^{\ell})\leq N'} 1\leq \sum_{\NN(\mf p^{k}\mf q^{\ell})=a,a\in\Pi_{2}(N')} 1\ll \pi_{2}(N)=o_{N\to\infty}(1).
    \end{split}
\end{equation}

\fi

For $i\in\N$, let  $\Pi_{i}(N)$ denote the set of integers  $x\in\{1,\dots,n\}$ with at most $i$ disinct factors. Let $\pi_{i}(N):=\vert\Pi_{i}(N)\vert$. It is known that (see \cite[Chapter II.6]{tenebaum_2015}, for example)
$$\pi_{i}(N)\sim C_{i}\frac{N(\log\log N)^{i-1}}{\log N}.$$
%(\WS{see for example https://math.stackexchange.com/questions/3568197/lower-bound-related-to-the-number-of-distinct-prime-numbers}).
Then it follows from  Lemma \ref{mv8}  that 
$$\sum_{\NN(\mf p^{k})\leq N'} 1=\sum_{a\in \Pi_{1}(N')}\sum_{\NN(\mf p^{k})=a} 1\ll \sum_{a\in \Pi_{1}(N')} 1=\pi_{1}(N')=O(\frac{N'}{\log N'}),$$
and
$$\sum_{\NN(\mf q^{\ell})\leq N'} \NN(\mf q^{\ell})^{-1}=\sum_{a\in \Pi_{1}(N')}\sum_{\NN(\mf q^{\ell})=a} a^{-1}\ll \sum_{a\in \Pi_{1}(N')} a^{-1}\ll \sum_{p\leq N', p \text{ is a prime}} p^{-1}=O(\log\log N'),$$
where the last inequality is a consequence of Merten's second theorem, see for example \cite[Chapter I.1, Theorem~9]{tenebaum_2015}. 
On the other hand, note that
%$$\sum_{\NN(\mf p^{k}\mf q^{\ell})\leq N', \mf p\neq \mf q} 1=\sum_{a\in \Pi_{2}(N')}\sum_{\NN(\mf p^{k}q^{\ell})=a, \mf p\neq \mf q} 1\ll \sum_{a\in \Pi_{2}(N')} 1=\pi_{2}(N')=O(\frac{N'\log\log N'}{\log N'}),$$
\begin{equation}\nonumber
    \begin{split}
      &\qquad \sum_{\NN(\mf p^{k}\mf q^{\ell})\leq N', \mf p\neq \mf q} 1=\sum_{a,b\in \Pi_{1}(N'),k,\ell\in\N,a^{k}b^{\ell}\leq N'}\sum_{\NN(\mf p)=a, \NN(\mf q)=b}1
       \ll \sum_{a,b\in \Pi_{1}(N'),k,\ell\in\N,a^{k}b^{\ell}\leq N'} 1
       \\&\leq \pi_{2}(N')+\sum_{p^{r}\in \pi_{1}(N')} r
       \leq O(\frac{N'\log\log N'}{\log N'})+\sum_{r=1}^{\lceil\frac{\log N'}{\log 2}\rceil}r\pi_{1}({N'}^{1/r}),
    \end{split}
\end{equation}
and 
\begin{equation}\nonumber
    \begin{split}
      &\qquad \sum_{r=1}^{\lceil\frac{\log N'}{\log 2}\rceil}r\pi_{1}({N'}^{1/r})
      \ll \sum_{r=1}^{\lceil\frac{\log N'}{\log 2}\rceil}\frac{r^{2}{N'}^{1/r}}{\log N'}
      \leq \frac{N'}{\log N'}+\sum_{r=2}^{\lceil\frac{\log N'}{\log 2}\rceil}\frac{r^{2}\sqrt{N'}}{\log N'}
      \\&\leq \frac{N'}{\log N'}+(\log N')^{2}\sqrt{N'}
      \leq O(\frac{N'\log\log N'}{\log N'}). 
    \end{split}
\end{equation}
The claim follows by combining the above estimates together.
 










It follows from the claim that
\begin{equation}\label{ltk3}
    \begin{split}
       \E_{\NN(u) \leq N}\left| h((Qu+a))-A_{N}\right|^{2}
       \leq %(1+o_{N\to\infty}(1)) 
       O(B_{N'})
    \end{split}
\end{equation}
when $h$ is non-negative.

\fi

When $h$ is real valued, then we let $h_{\pm}$ be the additive function defined by $h_{\pm}(\mf p^{k}):=\max\{\pm h(\mf p^{k}),0\}$. Then 
$$\left| h((Qu+a))-A_{Q,N}\right|^{2}\leq 2\left| h_{+}((Qu+a))-A_{Q,N,+}\right|^{2}+2\left| h_{-}((Qu+a))-A_{Q,N,-}\right|^{2},$$
$B_{N'}=B_{N',+}+B_{N',-}$ and $C_{N'}=C_{N',+}+C_{N',-}$, where $A_{N,\pm}, B_{N,\pm}$ and $C_{N,\pm}$ are defined similar to $A_{N}, B_{N}$ and $C_{N}$ but with $h$ replaced by $h_{\pm}$. So it follows from (\ref{snnn}) that (\ref{ltk}) holds when  $h$ is real valued.

When $h$ is complex valued, (\ref{ltk}) also holds by considering the real and imaginary part separately. We are done.
\end{proof}



\iffalse

\begin{proof}
%%Let 
%$$c_{M,K,Q}:=\sum_{M<\NN(\mf p^{k})\leq N,\mf p\nmid Q}h(\mf p^{k})\NN(\mf p)^{-k}(1-\NN(\mf p)^{-1})$$ 
%It suffices to show that 
%$\E_{\NN(u) \leq N}h(Qu+a)\approx c_{M,K,Q}$ and $\E_{\NN(u) \leq N}\vert h(Qu+a)\vert^{2}\approx c_{M,K,Q}^{2}$.
% $B_N:=\{ u\in\OK : \NN(u) \leq N\}$.
%Then
We first assume that $h$ is non-negative.
 Note that
\begin{equation}\label{firstTK}
    \begin{split}
        &\quad \E_{\NN(u) \leq N}h((Qu+a))
        =\frac{1}{\gamma_{K} N}\sum_{\NN(\frac{u-a}{Q})\leq N, u\equiv a \mod Q}h((u))
        \\&=\frac{1}{\gamma_{K} N}\sum_{\NN(\frac{u-a}{Q})\leq N, u\equiv a \mod Q}\sum_{\mf p^{k}|| u} h(\mf p^{k})
        \\&=\frac{1}{\gamma_{K} N}\sum_{\NN(\frac{u-a}{Q})\leq N, u\equiv a \mod Q}\sum_{\mf p^{k}|| u, \mf p\nmid Q} h(\mf p^{k}).
  %      \\&=\frac{1}{\gamma_{K} N}\sum_{\NN(\mf p^{k})\leq N', \mf p\nmid Q}h(\mf p^{k})\sum_{\NN(\frac{u-a}{Q})\leq N, \mf p^{k}||u, u\equiv a \mod Q} 1,
     %   \\&=\frac{1}{N^2}\sum_{p^{k}\in  [0,QN+a]^2, p\nmid Q}h(p^{k})(\sum_{u\in [0,QN+a]^2, p^{k}|u, u\equiv a \mod Q} 1-\sum_{u\in [0,QN+a]^2, p^{k+1}|u, u\equiv a \mod Q} 1)
  %  \\&\approx \frac{1}{N^2}\sum_{K<p^{k}\leq QN+a, p\nmid Q}h(p^{k})(\frac{QN+a}{p^{k}Q}-\frac{QN+a}{p^{k+1}Q})
  %  \\& =\sum_{\NN(\mf p^{k})\leq \NN(Q)N, \mf p\nmid Q}h(\mf p^{k})\NN(\mf p)^{-k}(1-\NN(\mf p)^{-1})+\Theta_{1}+\Theta_{2},
    \end{split}
\end{equation}
%where the upper bound $2\NN(Q)(N+1)$ in the last line follows from the estimate that if $\mf p^{k}\vert u$ for some $u$ with $\NN(\frac{u-a}{Q})\leq N$, then
%$$\NN(\mf p^{k})\leq \NN(u)\leq 2(\NN(u-a)+\NN(a))\leq 2(\NN(Q)N+\NN(Q))=2\NN(Q)(N+1).$$
%
 %
 Note that if $\mf p^{k}\vert u$ for some $u$ with $\NN(\frac{u-a}{Q})\leq N$, then
$$\NN(\mf p^{k})\leq \NN(u)\leq 2(\NN(u-a)+\NN(a))\leq 2(\NN(Q)N+\NN(Q))=N'.$$ 
By Lemma~\ref{countingCRT},
\begin{equation}\label{firstTK2}
    \begin{split}
        &\quad 
        \frac{1}{\gamma_{K} N}\sum_{\NN(\frac{u-a}{Q})\leq N, u\equiv a \mod Q}\sum_{\mf p^{k}|| u, \mf p\nmid Q} h(\mf p^{k})
        \\&=\frac{1}{\gamma_{K} N}\sum_{\NN(\mf p^{k})\leq N', \mf p\nmid Q}h(\mf p^{k})\sum_{\NN(\frac{u-a}{Q})\leq N, \mf p^{k}||u, u\equiv a \mod Q} 1
    \\& =\frac{1}{\gamma_{K} N}\sum_{\NN(\mf p^{k})\leq N', \mf p\nmid Q}h(\mf p^{k})\Bigl(\sum_{\NN(\frac{u-a}{Q})\leq N, \mf p^{k}|u, u\equiv a \mod Q} 1-\sum_{\NN(\frac{u-a}{Q})\leq N, \mf p^{k+1}|u, u\equiv a \mod Q} 1\Bigr)
    \\&=\sum_{\NN(\mf p^{k})\leq N', \mf p\nmid Q}h(\mf p^{k})((\NN(\mf p)^{-k}-\NN(\mf p)^{-(k+1)})+O(\frac{1}{N})).
    \end{split}
\end{equation}
\WS{here we need the stronger error term $O(\frac{1}{N})$ not just $o_{N\to\infty}(1)$}{\color{red} Under (\ref{newieq}), the error term should be $O(\frac{1}{\sqrt{N\NN(\mf p^{k})}})$}
%Finally, 
%\begin{equation}\label{firstTK3}
%    \begin{split}
 %       &\quad \frac{1}{\gamma_{K} N}\sum_{N<\NN(\mf p^{k})\leq 2\NN(Q)(N+1), \mf p\nmid Q}h(\mf p^{k})\sum_{\NN(\frac{u-a}{Q})\leq N, \mf p^{k}||u, u\equiv a \mod Q} 1
%    \\& \leq \frac{1}{\gamma_{K} N}\sum_{N<\NN(\mf p^{k})\leq \NN(Q)N, \mf p\nmid Q}h(\mf p^{k})\cdot 6(2\NN(Q)+1),
%    \end{split}
%\end{equation}
%where the constant 6 comes from the fact that $K$ has at most 6 units by Lemma \ref{l32}.
Combining (\ref{firstTK}) and (\ref{firstTK2}), %and (\ref{firstTK3}), 
we have that 
\begin{equation}\nonumber
    \begin{split}
        \vert\E_{u\in\OK\colon\NN(u) \leq N}h((Qu+a))-A_{N'}\vert\leq \frac{C}{N}\cdot\sum_{\NN(\mf p^{k})\leq N', \mf p\nmid Q}\vert h(\mf p^{k})\vert
    \end{split}
\end{equation}
for some $C>0$ depending only on $K$ and $Q$.
Therefore,
\begin{equation}\label{firstTK0}
    \begin{split}
       &\qquad \vert A_{N}\cdot(\E_{u\in\OK\colon\NN(u) \leq N}h((Qu+a))-A_{N'})\vert
       \\&\ll \frac{1}{N}\sum_{\NN(\mf p^{k})\leq N', \mf p\nmid Q}\vert h(\mf p^{k})\vert\cdot \sum_{\NN(\mf p^{k})\leq N', \mf p\nmid Q}\vert h(\mf p^{k})\vert\cdot \NN(\mf p^{k})^{-1}
        \\&\leq \frac{1}{N}\sum_{\NN(\mf p^{k})\leq N', \NN(\mf p)>M}\vert h(\mf p^{k})\vert\cdot \sum_{\NN(\mf p^{k})\leq N', \NN(\mf p)>M}\vert h(\mf p^{k})\vert\cdot \NN(\mf p^{k})^{-1}
          \\&\leq \frac{1}{N}\Bigl(B_{N'}\cdot\sum_{\NN(\mf p^{k})\leq N'}\NN(\mf p^{k})\Bigr)^{1/2}\cdot\Bigl(B_{N'}\cdot\sum_{\NN(\mf p^{k})\leq N'}\NN(\mf p^{k})^{-1}\Bigr)^{1/2}
          \\&=B_{N'}\cdot \frac{1}{N}\Bigl(\sum_{\NN(\mf p^{k}),\NN(\mf q^{\ell})\leq N'}\NN(\mf p^{k})\NN(\mf q^{\ell})^{-1}\Bigr)^{1/2}.
          \end{split}
\end{equation}

{\color{red}

Under (\ref{newieq}), this term should be 
\begin{equation}\label{firstTK0}
    \begin{split}
       &\qquad \vert A_{N}\cdot(\E_{u\in\OK\colon\NN(u) \leq N}h((Qu+a))-A_{N'})\vert
       \\&\ll \frac{1}{\sqrt{N}}\sum_{\NN(\mf p^{k})\leq N', \mf p\nmid Q}\vert h(\mf p^{k})\vert \cdot \NN(\mf p^{k})^{-1/2}\cdot \sum_{\NN(\mf p^{k})\leq N', \mf p\nmid Q}\vert h(\mf p^{k})\vert\cdot \NN(\mf p^{k})^{-1}
        \\&\leq \frac{1}{\sqrt{N}}\sum_{\NN(\mf p^{k})\leq N', \NN(\mf p)>M}\vert h(\mf p^{k})\vert\cdot \NN(\mf p^{k})^{-1/2}\cdot \sum_{\NN(\mf p^{k})\leq N', \NN(\mf p)>M}\vert h(\mf p^{k})\vert\cdot \NN(\mf p^{k})^{-1}
          \\&\leq \frac{1}{\sqrt{N}}\Bigl(B_{N'}\cdot\sum_{\NN(\mf p^{k})\leq N'} 1\Bigr)^{1/2}\cdot\Bigl(B_{N'}\cdot\sum_{\NN(\mf p^{k})\leq N'}\NN(\mf p^{k})^{-1}\Bigr)^{1/2}
          \\&=B_{N'}\cdot \frac{1}{\sqrt{N}}\Bigl(\sum_{\NN(\mf p^{k}),\NN(\mf q^{\ell})\leq N'}\NN(\mf q^{\ell})^{-1}\Bigr)^{1/2}.
          \end{split}
\end{equation}


}

\iffalse
\WS{
This part needs to be written carefully.

(1) I think in the last two lines, if $\mf p^{k}|| u$, then $\NN(\mf p)$ could be larger than $\NN(QN)$.

(2) Also for the range of $\mf p^{k}$, I don't see how to replace $\NN(Q)N$ by $N$ (this is claimed in [KMPT] Lemma 2.5 but I don't see why)

(3) In the proof we missed $M$

(4) Lemma 5.1 does not provide control for error terms. 
}
\FM{If we use that $\mf p \nmid Q$ does that help with (2)? 

In any case, for the last term, can we use some sort of trivial bound+prime number theorem to get something that goes to $0$ as $N\to \infty$? I just checked that there is an analog of the prime number theorem for number fields, which might help us here.}
%where
%\[
%\Theta_1=\sum_{\NN(\mf p) \leq \NN(Q)N, \mf p \nmid Q} |h(\mf p^k)|^2 \left( \frac{1}{\NN(\mf p^k)}-\frac{1}{\NN(\mf p^{k+1})}\right).
%\]



Note that for the last line, we used Lemma~\ref{countingCRT}  and we factored out $\NN(\mf p^k)$ in the $\left( \frac{1}{\NN(\mf p^k)}-\frac{1}{\NN(\mf p^{k+1})}\right)$ bracket.

For the second term, we have

\fi

Now we estimate the last term of (\ref{ltk}). We have
\begin{equation}\label{secondTK}
    \begin{split}
        &\quad \E_{\NN(u) \leq N}\vert h((Qu+a))\vert^{2}
        =\frac{1}{\gamma_{K} N}\sum_{\NN(\frac{u-a}{Q})\leq N, u\equiv a \mod Q}\vert h((u))\vert^{2}
        \\&=\frac{1}{\gamma_{K} N}\sum_{\NN(\frac{u-a}{Q})\leq N, u\equiv a \mod Q}\left| \sum_{\mf p^{k}|| u} h(\mf p^{k})\right|^{2}
        \\&=\frac{1}{\gamma_{K} N}\sum_{\NN(\frac{u-a}{Q})\leq N, u\equiv a \mod Q}\left(\sum_{\mf p^{k}|| u} \vert h(\mf p^{k})\vert^{2}+\sum_{\mf p^{k},\mf q^{\ell}|| u, \mf p\neq \mf q} h(\mf p^{k})h(\mf q^{\ell})\right)
          \\&=\frac{1}{\gamma_{K} N}\sum_{\NN(\frac{u-a}{Q})\leq N, u\equiv a \mod Q}\left(\sum_{\mf p^{k}|| u, \mf p\nmid Q} \vert h(\mf p^{k})\vert^{2}+\sum_{\mf p^{k},\mf q^{\ell}|| u, \mf p\neq \mf q, \mf p,\mf q\nmid Q} h(\mf p^{k})h(\mf q^{\ell})\right).
    \end{split}
\end{equation}
The first term in the last line of \eqref{secondTK} can be dealt with in the same way as in  \eqref{firstTK2} %and \eqref{firstTK3} 
(with $h$ replaced by $h^{2}$), which is  
\begin{equation}\label{secondTK2}
    \begin{split}
        &\quad 
        \frac{1}{\gamma_{K} N}\sum_{\NN(\frac{u-a}{Q})\leq N, u\equiv a \mod Q}\sum_{\mf p^{k}|| u, \mf p\nmid Q} h(\mf p^{k})^{2}
 %   \\&=\sum_{\NN(\mf p^{k})\leq N', \mf p\nmid Q}h(\mf p^{k})^{2}((\NN(\mf p)^{-k}-\NN(\mf p)^{-(k+1)})+O(\frac{1}{N}))
 %   \\&\leq B_{N'}+\frac{C}{N}\sum_{\NN(\mf p^{k})\leq N', \mf p\nmid Q}h(\mf p^{k})^{2}
  %  \\&\leq B_{N'}+\frac{C}{N}\sum_{\NN(\mf p^{k})\leq N', \NN(\mf p)>M}h(\mf p^{k})^{2}
     \\&=\frac{1}{\gamma_{K} N}\sum_{\NN(\mf p^{k})\leq N', \mf p\nmid Q}h(\mf p^{k})^{2}\sum_{\NN(\frac{u-a}{Q})\leq N, \mf p^{k}||u, u\equiv a \mod Q} 1
      \\&\leq\sum_{\NN(\mf p^{k})\leq N', \mf p\nmid Q}h(\mf p^{k})^{2}\NN(\mf p)^{-k}
 \leq\sum_{\NN(\mf p^{k})\leq N', \NN(\mf p)>M}h(\mf p^{k})^{2}\NN(\mf p)^{-k}=B_{N'}
    \end{split}
\end{equation}
\WS{here we need to use the lower bound for Lemma \ref{countingCRT} (to do so we also need to assume that $u\neq 0$).} {\color{red}
Under (\ref{newieq}), this should be 
\begin{equation}\label{secondTK2}
    \begin{split}
        &\quad 
        \frac{1}{\gamma_{K} N}\sum_{\NN(\frac{u-a}{Q})\leq N, u\equiv a \mod Q}\sum_{\mf p^{k}|| u, \mf p\nmid Q} h(\mf p^{k})^{2}
 %   \\&=\sum_{\NN(\mf p^{k})\leq N', \mf p\nmid Q}h(\mf p^{k})^{2}((\NN(\mf p)^{-k}-\NN(\mf p)^{-(k+1)})+O(\frac{1}{N}))
 %   \\&\leq B_{N'}+\frac{C}{N}\sum_{\NN(\mf p^{k})\leq N', \mf p\nmid Q}h(\mf p^{k})^{2}
  %  \\&\leq B_{N'}+\frac{C}{N}\sum_{\NN(\mf p^{k})\leq N', \NN(\mf p)>M}h(\mf p^{k})^{2}
     \\&=\frac{1}{\gamma_{K} N}\sum_{\NN(\mf p^{k})\leq N', \mf p\nmid Q}h(\mf p^{k})^{2}\sum_{\NN(\frac{u-a}{Q})\leq N, \mf p^{k}||u, u\equiv a \mod Q} 1
      \\&=\sum_{\NN(\mf p^{k})\leq N', \mf p\nmid Q}h(\mf p^{k})^{2}\NN(\mf p)^{-k}+\Theta,
    \end{split}
\end{equation}
where 
$$\vert\Theta\vert\ll \frac{1}{\sqrt{N}}\sum_{\NN(\mf p^{k})\leq N', \mf p\nmid Q}h(\mf p^{k})^{2}\cdot \NN(\mf p^{k})^{-1/2}\ll \frac{1}{\sqrt{N}}\sum_{\NN(\mf p^{k})\leq N', \mf p\nmid Q}h(\mf p^{k})^{2}\cdot \NN(\mf p^{k})^{-1}\leq B_{N'}.$$
(this term probablity should be $o(B_{N'})$)
}
For the second term in \eqref{secondTK}, similar to (\ref{firstTK}), we also swap the order of the sums and rewrite it as
%\begin{equation}
 %   \begin{split}
  %      &\quad \frac{1}{(2N)^2}\sum_{p^{k}, q^l \in Q[\pm N]^2+a, p \neq q, p, q \in \Z[i]_+}\vert h(p^{k})\vert^{2}\sum_{u\leq QN+a, u\equiv a \mod Q, p^{k}\Vert u} 1 
   %     \\&\approx \frac{1}{N}\sum_{p^{k}\leq QN+a}\vert h(p^{k})\vert^{2}p^{-k}(1-p^{-1}) 
    %    \\&\leq \frac{1}{N}\sum_{p^{k}\leq QN+a}\vert h(p^{k})\vert^{2}p^{-k}
   % \end{split}
%\end{equation}
%(the RHS of (\cite{20}) is similar to this.)
% On the other hand, we have
\begin{equation}\nonumber
    \begin{split}
        &\quad \frac{1}{\gamma_{K} N}\sum_{\NN(\mf p^{k}\mf q^{\ell}) \leq N', \mf p\neq \mf q,\mf p,\mf q\nmid Q} h(\mf p^{k})h(\mf q^{\ell})\sum_{\NN(\frac{u-a}{Q})\leq N, u\equiv a \mod Q, \mf p^{k}||u, \mf q^{\ell}||u}1.
    \end{split}
\end{equation}
Similar to (\ref{firstTK2}), we have that
\begin{equation}\label{secondTK3}
    \begin{split}
        &\quad \frac{1}{\gamma_{K} N}\sum_{\NN(\mf p^{k}\mf q^{\ell}) \leq N', \mf p\neq \mf q,\mf p,\mf q\nmid Q} h(\mf p^{k})h(\mf q^{\ell})\sum_{\NN(\frac{u-a}{Q})\leq N, u\equiv a \mod Q, \mf p^{k}||u, \mf q^{\ell}||u}1
        \\&=\sum_{\NN(\mf p^{k}\mf q^{\ell})\leq N', \mf p\neq \mf q, \mf p,\mf q\nmid Q} h(\mf p^{k})h(\mf q^{\ell})\Bigl(\NN(\mf p^k)^{-1}\NN(\mf q^l)^{-1}\left(1-\frac{1}{\NN(\mf p)}\right) \left(1-\frac{1}{\NN(\mf q)}\right)+O(\frac{1}{N})\Bigr)
        \\&\leq A_{N'}^{2}+\frac{C}{N}\sum_{\NN(\mf p^{k}\mf q^{\ell})\leq N', \mf p\neq \mf q, \NN(\mf p),\NN(\mf q)>M} h(\mf p^{k})h(\mf q^{\ell})     %  \\&=\vert\sum_{K<p^{k}\leq QN+a, p\nmid Q}h(p^{k})p^{-k}(1-p^{-1})\vert^{2}+error,
        \\&=A_{N'}^{2}+\frac{C}{N}\sum_{\NN(\mf p^{k}\mf q^{\ell})\leq N', \mf p\neq \mf q, \NN(\mf p),\NN(\mf q)>M} \frac{h(\mf p^{k})}{\NN(\mf p^{k})^{1/2}}\cdot\frac{h(\mf q^{\ell})}{\NN(\mf q^{\ell})^{1/2}}\cdot \NN(\mf p^{k})^{1/2}\cdot \NN(\mf q^{\ell})^{1/2}
        \\&\leq A_{N'}^{2}+B_{N'}\cdot\frac{C}{N}\Bigl(\sum_{\NN(\mf p^{k}\mf q^{\ell})\leq N', \mf p\neq \mf q} \NN(\mf p^{k})\cdot \NN(\mf q^{\ell})\Bigr)^{1/2}.
            \end{split}
\end{equation}

{\color{red}

Under (\ref{newieq}), this should be 
\begin{equation}\label{secondTK3}
    \begin{split}
        &\quad \frac{1}{\gamma_{K} N}\sum_{\NN(\mf p^{k}\mf q^{\ell}) \leq N', \mf p\neq \mf q,\mf p,\mf q\nmid Q} h(\mf p^{k})h(\mf q^{\ell})\sum_{\NN(\frac{u-a}{Q})\leq N, u\equiv a \mod Q, \mf p^{k}||u, \mf q^{\ell}||u}1
        \\&=\sum_{\NN(\mf p^{k}\mf q^{\ell})\leq N', \mf p\neq \mf q, \mf p,\mf q\nmid Q} h(\mf p^{k})h(\mf q^{\ell})\Bigl(\NN(\mf p^k)^{-1}\NN(\mf q^l)^{-1}\left(1-\frac{1}{\NN(\mf p)}\right) \left(1-\frac{1}{\NN(\mf q)}\right)+O(\frac{1}{\sqrt{N\NN(\mf p^{k}\mf q^{\ell})}})\Bigr)
        \\&\leq A_{N'}^{2}+\frac{C}{\sqrt{N}}\sum_{\NN(\mf p^{k}\mf q^{\ell})\leq N', \mf p\neq \mf q, \NN(\mf p),\NN(\mf q)>M} h(\mf p^{k})h(\mf q^{\ell})\NN(\mf p^{k})^{-1/2}\NN(\mf q^{\ell})^{-1/2}     %  \\&=\vert\sum_{K<p^{k}\leq QN+a, p\nmid Q}h(p^{k})p^{-k}(1-p^{-1})\vert^{2}+error,
       % \\&=A_{N'}^{2}+\frac{C}{\sqrt{N}}\sum_{\NN(\mf p^{k}\mf q^{\ell})\leq N', \mf p\neq \mf q, \NN(\mf p),\NN(\mf q)>M} \frac{h(\mf p^{k})}{\NN(\mf p^{k})^{1/2}}\cdot\frac{h(\mf q^{\ell})}{\NN(\mf q^{\ell})^{1/2}}\cdot \NN(\mf p^{k})^{1/2}\cdot \NN(\mf q^{\ell})^{1/2}
        \\&\leq A_{N'}^{2}+B_{N'}\cdot\frac{C}{\sqrt{N}}\Bigl(\sum_{\NN(\mf p^{k}\mf q^{\ell})\leq N', \mf p\neq \mf q} 1\Bigr)^{1/2}.
            \end{split}
\end{equation}
}

Combining (\ref{secondTK}), (\ref{secondTK2}) and (\ref{secondTK3}), we have that
\begin{equation}\label{secondTK0}
    \begin{split}
        &\quad \E_{\NN(u) \leq N}\vert h((Qu+a))\vert^{2}
        \\&\leq B_{N'}+A_{N'}^{2}+\frac{C}{N}\sum_{\NN(\mf p^{k})\leq N', \NN(\mf p)>M}h(\mf p^{k})^{2}+B_{N'}\cdot\frac{C}{N}\Bigl(\sum_{\NN(\mf p^{k}\mf q^{\ell})\leq N', \mf p\neq \mf q} \NN(\mf p^{k})\cdot \NN(\mf q^{\ell})\Bigr)^{1/2}.
    \end{split}
\end{equation}   
Finally, note that 
 \begin{equation}\label{snnn2}
    \begin{split}
        &\quad A_{N'}-A_{N}
        =\sum_{N<\NN(\mf p^{k})\leq N',\mf p\nmid Q}h(\mf p^{k})\NN(\mf p)^{-k}(1-\NN(\mf p)^{-1})
        \\&\leq \frac{2}{N}\sum_{N<\NN(\mf p^{k})\leq N',\mf p\nmid Q}h(\mf p^{k})
        \leq \frac{2}{N}\sum_{N<\NN(\mf p^{k})\leq N',\mf \NN(\mf p)>M}h(\mf p^{k})
        \\&\leq B_{N'}^{1/2}\cdot \frac{2}{N}\Bigl(\sum_{N<\NN(\mf p^{k})\leq N'}\NN(\mf p^{k})\Bigr)^{1/2}.
        \end{split}
\end{equation}
Combining %(\ref{ltk2}), 
(\ref{firstTK0}), (\ref{secondTK0}) and (\ref{snnn2}), we have
%the left hand side of (\ref{ltk}) can be bounded by
%
% \begin{equation}\label{snnn}
%    \begin{split}
%        &\quad \E_{\NN(u) \leq N}\left| h((Qu+a))-A_{N'}\right|^{2}  
%        \\&=-A_{N'}^{2}+2A_{N'}(\E_{\NN(u) \leq N}h((Qu+a))-A_{N'})+\E_{\NN(u) \leq N}\vert h((Qu+a))\vert^{2}
 %       \\&\leq B_{N'}(1+\frac{1}{N}\Bigl(\sum_{\NN(\mf p^{k}),\NN(\mf q^{\ell})\leq N'}\NN(\mf p^{k})\NN(\mf q^{\ell})^{-1}\Bigr)^{1/2}+???+\frac{C}{N}\Bigl(\sum_{\NN(\mf p^{k}\mf q^{\ell})\leq N', \mf p\neq \mf q} \NN(\mf p^{k})\cdot \NN(\mf q^{\ell})\Bigr)^{1/2}).
 %       \end{split}
%\end{equation}
%
 \begin{equation}\label{snnn}
    \begin{split}
        &\quad \E_{\NN(u) \leq N}\left| h((Qu+a))-A_{N}\right|^{2}  
        \\&=(A_{N}-A_{N'})^{2}-2A_{N}(\E_{\NN(u) \leq N}h((Qu+a))-A_{N'})+\E_{\NN(u) \leq N}\vert h((Qu+a))\vert^{2}-A_{N'}^{2}
        \\&\leq B_{N'}(1+\frac{2C}{N}\Bigl(\sum_{\NN(\mf p^{k}),\NN(\mf q^{\ell})\leq N'}\NN(\mf p^{k})\NN(\mf q^{\ell})^{-1}\Bigr)^{1/2}+\frac{C}{N}\Bigl(\sum_{\NN(\mf p^{k}\mf q^{\ell})\leq N', \mf p\neq \mf q} \NN(\mf p^{k})\cdot \NN(\mf q^{\ell})\Bigr)^{1/2}).
        \end{split}
\end{equation}
Since
$$\lim_{N\to\infty}\frac{2C}{N}\Bigl(\sum_{\NN(\mf p^{k}),\NN(\mf q^{\ell})\leq N'}\NN(\mf p^{k})\NN(\mf q^{\ell})^{-1}\Bigr)^{1/2}+\frac{C}{N}\Bigl(\sum_{\NN(\mf p^{k}\mf q^{\ell})\leq N', \mf p\neq \mf q} \NN(\mf p^{k})\cdot \NN(\mf q^{\ell})\Bigr)^{1/2}=0,$$
\WS{give a proof. This is the function $\epsilon(x)$ on page 302 of the book "Introduction to analytic and Probabilistic number theory"}

{\color{red} Under (\ref{newieq}), the error term is 
$$CB_{N'}\cdot(\Bigl(\frac{1}{N}\sum_{\NN(\mf p^{k}),\NN(\mf q^{\ell})\leq N'}\NN(\mf q^{\ell})^{-1}\Bigr)^{1/2}+\Bigl(\frac{1}{N}\sum_{\NN(\mf p^{k}\mf q^{\ell})\leq N', \mf p\neq \mf q} 1\Bigr)^{1/2})+\Theta$$

%$$\lim_{N\to\infty}\frac{2C}{\sqrt{N}}\Bigl(\sum_{\NN(\mf p^{k}),\NN(\mf q^{\ell})\leq N'}\NN(\mf p^{k})\NN(\mf q^{\ell})^{-1}\Bigr)^{1/2}+\frac{C}{\sqrt{N}}\Bigl(\sum_{\NN(\mf p^{k}\mf q^{\ell})\leq N', \mf p\neq \mf q} \NN(\mf p^{k})\cdot \NN(\mf q^{\ell})\Bigr)^{1/2}+\Theta.$$
}


\WS{P.S. ChatGPT claims that the TK theorem for general number field was proved in "Probabilistic Methods in The Theory of Numbers (Kubilius)" in chapter 9. I have not confirmed this claim, but I uploaded the book (named "book K").}

  we have that 
\begin{equation}\label{ltk3}
    \begin{split}
       \E_{\NN(u) \leq N}\left| h((Qu+a))-A_{N}\right|^{2}
       \ll (1+o_{N\to\infty}(1)) B_{N'}
    \end{split}
\end{equation}
when $h$ is non-negative.

When $h$ is real valued, then we let $h_{\pm}$ be the additive function defined by $h_{\pm}(\mf p^{k}):=\max\{\pm h(\mf p^{k}),0\}$. Then 
$$\left| h((Qu+a))-A_{N}\right|^{2}\leq 2\left| h_{+}((Qu+a))-A_{N,+}\right|^{2}+2\left| h_{-}((Qu+a))-A_{N,-}\right|^{2}$$
and $B_{N'}=B_{N',+}+B_{N',-}$, where $A_{N,\pm}$ and $B_{N,\pm}$ are defined similar to $A_{N}$ and $B_{N}$ but with $h$ replaced by $h_{\pm}$. So it follows from (\ref{ltk3}) that (\ref{ltk}) holds when  $h$ is real valued.

When $h$ is complex valued, (\ref{ltk}) also holds by considering the real and imaginary part separately. We are done.
\end{proof}

\begin{proposition}[Tur\'an-Kubilius inequality]\label{TK}
    Let $K=\Q(\sqrt{-d})$ for some squarefree $d\in\N$, $h\colon\mf I(K)\to\C$ be an additive function, $a,Q\in\mathcal{O}_K^{\times}$ with $\NN(a)\leq \NN(Q)$ and $(a)$ being coprime to $(Q)$, and $M,N\in\N$. Suppose that $u\vert Q$ for all $u\in\OK$ with $\NN(u)\leq M$.
    Then \begin{equation}\label{ltk}
    \begin{split}
       \E_{u\in\OK\colon\NN(u) \leq N}\left| h((Qu+a))-A_{N}\right|^{2}
       \ll (1+o_{N\to\infty}(1)) B_{N},
    \end{split}
\end{equation}
where
$$A_{N}:=\sum_{\NN(\mf p^{k})\leq N,\mf p\nmid Q}h(\mf p^{k})\NN(\mf p)^{-k}(1-\NN(\mf p)^{-1}) \text{ and } B_{N}:=\sum_{\NN(\mf p)>M, k\geq 1}\vert h(\mf p^{k})\vert^{2}\NN(\mf p^{k})^{-1}.$$
\end{proposition}
\begin{proof}
%%Let 
%$$c_{M,K,Q}:=\sum_{M<\NN(\mf p^{k})\leq N,\mf p\nmid Q}h(\mf p^{k})\NN(\mf p)^{-k}(1-\NN(\mf p)^{-1})$$ 
%It suffices to show that 
%$\E_{\NN(u) \leq N}h(Qu+a)\approx c_{M,K,Q}$ and $\E_{\NN(u) \leq N}\vert h(Qu+a)\vert^{2}\approx c_{M,K,Q}^{2}$.
% $B_N:=\{ u\in\OK : \NN(u) \leq N\}$.
%Then
The left hand side of (\ref{ltk}) can be expanded as  
\begin{equation}\label{ltk2}
    \begin{split}
       A_{N}^{2}-2A_{N}\cdot \E_{\NN(u) \leq N}h((Qu+a))+\E_{\NN(u) \leq N}h((Qu+a))^{2}.
    \end{split}
\end{equation}
We first estimate the second term of (\ref{ltk}). Note that
\begin{equation}\label{firstTK}
    \begin{split}
        &\quad \E_{\NN(u) \leq N}h((Qu+a))
        =\frac{1}{\gamma_{K} N}\sum_{\NN(\frac{u-a}{Q})\leq N, u\equiv a \mod Q}h((u))
        \\&=\frac{1}{\gamma_{K} N}\sum_{\NN(\frac{u-a}{Q})\leq N, u\equiv a \mod Q}\sum_{\mf p^{k}|| u} h(\mf p^{k})
        \\&=\frac{1}{\gamma_{K} N}\sum_{\NN(\frac{u-a}{Q})\leq N, u\equiv a \mod Q}\sum_{\mf p^{k}|| u, \mf p\nmid Q} h(\mf p^{k})
        \\&=\frac{1}{\gamma_{K} N}\sum_{\NN(\mf p^{k})\leq 2\NN(Q)(N+1), \mf p\nmid Q}h(\mf p^{k})\sum_{\NN(\frac{u-a}{Q})\leq N, \mf p^{k}||u, u\equiv a \mod Q} 1,
     %   \\&=\frac{1}{N^2}\sum_{p^{k}\in  [0,QN+a]^2, p\nmid Q}h(p^{k})(\sum_{u\in [0,QN+a]^2, p^{k}|u, u\equiv a \mod Q} 1-\sum_{u\in [0,QN+a]^2, p^{k+1}|u, u\equiv a \mod Q} 1)
  %  \\&\approx \frac{1}{N^2}\sum_{K<p^{k}\leq QN+a, p\nmid Q}h(p^{k})(\frac{QN+a}{p^{k}Q}-\frac{QN+a}{p^{k+1}Q})
  %  \\& =\sum_{\NN(\mf p^{k})\leq \NN(Q)N, \mf p\nmid Q}h(\mf p^{k})\NN(\mf p)^{-k}(1-\NN(\mf p)^{-1})+\Theta_{1}+\Theta_{2},
    \end{split}
\end{equation}
where the upper bound $2\NN(Q)(N+1)$ in the last line follows from the estimate that if $\mf p^{k}\vert u$ for some $u$ with $\NN(\frac{u-a}{Q})\leq N$, then
$$\NN(\mf p^{k})\leq \NN(u)\leq 2(\NN(u-a)+\NN(a))\leq 2(\NN(Q)N+\NN(Q))=2\NN(Q)(N+1).$$



On the other hand, by Lemma~\ref{countingCRT},
\begin{equation}\label{firstTK2}
    \begin{split}
        &\quad \frac{1}{\gamma_{K} N}\sum_{\NN(\mf p^{k})\leq N, \mf p\nmid Q}h(\mf p^{k})\sum_{\NN(\frac{u-a}{Q})\leq N, \mf p^{k}||u, u\equiv a \mod Q} 1
    \\& =\frac{1}{\gamma_{K} N}\sum_{\NN(\mf p^{k})\leq N, \mf p\nmid Q}h(\mf p^{k})\Bigl(\sum_{\NN(\frac{u-a}{Q})\leq N, \mf p^{k}|u, u\equiv a \mod Q} 1-\sum_{\NN(\frac{u-a}{Q})\leq N, \mf p^{k+1}|u, u\equiv a \mod Q} 1\Bigr)
    \\&=\sum_{\NN(\mf p^{k})\leq N, \mf p\nmid Q}h(\mf p^{k})((\NN(\mf p)^{-k}-\NN(\mf p)^{-(k+1)})+O(\frac{1}{N})).
    \end{split}
\end{equation}
\WS{here we need the stronger error term $O(\frac{1}{N})$ not just $o_{N\to\infty}(1)$}
Finally, 
\begin{equation}\label{firstTK3}
    \begin{split}
        &\quad \frac{1}{\gamma_{K} N}\sum_{N<\NN(\mf p^{k})\leq 2\NN(Q)(N+1), \mf p\nmid Q}h(\mf p^{k})\sum_{\NN(\frac{u-a}{Q})\leq N, \mf p^{k}||u, u\equiv a \mod Q} 1
    \\& \leq \frac{1}{\gamma_{K} N}\sum_{N<\NN(\mf p^{k})\leq \NN(Q)N, \mf p\nmid Q}h(\mf p^{k})\cdot 6(2\NN(Q)+1),
    \end{split}
\end{equation}
where the constant 6 comes from the fact that $K$ has at most 6 units by Lemma \ref{l32}.
Combining (\ref{firstTK}), (\ref{firstTK2}) and (\ref{firstTK3}), we have that 
\begin{equation}\nonumber
    \begin{split}
        \vert\E_{u\in\OK\colon\NN(u) \leq N}h((Qu+a))-A_{N}\vert\leq \frac{C}{N}\cdot\sum_{\NN(\mf p^{k})\leq \NN(Q)N, \mf p\nmid Q}\vert h(\mf p^{k})\vert
    \end{split}
\end{equation}
for some $C>0$ depending only on $K$ and $Q$.
Therefore,
\begin{equation}\nonumber
    \begin{split}
       &\qquad \vert A_{N}\cdot(\E_{u\in\OK\colon\NN(u) \leq N}h((Qu+a))-A_{N})\vert
       \\&\ll \frac{1}{N}\sum_{\NN(\mf p^{k})\leq \NN(Q)N, \mf p\nmid Q}\vert h(\mf p^{k})\vert\cdot \sum_{\NN(\mf p^{k})\leq \NN(Q)N, \mf p\nmid Q}\vert h(\mf p^{k})\vert\cdot \NN(\mf p^{k})^{-1}
        \\&\leq \frac{1}{N}\sum_{\NN(\mf p^{k})\leq \NN(Q)N, \NN(\mf p)>M}\vert h(\mf p^{k})\vert\cdot \sum_{\NN(\mf p^{k})\leq \NN(Q)N, \NN(\mf p)>M}\vert h(\mf p^{k})\vert\cdot \NN(\mf p^{k})^{-1}
          \\&\leq \frac{1}{N}\Bigl(B_{N}\cdot\sum_{\NN(\mf p^{k})\leq \NN(Q)N}\NN(\mf p^{k})\Bigr)^{1/2}\cdot\Bigl(B_{N}\cdot\sum_{\NN(\mf p^{k})\leq \NN(Q)N}\NN(\mf p^{k})^{-1}\Bigr)^{1/2}
          \\&=B_{N}\cdot \frac{1}{N}\Bigl(\sum_{\NN(\mf p^{k}),\NN(\mf q^{\ell})\leq \NN(Q)N}\NN(\mf p^{k})\NN(\mf q^{\ell})^{-1}\Bigr)^{1/2}.
          \end{split}
\end{equation}

\iffalse
\WS{
This part needs to be written carefully.

(1) I think in the last two lines, if $\mf p^{k}|| u$, then $\NN(\mf p)$ could be larger than $\NN(QN)$.

(2) Also for the range of $\mf p^{k}$, I don't see how to replace $\NN(Q)N$ by $N$ (this is claimed in [KMPT] Lemma 2.5 but I don't see why)

(3) In the proof we missed $M$

(4) Lemma 5.1 does not provide control for error terms. 
}
\FM{If we use that $\mf p \nmid Q$ does that help with (2)? 

In any case, for the last term, can we use some sort of trivial bound+prime number theorem to get something that goes to $0$ as $N\to \infty$? I just checked that there is an analog of the prime number theorem for number fields, which might help us here.}
%where
%\[
%\Theta_1=\sum_{\NN(\mf p) \leq \NN(Q)N, \mf p \nmid Q} |h(\mf p^k)|^2 \left( \frac{1}{\NN(\mf p^k)}-\frac{1}{\NN(\mf p^{k+1})}\right).
%\]



Note that for the last line, we used Lemma~\ref{countingCRT}  and we factored out $\NN(\mf p^k)$ in the $\left( \frac{1}{\NN(\mf p^k)}-\frac{1}{\NN(\mf p^{k+1})}\right)$ bracket.

For the second term, we have

\fi

Now we estimate the last term of (\ref{ltk}). We have
\begin{equation}\label{secondTK}
    \begin{split}
        &\quad \E_{\NN(u) \leq N}\vert h((Qu+a))\vert^{2}
        =\frac{1}{\gamma_{K} N}\sum_{\NN(\frac{u-a}{Q})\leq N, u\equiv a \mod Q}\vert h((u))\vert^{2}
        \\&=\frac{1}{\gamma_{K} N}\sum_{\NN(\frac{u-a}{Q})\leq N, u\equiv a \mod Q}\left| \sum_{\mf p^{k}|| u} h(\mf p^{k})\right|^{2}
        \\&=\frac{1}{\gamma_{K} N}\sum_{\NN(\frac{u-a}{Q})\leq N, u\equiv a \mod Q}\left(\sum_{\mf p^{k}|| u} \vert h(\mf p^{k})\vert^{2}+\sum_{\mf p^{k},\mf q^{\ell}|| u, \mf p\neq \mf q} h(\mf p^{k})\overline{h}(\mf q^{\ell})\right). 
    \end{split}
\end{equation}
The first term in the last line of \eqref{secondTK} can be dealt with in the same way as in \eqref{firstTK}, \eqref{firstTK2} and \eqref{firstTK3} (with $h$ replaced by $h^{2}$), which is  
\begin{equation}\nonumber
    \begin{split}
        &\quad \Bigl\vert\frac{1}{\gamma_{K} N}\sum_{\NN(\frac{u-a}{Q})\leq N, u\equiv a \mod Q}\sum_{\mf p^{k}|| u} \vert h(\mf p^{k})\vert^{2}-\sum_{\NN(\mf p^{k})\leq N,\mf p\nmid Q}\vert h(\mf p^{k})\vert^{2}\NN(\mf p)^{-k}(1-\NN(\mf p)^{-1})\Bigr\vert
        \\&\leq \frac{C}{N}\cdot\sum_{\NN(\mf p^{k})\leq \NN(Q)N, \mf p\nmid Q}\vert h(\mf p^{k})\vert^{2}.
    \end{split}
\end{equation}
This implies that 
\begin{equation}\label{secondTK2}
    \begin{split}
        &\quad \Bigl\vert\frac{1}{\gamma_{K} N}\sum_{\NN(\frac{u-a}{Q})\leq N, u\equiv a \mod Q}\sum_{\mf p^{k}|| u} \vert h(\mf p^{k})\vert^{2}-B_{N}\Bigr\vert
        \\&\ll \frac{1}{N}\cdot\sum_{\NN(\mf p^{k})\leq \NN(Q)N, \mf p\nmid Q}\vert h(\mf p^{k})\vert^{2}+\sum_{\NN(\mf p^{k})\leq N,\mf p\nmid Q}\vert h(\mf p^{k})\vert^{2}\NN(\mf p)^{-(k+1)}
      \\&\leq \frac{1}{N}\cdot\sum_{\NN(\mf p^{k})\leq \NN(Q)N, \NN(\mf p)>M}\vert h(\mf p^{k})\vert^{2}+\sum_{\NN(\mf p^{k})\leq N,\NN(\mf p)>M}\vert h(\mf p^{k})\vert^{2}\NN(\mf p)^{-(k+1)}
      \\&\leq  
    \end{split}
\end{equation}
{\color{red} check error term}
For the second term in \eqref{secondTK}, similar to (\ref{firstTK}), we also swap the order of the sums and rewrite it as
%\begin{equation}
 %   \begin{split}
  %      &\quad \frac{1}{(2N)^2}\sum_{p^{k}, q^l \in Q[\pm N]^2+a, p \neq q, p, q \in \Z[i]_+}\vert h(p^{k})\vert^{2}\sum_{u\leq QN+a, u\equiv a \mod Q, p^{k}\Vert u} 1 
   %     \\&\approx \frac{1}{N}\sum_{p^{k}\leq QN+a}\vert h(p^{k})\vert^{2}p^{-k}(1-p^{-1}) 
    %    \\&\leq \frac{1}{N}\sum_{p^{k}\leq QN+a}\vert h(p^{k})\vert^{2}p^{-k}
   % \end{split}
%\end{equation}
%(the RHS of (\cite{20}) is similar to this.)
% On the other hand, we have
\begin{equation}
    \begin{split}
        &\quad \frac{1}{\gamma_{K} N}\sum_{\NN(\mf p^{k}),\NN(\mf q^{\ell}) \leq 2\NN(Q)(N+1), \mf p\neq \mf q} h(\mf p^{k})\overline{h}(\mf q^{\ell})\sum_{u\in QB_N+a, u\equiv a \mod Q, \mf p^{k}||u, \mf q^{\ell}||u}1.
    \end{split}
\end{equation}
Similar to (\ref{firstTK2}) and , we have that
\begin{equation}
    \begin{split}
        &\quad \frac{1}{\gamma_{K} N}\sum_{\NN(\mf p^{k}),\NN(\mf q^{\ell}) \leq N, \mf p\neq \mf q} h(\mf p^{k})\overline{h}(\mf q^{\ell})\sum_{u\in QB_N+a, u\equiv a \mod Q, \mf p^{k}||u, \mf q^{\ell}||u}1
        \\&=\sum_{\NN(\mf p^{k}),\NN(\mf q^{\ell})\leq \NN(Q)N, \mf p\neq \mf q} h(\mf p^{k})\overline{h}(\mf q^{\ell})\Bigl(\NN(\mf p^k)^{-1}\NN(\mf q^l)^{-1}\left(1-\frac{1}{\NN(\mf p)}\right) \left(1-\frac{1}{\NN(\mf q)}\right)+O(\frac{1}{N})\Bigr)
        \\&\leq \left|\sum_{\NN(\mf p) \leq N} h(\mf p^{k})\left( \frac{1}{\NN(\mf p^k)}-\frac{1}{\NN(\mf p^{k+1})}\right)\right|^2+\Theta_2+\epsilon(N),
      %  \\&=\vert\sum_{K<p^{k}\leq QN+a, p\nmid Q}h(p^{k})p^{-k}(1-p^{-1})\vert^{2}+error,
    \end{split}
\end{equation}

where 
\[
\Theta_2=\sum_{\NN(\mf p) \leq N, \mf p \nmid Q} |h(\mf p^k)|^2 \left( \frac{1}{\NN(\mf p^k)}-\frac{1}{\NN(\mf p^{k+1})}\right)^2
\]

{\color{red} under construction}

\

\

\


(essentially, exchanging the sums and using Lemma~\ref{countingCRT}). This yields an error term $\Theta_1$ of the form
\[
\Theta_1=\sum_{\NN(\mf p) \leq \NN(Q)N, \mf p \nmid Q} |h(\mf p^k)|^2 \left( \frac{1}{\NN(\mf p^k)}-\frac{1}{\NN(\mf p^{k+1})}\right)
\]
%Note that $\frac{1}{N}\sum_{u\leq QN+a, u\equiv a \mod Q}\sum_{p^{k}|| u} \vert h(p^{k})\vert^{2}$ can be bounded by
For the second term in \eqref{secondTK}, we also swap the order of the sums and rewrite it as
%\begin{equation}
 %   \begin{split}
  %      &\quad \frac{1}{(2N)^2}\sum_{p^{k}, q^l \in Q[\pm N]^2+a, p \neq q, p, q \in \Z[i]_+}\vert h(p^{k})\vert^{2}\sum_{u\leq QN+a, u\equiv a \mod Q, p^{k}\Vert u} 1 
   %     \\&\approx \frac{1}{N}\sum_{p^{k}\leq QN+a}\vert h(p^{k})\vert^{2}p^{-k}(1-p^{-1}) 
    %    \\&\leq \frac{1}{N}\sum_{p^{k}\leq QN+a}\vert h(p^{k})\vert^{2}p^{-k}
   % \end{split}
%\end{equation}
%(the RHS of (\cite{20}) is similar to this.)
% On the other hand, we have
\begin{equation}
    \begin{split}
        &\quad \frac{1}{\gamma_{K} N}\sum_{\NN(\mf p^{k}),\NN(\mf q^{\ell}) \leq \NN(Q)N, \mf p\neq \mf q} h(\mf p^{k})\overline{h}(\mf q^{\ell})\sum_{u\in QB_N+a, u\equiv a \mod Q, \mf p^{k}||u, \mf q^{\ell}||u}1
        \\&\approx \frac{1}{\gamma_{K} N}\sum_{\NN(\mf p^{k}),\NN(\mf q^{\ell})\leq \NN(Q)N, \mf p\neq \mf q} h(\mf p^{k})\overline{h}(\mf q^{\ell})\NN(\mf p^k)\NN(\mf q^l)\left(1-\frac{1}{\NN(\mf p)}\right) \left(1-\frac{1}{\NN(\mf q)}\right)
        \\&= \left|\sum_{\NN(\mf p) \leq \NN(Q)N} h(\mf p^{k})\left( \frac{1}{\NN(\mf p^k)}-\frac{1}{\NN(\mf p^{k+1})}\right)\right|^2+\Theta_2,
      %  \\&=\vert\sum_{K<p^{k}\leq QN+a, p\nmid Q}h(p^{k})p^{-k}(1-p^{-1})\vert^{2}+error,
    \end{split}
\end{equation}
where in the second step we, once again, used Lemma~\ref{countingCRT}, and $\Theta_2$ is another error term of the form
\[
\Theta_2=\sum_{\NN(\mf p) \leq \NN(Q)N, \mf p \nmid Q} |h(\mf p^k)|^2 \left( \frac{1}{\NN(\mf p^k)}-\frac{1}{\NN(\mf p^{k+1})}\right)^2
\]
\end{proof}

\fi

 \subsection{The concentration estimate}
In this section we state and prove the concentration estimate we need.
  Let 
    \begin{equation}\label{folnerMult}
    \Phi_{M}:=\left\{\prod_{\NN(\mathfrak{p})\leq M}\mathfrak{p}^{a_{\mathfrak{p}}}\colon M<a_{\mathfrak{p}}\leq 2M, \prod_{\NN(\mathfrak{p})\leq M}\mathfrak{p}^{a_{\mathfrak{p}}} \text{ is principal} \right\}.
       \end{equation} 
    Note that since we can identify principal ideals with elements of $\mathcal{O}_K$ by using their generators, the subsets $\Phi_M$ can be seen as subsets of $\mathcal{O}_K$. 

\begin{lemma}
    The family of sets $(\Phi_{M})_{M}$ forms a F\o lner sequence in $\mathcal{O}_K$.
\end{lemma}
\begin{proof}
    let $\{\mathfrak{p}_{1},\dots,\mf p_k, \dots\}$ be an ordering of all the prime ideals with non-decreasing norms, and suppose that the set of $\mathfrak{p}$ with $\NN(\mathfrak{p})\leq M$ is $\mathfrak{p}_{1},\dots,\mathfrak{p}_{b_{M}}$. 
%{\color{red}
Let $G$ denote the ideal class group of $K$. For any ideal $I$, let $v(I)\in G$ denote the ideal class containing $I$.  

%and suppose that $\mathfrak{p}_{i}$ belongs to the ideal class $h_{i}\in H$. Then $\prod_{i=1}^{n}\mf p_i^{a_i}$ is a principle ideal if and only if $\prod_{i=1}^{n}h_{i}^{a_{i}}=e_{H}$.

%}


    
 %   Let $s$ be the ideal class number of $\mathcal{O}_K$. We let $\vec{\varepsilon}=\vec{0}$ if $\mf p$ is principal, and otherwise, take it to be the element $\vec{\varepsilon} \neq \vec{0}$ of the finite abelian group (already identified with a group of the form $\Z/a_1\Z\times \dots \times \Z/a_k\Z$).  

%    We will now show the F\o lner property claimed for $(\Phi_M)_{M}$. 
Let $u \in \mathcal{O}_K$, and suppose that $(u)=\prod_{i=1}^m \mf p_i^{c_i}$.
    Let $M$ be large so that $m\leq b_{M}$. 
    %Then
  %  $$u\cdot\Phi_{K}=\left\{\prod_{i=1}^{b_{K}}\mathfrak{p}_{i}^{a'_{i}}\colon K+c_{i}<a'_{i}\leq 2K+c_{i}, \prod_{i=1}^{b_{K}}\mathfrak{p}_{i}^{a'_{i}} \text{ is principle} \right\}.$$
    %This is because if $I_{1}I_{2}=I_{3}$ for some ideals $I_{1},I_{2},I_{3}$ with $I_{1}$ being principle, then $I_{2}$ is principle if and only if $I_{3}$ is principle (since $I_{2}$ and $I_{3}$ are in the same ideal class). Therefore, we have that (let $c=\max_{i}c_{i}$)
z5   
%{\color{red}
Since $v((u))=\prod_{i=1}^{b_{K}}v(\mathfrak{p}_{i})^{c_{i}}=e_{G}$, we have
\begin{equation*}
    \begin{split}
     &   u\cdot\Phi_{M}=\left\{\prod_{i=1}^{b_{M}}\mathfrak{p}^{a_{i}+c_{i}}_{i}\colon M<a_{i}\leq 2M, \prod_{i=1}^{b_{M}}v(\mathfrak{p}_{i})^{a_{i}+c_{i}}=e_{G}\right\}
     \\&\quad\;\;\;\;\quad=\left\{\prod_{i=1}^{b_{M}}\mathfrak{p}^{a'_{i}}_{i}\colon M+c_{i}<a'_{i}\leq 2M+c_{i}, \prod_{i=1}^{b_{M}}v(\mathfrak{p}_{i})^{a'_{i}}=e_{G}\right\},
    \end{split}
\end{equation*}
%}
   %  $$u\cdot\Phi_{K}=\left\{\prod_{i=1}^{b_{K}}\mathfrak{p}^{a_{i}+c_{i}}_{i}\colon K<a_{i}\leq 2K, \text{ principal }\right\}=\left\{\prod_{i=1}^{b_{K}}\mathfrak{p}^{a'_{i}}_{i}\colon K+c_{i}<a'_{i}\leq 2K+c_{i}, \text{ principal } \right\},$$
   where $c_{i}=0$ for $i>m$. %, and we used the fact that the product of two principal ideals is again principal.
   Then $\Phi_{M}$ is nonempty as its cardinality is approximately $M^{b_{M}}/\vert G\vert$, and we have that
    $$\frac{\vert\Phi_{M}\Delta (u\cdot \Phi_{M})\vert}{\vert\Phi_{M}\vert}\lesssim\frac{\vert G\vert c^{m}(M+c)^{b_{M}-1}}{M^{b_{M}}}\to 0$$
    %$$\frac{\vert\Phi_{K}\Delta u\cdot \Phi_{K}\vert}{\vert\Phi_{K}\vert}\approx\frac{smc(K+c)^{b_{K}-1}}{K^{b_{K}}}\to 0$$
    as $M\to\infty$,  where $c=2\cdot\max_{1\leq i\leq m}c_{i}$. So $(\Phi_{M})_M$ is a F\o lner sequence. %In particular, $\Phi_{K}$ is nonempty as its cardinality is approximately $K^{b_{K}}/s$.
\end{proof}


Let $f,g\colon\mf  I(K)\to\C$ be multiplicative functions. The truncated distance between $f$ and $g$ is given by:
$$\D(f,g;M_{1},M_{2}):=\sum_{M_{1}<\NN(\mf p)\leq M_{2}}\frac{1}{\NN(\mf p)}(1-\text{Re}(f(\mf p)\overline{g}(\mf p))).$$


%With the help of Lemma~\ref{countingCRT}, we can now prove the main result of this section; the following concentration estimate:

We are now ready to state the main concentration estimate of the paper:

\begin{proposition}\label{p25}
    Let $K=\Q(\sqrt{-d})$ for some squarefree $d\in\N$ and $f\colon \mathcal{O}_K^{\times}\to\S$ be a multiplicative function that is trivial on units.\footnote{We have already discussed in Lemma~\ref{tv} that this can be assumed without loss of generality.}
    Suppose that there exist a modified Dirichlet character $\chi$ of period $I$, some $\tau\in\R$ and  some extension $\tilde{f'}$ of the function 
    $$f'(u):=f(u)\overline{\chi}(u)\NN(u)^{-i\tau}$$
    with
    $\D(\tilde{f'}, 1)<\infty$. %for some modified Dirichlet character $\chi$, some extension $\tilde{f\chi}$ of $f\chi$  and some $\tau\in\R$. 
  %  $f\sim \chi\cdot n^{i\tau}$ for some Dirichlet character $\chi$ with period $I$ and some $\tau\in\R$. 
  %  Put %{\color{red}
 %   $$f'(u):=f(u)\overline{\chi}(u)\NN(u)^{-i\tau}.$$
  %  Then $f'$ is a multiplicative function taking values in $\S$ that is still trivial on units.
  %  Let $\tilde{f'}$ denote an extension of $f'$ to ideal domained multiplicative functions.
  Let also $a \in \mathcal{O}_K^{\times}$ with $\NN(a) \leq \NN(Q)$ and $(a)$ being coprime to $(Q)$.
  %
 % $(a,Q)=1$. \WS{does it mean they have no comment integer divider or ideal divider?} \FM{Should be that there is no non-unit $t \neq 0$ that divides both $a, Q$.}\WS{On the equality after (46), we need $Q$ and $a$ to have no comment ideal factors!} \FM{Whoops, then, I guess we will have to assume that. I think that in any case we only use it for $a$ a unit, so it should be ok.}
    %}
 %   By abuse of notation, we will still denote by $f\cdot \overline{\chi}$ any extension of $f\cdot \overline{\chi}$ to ideal domained multiplicative functions.
%    $$\Phi_{K}:=\left\{\prod_{\NN(p)\leq K}p^{a_{p}}\colon K<a_{p}\leq 2K\right\}$$
   Suppose that $M$ is large enough so that $\Phi_M \subseteq I$. %Then, letting $Q$ be an element of smallest possible norm that belongs to $\prod_{\NN(\mf p)\leq K} \mf p$, {\color{red} I think we need $Q$ to range over $\Phi_{K}$}
 %  {\color{blue} This shouldn't be the case: we pick $Q$ in this way to ensure that (at least in terms of $\N$ it will be divisible enough, so it is a multiple of the period of the Dirichlet character $\chi$. What we are doing should be analogous for $\mathcal{O}$
 %  } 
   We have 
   %{\color{red}
\begin{equation}\nonumber
    \begin{split}
       & \limsup_{N\to\infty}\max_{Q\in \Phi_{M}, a\in\mathcal{O}_K^{\times},\atop \NN(a)\leq \NN(Q), (a,Q)=1}\E_{ \NN(u)\leq N}\vert f(Qu+a)-\chi(a)\NN(Qu)^{i\tau}\exp(F(\tilde{f'},M,N;Q))\vert
     %  \\&\ll \D(f,\chi\cdot \NN(u)^{i\tau};K,\infty)+K^{-1/2},
     \\&\quad\quad\quad\quad\quad\quad\quad\quad\quad\quad\quad\quad\quad\quad\quad\quad\quad\quad\quad\quad\quad\quad\quad\ll \D(\tilde{f'},1;M,\infty)+M^{-1/2},
    \end{split}
\end{equation}
%}
%Question: (1) do we need $(Q,a)=1$? (2) At some point we need to use the fact that $Q\in\Phi_{K}\subseteq I$. Make it explicit when we used it.}
%for any $a\in\mathcal{O}_K^{\times}$ with $\NN(a)\leq \NN(Q)$ and $a \mod I\in (\mathcal{O}_K/I)^{\times}$,
where  the implicit constant depends only on $K$ and
$$F(\tilde{f'},M,N;Q):=\sum_{M<\NN(\mf p)\leq N, \mf p \nmid Q}\frac{1}{\NN(\mf p)}(1-\text{Re}(\tilde{f'}(\mf p))).$$
 %{\color{red} Note here $Q$ are integers in $\Z[i]$ not $\Z$, and $2N^{2}$ is the norm of $N+iN$.}
 \end{proposition}

%{\color{red} I changed the statement of this proposition, but we need to check the proof again. In particular check if it works for modified Dirichlet characters.}
 
\begin{proof} The proof is based on \cite[Lemma~2.5]{KMPT}. %{\color{blue} Here in the proof of Lemma 2.5 that we mention, the first step is, interestingly enough, to change $\chi$ for its modified non-zero version. Even if we do not write it as explicitly, the result of the first step is to assume that $\tau=0$ and $\chi=1$, so the concentration estimate will definitely be fine with the modification we had in mind.} Write $f(u)=f'(u)\chi(u)\NN(u)^{i\tau}$. Then $f'$ is a multiplicative function taking values in $\U$.
%
%
Note that if $Q \in I$, then
    $$f(Qu+a)=f'(Qu+a)\chi(a)\NN(Qu+a)^{i\tau}=f'(Qu+a)\chi(a)\NN(Qu)^{i\tau}+O_{\tau}(\log (\NN(1+u^{-1})).$$
    because $|e^{ix}-e^{iy}| \leq |x-y|$ for real numbers $x, y$ and $\ln \frac{\NN(Qu+a)}{\NN(Qu)}$ converges to 0 as $\NN(u) \to \infty$. 
    %\WS{where does the error term $O_{\tau}(\NN(u)^{-1})$ come from?}
   %\WS{So here we need the stronger statement that 
    %$$\lim_{\NN(u)\to\infty} \NN(u)\ln \frac{\NN(Qu+a)}{\NN(Qu)}<\infty.$$
    %}
    On the other hand, since  $\log \NN(1+u^{-1}) \to 0$ as $u \to \infty$ in the sense of leaving compact sets, we have that 
    $\lim_{N\to\infty}\E_{\NN(u) \leq N}\log \NN(1+u^{-1})=0$.
  So it suffices to show that 
\begin{equation}
    \begin{split}
       & \limsup_{N\to\infty}\max_{Q\in \Phi_{M}, a\in\mathcal{O}_K^{\times},\atop \NN(a)\leq \NN(Q), (a,Q)=1}\E_{ \NN(u)\leq N}\vert f'(Qu+a)-\exp(F(\tilde{f'},M,N;Q))\vert
     %  \\&\ll \D(f,\chi\cdot \NN(u)^{i\tau};K,\infty)+K^{-1/2},
     \\&\quad\quad\quad\quad\quad\quad\quad\quad\quad\quad\quad\quad\quad\quad\quad\quad\quad\quad\quad\quad\quad\quad\quad\ll \D(\tilde{f'},1;M,\infty)+M^{-1/2},
    \end{split}
\end{equation}
%where 
%\[
%F_{N}(f',K)=\sum_{\NN(\mf p)\leq N, \mf p\nmid Q}\frac{1}{\NN(\mf p)}(f'(\mf p)-1).
%\]

\iffalse
{\color{red} I think this paragraph is not needed.}
Indeed, we are estimating sums of differences of $F_N(f,Q)-F_N(f',Q)$ which can be bounded by the sum of the reciprocals of norms of primes. These grow like $\log \log N$, so when taking an exponential, we get something that grows like $\log N$. Since we need to divide by $N$ from the expectation, the assertion is true.
\\ \\

\fi

Let $h: \mf I(K)\to\C$ be the additive function whose values on powers of prime ideals are given by $h(\mf p^{k}):=\tilde{f'}(\mf p^{k})-1$. Since $z=e^{z-1}+O(\vert z-1\vert^{2})$, we have
$$f'(Qu+a)=\prod_{\mf p^{k}\Vert Qu+a, \NN(\mf p)>M} \tilde{f'}(\mf p^{k})=\prod_{\mf p^{k}\Vert Qu+a, \NN(\mf p)>M}\Bigl(\exp(h(\mf p^{k}))+O(\vert h(\mf p^{k})\vert^{2})\Bigr),$$
where the first equality follows from the unique factorization of ideals and the fact that $\mf p\vert Qu+a\Rightarrow \NN(\mf p)\geq M$, and $p^{k}\Vert Qu+a$ means that $k\in\N$ is the largest natural number for which $Qu+a\in{\mf p}^{k}$.
By the inequality $\left|\prod_{i=1}^{k}z_{i}-\prod_{i=1}^{k}w_{i}\right|\leq\sum_{i=1}^{k}\vert z_{i}-w_{i}\vert$ for  $z_{i},w_{i}\in\U$, which can be proved by induction, we have that
\begin{equation}\label{ffd}
    \begin{split}
       f'(Qu+a)=\exp(h(Qu+a))+O\left(\sum_{\mf p^{k}\Vert Qu+a, \NN(\mf p)>M}\vert h(\mf p^{k})\vert^{2}\right).
    \end{split}
\end{equation}
Also
note that
\begin{equation}\label{ddff}
    \begin{split}
        &\quad F(\tilde{f'},M,N;Q)=\sum_{M<\NN(\mf p^{k})\leq N, \mf p \nmid Q}h(\mf p^{k})\NN(\mf p)^{-k}(1-\NN(\mf p)^{-1})+O(M^{-1})\\&\quad\quad\quad\quad\quad\quad\quad\;:=\mu_{h,M,N}+O(M^{-1}).
    \end{split}
\end{equation}
Indeed, the term with $k=1$ can be compared to the actual definition of $F(\tilde{f'},M,N;Q)$, giving us an error bounded above by 
\begin{equation}\label{ddff3}
    \begin{split}
        \sum_{M<\NN(\mf p)\leq N} \sum_{k=2}^{\infty}\NN(\mf p)^{-k}=\sum_{M<\NN(\mf p)\leq N} \frac{O(1)}{\NN(p)^{2}}=O(1/M).
    \end{split}
\end{equation}
%
%$\sum_{M<\NN(\mf p)\leq N} \NN(\mf p)^{-2}\ll 1/M$. 
%{\color{red} check this error term}. 
%On the other hand, the terms with $k \geq 2$ can be trivially bounded above by a similar sum, which is again $O(1/M)$. 

Combining (\ref{ffd}), (\ref{ddff}) together with the fact that 
  $\vert e^{w}-e^{z}\vert\ll \vert w-z\vert$, we have 
\begin{equation}\label{firstO1Kbound}
    \begin{split}
        &\quad\E_{\NN(u) \leq N}\vert f'(Qu+a)-\exp(F(\tilde{f'},M,N;Q))\vert
        \\&\ll \E_{\NN(u) \leq N}\vert \exp(h(Qu+a))-\exp(F(\tilde{f'},M,N;Q))\vert+\E_{\NN(u) \leq N}\sum_{\mf p^{k}\Vert Qu+a, \NN(\mf p)>M}\vert h(\mf p^{k})\vert^{2}
        \\&\ll \E_{\NN(u) \leq N}\vert h(Qu+a)-F(\tilde{f'},M,N;Q)\vert+\E_{\NN(u) \leq N}\sum_{\mf p^{k}\Vert Qu+a, \NN(\mf p)>M}\vert h(\mf p^{k})\vert^{2}
        \\&=\E_{\NN(u) \leq N}\vert h(Qu+a)-\mu_{h,M,N}\vert+\E_{\NN(u) \leq N}\sum_{\mf p^{k}\Vert Qu+a, \NN(\mf p)>M}\vert h(\mf p^{k})\vert^{2}+O(M^{-1}).
    \end{split}
\end{equation}

For the second term, we have
\begin{equation}\label{204}
    \begin{split}
        &\quad \E_{\NN(u) \leq N}\sum_{\mf p^{k}\Vert Qu+a, \NN(\mf p)>M}\vert h(\mf p^{k})\vert^{2}
        =\sum_{\NN(\mf p)>M, k\geq 1}\vert h(\mf p^{k})\vert^{2}\frac{1}{\gamma_{K} N}\sum_{\NN(u) \leq N}\mathds{1}_{\mf p^{k}\Vert Qu+a}
        \\&=\sum_{\NN(\mf p)>M, k\geq 1}\vert h(\mf p^{k})\vert^{2}\NN(\mf p^{k})^{-1}+o_{N\to\infty}(1) \text{ (by Lemma \ref{countingCRT})}.
    \end{split}
\end{equation}


%Indeed, this last step is a count of the size of an ``AP'' with a certain step. The bound is fairly trivial over $\N$. If we move to $\mathcal{O}_K^{\times}$ we simply make use of Lemma \ref{mv7}, which counts of elements of certain norms to deduce that the same bound will hold up to a certain absolute constant depending only on the number field. 

For the first term, by Theorem \ref{TK},
we have 
\begin{equation}\label{20}
    \begin{split}
        &\quad \E_{\NN(u) \leq N}\vert h(Qu+a)-\mu_{h,M,N}\vert 
        \\&=\E_{\NN(u) \leq N}\left| h(Qu+a)-\sum_{M<\NN(\mf p^{k})\leq N, \mf p\nmid Q}h(\mf p^{k})\NN(\mf p)^{-k}(1-\NN(\mf p)^{-1})\right|
        \\&\leq \left(\E_{\NN(u) \leq N}\left| h(Qu+a)-\sum_{M<\NN(\mf p^{k})\leq N,\mf p\nmid Q}h(\mf p^{k})\NN(\mf p)^{-k}(1-\NN(\mf p)^{-1})\right|^{2}\right)^{1/2}
        \\&\leq \left((2+o_{M;N\to\infty}(1))\cdot\sum_{\NN(\mf p)>M, k\geq 1}\vert h(\mf p^{k})\vert^{2}\NN(\mf p^{k})^{-1}+\right. 
        \\&\quad\quad\quad\quad\quad\quad\quad\quad\quad\;\quad \left. O(1)\cdot \frac{1}{\sqrt{N}}\sum_{\NN(\mf p^{k})\leq N', \NN(\mf p)>M}h(\mf p^{k})^{2}\cdot \NN(\mf p^{k})^{-1/2}\right)^{1/2}.
    \end{split}
\end{equation}



 
 
%{\color{red} (1): to get the last step, we need a Turan-Kubilius inequality for $\Z[i]$. (2): also I don't see why in [FKM] this step follows from Turan-Kubilius (the one on wikipedia). I see that the last inequality holds if the term $h(Qu+a)$ were $h(u)$, but how did they pass to $Qu+a$? One thing to look at from (5.2) to (5.8) of [FKM2], this is the place where they replace $Qu+a$ by $u$.}
%{\color{red} reference for the theorem: Introduction To Analytic And Probabilistic Number Theory(Tenenbaum), page 319 of the pdf file.}

 
%In order to fully justify the last inequality of (\ref{20}), we need to show that

Note that by the Cauchy-Schwarz inequality and Parts (i) and (ii) of Lemma~\ref{l2},
\begin{equation}\label{208}
    \begin{split}
      &\qquad\frac{1}{\sqrt{N}}\sum_{\NN(\mf p^{k})\leq N', \NN(\mf p)>M}\vert h(\mf p^{k})\vert^{2}\cdot \NN(\mf p^{k})^{-1/2}
    \leq \frac{4}{\sqrt{N}}\sum_{\NN(\mf p^{k})\leq N'} \NN(\mf p^{k})^{-1/2}
       \\ & \quad\quad\quad\quad\quad\quad\quad\quad\quad\quad\quad\leq \frac{4}{\sqrt{N}}\left(\sum_{\NN(\mf p^{k})\leq N'} \NN(\mf p^{k})^{-1}\right)^{1/2}\cdot \left(\sum_{\NN(\mf p^{k})\leq N'} 1\right)^{1/2}
       \\&\quad\quad\quad\quad\quad\quad\quad\quad\quad\quad\;\;\;=O\left(\frac{\log\log N'}{\log N'}\right)^{1/2}=o_{M;N\to\infty}(1).
    \end{split}
\end{equation}
 So now it remains to control $\sum_{\NN(\mf p)>M, k\geq 1}\vert h(\mf p^{k})\vert^{2}\NN(\mf p^{k})^{-1}$. Similarly to (\ref{ddff3}), the trivial bound for $|h(\mf p^k)|$ gives
\begin{equation}\label{201}
    \begin{split}
       \sum_{\NN(\mf p)>M, k\geq 2}\vert h(\mf p^{k})\vert^{2}\NN(\mf p^{k})^{-1}=O(M^{-1}).
    \end{split}
\end{equation}
 %Thus, we see that it is enough to bound
%\[
% \sum_{p>K, k \geq 2} \frac{\text{const}}{p^k}=\sum_{p>K} \frac{1}{p^2} \frac{p}{p-1} \ll \sum_{p>K} \frac{1}{p^2} \ll K^{-1},
%\]
 %as desired.
 On the other hand,
\begin{equation}\label{202}
    \begin{split}
\sum_{\NN(\mf p)>M}\vert h(\mf p)\vert^{2}\NN(\mf p)^{-1}=\sum_{\NN(\mf p)>M}2(1-\text{Re}\tilde{f'}( \mf p))\NN(\mf p)^{-1}=2\D(\tilde{f'},1;M,\infty),    
\end{split}
\end{equation}
The conclusion follows by
combining (\ref{firstO1Kbound}), (\ref{204}), (\ref{20}), (\ref{208}), (\ref{201}) and (\ref{202}).
\end{proof}

\iffalse

\begin{equation}
    \begin{split}
        &\quad \E_{\NN(u) \leq N}\left| h(Qu+a)-\sum_{\NN(\mf p^{k})\leq N,\mf p\nmid Q}h(\mf p^{k})\NN(\mf p)^{-k}(1-\NN(\mf p)^{-1})\right|^{2}
        \\&\ll (1+o_{N\to\infty}(1)) \sum_{\NN(\mf p)>M, k\geq 1}\vert h(\mf p^{k})\vert^{2}\NN(\mf p^{k})^{-1}.
    \end{split}
\end{equation}
In order to do that, as in the proof of the Tur\'an-Kubilius theorem, we let 
$$c_{M,K,Q}:=\sum_{M<\NN(\mf p^{k})\leq N,\mf p\nmid Q}h(\mf p^{k})\NN(\mf p)^{-k}(1-\NN(\mf p)^{-1}).$$ It suffices to show that 
$\E_{\NN(u) \leq N}h(Qu+a)\approx c_{M,K,Q}$ and $\E_{\NN(u) \leq N}\vert h(Qu+a)\vert^{2}\approx c_{M,K,Q}^{2}$.
First observe that, setting $B_N:=\{ u : \NN(u) \leq N\}$, it is
\begin{equation}\label{firstTK}
    \begin{split}
        &\quad \E_{\NN(u) \leq N}h(Qu+a)
        =\frac{1}{\gamma_{K} N}\sum_{u\in QB_N+a, u\equiv a \mod Q}h(u)
        \\&=\frac{1}{\gamma_{K} N}\sum_{u\in QB_N+a, u\equiv a \mod Q}\sum_{\mf p^{k}|| u} h(\mf p^{k})
        \\&=\frac{1}{\gamma_{K} N}\sum_{u\in QB_N+a, u\equiv a \mod Q}\sum_{\mf p^{k}|| u, \mf p\nmid Q} h(\mf p^{k})
        \\&=\frac{1}{\gamma_{K} N}\sum_{\NN(\mf p^{k})\leq \NN(Q)N, \mf p\nmid Q}h(\mf p^{k})\sum_{u\in QB_N+a, \mf p^{k}||u, u\equiv a \mod Q} 1
     %   \\&=\frac{1}{N^2}\sum_{p^{k}\in  [0,QN+a]^2, p\nmid Q}h(p^{k})(\sum_{u\in [0,QN+a]^2, p^{k}|u, u\equiv a \mod Q} 1-\sum_{u\in [0,QN+a]^2, p^{k+1}|u, u\equiv a \mod Q} 1)
  %  \\&\approx \frac{1}{N^2}\sum_{K<p^{k}\leq QN+a, p\nmid Q}h(p^{k})(\frac{QN+a}{p^{k}Q}-\frac{QN+a}{p^{k+1}Q})
    \\& \approx \sum_{\NN(\mf p^{k})\leq \NN(Q)N, \mf p\nmid Q}h(\mf p^{k})\NN(\mf p)^{-k}(1-\NN(\mf p)^{-1})
    \end{split}
\end{equation}
Note that for the last line, we used Lemma~\ref{countingCRT} and we factored out $\NN(\mf p^k)$ in the $\left( \frac{1}{\NN(\mf p^k)}-\frac{1}{\NN(\mf p^{k+1})}\right)$ bracket.

For the second term, we have

\begin{equation}\label{secondTK}
    \begin{split}
        &\quad \E_{\NN(u) \leq N}\vert h(Qu+a)\vert^{2}
        =\frac{1}{\gamma_{K} N}\sum_{u\in  QB_N+a, u\equiv a \mod Q}\vert h(u)\vert^{2}
        \\&=\frac{1}{\gamma_{K} N}\sum_{u\in  QB_N+a, u\equiv a \mod Q}\left| \sum_{\mf p^{k}|| u} h(\mf p^{k})\right|^{2}
        \\&=\frac{1}{\gamma_{K} N}\sum_{u\in  QB_N+a, u\equiv a \mod Q}\left(\sum_{\mf p^{k}|| u} \vert h(\mf p^{k})\vert^{2}+\sum_{\mf p^{k},\mf q^{\ell}|| u, \mf p\neq \mf q} h(\mf p^{k})\overline{h}(\mf q^{\ell})\right). 
    \end{split}
\end{equation}
The first term in the last line of \eqref{secondTK} can be dealt with in the same way as in \eqref{firstTK} (essentially, exchanging the sums and using Lemma~\ref{countingCRT}). This yields an error term $\Theta_1$ of the form
\[
\Theta_1=\sum_{\NN(\mf p) \leq \NN(Q)N, \mf p \nmid Q} |h(\mf p^k)|^2 \left( \frac{1}{\NN(\mf p^k)}-\frac{1}{\NN(\mf p^{k+1})}\right)
\]
%Note that $\frac{1}{N}\sum_{u\leq QN+a, u\equiv a \mod Q}\sum_{p^{k}|| u} \vert h(p^{k})\vert^{2}$ can be bounded by
For the second term in \eqref{secondTK}, we also swap the order of the sums and rewrite it as
%\begin{equation}
 %   \begin{split}
  %      &\quad \frac{1}{(2N)^2}\sum_{p^{k}, q^l \in Q[\pm N]^2+a, p \neq q, p, q \in \Z[i]_+}\vert h(p^{k})\vert^{2}\sum_{u\leq QN+a, u\equiv a \mod Q, p^{k}\Vert u} 1 
   %     \\&\approx \frac{1}{N}\sum_{p^{k}\leq QN+a}\vert h(p^{k})\vert^{2}p^{-k}(1-p^{-1}) 
    %    \\&\leq \frac{1}{N}\sum_{p^{k}\leq QN+a}\vert h(p^{k})\vert^{2}p^{-k}
   % \end{split}
%\end{equation}
%(the RHS of (\cite{20}) is similar to this.)
% On the other hand, we have
\begin{equation}
    \begin{split}
        &\quad \frac{1}{\gamma_{K} N}\sum_{\NN(\mf p^{k}),\NN(\mf q^{\ell}) \leq \NN(Q)N, \mf p\neq \mf q} h(\mf p^{k})\overline{h}(\mf q^{\ell})\sum_{u\in QB_N+a, u\equiv a \mod Q, \mf p^{k}||u, \mf q^{\ell}||u}1
        \\&\approx \frac{1}{\gamma_{K} N}\sum_{\NN(\mf p^{k}),\NN(\mf q^{\ell})\leq \NN(Q)N, \mf p\neq \mf q} h(\mf p^{k})\overline{h}(\mf q^{\ell})\NN(\mf p^k)\NN(\mf q^l)\left(1-\frac{1}{\NN(\mf p)}\right) \left(1-\frac{1}{\NN(\mf q)}\right)
        \\&= \left|\sum_{\NN(\mf p) \leq \NN(Q)N} h(\mf p^{k})\left( \frac{1}{\NN(\mf p^k)}-\frac{1}{\NN(\mf p^{k+1})}\right)\right|^2+\Theta_2,
      %  \\&=\vert\sum_{K<p^{k}\leq QN+a, p\nmid Q}h(p^{k})p^{-k}(1-p^{-1})\vert^{2}+error,
    \end{split}
\end{equation}
where in the second step we, once again, used Lemma~\ref{countingCRT}, and $\Theta_2$ is another error term of the form
\[
\Theta_2=\sum_{\NN(\mf p) \leq \NN(Q)N, \mf p \nmid Q} |h(\mf p^k)|^2 \left( \frac{1}{\NN(\mf p^k)}-\frac{1}{\NN(\mf p^{k+1})}\right)^2
\]

 
 So now it remains to control $\sum_{\NN(\mf p)>K, k\geq 1}\vert h(\mf p^{k})\vert^{2}\NN(\mf p^{k})^{-1}$. As we saw before, the trivial bound for $|h(\mf p^k)|$ gives
 $$\sum_{\NN(\mf p)>M, k\geq 2}\vert h(\mf p^{k})\vert^{2}\NN(\mf p^{k})^{-1}\ll K^{-1}.$$
 %Thus, we see that it is enough to bound
%\[
% \sum_{p>K, k \geq 2} \frac{\text{const}}{p^k}=\sum_{p>K} \frac{1}{p^2} \frac{p}{p-1} \ll \sum_{p>K} \frac{1}{p^2} \ll K^{-1},
%\]
 %as desired.
 On the other hand,
 $$\sum_{\NN(\mf p)>M}\vert h(\mf p)\vert^{2}\NN(\mf p)^{-1}=\sum_{\NN(\mf p)>M}2(1-\text{Re}\tilde{f'}( \mf p))\NN(\mf p)^{-1}=2\D(\tilde{f'},1;M,\infty),$$
so putting everything together, we have that 
{\color{red} put everything together}

\fi

\section{Proofs of the recurrence results}\label{sec: proofmain}



%{\color{blue} I realized that in FKM, Pythagorean pairs of type 1 do not require concentration estimates. It might be a good idea to try to see if we can prove that part with our current knowledge in $\Z[i]$. This would correspond to Section 4 in [FKM] (not [FKM2]).}
%{\color{red} Yes the next step is to check if Section 4 in [FKM] works for the following sample case (say, take $a=b=1$ and $c=-1$):
In this section, we prove Propositions \ref{aperiodiclimit0},
\ref{zeroaveragesadeltabdelta}, and
\ref{fffss}.
  As we saw in Section~\ref{sec: proofstrat}, they will imply the partition regularity results promised in Section \ref{s1}.
%\begin{equation}\label{e14} 
 %   \liminf_{N \to \infty} \frac{1}{N^4}\sum_{n, m \in \Z[i] \cap [-N,N]^2} \int_{\mathcal{M}} f(\ell(m^2-n^2))\overline{f(\ell'mn)} \ d\sigma(f)>0.
%\end{equation}
%\begin{equation}\label{e15} 
 %   \liminf_{N \to \infty} \frac{1}{N^4}\sum_{n, m \in \Z[i] \cap [-N,N]^2} \int_{\mathcal{M}} f(\ell(m^2+n^2))\overline{f(\ell'mn)} \ d\sigma(f)>0.
%\end{equation}
%{\color{red} here please feel free to replace $[-N,N]^2$ by any convex set which is convenient.
%}

%}
Throughout this section, we assume that $K=\Q(\sqrt{-d})$ with $d\in\N$ squarefree.
We will make use of the next two simple lemmas:
\begin{lemma}\label{multav0}
Let $f: \mathcal{O}_K^{\times} \to \U$ be a non-trivial multiplicative function (that is, $f \neq 1$). Let $\Phi_M$ be a multiplicative F\o lner sequence in $\mathcal{O}_K^{\times}$. Then,
\[
\lim_{M \to \infty} \E_{u \in \Phi_M} f(u) = 0.
\]
\end{lemma}
\begin{proof}
Since $\U$ is compact, we consider limit points of the sequence $(\E_{u \in \Phi_M}f(u))_{M \in \N}$. Abusing notation, suppose that the limit of this sequence exists. 

Since $f\neq 1$, let $v \in \mathcal{O}_K^{\times}$ be such that $f(v)\neq 1$. Since $\Phi_M$ is a F\o lner sequence, we can write
\[
\lim_{M \to \infty} \E_{u \in \Phi_M} f(u) = \lim_{K \to \infty} \E_{u \in v\Phi_M} f(u) = \lim_{M \to \infty} \E_{u \in \Phi_M} f(vu) = f(v) \lim_{M \to \infty} \E_{u \in \Phi_M} f(u).
\]
Since $f(v) \neq 1$, this implies that the limit in question must be $0$. Therefore, $0$ is the only possible accumulation point for this sequence, so compactness of $\U$ completes the proof.
\end{proof}
\begin{lemma}\label{sequenceavfact}
Let $d\in\N$ be squarefree and $a: \OK^{\times} \to \U$ be a sequence. Let $l_1, l_2 \in \OK$ not both $0$. Suppose that for some $\varepsilon>0$ and some sequence $b: \N \to \U$ we have
\[
\limsup_{N \to \infty} \E_{\NN(u) \leq N} |a(u)-b(N)| \leq \varepsilon.
\]
Then,
\[
\limsup_{N \to \infty} \E_{\NN(n), \NN(m) \leq N} |a(l_1m+l_2n)-b(lN)| \ll_{K,l} \varepsilon,
\]
where $l:=2(\NN(l_1)+\NN(l_2))$.
\end{lemma}
\begin{proof}
We can estimate
\begin{equation}\label{sequenceavfactE1}
\E_{\NN(n), \NN(m) \leq N} |a(l_1m+l_2n)-b(lN)| \ll_{K,l} \frac{1}{N^2}\sum_{\NN(u) \leq lN} w_N(u) |a(u)-b(lN)|,
\end{equation}
%
%{\color{red}
%
%In our case, we want to have something like 
%\begin{equation}%\label{sequenceavfactE1}
%\E_{m,n\colon \NN(L_{j}(m,n))\leq N, 1\leq j\leq 4} |a(l_1m+l_2n)-L_{lN}| \ll \frac{1}{N^2}\sum_{\NN(u) \leq lN} w_N(u) |a(u)-L_{lN}|,
%\end{equation}
%and we are in good shape if we can have a good upper bound for $$w_N(u):=\left|\{(m,n): \NN(L_{j}(m,n))\leq N, 1\leq j\leq 4 \text{ and }l_1m+l_2n=u\}\right|$$
%}
%
where for $u \in \OK^{\times}$ we put $$w_N(u):=\left|\{(m,n)\in (\OK)^{2}: \NN(n), \NN(m) \leq N \text{ and }l_1m+l_2n=u\}\right|.$$ Since, given an $m$ with $\NN(m) \leq N$, there is only one solution to $l_1m+l_2n=u$ (as an equation in $n$), we see that $w_N(u) \ll_{K} N$. On the other hand, it is not hard to see that for the $u\in\OK$ for which $w_{N}(u)$ is non-empty has size at most $lN$. So the right hand side of \eqref{sequenceavfactE1} is bounded above   by
\[
\ll_{K,l}\E_{\NN(u) \leq lN} |a(u)-b(lN)|,
\]
so we are done, given our starting hypothesis.
\end{proof}
Before moving onto the main proof, we need a couple of preliminary results. Consider the sets on $\S^1$ given by
\[
I(\delta):=\{\exp(2\pi i\phi) : \phi \in (-\delta,\delta)\}.
\]
%\begin{remark}
%Throughout the rest of the section, and to guarantee the combinatorial result we promised in the Introduction, it will be important to ensure that we can pick elements $m,n$ from the averages that are distinct. This can be easily accomplished by demanding the stronger condition that $(m,n) \in S_q:=\left\{ (m,n): mn(m+\alpha n)(m+\beta n) \neq 0 \text{ and }\right.$ \\ $\left.\ell (m+\alpha n)(m+\beta n) \neq \ell' mn\right\}$. We will do so by introducing the relevant characteristic function, and noting that if $w: \mathcal{O}_K^{\times} \times \mathcal{O}_K^{\times} \to \C$ is a bounded function, then
%\[
%\left| \E_{\NN(m),\NN(n) \leq N} w(m,n)- \E_{\NN(m),\NN(n) \leq N} w(m,n) \mathds{1}_{S_q}(m,n)\right| \ll \]
%\[
%\E_{\NN(m),\NN(n) \leq N} \mathds{1}_{S_q}(m,n)\to 0, \text{ as } N \to \infty.
%\]
%(Indeed, note that the $L_i(m,n)$ not being $0$ part constitutes a set of $0$ density, and the quadratic part can also be easily seen to form a set of $0$ density, by considering the stronger weaker condition $\NN(\ell(m+\alpha n)(m+\beta n)) \neq \NN(\ell' mn)$, given that the norm of a quadratic imaginary field is non-negative).
%\end{remark}
%\WS{Do we really need this remark?}


\begin{lemma}\label{weights}
Let $0<\delta<\frac{1}{2}$ and consider the trapezoidal function $F_{\delta}: \S^1 \to [0,1]$ which is equal to $1$ on $I_{\delta/2}$ and $0$ outside of $I_{\delta}$. For $\ell,\ell' \in \mathcal{O}_K^{\times}$ and $\alpha,\beta \in \OK$, let
\begin{equation}\label{weight1}
w_{\delta}(m,n):= F_{\delta}\left(\NN\Bigl(\frac{\ell (m+\alpha n)(m+\beta n)}{\ell' mn}\Bigr)^i \right)\cdot \mathds{1}_{S_q}(m,n), \quad m, n \in \mathcal{O}_K^{\times}.
\end{equation}
%and
%\begin{equation}\label{weight1}
%\tilde{w}_{\delta}(m,n):= F_{\delta}((\NN(\ell(m^2+n^2))^i\cdot (\NN(\ell'mn))^{-i})\mathds{1}_{\NN(m) \neq \NN(n)}, \quad m, n \in \mathcal{O}^{\times}
%\end{equation}
Then,
\begin{equation}\label{weight2}
\lim_{N \to \infty} \E_{\NN(m), \NN(n) \leq N } w_{\delta}(m,n)>0.
%\text{ and } \lim_{N \to \infty} \E_{\NN(n), \NN(m) \leq N } \tilde{w}_{\delta}(m,n)>0,.
\end{equation}
\end{lemma}
\begin{proof}
%We will show that the limit in \eqref{weight1} is positive. Indeed, said limit can be rewritten (using homogeneity in all the relevant definitions) as
The left hand side of (\ref{weight2}) can be rewritten as 
\begin{equation}\label{weight3}
\lim_{N \to \infty} \E_{m',n'\in\frac{1}{N}\OK\colon\NN(m'), \NN(n') \leq 1}F_{\delta}\left(\NN\Bigl(\frac{\ell (m'+\alpha n')(m'+\beta n')}{\ell' m'n'}\Bigr)^i \right) \mathds{1}_{S_q}(m',n').
\end{equation}
%\[
%\lim_{N \to \infty} \E_{\NN(n), \NN(m) \leq N} F_{\delta}((\NN(\ell\cdot (m/N+\alpha n/N)\cdot (m/N+\beta n/N))^i\cdot (\NN(\ell'\cdot m/N\cdot n/N))^{-i}) \times
%\]
%\[
%\mathds{1}_{S_q}(m/N,n/N)
%\]

We extend the domain of $\NN$ from $K$ to $\C$ by continuity, since $K$ is dense in $\C$.
%Let $\NN: \C \to [0,\infty)$ be the norm in $\C$ given by
%$\NN(z):=z\sigma(z)$ for all $z\in\C$, where $\sigma$ is the non-trivial element of $\text{Gal}(K|\Q)$. \WS{I don't understand this extension} \FM{}
Let $B:=\{ w \in \C : \NN(w) \leq 1\}$. Define
  $$G_{\delta}(w,z):=F_{\delta}\left(\NN\Bigl(\frac{\ell (w+\alpha z)(w+\beta z)}{\ell' wz}\Bigr)^i \right)\cdot\mathds{1}_{B\times B, (w,z) \in S_q}(w,z).$$ 
  Then $G_{\delta}$ is continuous except for a set of measure zero with respect to the Lebesgue measure on $B \times B$, where $B:=\{ w \in \C : \NN(w) \leq 1\}$ -- here $(w,z) \in S_q$ is checked with the same definition, which also makes sense for $w,z \in \C$. Reinterpreting (\ref{weight3}) as Riemann sums, we may rewrite (\ref{weight3}) as   
\[
\int_{B}\int_{B} G_{\delta}(w,z) \ dwdz,
\]
where $dwdz$ is the appropriately normalized Lebesgue measure on $B \times B$. %Note that this is the case because $G_{\delta}$ only fails to be continuous on the set of points given by the set $S_q$, which is of Lebesgue measure zero (see the comments in the Remark above).

 

It remains to show that the integral of $G_{\delta}$ over $B \times B$ is positive. Since $G_{\delta}$ is non-negative, it suffices to show that $G_{\delta}$ is not identically $0$ outside the set of points where either $z$ or $w$ is equal to $0$; more particularly, we will show it is non-zero along a line of the form $w=az$, for some real number $a>1$. It follows that $G_{\delta}$ is non-zero for a small tube around said line (since $a>1$ implies that $w,z$ are distinct and non-zero), which implies the asserted positivity. 

In order to see that this is the case, observe that if we set $w=az$ for $a \in \R$ and $k \in \N$, the property $\NN(\ell(w+\alpha z)(w+\beta z))=\NN(\ell'wz)e^{2k\pi}$ will be satisfied, provided that we can solve the equation 
\[
\NN(\ell(a+\alpha)(a+\beta))=\NN(\ell'a)e^{2k\pi}.
\]
It is easy to see that if we move $\NN(a)$ to the left hand side of this equation, the function $H(a):=\NN((1+\beta a^{-1})(a+\alpha))$ is well defined for $a>0$, and because it is increasing, its range contains a ray of the form $[x_0,\infty)$. Now, simply pick $k \in \N$ such that $\frac{\NN(\ell')}{\NN(\ell)}e^{2k\pi}$ falls in this ray.

After raising such choices of norms of $w$ and $z$ to the $i$ and $-i$ powers respectively, we find that the function $F_{\delta}$ takes on values that are very close to $1$. Thus, by continuity, all points on a tubular neighbhorhood around it, will have value at least $1/2$, and also be non-zero, and such that $w \neq z$. This implies, as we argued, that $\int_B\int_B G_{\delta} \ dwdz>0$, as desired. 


%It is easy to check that one of the solutions is $a = \frac{\sqrt{4\NN(\ell)^2+\NN(\ell')^2}+\NN(\ell')}{2\NN(\ell)}$, which ensures that indeed, we can solve the equation and also have $a>1 $. Thus, on all points of the line $w=az$ inside the set $B \times B$, we have that $G_{\delta} \neq 0$, which in particular implies, as we argued, that $\int_B\int_B G_{\delta} \ dwdz>0.$ 

%For the second weight in \eqref{weight2} the argument we use is essentially the same, so we only give a brief outline.

%First, we consider the continuous function $\tilde{G}_{\delta}(w,z):=F_{\delta}( (\NN(\ell(w^2+z^2)))^i \cdot (\NN(\ell'\cdot w \cdot z))^{-i}\mathds{1}_{B \times B, \NN(w)\neq\NN(z)>0}(w,z)$. A similar argument used for the first weight, shows that the second limit in \eqref{weight2} is, in fact, equal to
%\[
%\int_{S}\int_{S} \tilde{G}_{\delta}(w,z) \ dwdz,
%\]
%where $dwdz$ is the same measure as for the first weight. Thus, in order to complete the proof, it only remains to show that $\tilde{G}_{\delta}$ is not identically zero. In order to show this, we take $b:=\NN(\ell'/\ell)e^{2k\pi}$ for some $k$ large enough so that $b>2$. Then, setting $a:=\frac{b+\sqrt{b^2-4}}{2}$ gives us an element such that
%\[
%\frac{\NN(\ell)}{\NN(\ell')}(a^2+1) = e^{2k\pi} a,
%\]
%where $a>1$. This being the case ensures that the argument of $\tilde{G}_{\delta}$ is equal to $1$ for non-zero values of $w, z$. As before, this implies that said function is positive on a small tube around the line $w=az$, so $\int_B\int_B \tilde{G}_{\delta} \ dwdz>0$ as we wished to show.
\end{proof}

\iffalse
Next we prove an analog of \cite[Proposition~4.1]{FKM1}: \WS{Is this proposition \ref{aperiodiclimit0}}
\begin{proposition}\label{aperiodic0}
Let $f: \mathcal{O}_K^{\times} \to \U$ be an aperiodic completely multiplicative function, let $\ell,\ell', Q,\alpha,\beta \in \mathcal{O}_K$ and $\delta>0$. Then, with the weight $w_{\delta}$  defined in Lemma~\ref{weights} we have
\[
\lim_{N \to \infty} \E_{\NN(m), \NN(n) \leq N} w_{\delta}(m,n) \cdot f(\ell \cdot (Qm+1+\alpha Qn)\cdot (Qm+1+\beta Qn) \overline{f(\ell' \cdot (Qm+1)Qn) } = 0.
\]
%and
%\[
%\lim_{N \to \infty} \E_{\NN(m), \NN(n) \leq N} \tilde{w}_{\delta}(m,n) \cdot f(\ell((Qm+1)^2+(Qn)^2))\cdot \overline{f(\ell'(Qm+1)(Qn))} = 0.
%\]
\end{proposition}
\fi
\begin{proof}[Proof of Proposition \ref{aperiodiclimit0}]
By the definition of $w_{\delta}(m,n)$, we can use the Stone-Weierstrass theorem on $\S^1$, so that together with linearity of the limit, the continuous function $F_{\delta}$ appearing in the weights can be replaced by a power $z^k$. This will simplify matters considerably, transforming the averages we need to show converge to $0$ into 
\[
\lim_{N \to \infty} \E_{\NN(m), \NN(n) \leq N} \mathds{1}_{S_q}(m,n)\NN\Bigl(\frac{\ell(m+\alpha n)(m+\beta n)}{\ell' mn}\Bigr)^{ki}  \times\]
\[
f(\ell(Qm+1+\alpha Qn)(Qm+1+\beta Qn))\cdot \overline{f(\ell'(Qm+1)(Qn))},
\]
for $k \in \Z$. Note that, with the same trick we used in Lemma~\ref{weights} we may as well replace $n, m$ by $Qn$ and $Qm$ respectively inside the norms, because the changes cancel out with the complex conjugate.

We next argue that, in fact, we can change $Qm$ by $Qm+1$ without affecting the averages. Indeed, for given $n \in \N$, it is not hard to see that 
$$\lim_{\NN(m) \to \infty} \log \frac{\NN(\ell(Qm+Q\alpha n)(Qm+Q\beta n))}{\NN(\ell(Qm+1+\alpha Qn)(Qm+1+\beta Qn))}=\lim_{\NN(m) \to \infty} \log \frac{\NN(\ell'(Q^2mn))}{\NN(\ell'(Qm+1)Qn)}=0.$$ So this change is valid. 

Therefore, it is enough to establish the convergence to $0$ of the averages
\[
\E_{\NN(m), \NN(n) \leq N} \mathds{1}_{S_q}(m,n)f_k(\ell(Qm+1+\alpha Qn)(Qm+1+\beta Qn))\cdot \overline{f_k(\ell'(Qm+1)(Qn))},
\]
where $f_k(n):=f(n) \cdot \NN(n)^{ik}$.
%{\color{red} Double check this, it should be correct}. {\color{blue}{Yep, that's what we want to show.}} 
Since $f$ is aperiodic, so are $f_k$ and $\overline{f_k}$ by Theorem \ref{characterization}. So the conclusion follows from Proposition~\ref{P: appendix}. %so using the analog of Frantzikinakis and Host's theorem for $\mathcal{O}_K$, namely, Proposition~\ref{P: appendix} we are done. 
%
%The argument for the second expression is exactly the same, we simply replace the minus with the plus. Note that we can still apply Proposition~\ref{P: appendix} because over $\mathcal{O}$ we have the factorization $(Qm+1)^2+(Qn)^2=(Qm+1+iQn)(Qm+1-iQn)$, which is needed for this application to be valid.
%{\color{red}
%Here is a point where we need the theorems to diverge between different rings of integers. 
%
%This simply means that it has to be stated for the correct equation which will be such that a parametrization of it will ahve one of its factors be of the form
%\[
%y= (m^2+dn^2)=(m+\sqrt{d}n)(m-\sqrt{d}n).
%\]
%Then, the argument employed will be the same, but the partition regularity result will be for a different equation.
%
%}
\end{proof}
%Using the weight described above, we now use a variation of \cite[Lemma~2.6]{FKM1}: {\color{blue} Is this Lemma 2.13?}
%\begin{lemma}
%    Let $f\colon \mathcal O_K^{\times}\to\S^{1}$ be a completely multiplicative function which is pretentious but not super pretentious. Let $\delta>0$. Then
%    \begin{equation} \label{e116}
 %   \lim_{M\to\infty}\E_{Q\in \Phi_{M}}\lim_{N \to \infty} \E_{\NN(m), \NN(n) \leq N} w_{\delta}(m,n)f(\ell(Qm+1+\alpha Qn)(Qm+1+\beta Qn))\overline{f(\ell'(Qm+1)Qn)}=0.
%\end{equation}
%and
%\begin{equation}\label{weight2withf}
%    \lim_{K\to\infty}\E_{Q\in \Phi_{K}}\lim_{N \to \infty} \E_{\NN(m), \NN(n) \leq N} \tilde{w}_{\delta}(m,n)f(\ell((Qm+1)^2+(Qn)^2))\overline{f(\ell'(Qm+1)Qn)}=0
%\end{equation}
%where $w_{\delta}$ is as defined in Lemma~\ref{weights}, and the limit exists because of Proposition~\ref{aperiodic0}.
%\end{lemma}



\begin{proof}[Proof of Proposition \ref{zeroaveragesadeltabdelta}]
%Let us see that \eqref{e116} holds.
%As we will point out later, the proof for \eqref{weight2withf} is very similar and only requires some small modifications.
By hypothesis on $f$,  $f(\e)=1$ for all units $\e$ and there exist some modified Dirichlet character $\chi$ of period $I$, some $\tau\in\R$, and some extension $\tilde{f'}$ of the function
$$f'(u):=f(u)\overline{\chi}(u)\NN(u)^{-i\tau}$$
such that 
  $\D(\tilde{f'}, 1)<\infty$.
Fix $\varepsilon>0$ and take $M_{0}=M_{0}(\varepsilon,f)$ so that 
\begin{equation}\label{tail}
    \sum_{\NN(\mf p)>M_{0}}\frac{1}{\NN( \mf p)}(1-\text{Re}\tilde{f'}(\mf p))+M_{0}^{-1/2}\leq \varepsilon
\end{equation}
and that the set $\Phi_{M_{0}}$ defined in (\ref{folnerMult}) is contained in $I$.  
By Proposition \ref{p25} and (\ref{tail}), %if $N$ is sufficiently large, then
$$\limsup_{N\to\infty}\E_{\NN(u)\leq N}\vert f(Qu+1)-\NN(Qu)^{i\tau}\exp(F(\tilde{f'},M,N;Q))\vert\ll\varepsilon$$
for every $M>M_{0}$ and $Q\in\Phi_{M}$.
For fixed  $\alpha \in \mathcal{O}_K$, it follows from Lemma~\ref{sequenceavfact}  that
\[
\limsup_{N\to\infty}\E_{\NN(m), \NN(n) \leq N}\left| f(Q(m+ \alpha n))+1)-\NN(Q(m+\alpha n))^{i\tau}\exp(F(\tilde{f'},M,lN;Q))\right|\ll\varepsilon,
\]
where $l_{\alpha}:=2(\NN(\alpha)+1)$.
%{\color{red} Note: in [FKM], for the $m-n$ case, the average was taken over those $m>n$ to ensure $m^{2}-n^{2}>0$.}
Combining the above two estimates together with the triangle inequality, %(since all the involved functions have absolute value less than one), 
we have that
 \begin{equation}\label{fweqfwf} 
    \begin{split}
&\limsup_{N\to\infty}\E_{\NN(m),\NN(n)\leq N}\Bigl\vert f\Bigl(\frac{(Qm+1+Q\alpha n)(Qm+1+Q\beta n)}{Qm+1}\Bigr) 
       \\&-\NN\Bigl(\frac{(Qm+Q\alpha n)(Qm+Q\beta n)}{Qm}\Bigr)^{i\tau}\exp(F(\tilde{f'},M,l_{\alpha}N;Q)+F(\tilde{f'},M,l_{\beta}N;Q)\\&\quad\quad\quad\quad\quad\quad\quad\quad\quad\quad\quad\quad\quad\quad\quad\quad\quad\quad\quad\quad\quad\quad\quad\quad\quad\quad\quad-F(\tilde{f'},M,N;Q))\Bigr\vert\ll\varepsilon. 
    \end{split}
\end{equation}
We now multiply both sides of (\ref{fweqfwf}) by $w_{\delta}(m,n)\cdot f(\frac{\ell}{\ell'n})\NN(Q)^{i\tau}$, which is bounded in absolute value by $1$, to obtain
 \begin{equation}\label{fpwepfw} 
    \begin{split}
       &\limsup_{N\to\infty}\E_{\NN(m), \NN(n) \leq N}  
         \Bigl\vert w_{\delta}(m,n)\cdot f\Bigl(\frac{\ell(Qm+1+Q\alpha n)(Qm+1+Q\beta n)}{\ell'(Qm+1)Qn}\Bigr)(f(Q)\NN(Q)^{-i\tau})
       \\&-w_{\delta}(m,n)\cdot f(\frac{\ell}{\ell'n}) \NN\left(\frac{(m+\alpha n)(m+\beta n)}{mn}\right)^{i\tau} 
       \cdot\exp(F(\tilde{f'},M,l_{\alpha}N;Q)+F(\tilde{f'},M,l_{\beta}N;Q)\\&\quad\quad\quad\quad\quad\quad\quad\quad\quad\quad\quad\quad\quad\quad\quad\quad\quad\quad\quad\quad\quad\quad\quad\quad\quad\quad\quad\quad\; -F(\tilde{f'},M,N;Q))\Bigr\vert\ll\varepsilon. 
    \end{split}
\end{equation}
Let
    $$\tilde{L}_{\delta}(Q):=\lim_{N \to \infty} \E_{\NN(m), \NN(n) \leq N}  w_{\delta}(m,n) \cdot f\Bigl(\frac{\ell(Qm+1+Q\alpha n)(Qm+1+Q\beta n)}{\ell'(Qm+1)Qn}\Bigr)(f(Q)\NN(Q)^{-i\tau}).$$
    Note that if $Q\in\Phi_{M}$, then $F(\tilde{f'},M,N;Q)$ is independent of $Q$ for all $N$. So  (\ref{fpwepfw}) implies that $|\tilde{L}_{\delta}(Q)-\tilde{L}_{\delta}(Q')|\ll \varepsilon$  for all $Q, Q' \in \Phi_M$. 
%
%    
%Since the second term of (\ref{fpwepfw}) is independent of $Q$, if $M>M_0$, we have that  $|\tilde{L}_{\delta}(Q)-\tilde{L}_{\delta}(Q')|\ll \varepsilon$  for all $Q, Q' \in \Phi_M$. 
%
%
Therefore, letting $M\to \infty$ we obtain
$$\limsup_{M\to\infty}\max_{Q,Q'\in\Phi_{M}}\vert\tilde{L}_{\delta}(Q)-\tilde{L}_{\delta}(Q')\vert=0.$$
For each $M\in\N$, pick some $Q_{M}\in\Phi_{M}$. Then
%the left hand side of (\ref{e116}) is equal to
$$\lim_{M \to \infty} \E_{Q \in \Phi_M} \lim_{N \to \infty} \E_{\NN(n),\NN(m) \leq N} A_{\delta}(f,Q; m,n)=\limsup_{M\to\infty}\tilde{L}_{\delta}(Q_{M})\E_{Q\in\Phi_M}\overline{f}(Q)\NN(Q)^{i\tau}.$$
Since $Q\mapsto \overline{f}(Q)Q^{i\tau}$ is a nontrivial multiplicative function, it follows from Lemma~\ref{multav0} that $\lim_{M \to \infty} \E_{Q\in\Phi_{M}}\overline{f}(Q)\NN(Q)^{i\tau}=0$ and we are done.
%{\color{red}
%As before, the main change that has to happen in this part of the argument is that of the parametrization, depending on the ring of integers under consideration, and repeat the proof, so it's probably fine to just have one paragraph.

%}
%The changes needed for \eqref{weight2withf} are very minor. Namely, all the expressions of the form $(Qm+1)^2-(Qn)^2$ get replaced by $(Qm+1)^2+(Qn)^2$. Then notice that Lemma~\ref{sequenceavfact} also implies that 
%\[
%\limsup_{N\to\infty}\E_{\NN(m), \NN(n) \leq N, \NN(m)\neq\NN(n)}\left| f(Q(m\pm in))+1)-\NN(Q(m\pm in))^{i\tau}\exp(F_{2N}(f,K))\right|\ll\varepsilon.
%\]
%We replace all the $m^2-n^2$ for $m^2+n^2$ in the other parts of the proof to conclude in the same manner.
\end{proof}

\iffalse
We also need the following variation of \cite[Lemma 2.8]{FKM1}. 
\begin{lemma}\label{lem: mainpos result}
    Let $\sigma$ be a Borel probability measure on $\mathcal{M}_{p}$ (the set of pretentious functions) with $\sigma(\{1\})>0$ and let $\mathcal{A}=\{(\NN(u)^{i\tau})_{u\in\mathcal O_K^{\times}}\colon \tau\in\R\}$.  There exist $\delta_{0},\rho_{0}>0$ dedpending only on $\sigma$ such that
    \begin{equation}\label{measurepositivity1}
    \liminf_{N\to\infty}\inf_{Q\in\mathcal O_K^{\times}}\text{Re}\left(\E_{\NN(m), \NN(n) \leq N}\int_{\mathcal{A}}w_{\delta}(m,n)f(\ell((Qm+1+\alpha Qn)(Qm+\beta Qn))\overline{f(\ell'(Qm+1)Qn)}\,d\sigma(f)\right)\geq \rho_{0}.
    \end{equation}
%    and
%    \begin{equation}\label{measurepositivity2}
 %   \liminf_{N\to\infty}\inf_{Q\in\mathcal O^{\times}}\text{Re}\left(\E_{\NN(m), \NN(n) \leq N}\int_{\mathcal{A}}\tilde{w}_{\delta}(m,n)f(\ell((Qm+1)^2+(Qn)^2))\overline{f(\ell'(Qm+1)Qn)}\,d\sigma(f)\right)\geq \rho_{0}
 %   \end{equation}
\end{lemma}
\fi
\begin{proof}[Proof of Proposition \ref{fffss}]
%Let us see why inequality \eqref{measurepositivity1} holds.
%since inequality \eqref{measurepositivity2} follows with the exact same methods, simply changing the minus for the plus throughout
%{\color{red}
%Here is a place where the changes with respect to the parametrization depending on $\mathcal O$ need to be effected.
%}


%. In particular, once the two inequalities are established, we can make sure the $\delta_0,\rho_0>0$ work for both without loss of generality.

Let $a:=\sigma(\{1\})>0$ and for $\delta>0$ set
    $$\mu_{\delta}:=\lim_{N\to\infty}\E_{\NN(m), \NN(n) \leq N}w_{\delta}(m,n).$$
By Lemma~\ref{weights}, we have that the limit exists and $\mu_{\delta}>0$. For $T>0$, denote
$$\mathcal{A}_{T}:=\{(\NN(u)^{i\tau})_{u\in\mathcal O_K^{\times}}\colon \tau\in [-T,T]\}\subseteq \mathcal{A},$$
which we claim is a closed Borel set. Indeed, we can write
\[
\mathcal{A}_T= \bigcap_{u \in \mathcal O_K^{\times}} \pi_u^{-1}\left(\{ \NN(u)^{i\tau} : \tau \in [-T,T] \}\right),
\]
where $\pi_u$ is the projection onto the $u$-th coordinate, with $u \in \mathcal O_K^{\times}$. Now, given $\tau \in \R$, the set $\{\NN(u)^{i\tau} : \tau \in [-T,T] \}$ is the image of the compact set $[-T,T]$ under the continuous map $\tau \mapsto\NN(u)^{-i\tau}$, which means it is compact, and since the target space is Hausdorff, it must be closed. The topology taken on $\mathcal{M}$ makes each of the $\pi_u$ continuous, and an aribtrary intersection of closed sets is closed. Thus, $\mathcal{A}_T$ is a Borel set.

By the monotone convergence theorem for sets, there exists $T_{0}=T_{0}(\sigma)>0$ such that $\sigma(\mathcal{A}\backslash\mathcal{A}_{T_{0}})\leq a/2$.
Since $\lim_{\NN(u)\to\infty}\sup_{Q\in \mathcal O_K^{\times}}\vert\log\NN(Qu+1)-\log\NN(Qu)\vert=0$,
we have that 
\begin{equation}\label{e78}
    \begin{split}
        &\lim_{N\to\infty}\sup_{f\in \mathcal{A}_{T_{0}},Q\in\mathcal O_K^{\times}}\E_{\NN(m), \NN(n) \leq N}w_{\delta}(m,n)\cdot\Bigl\vert f\Bigl(\frac{\ell(Qm+1+\alpha Qn)(Qm+1+\beta Qn)}{\ell'(Qm+1)Qn}\Bigr) 
        \\&\quad\quad\quad\quad\quad\quad\quad\quad\quad\quad\quad\quad\quad\quad\quad\quad\quad\quad\quad\quad\quad\quad\;\;\;\;-f\Bigl(\frac{\ell(m+\alpha n)(m+\beta n)}{\ell'mn}\Bigr)\Bigr\vert=0.
    \end{split}
\end{equation}
On the other hand, by the definition of $w_{\delta}$ above, the limit in $N$ after the supremum will only have contributions close to the value of $w_{\delta}$ when the value of $f\Bigl(\frac{\ell(m+\alpha n)(m+\beta n)}{\ell'mn}\Bigr)$ is very close to $1$. This will be independent of the choice of multiplicative function $f$ provided it is in the set $\mathcal{A}_{T_0}$, so it follows that
\begin{equation}\label{e79}
    \begin{split}
        \lim_{N\to\infty}\sup_{f\in \mathcal{A}_{T_{0}}} \left| \E_{\NN(m), \NN(n) \leq N}w_{\delta}(m,n)\cdot\Bigl( f\Bigl(\frac{\ell(m+\alpha n)(m+\beta n)}{\ell'mn}\Bigr)-1\Bigr)\right|=O(\delta).
    \end{split}
\end{equation}
 Combining (\ref{e78}) and (\ref{e79}), we have that 
Thus, taking the last limit in $\delta$ yields the equality that was claimed.
    \begin{equation}
    \begin{split}
       \lim_{\delta\to 0^{+}}\lim_{N\to\infty}\sup_{f\in \mathcal{A}_{T_{0}}} \left| \E_{\NN(m), \NN(n) \leq N}w_{\delta}(m,n)\cdot\Bigl( f\Bigl(\frac{\ell(Qm+1+\alpha Qn)(Qm+1+\beta Qn)}{\ell'(Qm+1)Qn}\Bigr)-1\Bigr)\right|=0.
    \end{split}
\end{equation}
So if $\delta_{0}$ is sufficiently small depending only on $T_{0}$, we have that
$$\liminf_{N\to\infty}\inf_{Q\in\mathcal O_K^{\times}}\text{Re}\left(\E_{\NN(m), \NN(n) \leq N}\int_{\mathcal{A}_{T_{0}}}w_{\delta}(m,n)f\Bigl(\frac{\ell(Qm+1+\alpha Qn)(Qm+1+\beta Qn)}{\ell'(Qm+1)Qn}\Bigr)\,d\sigma(f)\right)$$
is at least $\sigma(\mathcal{A}_{T_{0}})\mu_{\delta_{0}}\geq \sigma(\{1\})\mu_{\delta_{0}}=a\mu_{\delta_{0}}$.
On the other hand, 
$$\liminf_{N\to\infty}\sup_{Q\in\mathcal O_K^{\times}}\left|\E_{\NN(m), \NN(n) \leq N}\int_{\mathcal{A}\backslash\mathcal{A}_{T_{0}}}w_{\delta}(m,n)f\Bigl(\frac{\ell(Qm+1+\alpha Qn)(Qm+1+\beta Qn)}{\ell'(Qm+1)Qn}\Bigr)\,d\sigma(f)\right|$$
is bounded above by $a\mu_{\delta_{0}}/2$.  We are done by setting $\rho_{0}:=a\mu_{\delta_{0}}/2$.
\end{proof}
%We finally have proved all the intermediate results that allow us to show the following theorem, which we restate here:
%{\color{red}
%Change this into the correct equation - note that again, this will depend on the choice for $\mathcal  O$, and we will still get partition regularity.
%}

%\begin{theorem}
%Let $a, b, c \in \N$ be squares. Then, the equation
%\[
%ax^2\pm by^2 = \pm cz^2
%\]
%is partition regular in $x, y$ and $x, z$. More generally, let $a,b,c \in \Z[i]^{\times}$. Then, the equation
%\[
%a^2x^2+b^2y^2=c^2z^2
%\]
%is partition regular with respect to $x, y$ and $x, z$.
%\end{theorem}
%Note that we could also use this method to deal with more general quadratic forms that 
%{\color{red}
%\begin{remark}
%Write some comparison with the results in [FKM]: what they got, what we got, what can make sense over $\Z[i]^{\times}$ versus $\N$, etc.
%\end{remark}
%}
%\iffalse
\subsection{Further discussions}\label{SEC: further}
We end this article with a discussion of potential directions that can be studied next. More particularly, we will discuss the case of real quadratic fields, and also comment on the potential to develop quadratic concentration estimates (as a generalization of the linear concentration estimates that were obtained in Section~\ref{SECconcest}).


\subsubsection{Real quadratic fields}
The main issue with real quadratic fields is the fact that they have infinitely many units. This breaks down many of our arguments irreparably, especially because the balls one could hope to define with the norm (or even its absolute value) have infinitely many elements. This means that the connection with the generalized version of Hal\'asz's theorem, one of the key ingredients of this work, cannot be made.


A potential approach one may consider is to average along a suitable cut of the ball, where we only take a few representatives of each element times powers of the units. Next, it would be key to be able to apply a suitable version of Hal\'asz's theorem in this context, switching to boxes instead of balls, if necessary (and then showing that the two averaging methods can be compared). It might by useful to start by assuming that $f(\e)=1$ for all units $\e$, and then upgrade to an arbitrary $f$.


\iffalse
\begin{proposition}
    Let $K$ be a number field with $b_{1},\dots,b_{d}$ be an integral basis of $\mathcal{O}_{K}$, where $d=[K,\Q]$. Denote 
    $$B_{N}:=\{n_{1}b_{1}+\dots+n_{d}b_{d}\colon n_{1},\dots,n_{d}\in[N]\}.$$
    Let $\mathfrak{J}_{0}(K)$ denote the set of all the principle ideals of $K$. Let $f\colon K\to\C$ be a 1-bounded function vanishing on all the units. Then $f$ naturally lifts to a function $\tilde{f}\colon \mathfrak{J}_{0}(K)\to\C$ with $\tilde{f}((u))=f(u)$ for all $u\in \mathcal{O}_{K}$. If 
    $$\lim_{N\to\infty}\mathbb{E}_{I\in\mathfrak{J}_{0}(K),\NN(I)\leq N}\tilde{f}(I)=c,$$
    then 
    $$\lim_{N\to\infty}\mathbb{E}_{u\in B_{N}}f(u)=c.$$
\end{proposition}

Some remarks: 

(1) Our final goal is to replace the average on $B_{N}$ by averages along AP. But we can deal with the general case later.

(2) similar to what we did for the imaginary case, the average $\lim_{N\to\infty}\mathbb{E}_{I\in\mathfrak{J}_{0}(K),\NN(I)\leq N}\tilde{f}(I)$ is closely related to $\lim_{N\to\infty}\mathbb{E}_{I\in\mathfrak{J}(K),\NN(I)\leq N}\tilde{f}(I)$. If we know the information for the latter, then there is a high chance that we have information for the former.

(3) The case when $f$ does not vanish on units needs to be dealt with a different method.

(4) I think the same argument can be used to show that for Gaussian integers, averages longs boxes are the same as averages along balls. So please check if this argument is correct.

\begin{proof}
Denote $A_{N}:=\mathbb{E}_{I\in\mathfrak{J}_{0}(K),\NN(I)\leq N}\tilde{f}(I)$. Let $w(N)=\vert\{I\in\mathfrak{J}_{0}(K)\colon\NN(I)\leq N\}\vert$ 
     Fix $\e>0$ and suppose that 
    $\vert A_{N}-c\vert<\e,$
    for all $N\geq N_{0}$. Let  $h(I,N):=\vert I \cap B_{N}\vert$. 

\textbf{Claim 1.} there exists a function $h'\colon K\times \mathfrak{J}_{0}(K)\to\N$ such that 
$$\vert \frac{h(I,N)}{h'(\NN(I),N)}-1\vert<\delta(\NN(I),N)$$
for all $\NN(I)\leq C_{K}N^{d}$, where denoting $$\delta_{N}:=\frac{1}{N^{d}}\sum_{I\in\mathfrak{J}_{0}(K)\colon\NN(I)\leq C_{K}N^{d}}h'(\NN(I),N)\delta(\NN(I),N),$$
we have that $\lim_{N\to\infty}\delta_{N}=0$.


{\color{red} TBC.}

{\color{red} As is shown in the previous example, I think the value of $h'(\NN(I),N)$ can be estimated explicitly (in the previous example is was $2\log_{1+\sqrt{2}}(N/\NN(I)^{1/2})$). }

{\color{red} UPDATE:  this claim does not holds when there are finitely many units (i.e. when $h(I,N)$ does not diverge as $N\to\infty$). Let $K=\Q[i]$, then there exist  $z,z'\in\Z[i]$ with $\vert z\vert=\vert z'\vert$ and $z\in B_{N}, z'\notin B_{N}$. Then $h(z,N)-h(z',N)=1$, so $\delta$ can be as large as 1. But this claim might still hold when $h(I,N)$  diverges as $N\to\infty$, at least when $\NN(I)\leq C_{K}^{-1}N^{d}$ for sufficiently large $C_{K}$.}
Denote 
 
     If so,
    then 
    \begin{equation}\label{longeq}
        \begin{split}
            & \mathbb{E}_{u\in B_{N}}f(u)
            =\frac{1}{N^{d}}\sum_{I\in\mathfrak{J}_{0}(K)\colon\NN(I)\leq C_{K}N^{d}}h(I,N)\tilde{f}(I)
            \\&\frac{1}{N^{d}}\sum_{I\in\mathfrak{J}_{0}(K)\colon\NN(I)\leq C_{K}N^{d}}h'(\NN(I),N)\tilde{f}(I)+\delta_{N}
            \\&=\frac{1}{N^{d}}\sum_{k=1}^{C_{K}N^{d}}\sum_{I\in\mathfrak{J}_{0}(K)\colon\NN(I)=k}h'(k,N)\tilde{f}(I)+\delta_{N}
            \\&=\frac{1}{N^{d}}\sum_{k=1}^{C_{K}N^{d}}h'(k,N)(w_{k}A_{k}-w_{k-1}A_{k-1})+\delta_{N}
            \\&=\frac{1}{N^{d}}\Bigl(h'(C_{K}N^{d},N)w_{C_{K}N^{d}}A_{C_{K}N^{d}}+\sum_{k=1}^{C_{K}N^{d}-1}(h'(k+1,N)-h'(k,N))w_{k}A_{k}\Bigr)+\delta'_{N}
        \end{split}
    \end{equation}
where
$$\vert\delta'_{N}\vert\leq\frac{1}{N^{d}}\sum_{I\in\mathfrak{J}_{0}(K)\colon\NN(I)\leq C_{K}N^{d}}\delta(\NN(I),N)\vert\tilde{f}(I)\vert\leq \vert\delta_{N}\vert.$$
 Similarly, by applying (\ref{longeq}) with $f$ replaced by $c$, we have that 
   \begin{equation*}
        \begin{split}
            & c=\mathbb{E}_{u\in B_{N}}c
            \\&=\frac{1}{N^{d}}\Bigl(h'(C_{K}N^{d},N)w_{C_{K}N^{d}}c+\sum_{k=1}^{C_{K}N^{d}-1}(h'(k+1,N)-h'(k,N))w_{k}c\Bigr)+\delta''_{N}
        \end{split}
    \end{equation*}
    where
$\vert\delta''_{N}\vert\leq \vert\delta_{N}\vert.$
So
   \begin{equation*}
        \begin{split}
            & \vert\mathbb{E}_{u\in B_{N}}f(u)-c-\delta'_{N}+\delta''_{N}\vert
            \\&\leq\frac{1}{N^{d}}\Bigl(h'(C_{K}N^{d},N)w_{C_{K}N^{d}}\vert A_{C_{K}N^{d}}-c\vert+\sum_{k=1}^{C_{K}N^{d}-1}(h'(k+1,N)-h'(k,N))w_{k}\vert A_{k}-c\vert\Bigr)
            \\&\leq O_{N_{0}}(\frac{h'(N_{0}+1,N)}{N^{d}})+\frac{\e}{N^{d}}\Bigl(h'(C_{K}N^{d},N)w_{C_{K}N^{d}}+\sum_{k=1}^{C_{K}N^{d}-1}(h'(k+1,N)-h'(k,N))w_{k}\Bigr)
            \\&=O_{N_{0}}(\frac{h'(N_{0}+1,N)}{N^{d}})+\e(1+\delta'''_{N}),
        \end{split}
    \end{equation*}
    where $\vert\delta'''_{N}\vert\leq \vert\delta_{N}\vert.$
    and we used (\ref{longeq}) for $f=\e$ in the last step. 
    Since the first term  converges to 0 as $N\to\infty$, we are done.
\end{proof}



An idea for the case where there are inifnitely many units: assume that $\mathcal{O}=\Z[\sqrt{2}]$. We need to estimate the average of multiplicative functions $f$ over a box $B_{N}=\{m+\sqrt{2}n\colon m,n\in [N]\}$ intersected with an AP. To be gin with we ignore the AP and simply consider the average of  $f$ over $B_{N}$. {\color{red} Assume for the moment that $f(\e)=1$ for all units $\e$}. Then 
$$\mathbb{E}_{u\in B_{N}}f(u)=\frac{\sum_{\NN(I)\leq N}w_{N}(I)f(I)}{\sum_{\NN(I)\leq N}w_{N}(I)},$$
where the RHS are averages along ideals, and 
$w_{N}(I)$ is the set of $u\in B_{N}$ with $u\in I$. I claim that $\vert w_{N}(I)\vert\approx 2\log_{1+\sqrt{2}}(N/\NN(I)^{1/2})$. 



Proof of the claim: all the units in $\Z[\sqrt{2}]$ are of the form $\pm\epsilon_{0}^{k}$, where $\epsilon_{0}=1+\sqrt{2}$. Fix any point $m+n\sqrt{2}\in I$, then $\vert w_{N}(I)\vert$ is the number of $k$ such that $(m+n\sqrt{2})\epsilon_{0}^{k}\in B_{N}$. assume that $(m+n\sqrt{2})\epsilon_{0}^{k}=a+b\sqrt{2}$. Then 
$$(a,b)=(m,n)B^{k},$$
where $B$ is a $2\times 2$ matrix whose eigenvalues are $1\pm \sqrt{2}$. For each $I$, we may choose $m,n$ so that $m,n$ are both approximately $\NN(I)^{1/2}$ (check). Then $\vert w_{N}(I)\vert$ is approximately $2\log_{1+\sqrt{2}}(N/\NN(I)^{1/2})$.
%\subsection{Older stuff that might be helpful for this}

The next step is to show the following:
let $K$ be a number field and $\mathcal{O}_{K}$ be the ring of integers of $K$. Let $a,b\in\Z$.
If $\sqrt{a},\sqrt{b}\in \mathcal{O}_{K}$, then $ax^{2}+by^{2}=z^{2}$ is partition regular with respect to $x,y$.

In [Sun2], the condition was $\sqrt{a},\sqrt{b},\sqrt{a+b}\in \mathcal{O}_{K}$. Here one can think of $K$ as $\Q[\sqrt{d}]$ for some squarefree $d$.


Note that we have a parametrization 
$$x=k(\sqrt{b}m+n)(\sqrt{b}m-n), y=k2\sqrt{a}mn, z=k\sqrt{ab}(m^{2}+n^{2}).$$
So the question is reduced to showing that 
\begin{equation}\nonumber
    \lim_{N \to \infty} \E_{n,m \in R'_{N}}\int_{\mathcal{M}} f((\sqrt{b}m+n)(\sqrt{b}m-n)) \cdot \overline{f(2\sqrt{a}mn)} \ d\sigma(f)>0.
\end{equation}

{\color{blue}
So we need to eastmate
\begin{equation}\nonumber
    \liminf_{N\to\infty}\frac{1}{N^4}\sum_{m,n\in [-N,N]^{2}} f_{0}(m)f_{1}(n)f_{2}(m+\ell_{1}n)f_{3}(m+\ell_{2}n).
\end{equation}
it suffices to estimate (for some $C$ depending only on $K$ and $\ell_{1},\ell_{2}$)
$$\liminf_{N\to\infty}\frac{1}{N^4}\sum_{m,n\in [-N/C,N/C]^{2}} f_{0}(m)f_{1}(n)f_{2}(m+\ell_{1}n)f_{3}(m+\ell_{2}n)$$
as they only differ by a constant. 

We say that $z=a+\tau_{d}b\in \mathcal{O}_{K}$ is a \emph{primary element} if $\vert \epsilon z\vert\geq \vert z\vert$ for any unit $\epsilon$ of $K$, here $\vert\cdot\vert:\max\{a,b\}$. Let $R'_{N}$ be the set of all primary element with $\NN_{K}(u)\leq N^{2}$. 
It follows from ??? that $R'_{N}$ is a finite set.

The previous average can be rewritten as 
\begin{equation}\nonumber
\liminf_{N\to\infty}\frac{1}{N^4}\sum_{m,n\in [-N,N]^{2}} (\mathds{1}_{[-N/C,N/C]^{2}}f_{0})(m)(\mathds{1}_{[-N/C,N/C]^{2}}f_{1})(n)(\mathds{1}_{R'_{N}}f_{2})(m+\ell_{1}n)f_{3}(m+\ell_{2}n),
\end{equation}
which by Cauchy-Schwartz is bounded by $\Vert \mathds{1}_{R'_{N}}f_{2}\Vert_{U^{3}([N]^{2})}$.

We now wish to classify all of multiplicative functions $f_{2}$ such that $\lim_{N\to\infty}\Vert \mathds{1}_{R'_{N}}f_{2}\Vert_{U^{2}([N]^{2})}=0$ (and then show that for these functions $\lim_{N\to\infty}\Vert \mathds{1}_{R'_{N}}f_{2}\Vert_{U^{3}([N]^{2})}=0$, this part should follow from [Sun2] but needs to be checked). To study $\Vert \mathds{1}_{R'_{N}}f_{2}\Vert_{U^{2}([N]^{2})}$, it reduces to study 
$$\E_{u\in [N]^{2}}\mathds{1}_{R'_{N}}f_{2}(u)\chi(u)$$
for a multiplicative function $\chi$, which is reduced to study 
$$\E_{u\in [N]^{2}}\mathds{1}_{R'_{N}}g(u)$$
for all multiplicative function $g$. However, we claim that $\vert R'_{N}\vert$ is comparible to $N^{2}$, and so the average becomes 
$$\E_{u\in R'_{N}}g(u).$$

In order to handle the case when $K$ has infinitely many units, we say that $z\in \mathcal{O}_{K}$ is a \emph{primary element} if $\vert \epsilon z\vert\geq \vert z\vert$ for any unit $\epsilon$ of $K$, here $\vert\cdot\vert$ is the modulus of complex numbers, not $\NN_{K}(\cdot)$ {\color{red} maybe the correct definition is that $\vert a+\tau_{d}b\vert:=a^{2}+b^{2}$ or $\max\{a,b\}$.}, the norm of the field $K$. Let $R'_{N}$ be the set of all primary element with $\NN_{K}(u)\leq N$. This will be the set to work with


{\color{red}

Suppose that $\NN_{K}(a+\tau_{d}b)\leq N$. If $a+\tau_{d}b$ is primary, then we should have that $\vert a\vert, \vert b\vert\leq O_{d}(1)N$. For this, we have that if $\NN_{K}(a+\tau_{d}b), \NN_{K}(a'+\tau_{d}b')\leq N$ and if $a+\tau_{d}b$ and $a'+\tau_{d}b'$ are primary, then their convex combination has $\NN_{K}$ norm $\leq O_{d}(1)N$.

}

We can compare this average with 
\begin{equation}\nonumber
    \lim_{N \to \infty} \E_{n,m \in R_{N}}\int_{\mathcal{M}} f((\sqrt{b}m+n)(\sqrt{b}m-n)) \cdot \overline{f(2\sqrt{a}mn)} \ d\sigma(f)>0,
\end{equation}
where $R_{N}=\{a+\tau_{d}b\colon a,b\in [0,N]\}$.}
\fi
\subsubsection{Quadratic concentration estimates for rings of integers}
We conjecture that it is possible to develop quadratic concentration estimates in the spirit of \cite{FKM2}, but due to space considerations we have not attempted to develop them in the present paper. As a first step, one could try to prove them for $\Z[i]$ given that it is the closest to $\Z$, having a norm that is easy to work with, and still preserving the UFD property. Then, one can see if they can be extended to arbitrary quadratic imaginary fields, and from there try to move to general number fields (of course, the issue we raised in the previous subsubsection above of the existence of infinitely many units still needs to be circumvented).


\iffalse
%\section{Appendix: Seminorm estimates}
{\color{red} It might be a good idea to make this appendix into a useful seminorm comparison short section.}
The purpose of this Appendix is to provide the details for the following proposition:
\begin{proposition}
Let $s \in \N$ and $L_j(m,n)$, $j=1,\dots, s$ be linear forms with coefficients on $\Z[i]$ and suppose that $s=1$ or $s>1$ and the linear forms $L_1, L_j$ are linearly independent for $j=2,\dots, s$. For $j=1,\dots, s$, let $h_j: \Z[i] \to \C$. Suppose that $||h_1||_{U^{s'}[N]} \to 0$ as $N \to \infty$, where $s'=\max\{s-1,2\}$. Then,
\begin{equation}\label{aperiodicconvex0}
    \lim_{N \to \infty} \E_{\NN(m),\NN(n) \leq  N} \prod_{j=1}^s h_j(L_j(m,n))  = 0.
\end{equation}  



\end{proposition}
\begin{proof}
We can define Gowers seminorms for multiplicative functions using the averaging scheme afforded by the balls $B_N:=\{ n  \in \Z[i]^{\times} : \NN(n) \leq N\}$. Therefore, as a consequence of the Gowers-Cauchy-Schwartz inequality, we see that the left hand side of \eqref{aperiodicconvex0} is bounded above by $||h_1||_{U^{s'}(B_N)}$. Next, simply observe that by the non-negativity of the seminorms, we can bound these above by an appropriate seminorm of the form $||h_1||_{U^{s'}(\Z^2_{[\sqrt{N}]})}$ (in other words, the ball seminorms can be bounded above by a constant multiple of convenient box seminorms) 

Finally, by assumption, it follows that the last sequence goes to $0$ as $N \to \infty$, and we are done.
\iffalse
We can, and will assume, without loss of generality, that $|h_j| \leq 1$ for $j=1,\dots, s$. We can also further assume that the linear forms $L_j$ are pairwise linearly independent (over $\Q(i)$), for otherwise we can combine the two relevant $h_i, h_j$ (it must be that both $i,j>1$ in this case), so repeating this process as many times as needed, we end with pairwise independent $L_j$.

Next, we want to use Fourier analysis tools, so we need to reduce matters to cyclic groups. First, for $j=1,\dots, s$, we write $L_j(m,n)=\kappa_jm+\lambda_jn$ for $\kappa_j,\lambda_j \in \Z[i]$ not both zero. Note that \eqref{aperiodicconvex0} is implied by
\begin{equation}
\lim_{N \to \infty} \E_{n, m \in [\pm \sqrt{N}]^2} \mathds{1}_{K_N'}(n,m) \prod_{j=1}^s h_j(L_j(m,n)) = 0,
\end{equation}
where $K_N'$ is still a sequence of convex sets of $[\pm \sqrt{N}]^2$. Next, given $N \in \N$, take $\tilde{N}$ to be the smalles prime such that $\tilde{N}>2[\sqrt{N}]$, $\tilde{N}>2\NN(L_j(m,n))$ for $m, n \in [\pm \sqrt{N}]^2$ and $j=1,\dots, s$, and $\tilde{N}>\NN(\kappa_i\lambda_j-\lambda_i\kappa_j),$ for $i,j=1\dots, s$. Then, it follows that $\frac{\tilde{N}}{N}$ is bounded by a constant depending on the linear forms $L_1,\dots,L_s$ only.

We next put $\tilde{h}_j: \Z_{\tilde{N}}^2 \to \C$ so that it agrees with $h_j$ whenever $n,m \in [\pm [\tilde{N}/2]]^2$ and extend it by periodicity conveniently {\color{red} Check here exactly how to adapt it following [FH]; this definition may need some slight tweaking so that later arguments work out.}. Then, since given $(m,n) \in K_N'$ we have that $\NN(L_j(m,n))<\frac{\tilde{N}}{2}$, it is $h_j(L_j(m,n))=\tilde{h}_j(L_j(m,n))$, whence
\begin{equation}\label{movetoZN}
    \E_{m,n \in [\pm \sqrt{N}]^2} \mathds{1}_{K_N'} \prod_{j=1}^s h_j(L_j(m,n)) = \left(\frac{\tilde{N}}{N}\right)^2 \E_{n,m \in Z_{\tilde{N}}^2} \mathds{1}_{K_N'} \prod_{j=1}^s \tilde{h}_j(L_j(m,n)).
\end{equation}
From this point on, we will work with the right hand side of \eqref{movetoZN}, and assume that the forms $L_j$ and the functions $\tilde{h}_j$ are defined on $\Z_{\tilde{N}}^2$ under the natural identification with the quotient of $\Z[i]$. We will first show that the right hand side of \eqref{movetoZN} converges to $0$ as $N \to \infty$ under the assumption that
\begin{equation}\label{limitgowers0}
    \lim_{N \to \infty} ||\tilde{h}_j||_{U^{s'}(\Z_{\tilde{N}})} = 0
\end{equation}
and then verify that this assumption is satisfied.

We proceed by removing the cutoff function $\mathds{1}_{K_N'}(m,n)$. {\color{red} This will be very sketchy: need to check the original paper by Green and Tao where they lay out the arguments with more details}. Namely, we embed $(\Z_{\tilde{N}}^2)^2$ in the torus $\mathbb{T}^4$ and represent $K_{N}$ as the intersection of $(\Z_{\tilde{N}}^2)^2$ with a ``convex'' subset of $\mathbb{T}^4$. This convex set is then approximated by a sufficiently regular function on $\mathbb{T}^4$ and in turn, this is approximated by a polynomial with bounded coefficients and spectrum of bounded cardinality (with respect to $N$). We then only have to show 
\[
\lim_{N \to \infty} \max_{\eta, \xi \in Z_{\tilde{N}}^2} \left| \E_{m, n \in \Z_{\tilde{N}}^2} e\left( \frac{m\cdot\eta}{\tilde{N}^2}+\frac{\eta\cdot\xi}{\tilde{N}^2}\right) \prod_{j=1}^s \tilde{h}_h(L_j(m,n)) \right| = 0.
\]
\fi
\end{proof}

I am putting here something from the introduction dealing with seminorm properties that Wenbo wrote:

{\color{red}

A sample case:
Let $f_{i}\colon\mathbb{Z}^{2}\to\mathbb{C}$ be 1-bounded functions. Let $B_{N}\subseteq \mathbb{Z}^{2}$ be the ball of radius $N$.
Show that 
$$\frac{1}{N^4}\sum_{{\color{blue} m,n\in B_{N}}} f_{0}(m)f_{1}(m+n)f_{2}(m+2n)f_{3}(m+3n)$$
is bounded by $\Vert \mathds{1}_{B_{N}} f_{0}\Vert_{U^{3}([N]^{2})}$, or equivalently by $\Vert f_{0}\Vert_{U^{3}(B_{N})}$ (the localized Gowers norm).

One potential approach is that we change the blue part to the sum over all $m,n$ such that $m,m+n,m+2n,m+3n\in B_{N}$. If we can show that this set has positive density in $[N]^{2}$, then we can replace $f_{i}$ by $\mathds{1}_{B_{N}}f_{i}$.
}
%{\color{blue} Idea: Use the open mapping theorem with the identity on the $\ell^{\infty}$ spaces of functions with the two Gowers seminorms, one for balls and the usual one. Since the Cauchy-Schwarz inequality implies continuity, since they are both F-spaces the inverse must also be bounded, and the two seminorms are equivalent [check if this reasoning holds.]}
{\color{red}
the Cauchy-Schwarz inequality gives that $\lim_{N\to\infty}\Vert 1_{B_{N}}f\Vert_{U^{s}([N]^{2})}\ll \lim_{N\to\infty}\Vert f\Vert_{U^{s}([N]^{2})}$. How to show the converse direction?}
\medskip

%{\color{blue} The idea is that we consider the map $\text{id}:\ell^{\infty}(\Z[i],|||1_{B_N}f|||_{U^s([N]^2)}) \to \ell^{\infty}(\Z[i],|||f|||_{U^s([N]^2)})$. Then, by the previous observation, it is a bijective continuous map between two F-spaces. As such, I think (check if I'm not making a mistake here) that such a map must have a continuous inverse by the OMT, and so we would get some other similar inequality to the one in red above, with some constant that we wouldn't care about.
%}



{\color{red} Perhaps another way to overcome the difficulty is as follows:
In order to eastmate
\begin{equation}\label{bs}
    \liminf_{N\to\infty}\frac{1}{N^4}\sum_{m,n\in [-N,N]^{2}} f_{0}(m)f_{1}(m+n)f_{2}(m+2n)f_{3}(m+3n),
\end{equation}
it suffices to estimate 
$$\liminf_{N\to\infty}\frac{1}{N^4}\sum_{m,n\in [-N/10,N/10]^{2}} f_{0}(m)f_{1}(m+n)f_{2}(m+2n)f_{3}(m+3n)$$
as they only differ by a constant. But the later expression is the same as 
$$\liminf_{N\to\infty}\frac{1}{N^4}\sum_{m,n\in [-N,N]^{2}}\mathds{1}_{m,n\in [-N/10,N/10]^{2}} (\mathds{1}_{B_{N}}f_{0})(m)(\mathds{1}_{B_{N}}f_{1})(m+n)(\mathds{1}_{B_{N}}f_{2})(m+2n)(\mathds{1}_{B_{N}}f_{3})(m+3n).$$
The weighted average can be bounded by (check) $\liminf_{N\to\infty}\Vert \mathds{1}_{B_{N}}f_{0}\Vert_{U^{3}([N]^{2})}$, which is what we want. So it suffices to study (\ref{bs}) for non-aperiodic functions.

\

Similarly, if we want to estimate
\begin{equation}\label{bss}
    \liminf_{N\to\infty}\frac{1}{N^4}\sum_{m,n\in B_{N}} f_{0}(m)f_{1}(m+n)f_{2}(m+2n)f_{3}(m+3n),
\end{equation}
which is the same as 
$$\liminf_{N\to\infty}\frac{1}{N^4}\sum_{m,n\in [-N,N]^{2}}\mathds{1}_{B_{N}}(m)\mathds{1}_{B_{N}}(n) f_{0}(m)f_{1}(m+n)f_{2}(m+2n)f_{3}(m+3n),$$
or
$$\liminf_{N\to\infty}\frac{1}{N^4}\sum_{m,n\in [-N,N]^{2}}\mathds{1}_{B_{N}}(m)\mathds{1}_{B_{N}}(n) (\mathds{1}_{B_{10N}}f_{0})(m)(\mathds{1}_{B_{10N}}f_{1})(m+n)(\mathds{1}_{B_{10N}}f_{2})(m+2n)(\mathds{1}_{B_{10N}}f_{3})(m+3n),$$
it suffices to estimate
$$\liminf_{N\to\infty}\frac{1}{N^4}\sum_{m,n\in [-10N,10N]^{2}}\mathds{1}_{m,n\in B_{N}} (\mathds{1}_{B_{10N}}f_{0})(m)(\mathds{1}_{B_{10N}}f_{1})(m+n)(\mathds{1}_{B_{10N}}f_{2})(m+2n)(\mathds{1}_{B_{10N}}f_{3})(m+3n)$$
as they only differ by a constant. But the later weighted average can be bounded by (check) $\liminf_{N\to\infty}\Vert \mathds{1}_{B_{10N}}f_{0}\Vert_{U^{3}([10N]^{2})}$ which is what we want. So it suffices to study (\ref{bss}) for non-aperiodic functions.

So the conclusion is: we have many fexibilty the choose the domain of $m,n$, and they can always bounded by $\liminf_{N\to\infty}\Vert \mathds{1}_{B_{N}}f_{0}\Vert_{U^{3}([N]^{2})}$, which means the desired average is 0 for aperiodic functions. So it suffices to study the average for non-aperiodic functions, and we have the freedom to choose the most convenient domain of $m,n$.
}
%\fi
\fi
\appendix
\section{Seminorm estimates}\label{appapp}
The purpose of this appendix is to deal with the apparent disconnect between Halász's theorem (Theorem~\ref{C22intro}), which deals with averages along balls, and a result that is incredibly important to us: \cite[Theorem~1.12]{S23}, which works for averages along boxes, instead of balls.
%
%A key realization for us to be able to freely pass from balls to boxes as needed is the following: the average
%{\color{red}
%
%$$\lim_{N \to \infty} \E_{m,n\colon \NN(L_{j}(m,n)) \leq  N, 1\leq j\leq s} \prod_{j=1}^s h_j(L_j(m,n))  = 0.$$
%is a constant multiple of 
%$$\lim_{N \to \infty} \E_{m,n \in  [\pm CN]^{2}} \prod_{j=1}^s (\mathds{1}_{\NN(\cdot)\leq N}h_j)(L_j(m,n))  = 0$$
%for some $C$ depending on $L_{j}$.
%This average can be bounded by $\nnorm{\mathds{1}_{\NN(\cdot)\leq N}h_j}_{U^{s'}[\pm CN]^{2}}$. 
%}
%
%We need to be a bit careful with the fact that there are two potential definitions for aperiodic functions. The first one (call it \emph{type-1 aperodic}) is given by
%$$\sup_{P \text{ is an infinite AP}}\lim_{N\to\infty}\left|\E_{\NN(u)\leq N}\mathds{1}_P(u) f(u)\right|=0.$$
%On the other hand, the second (call it \emph{type-2 aperodic}) is given by
%$$\lim_{N\to\infty}\sup_{P \text{ is an infinite AP}}\left|\E_{u\in P,\NN(u)\leq N}h_{j}(u)\right|=0.$$
%{\color{red}
%$$\lim_{N\to\infty}\sup_{P \text{ is an infinite AP}}\left|\E_{\NN(u)\leq N}\mathds{1}_{P}(u)f(u)\right|=0.$$
%}
%We will later see that they are, in fact, equivalent, but we can now state the main result of this appendix:
To this end, we prove the following.
\begin{proposition}\label{P: appendix}
Let $K=\Q(\sqrt{-d})$ for some squarefree $d\in\N$, $s \in \N$ with $s \geq 3$ and $L_j(m,n)$, $j=1,\dots, s$ be linear forms with coefficients in $\OK$ such that the linear forms $L_1, L_j$ are linearly independent for $j=2,\dots, s$. For $j=1,\dots, s$, let $h_j: \OK \to \C$. If  $h_{1}$ is aperiodic, 
then
\begin{equation}\nonumber%\label{dwqefp}
   \lim_{N \to \infty} \E_{\NN(m),\NN(n) \leq  N} \prod_{j=1}^s h_j(L_j(m,n))  = 0.
\end{equation}
 \end{proposition}
%{\color{red} I will assume that $s\geq 3$. }
\iffalse
Clearly there exists $C>0$ depending only on $K,L_{1},\dots,L_{s}$ such that $\NN(L_{j}(m,n)) \leq  N^{2}$ for all $1\leq j\leq s$ implies that $m,n\in[\pm CN]^{2}$.%{\color{red} for this to be true I think we need $s\geq 2$}.  
Let $\tilde{N}$ be the smallest prime number which is greater than $10CN$.
Let $h_{j,N}\colon \left[\pm\frac{\tilde{N}-1}{2}\right]^{2}\to\C$ be the map given by $h_{j,N}(u):=\mathds{1}_{\NN(u)\leq N^{2}}\cdot h_{j}(u)$.
Then, for every $N \in \N$, the left hand side of (\ref{dwqefp}) is equal to (before taking limits)
    $$W_N\cdot \lim_{N \to \infty}\E_{m,n \in  \left[\pm\frac{\tilde{N}-1}{2}\right]^{2}} \prod_{j=1}^s h_{j,N}(L_j(m,n)),$$
where 
$$W_N:=\frac{\tilde{N}^{4}}{\vert\{(m,n)\colon \NN(L_{j}(m,n)) \leq  N^{2}, 1\leq j\leq s\}\vert}$$
is a finite positive real number depending only on $C$ and $N$. Indeed, it is easy to see that the numerator is between $N$ and $2^4N$ by Bertrand's postulate, whereas the denominator contains a small ball of radius a multiple of $N$, and is contained in one such ball. 
%{\color{red} give a proof}. %{\color{blue} I think one can show that there exists some $\lambda$ small enough such that the denominator is of order $\lambda N^4$, and the same is true by Bertrand's postulate for the numerator, so definitely $W>0$, only need to show that the limit actually exists.

%Suggestion: write all of this for a specific $N$ or $\tilde{N}$, and then take limsups, to avoid the problem of $W$ potentially not existing. It should still be okay for the rest of the argument, because clearly $0<W<\infty$ uniformly in $N$.
%}
It is possible that the limit of $W_N$ does not exist, but we can ensure that there exist real numbers $K_1,K_2$ such that
\[
0<K_1 \leq \liminf_{N \to \infty} W_N \leq \limsup_{N \to \infty } W_N \leq K_2 <\infty,
\]
which is enough for our purposes.

Let $\tilde{h}_{j,N}\colon \Z_{\tilde{N}}^{2}\to\C$ be the map given by $\tilde{h}_{j,N}:=h_{j,N}\circ\tau$, where $\tau\colon \Z_{\tilde{N}}^{2}\to \left[\pm\frac{\tilde{N}-1}{2}\right]^{2}$ is the natural embedding.
Then 
 the left hand side of (\ref{dwqefp}) is equal to $W_N$ times
    $$\E_{m,n \in  \Z_{\tilde{N}}^{2}} \prod_{j=1}^s \tilde{h}_{j,N}(\tilde{L}_j(m,n)),$$
    where $\tilde{L}_{j}\colon \Z_{\tilde{N}}^{4}\to \Z_{\tilde{N}}^{2}$ is the linear map induced by $L_{j}\colon \Z^{4}\to\Z^{2}$.
%Let $\tilde{h}_{j,N}\colon \Z_{\tilde{N}}^{2}\to\C$ be the map given by
%$\tilde{h}_{j,N}(x,y)=h_{j}(x+\tau_{d}y)$
%if $x,y\in [\pm CN]$ {\color{red} here to define it rigoruously, we can identify $\Z_{\tilde{N}}$ with the interval $\{-\frac{\tilde{N}-1}{2},\dots,\frac{\tilde{N}-1}{2}\}$} and $\NN(x+\tau_{d}y)\leq N$, and $\tilde{h}_{j,N}(x,y)=0$ otherwise. 
%Let $\tilde{h}_{j,N}(u):=\mathds{1}_{\NN(u)\leq N}\cdot h_{j}(u)$. Then
%the left hand side of (\ref{dwqefp}) is equal to 
%    $$  W\cdot \lim_{N \to \infty}\E_{m,n \in  \Z_{\tilde{N}}^{2}} \prod_{j=1}^s \tilde{h}_{j,N}(L_j(m,n)),$$
%where 
%$$W=\lim_{N\to\infty}\frac{\tilde{N}^{4}}{\vert\{(m,n)\colon \NN(L_{j}(m,n)) \leq  N, 1\leq j\leq s\}\vert}$$
%is a finite positive real number depending only on $C$ {\color{red} give a proof}. 
%%By a direct generalization of Lemma 10.7 and Theorem 2.1 of [FH], we have that 
%%$$\vert\E_{m,n \in  [\pm CN]^{2}} \prod_{j=1}^s \tilde{h}_{j,N}(L_j(m,n))\vert\leq C' \min_{1\leq j\leq s}{\norm{\tilde{h}_{j}}_{U^{s'}[\pm CN]^{2}}}^{1/3}+\frac{10}{\tilde{N}}$$
This average is known to be bounded by 
$\norm{\tilde{h}_{j,N}}_{U^{s-1}(\Z_{\tilde{N}}^{2})}$
{\color{red} give a reference, see for example
Theorem 3.1 of
T. Tao. A quantitative ergodic theory proof of Szemerédi’s theorem. Electron. J. Combin. 13
(2006), no.1, Research Paper 99, 49 pp.}

So to conclude the proof, it suffices to show that  $\lim_{N\to\infty}\norm{\tilde{h}_{j,N}}_{U^{s-1}(\Z_{\tilde{N}}^{2})}=0$.
\fi
\iffalse

For any function $f\colon \mathcal{O}_{K}\to\C$, let 
$$\Vert f\Vert_{U^{s}[\pm N]^{2}}:=\frac{1}{\Vert\mathds{1}_{[\pm N]^{2}}\Vert_{U^{s}(\Z_{N'}^{2})}}\Vert\mathds{1}_{[\pm N]^{2}}\cdot f\Vert_{U^{s}(\Z_{N'}^{2})}.$$


 

Similar to Lemma A.2 of [FH], this quantity is independent of the choice of $N'$ if $N'>10N$ and so it is well defined.
Similar to Lemma A.4 of  [FH], it suffices to show that 
$\lim_{N\to\infty}\norm{h_{j,N}}_{U^{s-1}([\pm\frac{\tilde{N}-1}{2}]^{2})}=0$.


Suppose on the countary that there exists arbitraily large $N$ such that 
$\norm{h_{j,N}}_{U^{s-1}([\pm\frac{\tilde{N}-1}{2}]^{2})}>\e$. Then by the inverse theorem, there exists an $(s-2)$-step nilsequence $\phi\colon [N]\to\C$ of complexity $O_{\e,s}(1)$ such that 

\fi

%{\color{blue}

%We wish to estimate the averages
%\begin{equation}\label{dwqefpfirst}
 %   \lim_{N \to \infty} \E_{\NN(m),\NN(n) \leq  N^{2}} \prod_{j=1}^s h_j(L_j(m,n))  = 0.
%\end{equation}

%Assume without loss of generality that $h_{1}$ is aperiodic.
The first step to proving Proposition \ref{P: appendix} is to bound the target average by the Gower's norm of a modified version of $h_{1}$. 
Clearly there exists $C_{1},C_{2}>0$ depending only on $K,L_{1},\dots,L_{s}$ such that
$$\NN(m),\NN(n)\leq N\Rightarrow \NN(L_{j}(m,n)) \leq  C_{1}N\Rightarrow m,n\in[\pm C_{2}\sqrt{N}]^{2},$$
where with a slight abuse of notations, we think of numbers $m,n\in\OK$ as a vector in $\Z^{2}$ (with $1,\tau_{d}$ being the basis).
Let $\tilde{N}$ be the smallest prime number which is greater than $10C_{2}\sqrt{N}$.
Let $h_{j,N}\colon \left[\pm\frac{\tilde{N}-1}{2}\right]^{2}\to\C$ be the map given by $h_{j,N}(u):=\mathds{1}_{\NN(u)\leq C_{1}N}\cdot h_{j}(u)$.
Then, for every $N \in \N$, %the left hand side of \eqref{dwqefp} is equal to (before taking limits)
\begin{equation}\label{fwewgb}
    \E_{\NN(m),\NN(n) \leq  N} \prod_{j=1}^s h_j(L_j(m,n))=W_N\cdot \E_{m,n \in  \left[\pm\frac{\tilde{N}-1}{2}\right]^{2}} \mathds{1}_{\NN(m),\NN(n)\leq N}\prod_{j=1}^s h_{j,N}(L_j(m,n)),
\end{equation}
where 
$$W_N:=\frac{\tilde{N}^{4}}{\vert\{(m,n)\in(\OK)^{2}\colon \NN(m),\NN(n) \leq  N\}\vert}$$
is a finite positive real number. %depending only on $C_{2}$ and $N$. Indeed, it is easy to see that the numerator is between $N$ and $2^4N$ by Bertrand's postulate, whereas the denominator contains a small ball of radius a multiple of $N$, and is contained in one such ball. 
%{\color{red} give a proof}. %{\color{blue} I think one can show that there exists some $\lambda$ small enough such that the denominator is of order $\lambda N^4$, and the same is true by Bertrand's postulate for the numerator, so definitely $W>0$, only need to show that the limit actually exists.
%
%Suggestion: write all of this for a specific $N$ or $\tilde{N}$, and then take limsups, to avoid the problem of $W$ potentially not existing. It should still be okay for the rest of the argument, because clearly $0<W<\infty$ uniformly in $N$.
%}
%It is possible that the limit of $W_N$ does not exist, but 
%
By Bertrand's postulate and Corollary \ref{lem: bounds for number of ideals of norm at most x},
  there exist  $0<K_1<K_2$ depending only on $K,L_{1},\dots,L_{s}$ such that
 \begin{equation}\label{k1k2}
    0<K_1 \leq \liminf_{N \to \infty} W_N \leq \limsup_{N \to \infty } W_N \leq K_2 <\infty.
    \end{equation}

%which is enough for our purposes.


Let $\tilde{h}_{j,N}\colon \Z_{\tilde{N}}^{2}\to\C$ be the map given by $\tilde{h}_{j,N}:=h_{j,N}\circ\tau$, where $\tau\colon \Z_{\tilde{N}}^{2}\to \left[\pm\frac{\tilde{N}-1}{2}\right]^{2}$ is the natural embedding.
Then it follows from
 (\ref{fwewgb}) that
    \begin{equation}\label{eq: averagetobebdd}
    \E_{\NN(m),\NN(n) \leq  N} \prod_{j=1}^s h_j(L_j(m,n))=W_{N}\cdot\E_{m,n \in  \Z_{\tilde{N}}^{2}} \mathds{1}_{\NN(\tau(m)),\NN(\tau(n))\leq N}\prod_{j=1}^s \tilde{h}_{j,N}(\tilde{L}_j(m,n)),
    \end{equation}
    where $\tilde{L}_{j}\colon \Z_{\tilde{N}}^{4}\to \Z_{\tilde{N}}^{2}$ is the linear map induced by $L_{j}\colon \Z^{4}\to\Z^{2}$.

%    Denote $\phi_{\tilde{N}}(m,n):=\mathds{1}_{\NN(\tau(m)),\NN(\tau(n))\leq N^{2}}$ to be a map from $\Z_{\tilde{N}}^{4}$ to $\R$. Then the above average is bounded by {\color{red} check}
    
 By (\ref{k1k2}), it follows from \cite[Proposition 7.1$'$]{GT10b} that the limit of (\ref{eq: averagetobebdd}) as $N\to\infty$ is zero if $\lim_{N\to\infty}\norm{\tilde{h}_{1,N}}_{U^{s-1}(\Z_{\tilde{N}}^{2})}=0$, where $\norm{\cdot}_{U^{s-1}(\Z_{\tilde{N}}^{2})}$ is the $(s-1)$-th Gowers norm on $\Z_{\tilde{N}}^{2}$ (see for example \cite[Definition B.1]{GT10b} for the definition). 
    The reasons said Proposition 7.1' can be applied are as follows. First, the pseudorandomness assumption is unneeded in our case, since our functions are bounded by $1$. Second, without loss of generality, we can assume that our linear forms are in the $s$-normal form, because the operation as described in \cite[Lemma~4.4]{GT10b}.
    
%$$\norm{\hat{\phi}_{\tilde{N}}}_{\ell^{1}(\Z^{4}_{\tilde{N}})}\cdot \max_{\xi\in \Z^{4}_{\tilde{N}}}\left| \E_{m,n \in  \Z_{\tilde{N}}^{2}} \exp\left(\frac{(m,n)\cdot \xi}{\tilde{N}}\right)\prod_{j=1}^s \tilde{h}_{j,N}(\tilde{L}_j(m,n))\right|.$$
%Fix $\xi$, there exists $\ell_{1},\dots,\ell_{s}\in\Z_{\tilde{N}}^{2}$ such that writing $h'_{j,N}(m):=\tilde{h}_{j,N}(m)\cdot\exp(\frac{m\cdot \ell_{j}}{\tilde{N}})$, we have
%\begin{equation}\label{eq: absorbedexp}
%\left| \E_{m,n \in  \Z_{\tilde{N}}^{2}} \exp\left(\frac{(m,n)\cdot \xi}{\tilde{N}}\right)\prod_{j=1}^s \tilde{h}_{j,N}(\tilde{L}_j(m,n))\right|=\max_{\xi\in \Z^{4}_{\tilde{N}}}\left| \E_{m,n \in  \Z_{\tilde{N}}^{2}}  \prod_{j=1}^s h'_{j,N}(\tilde{L}_j(m,n))\right|.
%\end{equation}
%Indeed, this relies on our ability to solve the equation
%\begin{equation}\label{eq: findingellj}
%\sum_{j=1}^{s}\tilde{L}_{j}(m,n)\cdot \ell_{j}=(m,n)\cdot \xi,
%\end{equation}
%which is possible since $s \geq 2$ and the $L_j$ are linearly independent. Thus, we must have terms involving both $n$ and $m,$ so finding the appropriate $\ell_j$ is then simply a matter of solving the system of equations afforded by \eqref{eq: findingellj}, which is always solvable.

%Now, the average in \eqref{eq: absorbedexp}
%is known to be bounded by 
%$
%{\color{red} give a reference, see for example
%Theorem 3.1 of
%T. Tao. A quantitative ergodic theory proof of Szemerédi’s theorem. Electron. J. Combin. 13
%(2006), no.1, Research Paper 99, 49 pp.} Which is the same as $\norm{\tilde{h}_{j,N}}_{U^{s-1}(\Z_{\tilde{N}}^{2})}$.
%On the other hand, one can show that $\Vert\hat{\phi}_{\tilde{N}}\Vert_{\ell^{1}(\Z^{4}_{\tilde{N}})}$ is approximately a constant {\color{red} check (in Lemma 10.7 of [FH], then first replace $\phi$ by a smooth function and then compute its $\ell^{1}$ norm).}
To conclude the proof of Proposition~\ref{P: appendix}, it suffices to show $\lim_{N\to\infty}\norm{\tilde{h}_{1,N}}_{U^{s-1}(\Z_{\tilde{N}}^{2})}=0$.
A similar estimate was obtained in \cite{S23} where $h_{j,N}(u)$ is replaced by a function of the form (say) $\mathds{1}_{u\in[\pm \sqrt{N}]^{2}}\cdot h_{j}(u)$ (see \cite[Theorem~1.12]{S23}). Our goal is to show that an argument similar to the ones used in Sections 7 and 8 of \cite{S23} can be applied to prove Proposition \ref{P: appendix}.

%}

%{\color{red} In the rest of proof, $\mathds{1}_{\NN(n)\leq N^{2}}$ needs to be replaced by  $\mathds{1}_{\NN(n)\leq (C_{1}N))^{2}}$ in the deifnition of $\tilde{h}_{j,N}$.}
In the rest of the proof we refer the reader to \cite{S23} for definitions.
Suppose on the countary that there exists an arbitraily large $N$ such that 
$\norm{\tilde{h}_{1,N}}_{U^{s-1}(\Z_{\tilde{N}}^{2})}>\e$. 
Then by the inverse theorem for Gowers norms (see for example \cite[Theorem~8.5]{S23}), there exist $\delta\gg_{\e,s} 1$ and an $(s-2)$-step $\tilde{N}$-periodic nilsequence $\phi_{N}\colon \Z_{\tilde{N}}\to\C$ of complexity $O_{\e,s}(1)$ such that 
$$\E_{u\in\Z_{\tilde{N}}^{2}}\tilde{h}_{1,N}(u)\phi_{N}(u)\gg \delta.$$
This is equivalent to saying that 
%$$\E_{n\in[\pm\frac{\tilde{N}-1}{2}]^{2}}\mathds{1}_{\NN(u)\leq N}h_{j}(u)\phi_{N}(u)\approx W'\cdot\E_{\NN(u)\leq N^{2}}h_{j}(u)\phi_{N}(u)\gg \delta,$$
%where 
%$$W':=\lim_{N\to\infty}\frac{\{u\colon\NN(u)\leq N^{2}\}}{\tilde{N}^{2}}$$
%is a finite positive real number depending only on $C$ {\color{red} give a proof}
$$\E_{u\in\left[\pm\frac{\tilde{N}-1}{2}\right]^{2}}\mathds{1}_{\NN(u)\leq C_1N}h_{1}(u)\phi_{N}(u)\gg \delta.$$

Before continuing the rest of the proof, we  need to generalize a few properties and results from \cite{S23} that will be key for the rest of our arguments.

First, recall that for $P\subseteq \mathcal{O}_{K}$ and $p,q\in\mathcal{O}_{K}$, we denote by 
$$I_{p,q}(P):=\{u\in \mathcal{O}_{K}\colon pu,qu\in P\}.$$
Moreover, we say that a set $P\subseteq \mathcal{O}_{K}$ is \emph{good} if for every pair of primes $p,q\in\mathcal{O}_{K}$, every (finite length) 2-dimensional arithmetic progression $P'$, and every line $\ell$ in $\mathcal{O}_{K}$, the set $I_{p,q}(P\cap P')\cap \ell$ is a (one-dimensional) arithmetic progression.
 
%We can now state the first result we will need.
We have the following generalization of \cite[Theorem~7.1]{S23}. 
%\WS{To do: we need to unify the format of citations, [1,Theorem 2] or Theorem 2 of [1]} 
\begin{proposition}\label{7.1s2}
    The analog of \cite[Theorem~7.1]{S23} holds if we replace $P$ by any good set (where the average is still taken along boxes).
\end{proposition}
\begin{proof}
The proof goes along the same lines as that of \cite[Theorem~7.1]{S23}. In order to avoid unnecessary repetition, we just wish to highlight the fact that the only key property of the set $P$ that we use is that it is a good set, that is: that it satisfies the intersection property with any line; everything else remains unchanged.    
\end{proof}
%{\color{red} I checked the proof and it seems that this is true. Please double check.}

As a consequence, we have the following generalization of \cite[Theorem 8.1]{S23}.
\begin{proposition}\label{8.1s2}
    The analog of \cite[Theorem 8.1]{S23} holds if we replace $P$ by any good set, and $P'$ by a set which is the intersection of $P$ and a 2-dimensional infinite arithmetic progression (where the average is still taken along boxes).
\end{proposition}
\begin{proof}
    The proof also goes along the same lines as that of \cite[Theorem~8.1]{S23}. In order to avoid needlessly repeating it, we just wish to mention that the only key property of the set $P$ that we use is that it is a good set, that is: that it satisfies the intersection property with any line; everything else remains unchanged.
\end{proof}
%{\color{red} I checked the proof and it seems that this is true (where we use Proposition \ref{7.1s2} to replace Theorem 7.1 of [Sun2]). Please double check.}
In order to actually use these propositions, we need to make sure that the balls we are working with are good sets, which is the content of the following lemma:
\begin{lemma}\label{7.1s23}
For every large enough $N \in \N$ and $C>0$, the set $B_{CN}:=\{u\colon\NN(u)\leq CN\}$ is good.
\end{lemma}
\begin{proof}
Given a pair of primes $p, q \in \OK$, a 2-dimensional arithmetic progression $P$, and a one-dimensional line $\ell \in \mathcal{O}$ we would like to show that the set 
\[
\{ u \in \OK : pu, qu \in B_{CN} \cap P\} \cap \ell
\]
is a one-dimensional arithmetic progression. Indeed, for quadratic imaginary extensions, multiplication always by a prime element always enlarges the norm, so since $\NN(pu)=\NN(p)\NN(u) \geq \NN(u)$, and the same for $qu$, it follows that the first condition of belonging to $B_{CN}$ is nothing but $u \in B_{cN}$ for some other constant $c>0$ which will depend on the norms of $p$ and $q$. For $B_{cN}$ to be non-trivial, we need $N$ to be large enough. 

On the other hand, using the characterization for arithmetic progressions (only for quadratic imaginary fields) coming from the last part of Section~\ref{ss5} we see that, $pu \in P$ if and only if $u \in P'$, some other arithmetic progression $P'$, and the same for $qu$. The intersection of these two arithmetic progressions is again an arithmetic progression, say $P'',$ so the statement becomes
\[
\{u \in \OK : u \in B_{cN} \cap P''\} \cap \ell
\]
is an arithmetic progression, which is now geometrically clear (the key properties are the convexity of the balls, and the fact that $P''$ is a 2D arithmetic progression).
\end{proof}
We now continue the proof. Since 
 $$\E_{u\in\left[\pm\frac{\tilde{N}-1}{2}\right]^{2}}\mathds{1}_{\NN(u)\leq C_1N}h_{1}(u)\phi_{N}(u)\gg \delta$$
 and since the set $B_{C_1N}=\{u\colon\NN(u)\leq C_1N\}$ is good, it follows from Proposition \ref{8.1s2} that 
 $$\left|\E_{u\in[\pm\frac{\tilde{N}-1}{2}]^{2}}\mathds{1}_{R_{N}}(u)h_{1}(u)\right|\gg_{\delta} 1$$
 for some $R_{N}$ which is the intersection of $B_{C_1N}$ and some infinite arithmetic progression $P_{N}$. Then 
  $$\left|\E_{u\in B_{N}}\mathds{1}_{R_{N}}(u)h_{1}(u)\right|\gg_{\delta} 1.$$
In conclusion, we have that 
$$\lim_{N\to\infty}\sup_{P\in\mathcal{AP}[\tau_{d}]}\left|\E_{\NN(u)\leq N}\mathds{1}_{P}(u)h_{1}(u)\right|\gg_{\delta} 1.$$
By Proposition \ref{prop: dircharequivalence}, this implies that $h_{1}$ is not aperiodic, a contradiction. This completes the proof of Proposition \ref{P: appendix}.


\iffalse

Assume that $\phi_{N}(u)=F_{N}(g_{N}(n)\Gamma)$ for some $(s-2)$-step nilmanifold $G/\Gamma$ of  complexity $O_{\e,s}(1)$, some polynomial sequence $g_{N}\colon\Z^{2}\to G$, and some $F\colon G/\Gamma\to\C$ of Lipschitz norm at most  $O_{\e,s}(1)$. Let $\omega\colon\N\to\R_{+}$ to be chosen later.
By \cite[Theorem~4.16]{S23}, there exists $M=O_{\e,s,\omega}(1)$ such that we may write 
$g_{N}=\e_{N}g'_{N}\gamma_{N}$ where $\e_{N}$ is $(M,N)$-smooth, $g'$ is totally $\omega(M)$-totally equidistributed on a subnilmanifold $X'=G'/\Gamma'$, and $\gamma$ is $M$-rational and $M$-periodic.




By the ball version of Lemma 2.20 of [Sun2] {\color{red} the ball version is in fact even simpler}, for any fintie collection $\mathcal{P}$ of prime numbers in $\mathcal{O}_{K}$ {\color{red} since there are only finitely many units, all elements in $\mathcal{O}_{K}$ are $O_{K}(1)$-regular and so  Lemma 2.20 of [Sun2] applies}, we have
$$\vert\E_{\NN(u)\leq N^{2}}h_{j}(u)\phi_{N}(u)\vert^{2}\ll_{C}\frac{1}{A_{\mathcal{P}^{2}}}\frac{C_{\mathcal{P}}(N)}{N^{2}}+\frac{1}{A_{\mathcal{P}}}+\frac{1}{A_{\mathcal{P}^{2}}}+o_{C,\mathcal{P}}(1),$$
where
$$C_{\mathcal{P}}(N):=\sum_{p,q\in\mathcal{P}, p\neq q}\vert\sum_{\NN(qu),\NN(pu)\leq N^{2}}\phi_{N}(pu)\overline{\phi}_{N}(qu)\vert$$
and 
$$\mathcal{A}_{P}:=\sum_{p\in\mathcal{P}}\frac{1}{\NN(p)}.$$

On the other hand, for fixed $p,q$, 
\begin{equation}
    \begin{split}
       &\qquad 
       \lim_{N\to\infty}\E_{\NN(qu),\NN(pu)\leq N^{2}}\phi_{N}(pu)\overline{\phi}_{N}(qu)
       \\&=\lim_{N\to\infty}\E_{\NN(u)\leq N^{2}}\phi_{N}(pu)\overline{\phi}_{N}(qu)
       \\&=\lim_{N\to\infty}\E_{\NN(u)\leq N^{2}}F_{N}(\epsilon_{N}(pu)g_{N}(pu)\gamma_{N}(pu)\Gamma)\overline{F}_{N}(\epsilon_{N}(qu)g_{N}(qu)\gamma_{N}(qu)\Gamma)
   %     \\&=\lim_{N\to\infty}\sum_{r\in [M]^{2}}\E_{\NN(Mu+r)\leq N^{2}}F_{N}(\epsilon_{N}(p(Mu+r)))g_{N}(p(Mu+r))\gamma_{N}(pr)\Gamma)\overline{F}_{N}(\epsilon_{N}(q(Mu+r))g_{N}(q(Mu+r)))\gamma_{N}(qr)\Gamma)
    \end{split}
\end{equation}
 
\fi



%For the relation between the two type of functions, we have
%\begin{itemize}
%    \item Type 2 implies Type 1.
%    \item We have the Gowers norm control (Proposition 2.1) for type 2 functions.
%    \item We have proved the aperiodic-pretentious dichotomy for Type 1 functions.
%    \item Type 1 is equivalent to Type 2 in the 1 dimensional case (this is Lemma A.1 of [Sun2])
%\end{itemize}

%{\color{red} So in order to complete the proof, we need to show either (i) Type 1 is equivalent to Type 2, i.e. Lemma A.1 of [Sun2] holds in the multi-dimensional case; or (ii) the aperiodic-pretentious dichotomy holds for Type 2 functions; or (iii) the Gowers norm control (Proposition 2.1) holds for type 1 functions.}
%Lastly, we give a proof of the equivalence between the two types of aperiodic functions we described above.
 


%\section{Graveyard}




%\subsection{Basic facts on number fields and their rings of integers}
%https://math.stackexchange.com/questions/1222002/finite-group-of-units

%It follows from here that

%https://math.stackexchange.com/questions/2732683/units-in-complex-quadratic-fields-mathbb-q-sqrt-d-with-d-equiv-1-pmod-4


%If $K=\mathbb{Q}[\sqrt{d}]$ for $d<0$, then 

%(1) if $d=-1$, then there are only 4 units (which is the Gaussian case)

%(2) if $d=-3$, then there are only 6 units. This case should be similar to the Gaussian case (and this is also a case of UFD)

%(3) for all other $d<0$, the only units are $\pm 1$ and so $R'_{N}$ is simply $[N]^{2}$ and we are fine in this step.

%(4) if $d>0$, there are infitely many units, however, the Dirichlet Unit Theorem 7.4 on page 42 of [BookN] (check if my computation is correct), and we have that all the unites are of the form $\pm \e^{r}$ for some unit $\e$ of $K$. 

%In the first 3 cases, we are in good shape at this step. We now consider (4). For example, assume that $K=\Q[\sqrt{2}]$. In this case all the units are of the form $\pm (1+\sqrt{2})^{r}$. Let $A=\begin{bmatrix}
%1 & 2\\
%1 & 1
%\end{bmatrix}$. Then if we think of $a+b\sqrt{2}$ as an element $\begin{bmatrix}
%a\\
%b
%\end{bmatrix}\in\Z^{2}$, then the map $a+b\sqrt{2}\mapsto \alpha(a+b\sqrt{2})$ is the map $\begin{bmatrix}
%a\\
%b
%\end{bmatrix}\mapsto A\begin{bmatrix}
%a\\
%b
%\end{bmatrix}$. This matrix has eigenvalues $1\pm \sqrt{2}$ with eigenvectors $\begin{bmatrix}
%\pm\sqrt{2}\\
%1
%\end{bmatrix}$.
%In order to handle the case when $K$ is not an UFD, here is a conjecture:

%\begin{conjecture}
%    For any prime ideal $I$ of $\mathcal{O}_{K}$, there exist $k\in\N$ and prime numbers $p_{1},\dots,p_{\ell},q_{1},\dots,q_{\ell'}\in \mathcal{O}_{K}$ (a prime number is a number that is not the trivial product of two numbers), such that $(q_{1})\dots(q_{\ell'})I^{k}=(p_{1})\dots(p_{\ell})$. 
%\end{conjecture}

%If the conjecture holds, then we can define the value $f(I)$ of a multiplicative function on an ideal $I$ as any $k$-th root of 
%$f(p_{1}\dots p_{\ell}/q_{1}\dots q_{\ell'})$, where $k$ is the smallest integer satisfying the above conjecture. This is a well defined multipicative function and we can apply Halasz to it.

%For example, in $\Z[\sqrt{-5}]$, we have $6=2\cdot 3=(1+\sqrt{-5})(1-\sqrt{-5})$, where $2,3,1\pm\sqrt{-5}$ are prime numbers, but the underlying prime ideals are $I_{1}=(2,1-\sqrt{-5})$, $I_{2}=(3,1-\sqrt{-5})$, $I_{3}=(3,1+\sqrt{-5})$. In fact,
%$$(2)=I_{1}^{2}, (3)=I_{2}I_{3}, (1+\sqrt{-5})=I_{1}I_{3}, (1-\sqrt{-5})=I_{1}I_{2}.$$
%We then define
%$f(I_{1})^{2}:=f(2)$ or $f(I_{1})^{2}:=f(1+\sqrt{-5})f(1-\sqrt{-5})/f(3)$, which are the same since $f$ is multiplicative. Similarly, we may define $f(I_{2})^{2}:=f(3)f(1-\sqrt{-5})/f(1+\sqrt{-5})$....

%{\color{red}



%  Let $a_{1},\dots,a_{k},b_{1},\dots,b_{k'}\in\mathcal{O}_{K}$. If 
%$$\{c_{1}a_{1}+\dots+c_{k}a_{k}\colon c_{i}\in\Z\}=\{d_{1}b_{1}+\dots+d_{k'}b_{k'}\colon d_{i}\in\Z\},$$
%then 
%$$\{\tau_{d}c'_{1}a_{1}+\dots+\tau_{d}c'_{k}a_{k}\colon c'_{i}\in\Z\}=\{\tau_{d}d'_{1}b_{1}+\dots+\tau_{d}d'_{k'}b_{k'}\colon d'_{i}\in\Z\},$$
%and so 
%$$\{C_{1}a_{1}+\dots+C_{k}a_{k}\colon C_{i}=c_{i}+\tau_{d}c'_{i}\in\mathcal{O}_{K}\}=\{D_{1}b_{1}+\dots+D_{k'}b_{k'}\colon D_{i}=d_{i}+\tau_{d}d'_{i}\in\mathcal{O}_{K}\}.$$

%On the other hand, for every ideal $I=\langle a_{1},\dots,a_{k}\rangle$ in $\mathcal{O}_{K}$, by the rank of abelian group, there exists $b_{1},\dots,b_{D}\in\mathcal{O}_{K}$ so that 
%$$\{c_{1}a_{1}+\dots+c_{k}a_{k}\colon c_{i}\in\Z\}=\{d_{1}b_{1}+\dots+d_{D}b_{D}\colon d_{i}\in\Z\}.$$
%This implies that every ideal of $\mathcal{O}_{K}$ is generated by at most $D:=[K,\Q]$ elements.

%{\color{blue} Update: it follows from Exercise 5 of [BookN] on page 23 that every ideal in a Dedekind domain is generated by at most 2 elements.}

%\textbf{Claim.} For every prime number $p$ in $\mathbb{Z}[\sqrt{-5}]$ (meaning that $p\neq rs$ for non-unit $r,s\in \mathbb{Z}[\sqrt{-5}]$), $(p)$ is the product of at most 2 prime ideals.

%To see this, assume $(p)=J_{1}J_{2}J_{3}$. Since  $\mathbb{Z}[\sqrt{-5}]$ has only two ideal classes, we may assume that $J_{1}$ and $J_{2}$ are in the same class. In both cases, we may write $J_{1}=(a_{1})J$ and  $J_{2}=(a_{2})J$ for some $J$ with $J^{2}=(b)$. Then $(p)=(a_{1}a_{2}b)J_{3}$, this should be impossible if $p$ is a prime number.

 
 
%}

%{\color{blue} We consider the following question: suppose that $f\colon K\to\S$ is a completely multiplicative function. Can we extend it to a completely multiplicative function $\tilde{f}$ whose domain is the set of all ideals of $K$?

%An example is to consider the case when $K=\Q[\sqrt{-5}]$. In this case we know that the class number is 2, meaning that there exists a non-principle ideal $I$ such that all the fractional ideals are of the form $(p)/(q)$ or $(p)I/(q)$ for some $p,q\in\Z[\sqrt{-5}]$. Therefore, we see that $\tilde{f}$ is completely determined by its value at $I$. 

%On the other hand, since in this case the class group is simply $\Z/2\Z$, we have that $I^{2}=(x)$ for some $x\in\Z[\sqrt{-5}]$. We define $\tilde{f}(I)$ to be any of the two numbers whose square is $f(x)$. 
%We now define $\tilde{f}$ as follows.

%Define $\tilde{f}((x)):=f(x)$ for principle ideals and $\tilde{f}((x)I):=f(x)\tilde{f}(I)$ for non-principle ideals.

%We claim that this $\tilde{f}$ is a well defined multiplicative function.
%To see the well definedness, assume that 
%$$I_{1}\dots I_{k}J_{1}\dots J_{\ell}=I'_{1}\dots I'_{k'}J'_{1}\dots J'_{\ell'}$$
%for some $I_{i}=(a_{i})$, $I'_{i}=(a'_{i})$, $J_{i}=(b_{i})I$, $J'_{i}=(b'_{i})I$. Then we must have $2\vert \ell-\ell'$ and that $x^{\ell/2}\prod_{i} a_{i}\prod_{j}b_{j}=x^{\ell'}\prod_{i} a_{i}\prod_{j}b_{j}$.
%If $\ell,\ell'$ are even, then
%\begin{equation*}
%    \begin{split}
%        & \tilde{f}(I_{1}\dots I_{k}J_{1}\dots J_{\ell})=\tilde{f}((x^{\ell/2}\prod_{i} a_{i}\prod_{j}b_{j}))
%        =f(x^{\ell/2}\prod_{i} a_{i}\prod_{j}b_{j})
%        \\&=f(x^{\ell'/2}\prod_{i'} a'_{i'}\prod_{j'}b_{j'})=\tilde{f}((x^{\ell'/2}\prod_{i'} a_{i'}\prod_{j'}b_{j'}))
%        =\tilde{f}(I'_{1}\dots I'_{k'}J'_{1}\dots J'_{\ell'}).
%    \end{split}
%\end{equation*}
%If $\ell,\ell'$ are odd, then
%\begin{equation*}
%    \begin{split}
%        & \tilde{f}(I_{1}\dots I_{k}J_{1}\dots J_{\ell})=\tilde{f}((x^{(\ell-1)/2}\prod_{i} a_{i}\prod_{j}b_{j})I)
%        =f(x^{(\ell-1)/2}\prod_{i} a_{i}\prod_{j}b_{j})\tilde{f}(I)
%        \\&=f(x^{(\ell'-1)/2}\prod_{i'} a'_{i'}\prod_{j'}b_{j'})\tilde{f}(I)=\tilde{f}((x^{(\ell'-1)/2}\prod_{i'} a_{i'}\prod_{j'}b_{j'}))
%        =\tilde{f}(I'_{1}\dots I'_{k'}J'_{1}\dots J'_{\ell'}).
%    \end{split}
%\end{equation*}
%So the value of $\tilde{f}$ is independent of the choice of the decomposition. It is also not hard to show that $\tilde{f}$ is completely multiplicative. We are done.  

%\textbf{Update:} I learned that the ideal class group of every number field is a finite abelian group (need to find a reference), therefore I think it is very likely that the above method applies to every number field (i.e. multiplicative functions on number fields can also be extended to a ideal-domained multiplicative function)!
%}
%{\color{red} We still need to define pretentious functions on more general number fields. If the norm takes on negative values, $\tau \mapsto \NN(u)^{i\tau}$ is not well defined. What could be an alternative?}
%}

%\begin{proposition}
%    Let $K$ be a number field and $\mathfrak{Q}(K)$ denote the set of all fractional ideals of $K$. For any completely multiplicative function $f\colon K^{\times}\to\S$ with  $f(\e)=1$ for all units $\e$, there exists a completely multiplicative function $\tilde{f}\colon \mathfrak{Q}(K)\to\S$ such that $\tilde{f}((x))=f(x)$ for all $x\in K^{\times}$. \
%\end{proposition}
%\begin{proof}
%    Let $G$ be the ideal class group of $K$. By ???, $G$ is a finite abelian group. By the structure theorem, there exist $k\in\N$, 
%    and generators $g_{1},\dots,g_{k}$ of $G$ of orders $d_{1},\dots,d_{k}\in\N$.
%    For each $1\leq i\leq k$, let $I_{i}\in \mathfrak{Q}(K)$ be a representative of $g_{i}$. Then it follows that $I_{i}^{d_{i}}=(x_{i})/(y_{i})$ for some $x_{i},y_{i}\in (\mathcal{O}_{K})^{\times}$. Let $a_{i}$ be any element in $\S$ such that $a_{i}^{d_{i}}=f(x_{i})/f(y_{i})$. We then define
%    $$\tilde{f}\left(\frac{(x)}{(y)}I_{1}^{b_{1}}\dots I_{k}^{b_{k}}\right):=f(x/y)a_{1}^{b_{1}}\dots a_{k}^{b_{k}}$$
%for all $x,y\in (\mathcal{O}_K)^{\times}$ and $0\leq b_{i}\leq d_{i}-1$. This defines a function on $\mathfrak{Q}(K)$, since by the structure of the ideal class group, every element of $\mathfrak{Q}(K)$ can be written in the form $II_{1}^{b_{1}}\dots I_{k}^{b_{k}}$ for some $I$ of the form $\frac{(x)}{(y)}$ for some $x,y\in (\mathcal{O}_K)^{\times}$ and for some $0\leq b_{i}\leq d_{i}-1$ in a unique way (the representation $I=\frac{(x)}{(y)}$ is not unique, but the value $\tilde{f}(I):=f(x/y)$ is unique since $f(\e)=1$ for all units $\e$).

%It is clear that $\tilde{f}((x))=f(x)$ for all $x\in K^{\times}$. To show that $\tilde{f}$ is multiplicative, it suffices to show the following claim.

 %\textbf{Claim.} For all $x,y\in (\mathcal{O}_K)^{\times}$ and $b_{i}\in\N$, we have that $$\tilde{f}\left(\frac{(x)}{(y)}I_{1}^{b_{1}}\dots I_{k}^{b_{k}}\right)=\tilde{f}((x))\tilde{f}((y))^{-1}\tilde{f}(I_{1})^{b_{1}}\dots \tilde{f}(I_{k})^{b_{k}}.$$

%Indeed, assume that $b_{i}=c_{i}d_{i}+r_{i}$ for some $c_{i}\in\N$ and $0\leq r_{i}\leq d_{i}-1$. Then 
%$$\frac{(x)}{(y)}I_{1}^{b_{1}}\dots I_{k}^{b_{k}}=\frac{(xx_{1}^{c_{1}}\dots x_{k}^{c_{k}})}{(yy_{1}^{c_{1}}\dots y_{k}^{c_{k}})}I_{1}^{r_{1}}\dots I_{k}^{r_{k}}.$$
%So 
%\begin{equation*}
%    \begin{split}
%        &\qquad\tilde{f}\left(\frac{(x)}{(y)}I_{1}^{b_{1}}\dots I_{k}^{b_{k}}\right)=\tilde{f}\left(\frac{(xx_{1}^{c_{1}}\dots x_{k}^{c_{k}})}{(yy_{1}^{c_{1}}\dots y_{k}^{c_{k}})}I_{1}^{r_{1}}\dots I_{k}^{r_{k}}\right)
%         =f\left(\frac{xx_{1}^{c_{1}}\dots x_{k}^{c_{k}}}{yy_{1}^{c_{1}}\dots y_{k}^{c_{k}}}\right)a_{1}^{r_{1}}\dots a_{k}^{r_{k}}
%         \\&=f(x/y)a_{1}^{d_{1}c_{1}+r_{1}}\dots a_{k}^{d_{k}c_{k}+r_{k}}
%         =\tilde{f}((x))\tilde{f}((y))^{-1}\tilde{f}(I_{1})^{b_{1}}\dots \tilde{f}(I_{k})^{b_{k}}.
%    \end{split}
%\end{equation*}
%We are done.
%\end{proof}


%Some thoughts for the case $K=\Q[\sqrt{-5}]$. In this case, we know that the ideal class number is 2. Every fractional ideal of $K$ is either a principle fractional ideal, or the product of a principle fractional ideal and $I=\langle 2, 1-\sqrt{-5}\rangle$. Let $\mathfrak{I}(K)$ denote the set of ideals in $\mathcal{O}_{K}$. Consder 
%$$\lim_{N\to\infty} \E_{\mathfrak{v}\in \mathfrak{I}(K), \NN_{K}(\mathfrak{v})\leq N}\tilde{f}(\mathfrak{v}).$$

%\begin{lemma}
%    All the ideals are either of the form $(a)$ for some $a\in\Z[\sqrt{-5}]$ or of the form $\frac{(x+y\sqrt{-5})}{(2)}I$ for some $x,y\in\Z$ which are both even or both odd.
%\end{lemma}
%\begin{proof}
%    Assume that $\frac{(p)}{(q)}I$ is an ideal. We may assume WLOG that $q\in\Z$. Then $\NN(q)\vert 2\NN(p)$. If $4\vert \NN(p),\NN(q)$, then $p/2,q/2\in \mathcal{O}_{K}$ as their norms are integers, and $\frac{(p)}{(q)}I=\frac{(p/2)}{(q/2)}I$. So we may assume WLOG that at most one of $\NN(p),\NN(q)$ is divisible by 4. If $\NN(q)$ is even, then since $q\in\Z$, we have that $4\vert\NN(q)$ and so $4\nmid\NN(p)$. Let $q'=q/2\in\Z$, then $\NN(q)\vert 2\NN(p)$ implies that $2\NN(q')\vert \NN(p)$. In particular $\NN(p/q')\in\Z$ and so $p':=p/q'\in\mathcal{O}_{K}$. So $\frac{(p)}{(q)}I=\frac{(p')}{(2)}I$. Since its norm is an integer, we have that $\NN(p')\in 2\Z$, which implies that $p'=x+y\sqrt{-5}$ with $2\vert x+y$.

%    Now assume that $\NN(q)$ is odd. Then $\NN(q)\vert 2\NN(p)$ implies that $\NN(q)\vert \NN(p)$ and so $\NN(p/q)\in\Z$, this means that $p/q\in\mathcal{O}_{K}$ and so $\frac{(p)}{(q)}I=(p/q)I=\frac{(2p/q)}{(2)}I$.
%\end{proof}


%By this lemma, we we decompose the ideals in type 1 ideals $\mathfrak{I}_{1}(K)$ and type 2 ideals $\mathfrak{I}_{2}(K)$.
%Then

%\begin{equation*}
%    \begin{split}
%        &\qquad\lim_{N\to\infty} \E_{\mathfrak{v}\in \mathfrak{I}(K), \NN_{K}(\mathfrak{v})\leq N}\tilde{f}(\mathfrak{v})
%        \\&=\sum_{i=1}^{2}\lim_{N\to\infty} \E_{\mathfrak{v}\in \mathfrak{I}(K), \NN_{K}(\mathfrak{v})\leq N}\mathds{1}_{\mathfrak{v}\in \mathfrak{I}_{i}(K)}\tilde{f}(\mathfrak{v})
%        \\&=\lim_{N\to\infty} \frac{1}{\vert\{\mathfrak{v}\in \mathfrak{I}(K)\colon \NN_{K}(\mathfrak{v})\leq N\}\vert}\sum_{u\in\mathcal{O}_{K},\NN_{K}(u)\leq N}\tilde{f}((u))
%        \\&+\lim_{N\to\infty} \frac{1}{\vert\{\mathfrak{v}\in \mathfrak{I}(K)\colon \NN_{K}(\mathfrak{v})\leq N\}\vert}\sum_{2\vert x+y, \NN_{K}(x+y\sqrt{-5})\leq 2N}\tilde{f}(\frac{(x+y\sqrt{-5})}{(2)}I)
%        \\&=\frac{1}{2}\lim_{N\to\infty} \E_{u\in\mathcal{O}_{K},\NN_{K}(u)\leq N}\tilde{f}((u))
%        +\frac{1}{2}\lim_{N\to\infty} \E_{2\vert x+y, \NN_{K}(x+y\sqrt{-5})\leq 2N}\tilde{f}(\frac{(x+y\sqrt{-5})}{(2)}I) {\color{red} check}.
%    \end{split}
%\end{equation*}


% The first part is an average along balls over integers, the second part is an average along balls under some arithmetic progressions.

%{\color{red} To overcome this difficulty, if we can show that the average $\lim_{N\to\infty} \E_{\mathfrak{v}\in \mathfrak{I}(K), \NN_{K}(\mathfrak{v})\leq N}\tilde{f}(\mathfrak{v})$ is zero for all the possible extensions $\tilde{f}$ of $f$, then this will cancel the second term and gives us that $\lim_{N\to\infty} \E_{u\in\mathcal{O}_{K},\NN_{K}(u)\leq N}\tilde{f}((u))=0$. (Details to be added)

% \textbf{Claim.} The followings are equivalent (for imaginary quadratic fields):

 % (1) The average of $f$ along every AP is 0.
  
 %(2) For any Dirichlet chracter $\chi$, the average of $f\chi$ is 0.

 %(3) For any Dirichlet chracter $\chi$ and any extension $\tilde{f\chi}$ of $f\chi$, the average of $\tilde{f\chi}$ (along ideals) is 0.

 %\

 %We proved before that $(1)\Leftrightarrow (2)$. If (3) holds, then the previous argument shows that $\lim_{N\to\infty} \E_{u\in\mathcal{O}_{K},\NN_{K}(u)\leq N}\tilde{f\chi}((u))=0$ for any $\chi$ and any extension. But this is just the average of $f\chi$, which implies (2).

 %If (2) holds, then the previous equation shows that (3) holds.

 %Therefore, in order to classify all the aperiodic function, it suffices to classify all the functions with $\lim_{N\to\infty} \E_{\mathfrak{v}\in \mathfrak{I}(K), \NN_{K}(\mathfrak{v})\leq N}\tilde{f}(\mathfrak{v})=0$.
 
 %}


\iffalse

\begin{proposition}
    Let $f\colon \mathfrak{Q}(K)\to\S$ be a completely multiplicative function. Let $\mathfrak{I}(K)$ denote the set of ideals in $\mathcal{O}_{K}$. If $K$ has finitely many units, then
    $$\lim_{N\to\infty} \E_{\mathfrak{v}\in \mathfrak{I}(K), \NN_{K}(\mathfrak{v})\leq N}f(\mathfrak{v})=0 \Leftrightarrow \lim_{N\to\infty} \E_{u\in\mathcal{O}_{K}, \NN_{K}(u)\leq N}f((u))=0.$$
\end{proposition}
\begin{proof}

\iffalse

 Let $G$ be the ideal class group of $K$. By ???, $G$ is a finite abelian group. Then 
   $$\lim_{N\to\infty} \E_{\mathfrak{v}\in \mathfrak{I}(K), \NN_{K}(\mathfrak{v})\leq N}f(\mathfrak{v})=\sum_{g\in G}\lim_{N\to\infty} \E_{\mathfrak{v}\in \mathfrak{I}(K), \NN_{K}(\mathfrak{v})\leq N}\mathds{1}_{\mathfrak{v}\in \mathcal{C}_{g}}f(\mathfrak{v}),$$
where $\mathcal{C}_{g}$ is the ideal class of $g$.
By ???, when $N\to\infty$, the number of ideals with norm at most $N$ from any fixed ideal class are approximately the same. So
$$\lim_{N\to\infty} \E_{\mathfrak{v}\in \mathfrak{I}(K), \NN_{K}(\mathfrak{v})\leq N}f(\mathfrak{v})=\E_{g\in G}\lim_{N\to\infty} \E_{\mathfrak{v}\in \mathcal{C}_{g}, \NN_{K}(\mathfrak{v})\leq N}f(\mathfrak{v}).$$

\textbf{Claim.} for every $g\in G$, there exists  $\mathfrak{u}_{g}\in\mathcal{C}_{g}$ such that 
$$\mathcal{C}_{g}=\{(p)\mathfrak{u}_{g}\colon p\in\mathcal{O}_{K}^{\times}\}.$$


\fi

 


 Let $G$ be the ideal class group of $K$. By ???, $G$ is a finite abelian group. By the structure theorem, there exist $k\in\N$, 
    and generators $g_{1},\dots,g_{k}$ of $G$ of orders $d_{1},\dots,d_{k}\in\N$.
    For each $1\leq i\leq k$, let $I_{i}\in \mathfrak{Q}(K)$ be a representative of $g_{i}$. Then it follows that $I_{i}^{d_{i}}=(x_{i})/(y_{i})$ for some $x_{i},y_{i}\in (\mathcal{O}_{K})^{\times}$. 

   For any $g=(v_{1},\dots,v_{k})\in \{0,\dots,d_{1}-1\}\times\dots\times\{0,\dots,d_{k}-1\}$, let $\mathcal{C}_{g}$ be the ideal class of $J_{g}:=I_{1}^{v_{1}}\dots I_{k}^{v_{k}}$. Then 
   $$\lim_{N\to\infty} \E_{\mathfrak{v}\in \mathfrak{I}(K), \NN_{K}(\mathfrak{v})\leq N}f(\mathfrak{v})=\sum_{g\in G}\lim_{N\to\infty} \E_{\mathfrak{v}\in \mathfrak{I}(K), \NN_{K}(\mathfrak{v})\leq N}\mathds{1}_{\mathfrak{v}\in \mathcal{C}_{g}}f(\mathfrak{v}),$$
where $\mathcal{C}_{g}$ is the ideal class of $g$.

   For each $g$, it is clear that there exist $p_{g},q_{g}\in\mathcal{O}_{K}^{\times}$ such that 
   $$\mathcal{C}_{g}=\{\frac{(p_{g}r)}{(q_{g})}I_{g}\colon r\in\mathcal{O}_{K}^{\times}\}$$
   {\color{red} I think this is not correct....}. So
   \begin{equation*}
    \begin{split}
        &\qquad\lim_{N\to\infty} \E_{\mathfrak{v}\in \mathfrak{I}(K), \NN_{K}(\mathfrak{v})\leq N}\mathds{1}_{\mathfrak{v}\in \mathcal{C}_{g}}f(\mathfrak{v})
        \\&=\lim_{N\to\infty}\frac{\vert\{r\in\mathcal{O}_{K}^{\times}\colon \NN_{K}(\frac{(p_{g}r)}{(q_{g})}I_{g})\leq N\}\vert}{\vert\{\mathfrak{v}\in \mathfrak{I}(K)\colon \NN_{K}(\mathfrak{v})\leq N\}\vert}\E_{r\in\mathcal{O}_{K}^{\times}\colon \NN_{K}(r)\leq N/\NN(\frac{(p_{g})}{(q_{g})}I_{g})}f(\frac{(p_{g}r)}{(q_{g})}I_{g})
        \\&=\lim_{N\to\infty}\frac{\vert\{r\in\mathcal{O}_{K}^{\times}\colon \NN_{K}(r)\leq N/\NN(\frac{(p_{g})}{(q_{g})}I_{g})\}\vert}{\vert\{\mathfrak{v}\in \mathfrak{I}(K)\colon \NN_{K}(\mathfrak{v})\leq N\}\vert}\E_{r\in\mathcal{O}_{K}^{\times}\colon \NN_{K}(r)\leq N/\NN(\frac{(p_{g})}{(q_{g})}I_{g})}f(\frac{(p_{g}r)}{(q_{g})}I_{g}).
    \end{split}
\end{equation*}

Since the set of $r$ with $\NN(r)\leq N$ is approximately $CN$ plus an error term for some $C$ depending only on $K$, the above average is 
   \begin{equation*}
    \begin{split}
        &\qquad \lim_{N\to\infty}\frac{N/\NN(\frac{(p_{g})}{(q_{g})}I_{g})}{\sum_{g'\in G}N/\NN(\frac{(p_{g'})}{(q_{g'})}I_{g'})}\E_{r\in\mathcal{O}_{K}^{\times}\colon \NN_{K}(r)\leq N/\NN(\frac{(p_{g})}{(q_{g})}I_{g})}f(\frac{(p_{g}r)}{(q_{g})}I_{g})
        \\&=\lim_{N\to\infty}\frac{1/\NN(\frac{(p_{g})}{(q_{g})}I_{g})}{\sum_{g'\in G}1/\NN(\frac{(p_{g'})}{(q_{g'})}I_{g'})}\E_{r\in\mathcal{O}_{K}^{\times}\colon \NN_{K}(r)\leq N/\NN(\frac{(p_{g})}{(q_{g})}I_{g})}f(\frac{(p_{g}r)}{(q_{g})}I_{g})
         \\&=\frac{1/\NN(\frac{(p_{g})}{(q_{g})}I_{g})}{\sum_{g'\in G}1/\NN(\frac{(p_{g'})}{(q_{g'})}I_{g'})}\lim_{N\to\infty}\E_{r\in\mathcal{O}_{K}^{\times}\colon \NN_{K}(r)\leq N}f(\frac{(p_{g}r)}{(q_{g})}I_{g})
          \\&= \frac{f(\frac{(p_{g})}{(q_{g})}I_{g})/\NN(\frac{(p_{g})}{(q_{g})}I_{g})}{\sum_{g'\in G}1/\NN(\frac{(p_{g'})}{(q_{g'})}I_{g'})}\lim_{N\to\infty}\E_{r\in\mathcal{O}_{K}^{\times}\colon \NN_{K}(r)\leq N}f((r)).
    \end{split}
\end{equation*}
Therefore, 
 $$\lim_{N\to\infty} \E_{\mathfrak{v}\in \mathfrak{I}(K), \NN_{K}(\mathfrak{v})\leq N}f(\mathfrak{v})=M\lim_{N\to\infty} \E_{u\in\mathcal{O}_{K}, \NN_{K}(u)\leq N}f((u)),$$
 where 
 $$M:=\frac{\sum_{g\in G}f(\frac{(p_{g})}{(q_{g})}I_{g})/\NN(\frac{(p_{g})}{(q_{g})}I_{g})}{\sum_{g\in G}1/\NN(\frac{(p_{g})}{(q_{g})}I_{g})}.$$

 So it suffices to show that $M\neq 0$. {\color{red} I need to think a little more on this, perhaps we need to be careful in choosing $p_{g}$ and $q_{g}$. However, for some simple cases I can prove it. Consider the case where $K=\Q[\sqrt{-5}]$, whose ideal class group is $\Z/2\Z$. In this case, the numerator of $M$ is simply $1+f(I)/\NN(I)$, where $I$ is a fractional ideal such that every ideal of $\Q[\sqrt{-5}]$ is either principle or is the product of a principle ideal anf $I$. In this case, I think $I=(2,1-\sqrt{-5})$ (proof given in the next paragraph). Since $\vert f(I)\vert/\vert\NN(I)\vert\leq 1/2$, we have that  $1+f(I)/\NN(I)\neq 0$ and we are done.
 
 To see that $I=(2,1-\sqrt{-5})$, suppose that $\frac{(p)}{(q)}I$ is an ideal. Then its norm is an integer, so $\NN(I)p^{2}/q^{2}=2p^{2}/q^{2}\in\Z$. If $2\vert q$, then we also have that $w\vert p$. So we may assume WLOG $2\nmid q$. Then $2p^{2}/q^{2}\in\Z$ implies that $p^{2}/q^{2}\in\Z$ and that $p/q\in\Z$. So all the ideals in the ideal class of $I$ is the product of a principle ideal and $I$.


 \


 On the other hand, the following weaker property holds (which maybe is also good for our purposes):
 
 \begin{proposition}
     For any completely multiplicative function $f\colon K^{\times}\to\S$ with  $f(\e)=1$ for all units $\e$, if $K$ has finitely many units, there exists a completely multiplicative function $\tilde{f}\colon \mathfrak{Q}(K)\to\S$ such that $\tilde{f}((x))=f(x)$ for all $x\in K^{\times}$, and that   $$\lim_{N\to\infty} \E_{\mathfrak{v}\in \mathfrak{I}(K), \NN_{K}(\mathfrak{v})\leq N}\tilde{f}(\mathfrak{v})=0 \Leftrightarrow \lim_{N\to\infty} \E_{u\in\mathcal{O}_{K}, \NN_{K}(u)\leq N}\tilde{f}((u))=0.$$
 \end{proposition}
   

 To see this, let $z_{i}:=\tilde{f}(I_{(0,\dots,0,1,0,\dots,0)})$, where 1 is the $i$-th coordinate. Then the numerator of $M$ given by $\sum_{g\in G}\tilde{f}(\frac{(p_{g})}{(q_{g})}I_{g})/\NN(\frac{(p_{g})}{(q_{g})}I_{g})$ can be written as a polynomial $P$ of $z_{1},\dots,z_{k}$ of multi-degree $d_{1}-1,\dots,d_{k}-1$.
 Fix $1\leq i\leq k$. Recall that we defined $\tilde{f}(I_{i})$ to be an element whose $d_{i}$-th power is $f(x_{i}/y_{i})$ (see Proposition 1.5).
 If $P$ is not independent of the $i$-th coordinate, then for fixed $z_{j}, 1\leq j\leq k, j\neq i$, there are at most $d_{i}-1$ values of $z_{i}$ which make $P$ vanishes. So we may pick  $\tilde{f}(I_{i})$ to be a $d_{i}$-th root of $f(x_{i}/y_{i})$ so that $P$ and thus $M$ is non-zero. So we are done.

 If  $P$ is  independent of all the coordinate, then $P$ is a constant and clearly
 $M=\tilde{f}(\frac{(p_{0})}{(q_{0})}I_{0})/\NN(\frac{(p_{0})}{(q_{0})}I_{0})\neq 0$. So $M\neq  0$ and we are done.
 }
\end{proof}

\fi


%\subsection{Extensions of multiplicative functions}



%Therefore, our task is redued to find a classification for all the ideal domain completely multiplicative function $g\colon\mathfrak{I}(K)\to\S$ such that 
%$$\lim_{N\to\infty}\E_{\mathfrak{v}\in \mathfrak{I}(K)\colon \NN(\mathfrak{v})\leq N}g(\mathfrak{v})=0.$$
 
%To do so, for any multiplicative functions $f,g\colon\mathfrak{Q}(K)\to\C$, define
%$$\D(f,g):=\sum_{\mathfrak{p} \text{ is a prime ideal of } \mathcal{O}_{K}}\frac{1}{\NN(\mathfrak{p})}(1-\text{Re}f(\mathfrak{p})\overline{g}(\mathfrak{p})).$$
%For $f,g\colon K\to\C$, define
%$$\D_{0}(f,g):=\sum_{p \text{ is a prime element of } \mathcal{O}_{K}}\frac{1}{\NN(p)}(1-\text{Re}f(p)\overline{g}(p)).$$
%Here a \emph{prime element} is an element in $\mathcal{O}_{K}$ that can be be written as the product of two non-units in $\mathcal{O}_{K}$.

 
 

\iffalse

We say that a function $f\colon K\to\S$ is \emph{aperiodic} condition (i) of Proposition \ref{prop: dircharequivalence} holds.
With all the ingredients above, we have the following 
\begin{theorem}
    Let $K$ be an imaginary quadratic number field. For any completely multiplicative function $f\colon \mathcal{O}_K^{\times}\to\S^1$ with $f(\e)=1$ for all units of $\e$, either $f$ is aperiodic or $\D(\tilde{f\chi},\NN(\mathfrak{v})^{i\tau})<\infty$ for some Dirichlet character $\chi$, some extension $\tilde{f\chi}$ of $f\chi$, and some $\tau\in\R$.
    %$\D_{0}(f\chi,\NN(u)^{i\tau})<\infty$ for some Dirichlet character $\chi$ and some $\tau\in\R$. {\color{red} Note: this is not a dichotomy, as I think the two cases may overlap for non-PID. But I hope this is good enough for our purposes.}
\end{theorem}
\begin{proof}
Suppose that $f$ is a completely multiplicative function which is not aperiodic.  
    It follows from Proposition~\ref{prop: dircharequivalence} that there exists a Dirichlet character $\chi\colon \OK\to\C$ and some extension $\tilde{f\chi}$ of $f\chi$, we have that 
        $$\lim_{N\to\infty}\E_{\mathfrak{v}\in \mathfrak{I}(K)\colon \NN(\mathfrak{v})\leq N}\tilde{f\chi}(\mathfrak{v})=0.$$
Since $g$ is completely multiplicative, case (iii) of Conjecture \ref{c27k} cannot happen, and thus we have that $\D(\tilde{f\chi},\NN(\mathfrak{v})^{i\tau})<\infty$ for some $\tau\in\R$. %Restricting to the sum over principle prime ideals in $\D(\tilde{f\chi},\NN(\mathfrak{v})^{i\tau})$, this implies that
%  $\D_{0}(f\chi,\NN(u)^{i\tau})\leq \D(\tilde{f\chi},\NN(\mathfrak{v})^{i\tau})<\infty$.    {\color{red} Here I am not using the full strength of the fact that $\D(\tilde{f\chi},\NN(\mathfrak{v})^{i\tau})<\infty$, but I hope this is good enough for our purposes.}
\end{proof}

\fi
%{\color{red} in the references, sometimes the first name is wrriten in full while sometimes it is not}

\bibliographystyle{abbrv}
\bibliography{library}


%}
 
\end{document}
